%% Generated by Sphinx.
\def\sphinxdocclass{report}
\documentclass[letterpaper,10pt,english,openany,oneside]{sphinxmanual}
\ifdefined\pdfpxdimen
   \let\sphinxpxdimen\pdfpxdimen\else\newdimen\sphinxpxdimen
\fi \sphinxpxdimen=.75bp\relax
\ifdefined\pdfimageresolution
    \pdfimageresolution= \numexpr \dimexpr1in\relax/\sphinxpxdimen\relax
\fi
%% let collapsible pdf bookmarks panel have high depth per default
\PassOptionsToPackage{bookmarksdepth=5}{hyperref}

\PassOptionsToPackage{booktabs}{sphinx}
\PassOptionsToPackage{colorrows}{sphinx}

\PassOptionsToPackage{warn}{textcomp}
\usepackage[utf8]{inputenc}
\ifdefined\DeclareUnicodeCharacter
% support both utf8 and utf8x syntaxes
  \ifdefined\DeclareUnicodeCharacterAsOptional
    \def\sphinxDUC#1{\DeclareUnicodeCharacter{"#1}}
  \else
    \let\sphinxDUC\DeclareUnicodeCharacter
  \fi
  \sphinxDUC{00A0}{\nobreakspace}
  \sphinxDUC{2500}{\sphinxunichar{2500}}
  \sphinxDUC{2502}{\sphinxunichar{2502}}
  \sphinxDUC{2514}{\sphinxunichar{2514}}
  \sphinxDUC{251C}{\sphinxunichar{251C}}
  \sphinxDUC{2572}{\textbackslash}
\fi
\usepackage{cmap}
\usepackage[OT1]{fontenc}
\usepackage{amsmath,amssymb,amstext}
\usepackage{babel}





\usepackage[Bjarne]{fncychap}
\usepackage{sphinx}

\fvset{fontsize=auto}
\usepackage{geometry}


% Include hyperref last.
\usepackage{hyperref}
% Fix anchor placement for figures with captions.
\usepackage{hypcap}% it must be loaded after hyperref.
% Set up styles of URL: it should be placed after hyperref.
\urlstyle{same}

\addto\captionsenglish{\renewcommand{\contentsname}{Part I. Introduction}}

\usepackage{sphinxmessages}
\setcounter{tocdepth}{1}


\usepackage{amssymb}
\DeclareUnicodeCharacter{2713}{\checkmark}
\DeclareUnicodeCharacter{039B}{\ensuremath{\Lambda}}
\DeclareUnicodeCharacter{03C3}{\ensuremath{\sigma}}
\DeclareUnicodeCharacter{2264}{\ensuremath{\le}}
\DeclareUnicodeCharacter{2212}{\ensuremath{-}}
\DeclareUnicodeCharacter{207A}{\textsuperscript{+}}
\DeclareUnicodeCharacter{207B}{\textsuperscript{\ensuremath{-}}}
\DeclareUnicodeCharacter{00BB}{\ensuremath{\gg}}


\title{Priroda Unofficial Manual}
\date{Oct 17, 2026}
\release{}
\author{Program by Dimitri Laikov}
\newcommand{\sphinxlogo}{\vbox{}}
\renewcommand{\releasename}{}
\makeindex
\begin{document}

\ifdefined\shorthandoff
  \ifnum\catcode`\=\string=\active\shorthandoff{=}\fi
  \ifnum\catcode`\"=\active\shorthandoff{"}\fi
\fi

\pagestyle{empty}
\sphinxmaketitle
\pagestyle{plain}
\sphinxtableofcontents
\pagestyle{normal}
\phantomsection\label{\detokenize{index::doc}}


\sphinxAtStartPar
An unofficial user manual for Priroda, the quantum\sphinxhyphen{}chemistry program by Dimitri Laikov.

\begin{sphinxadmonition}{important}{Important:}
\sphinxAtStartPar
This manual is not an official Priroda document. It was not written or reviewed by the developer. It is compiled
from the program’s own files, the short keyword list that ships with it, and user reports from forums and user
pages; {\hyperref[\detokenize{sources:sources}]{\sphinxcrossref{\DUrole{std}{\DUrole{std-ref}{Sources}}}}} lists them all. When this manual and the program disagree, the program is right.
\end{sphinxadmonition}

\begin{sphinxadmonition}{warning}{Warning:}
\sphinxAtStartPar
Work in progress. Unless an entry says otherwise, nothing here has yet been checked against a running Priroda binary.
Each keyword entry gives its sources and its verification status.
\end{sphinxadmonition}

\sphinxstepscope


\chapter{About Priroda}
\label{\detokenize{start/about:about-priroda}}\label{\detokenize{start/about::doc}}
\sphinxAtStartPar
Priroda (Russian for “nature”) is a quantum\sphinxhyphen{}chemistry program for molecules written by Dimitri N. Laikov.
It began as a fast program for GGA density\sphinxhyphen{}functional calculations and grew to cover hybrid DFT, RI\sphinxhyphen{}MP2,
coupled\sphinxhyphen{}cluster methods, TDDFT, a scalar\sphinxhyphen{}relativistic Hamiltonian, NMR, and Laikov’s own semiempirical model.



\sphinxAtStartPar
The program is closed\sphinxhyphen{}source and has no complete official manual. Priroda’s own documentation is a short keyword
list, the “Input Summary” (\sphinxcode{\sphinxupquote{input.txt}}), which ships with the program. This manual collects that text, the example
files on the official site, and twenty years of user experience from Russian\sphinxhyphen{}language forums. It is not an official
document: the developer has not written or reviewed it. {\hyperref[\detokenize{sources:sources}]{\sphinxcrossref{\DUrole{std}{\DUrole{std-ref}{Sources}}}}} lists what it is based on. To download the
program, see {\hyperref[\detokenize{start/getting-started:obtaining-priroda}]{\sphinxcrossref{\DUrole{std}{\DUrole{std-ref}{Obtaining Priroda}}}}}.




\section{Main features}
\label{\detokenize{start/about:main-features}}\begin{description}
\sphinxlineitem{Density\sphinxhyphen{}functional theory}
\sphinxAtStartPar
GGA DFT in which the Coulomb and exchange\textendash{}correlation terms are evaluated with an auxiliary (fitting) basis.
This is the source of Priroda’s speed. Every orbital basis comes with a matched fitting basis.
GGA DFT has analytic second derivatives.




\sphinxAtStartPar
Hybrid functionals with exact exchange evaluated by the resolution of the identity, under
{\hyperref[\detokenize{reference/control:control-theory}]{\sphinxcrossref{\DUrole{std}{\DUrole{std-ref}{theory=ridft}}}}}.



\sphinxlineitem{Correlated methods}
\sphinxAtStartPar
RI\sphinxhyphen{}MP2 and, as undocumented or less\sphinxhyphen{}used options, conventional MP2\textendash{}MP4, CCSD and CCSD(T). Gradients exist for
MP2, CCSD and CCSD(T).



\sphinxlineitem{Excited states}
\sphinxAtStartPar
GGA TDDFT with excited\sphinxhyphen{}state gradients, so excited\sphinxhyphen{}state geometries can be optimized.



\sphinxlineitem{Relativistic calculations}
\sphinxAtStartPar
A scalar\sphinxhyphen{}relativistic Hamiltonian ({\hyperref[\detokenize{reference/control:control-four}]{\sphinxcrossref{\DUrole{std}{\DUrole{std-ref}{four=1}}}}}) with all\sphinxhyphen{}electron basis sets and a finite\sphinxhyphen{}nucleus
model. The user FAQ of one computing centre names this, together with basis sets for the whole periodic table, as
what sets Priroda apart. Effective core potentials (SBK) are also available.



\sphinxlineitem{Geometry and dynamics}
\sphinxAtStartPar
Minimization in redundant internal coordinates, transition\sphinxhyphen{}state search, IRC, relaxed scans, Hessians with
vibrational analysis and thermochemistry, and molecular dynamics with a thermostat.



\sphinxlineitem{Properties}
\sphinxAtStartPar
NMR shielding tensors and chemical shifts (also at points in space, for NICS), Hirshfeld and Mulliken charges, bond
orders, ESR parameters, critical points of the electron density, Raman intensities, and orbitals and normal modes
for Molden.



\sphinxlineitem{Semiempirical model}
\sphinxAtStartPar
Laikov’s parametrized “QM” model ({\hyperref[\detokenize{reference/control:control-theory}]{\sphinxcrossref{\DUrole{std}{\DUrole{std-ref}{theory=qm\_n3}}}}}), used for systems of thousands of atoms.



\sphinxlineitem{Parallel execution}
\sphinxAtStartPar
MPI parallelism. In practice it scales well within one node and poorly beyond it; see {\hyperref[\detokenize{start/running::doc}]{\sphinxcrossref{\DUrole{doc}{Installing and running}}}}.



\end{description}

\sphinxAtStartPar
Experienced users report that Priroda is markedly faster than general\sphinxhyphen{}purpose programs for routine GGA work. One
2018 comparison found it about three times faster per optimization step than ORCA 4.0 (PBE/cc\sphinxhyphen{}pVDZ, one core).




\section{Capability matrix}
\label{\detokenize{start/about:capability-matrix}}
\sphinxAtStartPar
The table shows which \sphinxcode{\sphinxupquote{task=}} values are known to work with each \sphinxcode{\sphinxupquote{theory=}} value. It is assembled from example inputs,
user reports and the test runs on version 25, not from a statement by the developer.
\begin{itemize}
\item {} 
\sphinxAtStartPar
\(\checkmark\): shown to work in an example input, a user report or a test run.

\item {} 
\sphinxAtStartPar
\textendash{}: not available, according to the sources.

\item {} 
\sphinxAtStartPar
?: unverified. No source shows it working or failing.

\end{itemize}


\begin{savenotes}\sphinxattablestart
\sphinxthistablewithglobalstyle
\centering
\begin{tabulary}{\linewidth}[t]{TTTTTTT}
\sphinxtoprule
\sphinxstyletheadfamily 
\sphinxAtStartPar
Task
&\sphinxstyletheadfamily 
\sphinxAtStartPar
\sphinxcode{\sphinxupquote{dft}}
&\sphinxstyletheadfamily 
\sphinxAtStartPar
\sphinxcode{\sphinxupquote{ridft}}
&\sphinxstyletheadfamily 
\sphinxAtStartPar
\sphinxcode{\sphinxupquote{tddft}}
&\sphinxstyletheadfamily 
\sphinxAtStartPar
\sphinxcode{\sphinxupquote{rimp2}}
&\sphinxstyletheadfamily 
\sphinxAtStartPar
\sphinxcode{\sphinxupquote{ccsd}} / \sphinxcode{\sphinxupquote{ccsd4t}}
&\sphinxstyletheadfamily 
\sphinxAtStartPar
\sphinxcode{\sphinxupquote{qm\_n3}}
\\
\sphinxmidrule
\sphinxtableatstartofbodyhook
\sphinxAtStartPar
\sphinxcode{\sphinxupquote{energy}}
&
\sphinxAtStartPar
\(\checkmark\)
&
\sphinxAtStartPar
\(\checkmark\)
&
\sphinxAtStartPar
\(\checkmark\)
&
\sphinxAtStartPar
\(\checkmark\)
&
\sphinxAtStartPar
\(\checkmark\)
&
\sphinxAtStartPar
\(\checkmark\)
\\
\sphinxhline
\sphinxAtStartPar
\sphinxcode{\sphinxupquote{gradient}}
&
\sphinxAtStartPar
\(\checkmark\)
&
\sphinxAtStartPar
\(\checkmark\)
&
\sphinxAtStartPar
\(\checkmark\)
&
\sphinxAtStartPar
\(\checkmark\)
&
\sphinxAtStartPar
\(\checkmark\) (1)
&
\sphinxAtStartPar
\(\checkmark\)
\\
\sphinxhline
\sphinxAtStartPar
\sphinxcode{\sphinxupquote{hessian}}
&
\sphinxAtStartPar
\(\checkmark\) (2)
&
\sphinxAtStartPar
\(\checkmark\) (3)
&
\sphinxAtStartPar
?
&
\sphinxAtStartPar
\(\checkmark\) (3)
&
\sphinxAtStartPar
\(\checkmark\) (3)
&
\sphinxAtStartPar
?
\\
\sphinxhline
\sphinxAtStartPar
\sphinxcode{\sphinxupquote{optimize}}
&
\sphinxAtStartPar
\(\checkmark\)
&
\sphinxAtStartPar
\(\checkmark\)
&
\sphinxAtStartPar
\(\checkmark\) (4)
&
\sphinxAtStartPar
\(\checkmark\)
&
\sphinxAtStartPar
\(\checkmark\) (1)
&
\sphinxAtStartPar
\(\checkmark\)
\\
\sphinxhline
\sphinxAtStartPar
TS (\sphinxcode{\sphinxupquote{saddle=1}})
&
\sphinxAtStartPar
\(\checkmark\)
&
\sphinxAtStartPar
?
&
\sphinxAtStartPar
?
&
\sphinxAtStartPar
?
&
\sphinxAtStartPar
?
&
\sphinxAtStartPar
?
\\
\sphinxhline
\sphinxAtStartPar
\sphinxcode{\sphinxupquote{irc}}
&
\sphinxAtStartPar
\(\checkmark\)
&
\sphinxAtStartPar
?
&
\sphinxAtStartPar
?
&
\sphinxAtStartPar
?
&
\sphinxAtStartPar
?
&
\sphinxAtStartPar
?
\\
\sphinxhline
\sphinxAtStartPar
\sphinxcode{\sphinxupquote{scan}}
&
\sphinxAtStartPar
\(\checkmark\)
&
\sphinxAtStartPar
?
&
\sphinxAtStartPar
?
&
\sphinxAtStartPar
?
&
\sphinxAtStartPar
?
&
\sphinxAtStartPar
?
\\
\sphinxhline
\sphinxAtStartPar
\sphinxcode{\sphinxupquote{nmr}}
&
\sphinxAtStartPar
\(\checkmark\) (5)
&
\sphinxAtStartPar
\textendash{}
&
\sphinxAtStartPar
\textendash{}
&
\sphinxAtStartPar
\textendash{}
&
\sphinxAtStartPar
\textendash{}
&
\sphinxAtStartPar
\textendash{}
\\
\sphinxhline
\sphinxAtStartPar
TDDFT excited states
&
\sphinxAtStartPar
\textendash{}
&
\sphinxAtStartPar
\textendash{} (6)
&
\sphinxAtStartPar
\(\checkmark\)
&
\sphinxAtStartPar
\textendash{}
&
\sphinxAtStartPar
\textendash{} (7)
&
\sphinxAtStartPar
\textendash{}
\\
\sphinxhline
\sphinxAtStartPar
MD (\sphinxcode{\sphinxupquote{thermostat}})
&
\sphinxAtStartPar
?
&
\sphinxAtStartPar
?
&
\sphinxAtStartPar
?
&
\sphinxAtStartPar
?
&
\sphinxAtStartPar
?
&
\sphinxAtStartPar
\(\checkmark\)
\\
\sphinxbottomrule
\end{tabulary}
\sphinxtableafterendhook\par
\sphinxattableend\end{savenotes}



\sphinxAtStartPar
Notes to the table:
\begin{enumerate}
\sphinxsetlistlabels{\arabic}{enumi}{enumii}{}{.}%
\item {} 
\sphinxAtStartPar
A CCSD optimization is shown in a user input. CCSD(T) gradients are implied by the developer’s 2011 note on a bug
fix that names the CCSD(T) gradient. 

\item {} 
\sphinxAtStartPar
Analytic second derivatives. 

\item {} 
\sphinxAtStartPar
Computed by finite differences of the gradient (\sphinxcode{\sphinxupquote{NUMERICAL COMPUTATION OF SECOND DERIVATIVES}}), which takes
many gradient calculations. Shown for \sphinxcode{\sphinxupquote{ridft}} and \sphinxcode{\sphinxupquote{rimp2}} in version 25 tests; users report CCSD and CCSD(T)
Hessians. 

\item {} 
\sphinxAtStartPar
Excited\sphinxhyphen{}state optimization works, but Priroda has no root following, so the state can switch during the
optimization. 

\item {} 
\sphinxAtStartPar
NMR works only with \sphinxcode{\sphinxupquote{theory=dft}} and a non\sphinxhyphen{}hybrid functional; with \sphinxcode{\sphinxupquote{ridft}}, version 25 silently computes only the
energy. Users report that relativistic (\sphinxcode{\sphinxupquote{four=1}}) NMR runs hang. 

\item {} 
\sphinxAtStartPar
There is no hybrid TDDFT. See {\hyperref[\detokenize{start/about:what-priroda-does-not-do}]{\sphinxcrossref{\DUrole{std}{\DUrole{std-ref}{What Priroda does not do}}}}}.

\item {} 
\sphinxAtStartPar
Excited states at the coupled\sphinxhyphen{}cluster level use a different \sphinxcode{\sphinxupquote{theory=}} value, \sphinxcode{\sphinxupquote{ccsdsd}} (linear\sphinxhyphen{}response CCSD, RHF
only), whose status is undocumented. 

\end{enumerate}

\sphinxAtStartPar
The scalar\sphinxhyphen{}relativistic option \sphinxcode{\sphinxupquote{four=1}} has been used with \sphinxcode{\sphinxupquote{dft}}, \sphinxcode{\sphinxupquote{ridft}} and \sphinxcode{\sphinxupquote{rimp2}}. The full list of \sphinxcode{\sphinxupquote{theory=}}
values, including the undocumented ones, is in {\hyperref[\detokenize{reference/control:control-theory}]{\sphinxcrossref{\DUrole{std}{\DUrole{std-ref}{\$control theory}}}}}.



\begin{sphinxadmonition}{warning}{Warning:}
\sphinxAtStartPar
Hybrid functionals work only with \sphinxcode{\sphinxupquote{theory=ridft}}. With \sphinxcode{\sphinxupquote{theory=dft}}, older versions silently drop the exact
exchange, and newer versions, including version 25, stop right after printing the basis sets with no error message
in the output. With \sphinxcode{\sphinxupquote{theory=tddft}}, version 25 still runs a hybrid without its exact exchange and gives wrong
energies with no warning. See {\hyperref[\detokenize{reference/control:control-theory}]{\sphinxcrossref{\DUrole{std}{\DUrole{std-ref}{\$control theory}}}}}.
\end{sphinxadmonition}




\section{What Priroda does not do}
\label{\detokenize{start/about:what-priroda-does-not-do}}\label{\detokenize{start/about:id1}}
\sphinxAtStartPar
Check this list before planning a project around Priroda.
\begin{description}
\sphinxlineitem{No use of point\sphinxhyphen{}group symmetry in older versions}
\sphinxAtStartPar
Older versions always run in C1. A symmetric starting structure stays symmetric during optimization, so a
Hessian is needed to show whether it is a minimum. Newer outputs (v13 onwards) print a detected point group, and
the official 2023\textendash{}2025 example inputs use \sphinxcode{\sphinxupquote{sym=1}} in \sphinxcode{\sphinxupquote{\$molecule}} and \sphinxcode{\sphinxupquote{\$optimize}}. These options are not documented.
In version 25 tests, \sphinxcode{\sphinxupquote{\$molecule sym=1,0.01}} symmetrized the structure and \sphinxcode{\sphinxupquote{\$optimize sym=1}} kept the symmetry
during the optimization; each option alone does only half of this. See {\hyperref[\detokenize{reference/molecule:molecule-sym}]{\sphinxcrossref{\DUrole{std}{\DUrole{std-ref}{\$molecule sym}}}}} and
{\hyperref[\detokenize{reference/optimize:optimize-symmetry}]{\sphinxcrossref{\DUrole{std}{\DUrole{std-ref}{\$optimize symmetry}}}}}.



\sphinxlineitem{No continuum solvation}
\sphinxAtStartPar
No version has a continuum solvent model. Solvent can only be included as explicit molecules.



\sphinxlineitem{No D3 or other empirical dispersion correction}
\sphinxAtStartPar
The developer has declined to add one (relayed from correspondence). Users add Grimme’s D3 a posteriori with the
external \sphinxcode{\sphinxupquote{dftd3}} program. The 2023\textendash{}2025 functionals \sphinxcode{\sphinxupquote{LPBEVV}} and \sphinxcode{\sphinxupquote{LPBEVVV}} may include dispersion, but this is
inferred from their test systems and not documented.



\sphinxlineitem{No periodic systems}
\sphinxAtStartPar
Priroda handles molecules and clusters only.



\sphinxlineitem{No dummy atoms, so no counterpoise (BSSE) correction}
\sphinxAtStartPar
Ghost centres cannot be placed in \sphinxcode{\sphinxupquote{\$molecule}}. For NICS, give the probe points in
{\hyperref[\detokenize{reference/nmr:nmr-points}]{\sphinxcrossref{\DUrole{std}{\DUrole{std-ref}{\$nmr points}}}}} instead.



\sphinxlineitem{No hybrid TDDFT}
\sphinxAtStartPar
TDDFT is GGA only. A hybrid functional with \sphinxcode{\sphinxupquote{theory=tddft}} does not give hybrid TDDFT: version 25 drops the exact
exchange without a warning, and the energies are wrong. 



\end{description}

\sphinxAtStartPar
Other gaps reported by users:
\begin{itemize}
\item {} 
\sphinxAtStartPar
no spin\textendash{}orbit coupling, so no phosphorescence calculations; 

\item {} 
\sphinxAtStartPar
no multireference methods (CASSCF, MRCI, CASPT2) that users could get working; 

\item {} 
\sphinxAtStartPar
no freezing of atoms in Cartesian space; only internal coordinates can be fixed
({\hyperref[\detokenize{reference/optimize:optimize-fix}]{\sphinxcrossref{\DUrole{std}{\DUrole{std-ref}{\$optimize fix}}}}}); 

\item {} 
\sphinxAtStartPar
no general job chaining in one input: \sphinxcode{\sphinxupquote{task=optimize+hessian}} is the only combined task, and two \sphinxcode{\sphinxupquote{task=}} lines run
only the first. 

\end{itemize}


\section{Version history}
\label{\detokenize{start/about:version-history}}
\sphinxAtStartPar
Priroda numbers its versions with integers and prints the build date in the output banner, for example
\sphinxcode{\sphinxupquote{Priroda 13 (2013.01.10)}}. Energies differ slightly between versions, so keep every calculation you compare on one
version.




\begin{savenotes}\sphinxattablestart
\sphinxthistablewithglobalstyle
\centering
\begin{tabulary}{\linewidth}[t]{TTT}
\sphinxtoprule
\sphinxstyletheadfamily 
\sphinxAtStartPar
Version
&\sphinxstyletheadfamily 
\sphinxAtStartPar
Build dates seen
&\sphinxstyletheadfamily 
\sphinxAtStartPar
Notes
\\
\sphinxmidrule
\sphinxtableatstartofbodyhook
\sphinxAtStartPar
2
&
\sphinxAtStartPar
2002.01.14 (v2.10)
&
\sphinxAtStartPar
Banner \sphinxcode{\sphinxupquote{program Priroda version 2.10}}. Had \sphinxcode{\sphinxupquote{task=optimize+hessian}}.
\\
\sphinxhline
\sphinxAtStartPar
4
&
\sphinxAtStartPar
2004.09.14 (v4.07)
&
\sphinxAtStartPar
The finite\sphinxhyphen{}nucleus model changed between versions 3 and 4, so pre\sphinxhyphen{}v4 relativistic basis sets must not be used with v4 and later.
\\
\sphinxhline
\sphinxAtStartPar
5
&
\sphinxAtStartPar
—
&
\sphinxAtStartPar
Version described by the “Input Summary” mirrored on the KNC wiki. Introduced the systematic L basis sets. Hybrid RI\sphinxhyphen{}DFT was implemented in November 2004.
\\
\sphinxhline
\sphinxAtStartPar
6
&
\sphinxAtStartPar
2006.08.20
&
\sphinxAtStartPar
Long the most widely used version: 32\sphinxhyphen{} and 64\sphinxhyphen{}bit Linux builds with MPI, and a Windows build (\sphinxcode{\sphinxupquote{p6.exe}}). \sphinxcode{\sphinxupquote{optimize+hessian}} was removed. Parallel Hessians use at most 2 cores.
\\
\sphinxhline
\sphinxAtStartPar
7
&
\sphinxAtStartPar
—
&
\sphinxAtStartPar
Sent as object files for linking against a local MPI library.
\\
\sphinxhline
\sphinxAtStartPar
8
&
\sphinxAtStartPar
2008.09.05
&
\sphinxAtStartPar
Built on request for a Myrinet cluster.
\\
\sphinxhline
\sphinxAtStartPar
9
&
\sphinxAtStartPar
2009\textendash{}2011
&
\sphinxAtStartPar
\sphinxcode{\sphinxupquote{optimize+hessian}} returned. In mid\sphinxhyphen{}2011 the developer fixed a long\sphinxhyphen{}standing bug that in rare cases gave wrong gradients for MP2, MP4, CCSD, CCSD(T), CIS(2) and EOM\sphinxhyphen{}CCSD.
\\
\sphinxhline
\sphinxAtStartPar
10
&
\sphinxAtStartPar
2011.04.14
&
\sphinxAtStartPar
Linux build \sphinxcode{\sphinxupquote{p}} with a bundled \sphinxcode{\sphinxupquote{mpiexec}}.
\\
\sphinxhline
\sphinxAtStartPar
11
&
\sphinxAtStartPar
—
&
\sphinxAtStartPar
Contains the semiempirical QM model.
\\
\sphinxhline
\sphinxAtStartPar
13
&
\sphinxAtStartPar
2013.01.10, 2013.01.21, 2013.09.13
&
\sphinxAtStartPar
Builds behind the example outputs on the official site.
\\
\sphinxhline
\sphinxAtStartPar
14
&
\sphinxAtStartPar
2014.01.01
&
\sphinxAtStartPar
Input geometry echoed with a \sphinxcode{\sphinxupquote{\textless{}}} prefix (\sphinxcode{\sphinxupquote{\textless{}\$molecule}}). 8060\sphinxhyphen{}atom QM\sphinxhyphen{}model run on 32 processes.
\\
\sphinxhline
\sphinxAtStartPar
15
&
\sphinxAtStartPar
2015.09.14
&
\sphinxAtStartPar
Scalar\sphinxhyphen{}relativistic example (Pt complex).
\\
\sphinxhline
\sphinxAtStartPar
16, 17
&
\sphinxAtStartPar
2016\textendash{}2018
&
\sphinxAtStartPar
The Windows build runs in parallel from version 16. Version 17 runs Hessians on 4 or more cores. Around 2016 the developer removed limits on SCF memory allocation and on the system size for Hessians.
\\
\sphinxhline
\sphinxAtStartPar
23
&
\sphinxAtStartPar
2023.06.22
&
\sphinxAtStartPar
Newest numbered version reported by users (2024). New functionals \sphinxcode{\sphinxupquote{LPBEVV}} and \sphinxcode{\sphinxupquote{LPBEVVV}}.
\\
\sphinxhline
\sphinxAtStartPar
25
&
\sphinxAtStartPar
2025.02.28
&
\sphinxAtStartPar
Banner \sphinxcode{\sphinxupquote{Priroda 25 (2025.02.28)}}. The version used for the test runs behind this manual. Its date matches the \sphinxcode{\sphinxupquote{p4\_d}} build on the official site; \sphinxcode{\sphinxupquote{p\_l}} is dated 2025\sphinxhyphen{}04\sphinxhyphen{}14. The 2025 example inputs use \sphinxcode{\sphinxupquote{theory=DFT\_n4}}.
\\
\sphinxbottomrule
\end{tabulary}
\sphinxtableafterendhook\par
\sphinxattableend\end{savenotes}



\sphinxAtStartPar
Version numbers missing from the table (3, 12, 18\textendash{}22, 24) do not appear in any source. The Russian file \sphinxcode{\sphinxupquote{Dft\_ansi.txt}}
describes Priroda’s predecessor, which used different group names (\sphinxcode{\sphinxupquote{\$coordinates}}, \sphinxcode{\sphinxupquote{\$SCF}}, \sphinxcode{\sphinxupquote{\$search}}, \sphinxcode{\sphinxupquote{\$force}}).
Those names are not part of the current keyword list.




\section{How to cite}
\label{\detokenize{start/about:how-to-cite}}\label{\detokenize{start/about:id2}}
\sphinxAtStartPar
Priroda has no single program paper for the current version. Users cite the papers below; the first three are the
list given by the KNC wiki FAQ and repeated in the forums.

\sphinxAtStartPar
Program:
\begin{itemize}
\item {} 
\sphinxAtStartPar
D. N. Laikov, Yu. A. Ustynyuk. PRIRODA\sphinxhyphen{}04: a quantum\sphinxhyphen{}chemical program suite. New possibilities in the study of
molecular systems with the application of parallel computing. \sphinxstyleemphasis{Russ. Chem. Bull., Int. Ed.} \sphinxstylestrong{54} (2005) 820\textendash{}826.
Russian original: \sphinxstyleemphasis{Izv. Akad. Nauk, Ser. Khim.} 2005, No. 3, 804\textendash{}810. 

\end{itemize}

\sphinxAtStartPar
Density\sphinxhyphen{}fitting DFT (cite for any DFT calculation):
\begin{itemize}
\item {} 
\sphinxAtStartPar
D. N. Laikov. Fast evaluation of density functional exchange\sphinxhyphen{}correlation terms using the expansion of the electron
density in auxiliary basis sets. \sphinxstyleemphasis{Chem. Phys. Lett.} \sphinxstylestrong{281} (1997) 151\textendash{}156. 

\end{itemize}

\sphinxAtStartPar
L\sphinxhyphen{}family basis sets (\sphinxcode{\sphinxupquote{L1}}, \sphinxcode{\sphinxupquote{L2}}, \sphinxcode{\sphinxupquote{L11}}, …):
\begin{itemize}
\item {} 
\sphinxAtStartPar
D. N. Laikov. A new class of atomic basis functions for accurate electronic structure calculations of molecules.
\sphinxstyleemphasis{Chem. Phys. Lett.} \sphinxstylestrong{416} (2005) 116\textendash{}120. 

\end{itemize}

\sphinxAtStartPar
The \sphinxcode{\sphinxupquote{3z}} basis set, which has no journal paper:
\begin{itemize}
\item {} 
\sphinxAtStartPar
D. N. Laikov. “Development of an economical approach to the density functional calculation of molecules and its
application to complex chemical problems”. PhD dissertation, Moscow State University, 2000 (in Russian).



\end{itemize}

\sphinxAtStartPar
Scalar\sphinxhyphen{}relativistic DFT (\sphinxcode{\sphinxupquote{four=1}}):
\begin{itemize}
\item {} 
\sphinxAtStartPar
D. N. Laikov. An implementation of the scalar relativistic density functional theory for molecular calculations
with Gaussian basis sets. DFT2000 Conference, Menton, France, 2000. 

\end{itemize}

\sphinxAtStartPar
Semiempirical QM model (\sphinxcode{\sphinxupquote{theory=qm\_n*}}):
\begin{itemize}
\item {} 
\sphinxAtStartPar
D. N. Laikov. \sphinxstyleemphasis{J. Chem. Phys.} \sphinxstylestrong{135} (2011) 134120. 



\end{itemize}

\sphinxAtStartPar
For version 6, users cited the program itself as “D. N. Laikov, PRIRODA, Electronic Structure Code, Version 6, 2006”.
A 2014 forum suggestion adds Laikov’s home page, \sphinxurl{http://rad.chem.msu.ru/~laikov/}, as the reference for the current
version.



\sphinxAtStartPar
Also cite the original papers for the functional, and for the Hamiltonian and nuclear model if you use \sphinxcode{\sphinxupquote{four=1}}.
{\hyperref[\detokenize{reference/dft:dft-functional}]{\sphinxcrossref{\DUrole{std}{\DUrole{std-ref}{\$dft functional}}}}} and {\hyperref[\detokenize{reference/control:control-four}]{\sphinxcrossref{\DUrole{std}{\DUrole{std-ref}{\$control four}}}}} give them.

\sphinxstepscope


\chapter{Installing and running}
\label{\detokenize{start/running:installing-and-running}}\label{\detokenize{start/running::doc}}
\sphinxAtStartPar
Priroda has no installer. You copy an executable, the basis\sphinxhyphen{}set files you need and, for parallel runs, the MPI
launcher into a directory, write an input file, and start the program with the input and output file names as
arguments.


\section{Executables}
\label{\detokenize{start/running:executables}}

\subsection{Linux}
\label{\detokenize{start/running:linux}}
\sphinxAtStartPar
The official site lists these Linux files. Sizes and dates are those of the 2025 file index.




\begin{savenotes}\sphinxattablestart
\sphinxthistablewithglobalstyle
\centering
\begin{tabulary}{\linewidth}[t]{TTT}
\sphinxtoprule
\sphinxstyletheadfamily 
\sphinxAtStartPar
File
&\sphinxstyletheadfamily 
\sphinxAtStartPar
Date
&\sphinxstyletheadfamily 
\sphinxAtStartPar
Description
\\
\sphinxmidrule
\sphinxtableatstartofbodyhook
\sphinxAtStartPar
\sphinxcode{\sphinxupquote{p}}
&
\sphinxAtStartPar
2023\sphinxhyphen{}06\sphinxhyphen{}25
&
\sphinxAtStartPar
Main x86\sphinxhyphen{}64 executable, statically linked.
\\
\sphinxhline
\sphinxAtStartPar
\sphinxcode{\sphinxupquote{p4}}
&
\sphinxAtStartPar
2023\sphinxhyphen{}06\sphinxhyphen{}25
&
\sphinxAtStartPar
Build “for FMA/AVX2” processors.
\\
\sphinxhline
\sphinxAtStartPar
\sphinxcode{\sphinxupquote{p\_d}}
&
\sphinxAtStartPar
2021\sphinxhyphen{}11\sphinxhyphen{}12
&
\sphinxAtStartPar
Dynamically linked build.
\\
\sphinxhline
\sphinxAtStartPar
\sphinxcode{\sphinxupquote{p4\_d}}
&
\sphinxAtStartPar
2025\sphinxhyphen{}02\sphinxhyphen{}28
&
\sphinxAtStartPar
Dynamically linked FMA/AVX2 build.
\\
\sphinxhline
\sphinxAtStartPar
\sphinxcode{\sphinxupquote{p\_l}}
&
\sphinxAtStartPar
2025\sphinxhyphen{}04\sphinxhyphen{}14
&
\sphinxAtStartPar
Build for a particular cluster, started through the site’s Intel MPI wrapper.
\\
\sphinxhline
\sphinxAtStartPar
\sphinxcode{\sphinxupquote{p\_mvs1p1}}, \sphinxcode{\sphinxupquote{p\_mvs10p}}, \sphinxcode{\sphinxupquote{p\_mvs10q}}
&
\sphinxAtStartPar
2019\textendash{}2023
&
\sphinxAtStartPar
Builds for particular clusters.
\\
\sphinxhline
\sphinxAtStartPar
\sphinxcode{\sphinxupquote{p\_arm}}, \sphinxcode{\sphinxupquote{mpiexec\_arm}}
&
\sphinxAtStartPar
2015\sphinxhyphen{}09\sphinxhyphen{}23
&
\sphinxAtStartPar
ARM build and its MPI launcher.
\\
\sphinxhline
\sphinxAtStartPar
\sphinxcode{\sphinxupquote{mpiexec}}
&
\sphinxAtStartPar
2012\sphinxhyphen{}09\sphinxhyphen{}09
&
\sphinxAtStartPar
MPI launcher shipped with \sphinxcode{\sphinxupquote{p}}.
\\
\sphinxhline
\sphinxAtStartPar
\sphinxcode{\sphinxupquote{xm}}
&
\sphinxAtStartPar
2025\sphinxhyphen{}04\sphinxhyphen{}24
&
\sphinxAtStartPar
The developer’s molecule viewer. It takes a Priroda input file as its argument (\sphinxcode{\sphinxupquote{xm C60.qm.in}}).
\\
\sphinxbottomrule
\end{tabulary}
\sphinxtableafterendhook\par
\sphinxattableend\end{savenotes}

\sphinxAtStartPar
That \sphinxcode{\sphinxupquote{p}} is statically linked is inferred from the description of \sphinxcode{\sphinxupquote{p\_d}} as “dynamically linked”. The site does not say
which cluster each \sphinxcode{\sphinxupquote{p\_mvs*}} and \sphinxcode{\sphinxupquote{p\_l}} build is for.



\sphinxAtStartPar
Older distributions used other names: \sphinxcode{\sphinxupquote{p6}}, \sphinxcode{\sphinxupquote{p6\_32}} and \sphinxcode{\sphinxupquote{p6\_64}} for version 6, \sphinxcode{\sphinxupquote{p9}}, and \sphinxcode{\sphinxupquote{p1}} for version 11 at one
site. Scripts and forum posts written for those versions use these names. This manual writes \sphinxcode{\sphinxupquote{p}} throughout.



\sphinxAtStartPar
Some users run the Linux build under WSL on Windows. 


\subsection{Windows}
\label{\detokenize{start/running:windows}}
\sphinxAtStartPar
The Windows build is a Cygwin build for x86\sphinxhyphen{}64. It is in the \sphinxcode{\sphinxupquote{exe64/}} directory of the official site:




\begin{savenotes}\sphinxattablestart
\sphinxthistablewithglobalstyle
\centering
\begin{tabulary}{\linewidth}[t]{TTT}
\sphinxtoprule
\sphinxstyletheadfamily 
\sphinxAtStartPar
File
&\sphinxstyletheadfamily 
\sphinxAtStartPar
Date
&\sphinxstyletheadfamily 
\sphinxAtStartPar
Description
\\
\sphinxmidrule
\sphinxtableatstartofbodyhook
\sphinxAtStartPar
\sphinxcode{\sphinxupquote{p1.exe}}
&
\sphinxAtStartPar
2021\sphinxhyphen{}10\sphinxhyphen{}01
&
\sphinxAtStartPar
Serial executable.
\\
\sphinxhline
\sphinxAtStartPar
\sphinxcode{\sphinxupquote{p1\_190306.exe}}
&
\sphinxAtStartPar
2019\sphinxhyphen{}03\sphinxhyphen{}12
&
\sphinxAtStartPar
Older serial build.
\\
\sphinxhline
\sphinxAtStartPar
\sphinxcode{\sphinxupquote{p\_190306.exe}}
&
\sphinxAtStartPar
2019\sphinxhyphen{}03\sphinxhyphen{}12
&
\sphinxAtStartPar
Older parallel (MPI) build.
\\
\sphinxhline
\sphinxAtStartPar
\sphinxcode{\sphinxupquote{xm.exe}}
&
\sphinxAtStartPar
2019\sphinxhyphen{}03\sphinxhyphen{}08
&
\sphinxAtStartPar
Viewer.
\\
\sphinxhline
\sphinxAtStartPar
\sphinxcode{\sphinxupquote{cygwin1.dll}}
&
\sphinxAtStartPar
2021\sphinxhyphen{}10\sphinxhyphen{}01
&
\sphinxAtStartPar
Cygwin runtime library.
\\
\sphinxbottomrule
\end{tabulary}
\sphinxtableafterendhook\par
\sphinxattableend\end{savenotes}

\sphinxAtStartPar
For serial runs you need only \sphinxcode{\sphinxupquote{p1.exe}} and \sphinxcode{\sphinxupquote{cygwin1.dll}} in the same directory. A full Cygwin installation is not
required, and no administrator rights are needed. For parallel runs see {\hyperref[\detokenize{start/running:running-windows-parallel}]{\sphinxcrossref{\DUrole{std}{\DUrole{std-ref}{Parallel runs on Windows}}}}}.



\begin{sphinxadmonition}{note}{Note:}
\sphinxAtStartPar
In 2017 users referred to the parallel Windows build as \sphinxcode{\sphinxupquote{exe64/p.exe}}. The current index has no \sphinxcode{\sphinxupquote{p.exe}}; the newest
parallel build there is \sphinxcode{\sphinxupquote{p\_190306.exe}}.
\end{sphinxadmonition}



\sphinxAtStartPar
Old 32\sphinxhyphen{}bit Windows builds could not allocate more than about 1 GB of memory; \sphinxcode{\sphinxupquote{malloc(...) failed}} appeared above that.
Users report that the 64\sphinxhyphen{}bit build has no such limit.




\section{Running a calculation}
\label{\detokenize{start/running:running-a-calculation}}

\subsection{Serial runs}
\label{\detokenize{start/running:serial-runs}}
\sphinxAtStartPar
Give the input file first and the output file second:

\begin{sphinxVerbatim}[commandchars=\\\{\}]
p input.in output.out
\end{sphinxVerbatim}

\sphinxAtStartPar
On Windows:

\begin{sphinxVerbatim}[commandchars=\\\{\}]
p1.exe input.in output.out
\end{sphinxVerbatim}



\sphinxAtStartPar
On Linux the current directory is usually not in \sphinxcode{\sphinxupquote{PATH}}, so start the program as \sphinxcode{\sphinxupquote{./p}} or with its full path. If the
program cannot find the input or the basis file, give full paths everywhere: for the executable, the input, the
output, and the basis file inside the input.



\sphinxAtStartPar
Run in a directory you own and can write to. Keep file and directory names to Latin letters, digits and simple
punctuation, with no spaces or Cyrillic. This also applies to the computer name on Windows. File names are
case\sphinxhyphen{}sensitive on Linux, although keywords are not.



\sphinxAtStartPar
One user of an older version saw the complaint \sphinxcode{\sphinxupquote{0: output file \textquotesingle{}tasK.out\textquotesingle{}?}} when the output file name contained a
capital letter. Version 25 writes to an output file with capital letters in its name without complaint.



\sphinxAtStartPar
If no output file is given (\sphinxcode{\sphinxupquote{p input.in}}), the output goes to standard output. 


\subsection{Parallel runs on Linux}
\label{\detokenize{start/running:parallel-runs-on-linux}}
\sphinxAtStartPar
Priroda is parallelized with MPI. Start a parallel run through an MPI launcher. The official site shows the
\sphinxcode{\sphinxupquote{mpiexec}} that it ships next to \sphinxcode{\sphinxupquote{p}}:

\begin{sphinxVerbatim}[commandchars=\\\{\}]
./mpiexec \PYGZhy{}np 4 ./p input.in output.out
\end{sphinxVerbatim}



\sphinxAtStartPar
The output then starts with \sphinxcode{\sphinxupquote{parallel on 4 processors}} and one \sphinxcode{\sphinxupquote{node N:}} line per process. A two\sphinxhyphen{}process test gave
exactly the serial energy.



\sphinxAtStartPar
Older versions also accepted \sphinxcode{\sphinxupquote{p \sphinxhyphen{}np 4 input.in output.out}}, which users ran with versions 6 and 10. \sphinxstylestrong{Version 25 does
not}: it takes \sphinxcode{\sphinxupquote{\sphinxhyphen{}np}} as the input file name, prints \sphinxcode{\sphinxupquote{input file \textquotesingle{}\sphinxhyphen{}np\textquotesingle{}?.}} and stops. One site ran version 14 as
\sphinxcode{\sphinxupquote{mpiexec \sphinxhyphen{}np 4 p input.in}}, with the system’s MPI launcher.



\sphinxAtStartPar
Check the start of the output. A parallel run prints the number of processes and one line per process:

\begin{sphinxVerbatim}[commandchars=\\\{\}]
 parallel on  4 processors
 node  0: \PYGZsq{}linux\PYGZhy{}ldmm\PYGZsq{}
 node  1: \PYGZsq{}linux\PYGZhy{}ldmm\PYGZsq{}
 ...
\end{sphinxVerbatim}



\sphinxAtStartPar
The cluster builds are started through the cluster’s own MPI. The official example for \sphinxcode{\sphinxupquote{p\_l}} under SLURM is:

\begin{sphinxVerbatim}[commandchars=\\\{\}]
sbatch \PYGZhy{}p test \PYGZhy{}n 16 \PYGZhy{}t 6 impi \PYGZti{}/bin/p\PYGZus{}l file.\PYGZob{}in,out\PYGZcb{}
\end{sphinxVerbatim}

\sphinxAtStartPar
Here \sphinxcode{\sphinxupquote{impi}} is that cluster’s wrapper script for Intel MPI, not part of Priroda.




\subsection{Parallel runs on Windows}
\label{\detokenize{start/running:parallel-runs-on-windows}}\label{\detokenize{start/running:running-windows-parallel}}
\sphinxAtStartPar
The official site shows this command for a two\sphinxhyphen{}process run:

\begin{sphinxVerbatim}[commandchars=\\\{\}]
mpirun.exe \PYGZhy{}np 2 p.exe C60.qm.in C60.qm.w.2.out
\end{sphinxVerbatim}

\sphinxAtStartPar
The MPI runtime is not on the site. Users assembled it from Cygwin as follows:


\begin{enumerate}
\sphinxsetlistlabels{\arabic}{enumi}{enumii}{}{.}%
\item {} 
\sphinxAtStartPar
Download the 64\sphinxhyphen{}bit Cygwin installer (\sphinxcode{\sphinxupquote{setup\sphinxhyphen{}x86\_64.exe}} from \sphinxurl{https://cygwin.com}). Install the default packages
plus the \sphinxcode{\sphinxupquote{openmpi}} package.

\item {} 
\sphinxAtStartPar
Copy these files from the Cygwin installation into the directory with the Priroda executable:
\sphinxcode{\sphinxupquote{cygwin1.dll}}, \sphinxcode{\sphinxupquote{cyggcc\_s\sphinxhyphen{}seh\sphinxhyphen{}1.dll}}, \sphinxcode{\sphinxupquote{cygmpi\sphinxhyphen{}12.dll}}, \sphinxcode{\sphinxupquote{cygopen\sphinxhyphen{}pal\sphinxhyphen{}13.dll}}, \sphinxcode{\sphinxupquote{cygopen\sphinxhyphen{}rte\sphinxhyphen{}12.dll}}, \sphinxcode{\sphinxupquote{orterun.exe}}.
Newer OpenMPI packages have \sphinxcode{\sphinxupquote{cygopen\sphinxhyphen{}pal\sphinxhyphen{}40.dll}} and \sphinxcode{\sphinxupquote{cygopen\sphinxhyphen{}rte\sphinxhyphen{}40.dll}} instead of the \sphinxcode{\sphinxupquote{\sphinxhyphen{}13}}/\sphinxcode{\sphinxupquote{\sphinxhyphen{}12}} files; take
those. Take \sphinxcode{\sphinxupquote{cygwin1.dll}} from the same installation, because the other DLLs work only with matching versions.
After this, Cygwin itself is no longer needed.

\item {} 
\sphinxAtStartPar
Create a directory \sphinxcode{\sphinxupquote{dev}} one level above the directory you start Priroda from, with subdirectories \sphinxcode{\sphinxupquote{shm}} and
\sphinxcode{\sphinxupquote{mqueue}}. For Priroda in \sphinxcode{\sphinxupquote{C:\textbackslash{}Priroda}}, that means \sphinxcode{\sphinxupquote{C:\textbackslash{}dev\textbackslash{}shm}} and \sphinxcode{\sphinxupquote{C:\textbackslash{}dev\textbackslash{}mqueue}}. Without a Cygwin installation,
the Cygwin libraries take the parent of the start directory as the root \sphinxcode{\sphinxupquote{/}}.

\item {} 
\sphinxAtStartPar
If the program is in a system directory such as \sphinxcode{\sphinxupquote{C:\textbackslash{}Program Files}}, run it as administrator.

\item {} 
\sphinxAtStartPar
Set \sphinxcode{\sphinxupquote{TMPDIR}} and start the run from \sphinxcode{\sphinxupquote{cmd.exe}}:

\begin{sphinxVerbatim}[commandchars=\\\{\}]
cd /d \PYGZdq{}C:\PYGZbs{}Priroda\PYGZdq{}
set TMPDIR=\PYGZdq{}\PYGZpc{}TEMP\PYGZpc{}\PYGZdq{}
orterun.exe \PYGZhy{}np 2 p.exe \PYGZdq{}C:\PYGZbs{}Work\PYGZbs{}file.in\PYGZdq{} \PYGZdq{}C:\PYGZbs{}Work\PYGZbs{}file.out\PYGZdq{}
\end{sphinxVerbatim}

\end{enumerate}



\sphinxAtStartPar
Problems reported with this setup:
\begin{itemize}
\item {} 
\sphinxAtStartPar
The recipe needs a Priroda build that supports parallel runs. Version 6 (\sphinxcode{\sphinxupquote{p6.exe}}) failed with an OpenMPI
\sphinxcode{\sphinxupquote{agent\sphinxhyphen{}not\sphinxhyphen{}found}} error; users moved to version 17. 

\item {} 
\sphinxAtStartPar
On a computer connected by Wi\sphinxhyphen{}Fi that also has a wired network adapter, OpenMPI stops with
\sphinxcode{\sphinxupquote{tcp\_peer\_send\_blocking: send() to socket 12 failed: Transport endpoint is not connected (128)}}. Disable the wired
adapter. 

\item {} 
\sphinxAtStartPar
With the Ethernet cable unplugged, one user’s run used only two of the cores. 

\item {} 
\sphinxAtStartPar
\sphinxcode{\sphinxupquote{Warning: \textquotesingle{}/dev/shm\textquotesingle{} does not exists or is not a directory.}} means the Cygwin environment is missing. Create the
\sphinxcode{\sphinxupquote{dev}} directories (step 3), or start from a Cygwin shell. 

\end{itemize}

\begin{sphinxadmonition}{warning}{Warning:}
\sphinxAtStartPar
A user reported in 2021 that this recipe no longer worked on Windows 10 with a current Cygwin, and no fix was posted.
Expect to adapt it to the Cygwin and OpenMPI versions you have.
\end{sphinxadmonition}




\section{Memory, disk and scratch: \sphinxstyleliteralintitle{\sphinxupquote{\$system}}}
\label{\detokenize{start/running:memory-disk-and-scratch-system}}
\sphinxAtStartPar
The {\hyperref[\detokenize{reference/system:group-system}]{\sphinxcrossref{\DUrole{std}{\DUrole{std-ref}{\$system}}}}} group sets the resources. Every limit applies \sphinxstylestrong{per MPI process}.
\sphinxcode{\sphinxupquote{disk=10}} on four processes can use up to 40 GB.



\begin{sphinxVerbatim}[commandchars=\\\{\}]
\PYGZdl{}system memory=2048 disk=\PYGZhy{}2048 path=/scratch/user \PYGZdl{}end
\end{sphinxVerbatim}
\begin{description}
\sphinxlineitem{{\hyperref[\detokenize{reference/system:system-memory}]{\sphinxcrossref{\DUrole{std}{\DUrole{std-ref}{memory}}}}}}
\sphinxAtStartPar
Memory in MB. Two values set rank 0 and the other ranks separately: \sphinxcode{\sphinxupquote{mem=2048,256}} gives the master process
2048 MB and every other process 256 MB. In a parallel run, most of the memory and all of the disk storage are
used by rank 0, so give the other ranks less.



\sphinxlineitem{{\hyperref[\detokenize{reference/system:system-disk}]{\sphinxcrossref{\DUrole{std}{\DUrole{std-ref}{disk}}}}}}
\sphinxAtStartPar
A positive value is scratch disk in \sphinxstylestrong{GB}. The program opens one scratch file per gigabyte. A negative value is
scratch kept in memory, in \sphinxstylestrong{MB} (a RAM disk): \sphinxcode{\sphinxupquote{disk=\sphinxhyphen{}2048}} holds 2048 MB of scratch in RAM. A second value
applies to the other ranks, as for \sphinxcode{\sphinxupquote{memory}}; the official example inputs typically give them \sphinxcode{\sphinxupquote{\sphinxhyphen{}1}}.



\sphinxlineitem{{\hyperref[\detokenize{reference/system:system-path}]{\sphinxcrossref{\DUrole{std}{\DUrole{std-ref}{path}}}}}}
\sphinxAtStartPar
The scratch directory. A colon\sphinxhyphen{}separated list gives one directory per process, which pays off only when the
directories are on different physical drives. For several processes and one fast disk, a single directory is
enough.



\end{description}


\subsection{Keeping scratch in memory}
\label{\detokenize{start/running:keeping-scratch-in-memory}}
\sphinxAtStartPar
With a negative \sphinxcode{\sphinxupquote{disk}}, \sphinxcode{\sphinxupquote{memory}} plus |\sphinxcode{\sphinxupquote{disk}}| must fit into the physical RAM for each process, or the run stops at
once with \sphinxcode{\sphinxupquote{malloc(...) failed}}. One user’s setting for a 24 GB node is \sphinxcode{\sphinxupquote{memory=9000 disk=\sphinxhyphen{}14000}}. Keeping the scratch
in RAM is faster than any disk. In one RI\sphinxhyphen{}MP2 benchmark on 8 cores it cut the time from 86 minutes (scratch on disk)
to 12.4 minutes.




\subsection{Where scratch files go}
\label{\detokenize{start/running:where-scratch-files-go}}
\sphinxAtStartPar
The scratch directory is taken from \sphinxcode{\sphinxupquote{path=}}, then from the \sphinxcode{\sphinxupquote{TMPDIR}} environment variable. The files are named
\sphinxcode{\sphinxupquote{tmp\textless{}number\textgreater{}.\textless{}rank\textgreater{}}}, for example \sphinxcode{\sphinxupquote{/tmp/tmp2960.0}}. Priroda deletes them at the end of a normal run; after a crash,
delete the leftovers yourself.



\sphinxAtStartPar
When neither \sphinxcode{\sphinxupquote{path=}} nor \sphinxcode{\sphinxupquote{TMPDIR}} is set, version 25 puts scratch files in \sphinxcode{\sphinxupquote{/tmp}} (echo \sphinxcode{\sphinxupquote{Disk:   1g \textquotesingle{}/tmp\textquotesingle{}}}), as
version 6 did. \sphinxcode{\sphinxupquote{input.txt}} says the current working directory instead, and a 2020 user report the directory of the
input file; neither matches version 25. \sphinxcode{\sphinxupquote{/tmp}} is often small or shared, so set \sphinxcode{\sphinxupquote{path=}} or \sphinxcode{\sphinxupquote{TMPDIR}} for real jobs.



\sphinxAtStartPar
On Windows there is no \sphinxcode{\sphinxupquote{/tmp}}. Without \sphinxcode{\sphinxupquote{path=}} or \sphinxcode{\sphinxupquote{TMPDIR}}, older versions stopped after printing
\sphinxcode{\sphinxupquote{0: opening file /tmp/tmp2960.0}}. Use \sphinxcode{\sphinxupquote{\$system path=. \$end}} (the current directory) or set \sphinxcode{\sphinxupquote{TMPDIR}}.




\subsection{Disk speed}
\label{\detokenize{start/running:disk-speed}}
\begin{sphinxadmonition}{warning}{Warning:}
\sphinxAtStartPar
Never put scratch files on NFS or another network file system. On most clusters the home directories are on NFS;
use the node\sphinxhyphen{}local scratch disk instead.
\end{sphinxadmonition}



\sphinxAtStartPar
Priroda is sensitive to disk speed. Priroda’s own documentation (\sphinxcode{\sphinxupquote{input.txt}}) warns that for all methods except GGA DFT, performance
“will be degraded terribly” if more than one process reads and writes large files on the same physical drive. In
order of preference:


\begin{enumerate}
\sphinxsetlistlabels{\arabic}{enumi}{enumii}{}{.}%
\item {} 
\sphinxAtStartPar
keep the scratch in RAM (negative \sphinxcode{\sphinxupquote{disk}}) if it fits;

\item {} 
\sphinxAtStartPar
give each process its own physical drive with \sphinxcode{\sphinxupquote{path=a:b:c:d}};

\item {} 
\sphinxAtStartPar
put all processes on one fast local disk (SSD), accepting some loss.

\end{enumerate}

\sphinxAtStartPar
If the disk is the bottleneck for RI\sphinxhyphen{}MP2 or RI\sphinxhyphen{}DFT, {\hyperref[\detokenize{reference/integral:integral-direct}]{\sphinxcrossref{\DUrole{std}{\DUrole{std-ref}{\$integral direct=1}}}}} recomputes the
integrals instead of storing them. One user found it uses about a third of the memory but runs about ten times
longer in CPU time. In a disk\sphinxhyphen{}bound RI\sphinxhyphen{}MP2 optimization another user got 20 steps overnight with it instead of 2.



\sphinxAtStartPar
Users also recommend more RAM over faster disks. Scratch on a separate SSD helps, but RAM matters more.




\subsection{When the limits are too small}
\label{\detokenize{start/running:when-the-limits-are-too-small}}
\sphinxAtStartPar
Priroda does not fall back to a slower algorithm when memory or scratch runs out. The run stops with an error such as
\sphinxcode{\sphinxupquote{need … bytes.}} (the \sphinxcode{\sphinxupquote{memory}} limit is too small), \sphinxcode{\sphinxupquote{malloc(536870912) failed}} (the machine cannot supply the memory)
or \sphinxcode{\sphinxupquote{0: out of disk space (VNNNN), 184441541KB}} (scratch). The number in the disk
message is the amount needed. Raise the limit in \sphinxcode{\sphinxupquote{\$system}}, move to a node with more resources, or choose a cheaper
method (RI\sphinxhyphen{}MP2 instead of conventional MP2).



\sphinxAtStartPar
If you omit \sphinxcode{\sphinxupquote{\$system}}, version 6 printed \sphinxcode{\sphinxupquote{Memory: 128 (128) MB}} and \sphinxcode{\sphinxupquote{Disk: 1 GB \textquotesingle{}/tmp\textquotesingle{}}}. Version 25 prints
\sphinxcode{\sphinxupquote{Memory: 256m + 512m (256m + 512m)}} and no \sphinxcode{\sphinxupquote{Disk:}} line: 256 MB of memory and a 512 MB RAM disk per process.
That is enough for small molecules only, so always give \sphinxcode{\sphinxupquote{\$system}} for anything larger than a test.


\phantomsection\label{\detokenize{start/running:running-parallel-scaling}}
\sphinxAtStartPar
If a run fails for lack of memory or disk, see {\hyperref[\detokenize{troubleshooting/common-problems:troubleshooting-memory}]{\sphinxcrossref{\DUrole{std}{\DUrole{std-ref}{Memory, disk and scratch errors}}}}}.


\section{Parallel performance}
\label{\detokenize{start/running:parallel-performance}}
\sphinxAtStartPar
Users consistently report the following:


\begin{itemize}
\item {} 
\sphinxAtStartPar
\sphinxstylestrong{DFT scales well within one node, up to about 6\textendash{}8 cores.} Beyond 8 cores the gain is at most about 10\%, even for
large systems. One 38\sphinxhyphen{}step optimization took 101, 63 and 37 s on single nodes with 8, 12 and 24 cores. A small
PBE/L1 optimization with version 10 took 81.5, 42.5, 29.8 and 29.3 s on 1\textendash{}4 cores.



\item {} 
\sphinxAtStartPar
\sphinxstylestrong{Parallel results equal serial ones.} In version 25 tests on 2 processes, a water energy and a numerical
\sphinxcode{\sphinxupquote{ridft}} Hessian gave exactly the serial numbers; the Hessian took 4.9 s instead of 7.2 s.



\item {} 
\sphinxAtStartPar
\sphinxstylestrong{Running across nodes gives no gain for DFT.} Use one node per job.



\item {} 
\sphinxAtStartPar
\sphinxstylestrong{Memory per process does not drop} when you add processes, and in a parallel Hessian nearly all of the disk use is
on rank 0.



\item {} 
\sphinxAtStartPar
\sphinxstylestrong{RI\sphinxhyphen{}MP2 parallelizes better than DFT}, at least in memory. The v5 documentation limits useful parallelism to
Nproc ≤ min(N\_atoms, N\_occ) and recommends a number of processes that divides the number of occupied orbitals
nearly evenly. For N\_occ = 14 it rates 2, 7 or 14 processes as best, 3 or 5 as acceptable, 6 as poor and 13 as very
bad.



\item {} 
\sphinxAtStartPar
\sphinxstylestrong{Small RI jobs on many processes can hang.} RI\sphinxhyphen{}MP2 and RI\sphinxhyphen{}DFT hang with no message when a process receives no
work. The output then shows a line \sphinxcode{\sphinxupquote{0 blocks ... on node N}} after \sphinxcode{\sphinxupquote{integral estimates}}. Run small RI jobs on one or a
few cores.



\item {} 
\sphinxAtStartPar
\sphinxstylestrong{Several independent jobs often beat one parallel job.} On a 4\sphinxhyphen{}core machine, four single\sphinxhyphen{}core jobs ran about 30\%
faster in total than the same jobs run one after another on 4 cores each.



\end{itemize}

\sphinxAtStartPar
Results can differ in the last digits between runs on different numbers of processes. An unstable TS search was once
seen to follow different paths for different \sphinxcode{\sphinxupquote{\sphinxhyphen{}np}}.



\begin{sphinxadmonition}{note}{Note:}
\sphinxAtStartPar
Older versions scaled worse. In version 6 a parallel Hessian could use at most 2 cores and could hang on 4, and a
version 7 MP2 run hung on more than about 24 cores. Use version 17 or later for parallel work.
\end{sphinxadmonition}



\sphinxAtStartPar
Hangs, crashes and poor speed\sphinxhyphen{}up in parallel runs are covered in {\hyperref[\detokenize{troubleshooting/common-problems:troubleshooting-parallel}]{\sphinxcrossref{\DUrole{std}{\DUrole{std-ref}{Parallel and MPI problems}}}}}.


\section{Batch systems}
\label{\detokenize{start/running:batch-systems}}
\sphinxAtStartPar
Priroda needs nothing special from a batch system. The job script has to:
\begin{enumerate}
\sphinxsetlistlabels{\arabic}{enumi}{enumii}{}{.}%
\item {} 
\sphinxAtStartPar
request \sphinxstylestrong{one node} and the number of cores you will use (see {\hyperref[\detokenize{start/running:running-parallel-scaling}]{\sphinxcrossref{\DUrole{std}{\DUrole{std-ref}{parallel performance above}}}}});

\item {} 
\sphinxAtStartPar
point the scratch to the node’s local disk or to RAM, by exporting \sphinxcode{\sphinxupquote{TMPDIR}} or by \sphinxcode{\sphinxupquote{path=}} in \sphinxcode{\sphinxupquote{\$system}};

\item {} 
\sphinxAtStartPar
make the memory request cover \sphinxcode{\sphinxupquote{memory}} + |negative \sphinxcode{\sphinxupquote{disk}}| for every process;

\item {} 
\sphinxAtStartPar
start Priroda with the input and output file names.

\end{enumerate}



\sphinxAtStartPar
A SLURM script along these lines should work. Directory names and module setup are site\sphinxhyphen{}specific.

\begin{sphinxVerbatim}[commandchars=\\\{\}]
\PYGZsh{}!/bin/bash
\PYGZsh{}SBATCH \PYGZhy{}\PYGZhy{}job\PYGZhy{}name=mol
\PYGZsh{}SBATCH \PYGZhy{}\PYGZhy{}nodes=1
\PYGZsh{}SBATCH \PYGZhy{}\PYGZhy{}ntasks=8
\PYGZsh{}SBATCH \PYGZhy{}\PYGZhy{}time=24:00:00

PRIRODA=\PYGZdl{}HOME/priroda/p                      \PYGZsh{} full path to the executable
cd \PYGZdq{}\PYGZdl{}SLURM\PYGZus{}SUBMIT\PYGZus{}DIR\PYGZdq{}
export TMPDIR=/local/scratch/\PYGZdl{}SLURM\PYGZus{}JOB\PYGZus{}ID   \PYGZsh{} node\PYGZhy{}local disk; site\PYGZhy{}specific
mkdir \PYGZhy{}p \PYGZdq{}\PYGZdl{}TMPDIR\PYGZdq{}

\PYGZdl{}HOME/priroda/mpiexec \PYGZhy{}np \PYGZdq{}\PYGZdl{}SLURM\PYGZus{}NTASKS\PYGZdq{} \PYGZdq{}\PYGZdl{}PRIRODA\PYGZdq{} mol.in mol.out   \PYGZsh{} the mpiexec from the official site

rm \PYGZhy{}rf \PYGZdq{}\PYGZdl{}TMPDIR\PYGZdq{}
\end{sphinxVerbatim}



\sphinxAtStartPar
Under PBS/Torque the same steps read the number of cores from \sphinxcode{\sphinxupquote{\$PBS\_NODEFILE}}. PBS often provides a per\sphinxhyphen{}job local
scratch directory and deletes it after the job.



\begin{sphinxVerbatim}[commandchars=\\\{\}]
\PYGZsh{}!/bin/bash
\PYGZsh{}PBS \PYGZhy{}N mol
\PYGZsh{}PBS \PYGZhy{}l nodes=1:ppn=8,walltime=24:00:00

cd \PYGZdq{}\PYGZdl{}PBS\PYGZus{}O\PYGZus{}WORKDIR\PYGZdq{}
NP=\PYGZdl{}(wc \PYGZhy{}l \PYGZlt{} \PYGZdq{}\PYGZdl{}PBS\PYGZus{}NODEFILE\PYGZdq{})
export TMPDIR=/local/scratch/\PYGZdl{}PBS\PYGZus{}JOBID      \PYGZsh{} site\PYGZhy{}specific
mkdir \PYGZhy{}p \PYGZdq{}\PYGZdl{}TMPDIR\PYGZdq{}

\PYGZdl{}HOME/priroda/mpiexec \PYGZhy{}np \PYGZdq{}\PYGZdl{}NP\PYGZdq{} \PYGZdl{}HOME/priroda/p mol.in mol.out

rm \PYGZhy{}rf \PYGZdq{}\PYGZdl{}TMPDIR\PYGZdq{}
\end{sphinxVerbatim}



\sphinxAtStartPar
If \sphinxcode{\sphinxupquote{/tmp}} on the node is small or has a quota, a large job dies writing a scratch file, for example
\sphinxcode{\sphinxupquote{0: write file /tmp/tmp9633.11 (leAX): ...}}. Point the scratch at the cluster’s scratch file system or at RAM.



\begin{sphinxadmonition}{note}{Note:}
\sphinxAtStartPar
\sphinxstylestrong{Historical: the KNC wrappers.} The Kazan Scientific Centre clusters (2006\textendash{}2012) provided the wrappers
\sphinxcode{\sphinxupquote{pbspriroda5}}, \sphinxcode{\sphinxupquote{pbspriroda5\sphinxhyphen{}mpi}} and \sphinxcode{\sphinxupquote{pbspriroda6}}, called as \sphinxcode{\sphinxupquote{pbspriroda6 file\_name \sphinxhyphen{}time hh:mm:ss \sphinxhyphen{}queue queue \sphinxhyphen{}mem MB \sphinxhyphen{}disk size \sphinxhyphen{}procs 1:ppn=N}}. The wrappers wrote the \sphinxcode{\sphinxupquote{\$system}} group themselves from \sphinxcode{\sphinxupquote{\sphinxhyphen{}mem}} and \sphinxcode{\sphinxupquote{\sphinxhyphen{}disk}}, so
inputs submitted through them had to omit \sphinxcode{\sphinxupquote{\$system}}. Forum posts and old FAQs that mention these commands refer to
that site only.
\end{sphinxadmonition}




\section{Input file requirements}
\label{\detokenize{start/running:input-file-requirements}}
\sphinxAtStartPar
The input is a plain text file. Write it in a plain\sphinxhyphen{}text editor, not a word processor.
\begin{itemize}
\item {} 
\sphinxAtStartPar
\sphinxstylestrong{Plain ASCII.} Keywords, values and file names should contain ASCII characters only. Do not save the file as
Unicode.



\item {} 
\sphinxAtStartPar
\sphinxstylestrong{Unix line endings (LF).} This holds on Windows too. Saving from Word as “plain text” still leaves Windows (CRLF)
line endings, and some editors need to be told explicitly to write Unix line endings.



\item {} 
\sphinxAtStartPar
\sphinxstylestrong{Straight quotes.} Write \sphinxcode{\sphinxupquote{basis="3z.bas"}} with the ASCII double quote. Typographic quotes, as in \sphinxcode{\sphinxupquote{basis=»3z.bas»}},
break the basis\sphinxhyphen{}file name.



\end{itemize}

\sphinxAtStartPar
An input with CRLF line endings fails in \sphinxcode{\sphinxupquote{\$molecule}}, because the carriage return becomes part of the keyword. The
output stops after a line like this, where the keyword \sphinxcode{\sphinxupquote{z\sphinxhyphen{}matrix}} or \sphinxcode{\sphinxupquote{cartesian}} is shown garbled by the carriage
return:

\begin{sphinxVerbatim}[commandchars=\\\{\}]
 molecule input: \PYGZsq{}\PYGZsq{}
\PYGZsq{} ?matrix
\end{sphinxVerbatim}

\sphinxAtStartPar
Variants seen are \sphinxcode{\sphinxupquote{?rtesian}} for \sphinxcode{\sphinxupquote{cartesian}}, \sphinxcode{\sphinxupquote{\textquotesingle{}z\sphinxhyphen{}matrix\textquotesingle{} ?}}, and \sphinxcode{\sphinxupquote{\textquotesingle{}cartesian}} followed by \sphinxcode{\sphinxupquote{\textquotesingle{} ?}} on the next line. A
less obvious case: an input edited in Word ran \sphinxcode{\sphinxupquote{task=scan}} as a plain optimization.



\sphinxAtStartPar
To check and fix the line endings on Linux:

\begin{sphinxVerbatim}[commandchars=\\\{\}]
file mol.in        \PYGZsh{} reports \PYGZdq{}with CRLF line terminators\PYGZdq{} if the file has Windows endings
dos2unix mol.in
\end{sphinxVerbatim}

\sphinxAtStartPar
Running \sphinxcode{\sphinxupquote{dos2unix}} on a file that already has Unix line endings does no harm. On Windows, Notepad++ and the FAR
Manager editor (Shift+F2, then UNIX (LF)) can convert line endings.



\sphinxAtStartPar
For the full list of input\sphinxhyphen{}reading failures and their symptoms, see {\hyperref[\detokenize{troubleshooting/common-problems:troubleshooting-input}]{\sphinxcrossref{\DUrole{std}{\DUrole{std-ref}{The input is not read, or is read wrongly}}}}}.


\section{Basis\sphinxhyphen{}set files}
\label{\detokenize{start/running:basis-set-files}}
\sphinxAtStartPar
The basis sets, their auxiliary (fitting) sets and the ECPs are read from a separate file named by
{\hyperref[\detokenize{reference/control:control-basis}]{\sphinxcrossref{\DUrole{std}{\DUrole{std-ref}{\$control basis}}}}}. The set itself is chosen with {\hyperref[\detokenize{reference/molecule:molecule-set}]{\sphinxcrossref{\DUrole{std}{\DUrole{std-ref}{set=}}}}} in \sphinxcode{\sphinxupquote{\$molecule}}; without
it, the default set of the file is used. Without \sphinxcode{\sphinxupquote{basis=}}, the basis data are read from the input file itself.



\sphinxAtStartPar
Give the basis file as a \sphinxstylestrong{full path}:

\begin{sphinxVerbatim}[commandchars=\\\{\}]
\PYGZdl{}control task=optimize basis=/home/user/priroda/basis/basis.in \PYGZdl{}end
\end{sphinxVerbatim}

\sphinxAtStartPar
A relative name is looked up relative to the directory the program runs in; in version 25, \sphinxcode{\sphinxupquote{basis=basis.in}} works
when the file is in that directory. One user’s parallel run stopped after
the node list with \sphinxcode{\sphinxupquote{application called MPI\_Abort(MPI\_COMM\_WORLD, 1) \sphinxhyphen{} process 0}} and worked once the basis file was
given with its full path.


\begin{itemize}
\item {} 
\sphinxAtStartPar
Put the name in double quotes if the path contains special characters: \sphinxcode{\sphinxupquote{basis="/data/basis sets/basis.in"}}. Two
single quotes do not work; \sphinxcode{\sphinxupquote{basis=\textquotesingle{}\textquotesingle{}3z.bas\textquotesingle{}\textquotesingle{}}} fails with \sphinxcode{\sphinxupquote{error opening file \textquotesingle{}\textquotesingle{}3z.bas\textquotesingle{}\textquotesingle{}}}.



\item {} 
\sphinxAtStartPar
Windows paths work: \sphinxcode{\sphinxupquote{basis=D:\textbackslash{}NMR\textbackslash{}Nature\textbackslash{}basis\textbackslash{}6\sphinxhyphen{}31+G(d)}}.



\item {} 
\sphinxAtStartPar
The relativistic file \sphinxcode{\sphinxupquote{basis4.in}} needs \sphinxcode{\sphinxupquote{four=1}} in \sphinxcode{\sphinxupquote{\$control}}. Without it the basis input fails.



\end{itemize}

\sphinxAtStartPar
Which basis files exist and which sets they contain is described with {\hyperref[\detokenize{reference/control:control-basis}]{\sphinxcrossref{\DUrole{std}{\DUrole{std-ref}{\$control basis}}}}}.

\sphinxstepscope


\chapter{Getting started}
\label{\detokenize{start/getting-started:getting-started}}\label{\detokenize{start/getting-started:start-getting-started}}\label{\detokenize{start/getting-started::doc}}
\sphinxAtStartPar
This chapter takes you through four things: getting the program, a single\sphinxhyphen{}point energy of water, a geometry
optimization of ammonia followed by a Hessian that proves the structure is a minimum, and the parts of the output
you need to read along the way. {\hyperref[\detokenize{start/running::doc}]{\sphinxcrossref{\DUrole{doc}{Installing and running}}}} covers running the program in more detail.

\sphinxAtStartPar
The example inputs were run on Priroda 25 (2025.02.28), and most output excerpts are copied from those runs. A few
excerpts come from outputs posted by users of older versions; the text says so where it matters.




\section{Obtaining Priroda}
\label{\detokenize{start/getting-started:obtaining-priroda}}\label{\detokenize{start/getting-started:id1}}
\sphinxAtStartPar
Priroda is free of charge. Everything you need to run it can be downloaded from the official site,
\sphinxurl{https://rad.chem.msu.ru/~laikov/p\_w/}: the executables, the basis\sphinxhyphen{}set files and the helper scripts. Put the
executable and \sphinxcode{\sphinxupquote{basis.in}} in a directory of your own. The main Linux executable, \sphinxcode{\sphinxupquote{p}}, is statically linked, so a
serial run needs no other libraries. The page is a bare file index with no documentation, license text or contact information. It holds:
\begin{itemize}
\item {} 
\sphinxAtStartPar
Linux x86\sphinxhyphen{}64 executables (\sphinxcode{\sphinxupquote{p}}, \sphinxcode{\sphinxupquote{p4}}, \sphinxcode{\sphinxupquote{p\_d}}, \sphinxcode{\sphinxupquote{p4\_d}}), builds for particular clusters (\sphinxcode{\sphinxupquote{p\_l}}, \sphinxcode{\sphinxupquote{p\_mvs1p1}}, \sphinxcode{\sphinxupquote{p\_mvs10p}},
\sphinxcode{\sphinxupquote{p\_mvs10q}}), an ARM build (\sphinxcode{\sphinxupquote{p\_arm}}), a bundled \sphinxcode{\sphinxupquote{mpiexec}}, and the molecule viewer \sphinxcode{\sphinxupquote{xm}};

\item {} 
\sphinxAtStartPar
Windows (Cygwin, x86\sphinxhyphen{}64) executables in \sphinxcode{\sphinxupquote{exe64/}};

\item {} 
\sphinxAtStartPar
example inputs and outputs;

\item {} 
\sphinxAtStartPar
basis\sphinxhyphen{}set and parameter files (\sphinxcode{\sphinxupquote{basis.in}}, \sphinxcode{\sphinxupquote{basis4.in}}, \sphinxcode{\sphinxupquote{b2*.in}}, \sphinxcode{\sphinxupquote{b4s*.in}}, \sphinxcode{\sphinxupquote{3z.bas}}, \sphinxcode{\sphinxupquote{sbk.bas}}, \sphinxcode{\sphinxupquote{qm.in}});

\item {} 
\sphinxAtStartPar
small shell scripts: \sphinxcode{\sphinxupquote{ge}} (append the \sphinxcode{\sphinxupquote{\$Energy}} block from an output to the input), \sphinxcode{\sphinxupquote{gm}} (replace \sphinxcode{\sphinxupquote{\$molecule}} with
the optimized geometry) and \sphinxcode{\sphinxupquote{grx}} (join two IRC halves).

\end{itemize}



\sphinxAtStartPar
Which executable to choose is described in {\hyperref[\detokenize{start/running::doc}]{\sphinxcrossref{\DUrole{doc}{Installing and running}}}}. For the examples in this chapter you need one Linux
executable (for example \sphinxcode{\sphinxupquote{p}}) and \sphinxcode{\sphinxupquote{basis.in}}.

\sphinxAtStartPar
Users have also obtained the program by writing to the developer, Dimitri Laikov, by e\sphinxhyphen{}mail. Users reported the
address \sphinxcode{\sphinxupquote{dimitri\_laikov@mail.ru}} from 2010 to 2018; earlier it was \sphinxcode{\sphinxupquote{laikov@physto.se}}. Users note that he does not always answer quickly.



\begin{sphinxadmonition}{note}{Note:}
\sphinxAtStartPar
The legal status is not stated on the official site. A 2009 presentation relayed on a forum said that Priroda is
registered with Rospatent and distributed through the company OOO “Dimonta”. A user reported in 2018 that the developer
had said, a couple of years earlier, that there was no licence agreement. Users customarily ask the developer’s
permission before using the program.
\end{sphinxadmonition}



\sphinxAtStartPar
The official site has hosted builds with known defects. In 2016 it carried a note that the \sphinxcode{\sphinxupquote{w64/*.exe}} builds (an
attempt at 64\sphinxhyphen{}bit Cygwin) had a DFT/ECP bug. If a build behaves strangely, ask the developer whether it is the right one.




\section{Before you start}
\label{\detokenize{start/getting-started:before-you-start}}
\sphinxAtStartPar
Four practical points save most first\sphinxhyphen{}time failures.
\begin{itemize}
\item {} 
\sphinxAtStartPar
\sphinxstylestrong{Plain text with Unix line endings.} Write inputs as plain ASCII with LF line endings and straight quotes. A file
saved with Windows (CRLF) line endings gives errors such as \sphinxcode{\sphinxupquote{molecule input: \textquotesingle{}\textquotesingle{}}} or \sphinxcode{\sphinxupquote{?rtesian}}, or makes
\sphinxcode{\sphinxupquote{task=scan}} silently run an optimization. Convert such files with \sphinxcode{\sphinxupquote{dos2unix}}.



\item {} 
\sphinxAtStartPar
\sphinxstylestrong{Give the basis file by its full path.} The {\hyperref[\detokenize{reference/control:control-basis}]{\sphinxcrossref{\DUrole{std}{\DUrole{std-ref}{basis}}}}} keyword names the basis\sphinxhyphen{}set file.
A relative name works in version 25 when the file is in the directory you run from, but relative names have
failed in older versions and in parallel runs, so an absolute path is the safe choice. The examples below use
\sphinxcode{\sphinxupquote{/home/user/priroda/basis.in}}; replace it with the location of \sphinxcode{\sphinxupquote{basis.in}} on your machine.



\item {} 
\sphinxAtStartPar
\sphinxstylestrong{Atomic numbers are integers.} Each atom line in \sphinxcode{\sphinxupquote{\$molecule}} starts with the atomic number written as an integer
(\sphinxcode{\sphinxupquote{8}}, not \sphinxcode{\sphinxupquote{8.0}} and not \sphinxcode{\sphinxupquote{O}}). A real number such as \sphinxcode{\sphinxupquote{8.0}} gives no error: the program silently reads a different
molecule. 

\item {} 
\sphinxAtStartPar
\sphinxstylestrong{Running.} A serial run is \sphinxcode{\sphinxupquote{p input output}}. Parallel runs are described in {\hyperref[\detokenize{start/running::doc}]{\sphinxcrossref{\DUrole{doc}{Installing and running}}}}.



\end{itemize}


\section{A single\sphinxhyphen{}point energy}
\label{\detokenize{start/getting-started:a-single-point-energy}}

\subsection{The input}
\label{\detokenize{start/getting-started:the-input}}
\sphinxAtStartPar
Save the following as \sphinxcode{\sphinxupquote{water.in}}:

\begin{sphinxVerbatim}[commandchars=\\\{\}]
\PYGZdl{}system mem=128 disk=1 \PYGZdl{}end
\PYGZdl{}control
 task=energy
 theory=dft
 basis=/home/user/priroda/basis.in
 print=+charges+bonds
\PYGZdl{}end
\PYGZdl{}dft functional=PBE \PYGZdl{}end
\PYGZdl{}molecule
 charge=0 mult=1
 cartesian
 set=L1
  8    0.000000    0.000000    0.117300
  1    0.000000    0.757200   \PYGZhy{}0.469200
  1    0.000000   \PYGZhy{}0.757200   \PYGZhy{}0.469200
\PYGZdl{}end
\end{sphinxVerbatim}



\sphinxAtStartPar
Group by group:
\begin{itemize}
\item {} 
\sphinxAtStartPar
{\hyperref[\detokenize{reference/system:group-system}]{\sphinxcrossref{\DUrole{std}{\DUrole{std-ref}{\$system}}}}} sets the memory limit in MB and the scratch space. Here \sphinxcode{\sphinxupquote{disk=1}} means 1 GB of disk
for temporary files; a negative value would mean a RAM disk of that many MB. These values are per process.



\item {} 
\sphinxAtStartPar
{\hyperref[\detokenize{reference/control:group-control}]{\sphinxcrossref{\DUrole{std}{\DUrole{std-ref}{\$control}}}}} sets the job type and the method. {\hyperref[\detokenize{reference/control:control-task}]{\sphinxcrossref{\DUrole{std}{\DUrole{std-ref}{task=energy}}}}} asks for the
energy at the given geometry. {\hyperref[\detokenize{reference/control:control-theory}]{\sphinxcrossref{\DUrole{std}{\DUrole{std-ref}{theory=dft}}}}} is GGA DFT with density fitting, the program’s
standard method. \sphinxcode{\sphinxupquote{print=+charges+bonds}} adds Hirshfeld charges and Mulliken populations to the output (see
{\hyperref[\detokenize{start/getting-started:qs-charges}]{\sphinxcrossref{\DUrole{std}{\DUrole{std-ref}{Charges and bond orders}}}}} below).

\item {} 
\sphinxAtStartPar
{\hyperref[\detokenize{reference/dft:group-dft}]{\sphinxcrossref{\DUrole{std}{\DUrole{std-ref}{\$dft}}}}} selects the {\hyperref[\detokenize{reference/dft:dft-functional}]{\sphinxcrossref{\DUrole{std}{\DUrole{std-ref}{functional}}}}}. PBE is already the default, so this line
only makes the choice visible. 

\item {} 
\sphinxAtStartPar
{\hyperref[\detokenize{reference/molecule:group-molecule}]{\sphinxcrossref{\DUrole{std}{\DUrole{std-ref}{\$molecule}}}}} holds the {\hyperref[\detokenize{reference/molecule:molecule-charge}]{\sphinxcrossref{\DUrole{std}{\DUrole{std-ref}{charge}}}}}, the
{\hyperref[\detokenize{reference/molecule:molecule-mult}]{\sphinxcrossref{\DUrole{std}{\DUrole{std-ref}{multiplicity}}}}} (2S+1), the geometry, and the basis set. The geometry is given in Cartesian coordinates (Å), one line per atom: atomic number, then x, y,
z. {\hyperref[\detokenize{reference/molecule:molecule-set}]{\sphinxcrossref{\DUrole{std}{\DUrole{std-ref}{set=L1}}}}} picks the set named \sphinxcode{\sphinxupquote{L1}} from \sphinxcode{\sphinxupquote{basis.in}}. \sphinxcode{\sphinxupquote{L1}} is also the default
set of that file, and users compare it to cc\sphinxhyphen{}pVDZ in size. 

\end{itemize}

\begin{sphinxadmonition}{note}{Note:}
\sphinxAtStartPar
Hybrid functionals such as \sphinxcode{\sphinxupquote{PBE1}} or \sphinxcode{\sphinxupquote{B3LYP}} need {\hyperref[\detokenize{reference/control:control-theory}]{\sphinxcrossref{\DUrole{std}{\DUrole{std-ref}{theory=ridft}}}}}. With \sphinxcode{\sphinxupquote{theory=dft}}, older
versions silently drop the exact exchange, and newer ones abort with no error message. With \sphinxcode{\sphinxupquote{theory=tddft}},
version 25 still silently drops it. See {\hyperref[\detokenize{reference/control:control-theory}]{\sphinxcrossref{\DUrole{std}{\DUrole{std-ref}{\$control theory}}}}}.


\end{sphinxadmonition}

\sphinxAtStartPar
Run it:

\begin{sphinxVerbatim}[commandchars=\\\{\}]
p\PYG{+w}{ }water.in\PYG{+w}{ }water.out
\end{sphinxVerbatim}


\subsection{What to look for in the output}
\label{\detokenize{start/getting-started:what-to-look-for-in-the-output}}
\sphinxAtStartPar
The excerpts below are from this input run on Priroda 25. The output starts with the version banner,
\sphinxcode{\sphinxupquote{Priroda 25 (2025.02.28)}}, and the memory and disk limits. The line \sphinxcode{\sphinxupquote{molecule input: \textquotesingle{}\textquotesingle{}}} follows in every run; it is
not an error. Then come the coordinates, the \sphinxcode{\sphinxupquote{formula:}} line and a table of interatomic distances. Check the
formula: it is the quickest way to catch a misread \sphinxcode{\sphinxupquote{\$molecule}} group.



\sphinxAtStartPar
Next, a summary box repeats what the program understood from your input. Check that the charge, multiplicity,
method and functional are what you meant:

\begin{sphinxVerbatim}[commandchars=\\\{\}]
 +\PYGZhy{}\PYGZhy{}\PYGZhy{}\PYGZhy{}\PYGZhy{}\PYGZhy{}\PYGZhy{}\PYGZhy{}\PYGZhy{}\PYGZhy{}\PYGZhy{}\PYGZhy{}\PYGZhy{}\PYGZhy{}\PYGZhy{}\PYGZhy{}\PYGZhy{}\PYGZhy{}\PYGZhy{}\PYGZhy{}\PYGZhy{}\PYGZhy{}\PYGZhy{}\PYGZhy{}\PYGZhy{}\PYGZhy{}\PYGZhy{}\PYGZhy{}\PYGZhy{}\PYGZhy{}\PYGZhy{}\PYGZhy{}\PYGZhy{}\PYGZhy{}\PYGZhy{}\PYGZhy{}\PYGZhy{}\PYGZhy{}\PYGZhy{}
 | Number of atoms                      3
 | Charge of molecule                   0
 | Spin multiplicity                    1
 | Number of alpha electrons            5
 | Number of beta  electrons            5
 | Number of core  electrons            1
 | Theoretical Method                 DFT
 | Approximation to E(xc)             PBE
 | Dimension of wavefunction basis     24
 | Dimension of   auxiliary  basis     50
 | Hamiltonian: non\PYGZhy{}relativistic
 +\PYGZhy{}\PYGZhy{}\PYGZhy{}\PYGZhy{}\PYGZhy{}\PYGZhy{}\PYGZhy{}\PYGZhy{}\PYGZhy{}\PYGZhy{}\PYGZhy{}\PYGZhy{}\PYGZhy{}\PYGZhy{}\PYGZhy{}\PYGZhy{}\PYGZhy{}\PYGZhy{}\PYGZhy{}\PYGZhy{}\PYGZhy{}\PYGZhy{}\PYGZhy{}\PYGZhy{}\PYGZhy{}\PYGZhy{}\PYGZhy{}\PYGZhy{}\PYGZhy{}\PYGZhy{}\PYGZhy{}\PYGZhy{}\PYGZhy{}\PYGZhy{}\PYGZhy{}\PYGZhy{}\PYGZhy{}\PYGZhy{}\PYGZhy{}
\end{sphinxVerbatim}



\sphinxAtStartPar
The electron counts are per spin. “Number of core electrons” counts core orbitals (here the oxygen 1s); it matters
only for correlated methods.



\sphinxAtStartPar
Two lines then echo the SCF and grid settings, including defaults you did not set:

\begin{sphinxVerbatim}[commandchars=\\\{\}]
 scf: conv = 1.0e\PYGZhy{}06, 1.0e\PYGZhy{}03; iter=32,16; xmax = 0.25000; 0s; n1=12
 grid: accuracy=1.0e\PYGZhy{}08,1.0e\PYGZhy{}05, cut=1.0e\PYGZhy{}10,1.0e\PYGZhy{}07, t=\PYGZhy{}1
\end{sphinxVerbatim}



\sphinxAtStartPar
Next comes the SCF iteration table (columns for the energy, the energy change and the gradient, with indented
micro\sphinxhyphen{}iterations). The grid is enlarged partway through, which is why the energy jumps once. A successful SCF ends
with the line \sphinxcode{\sphinxupquote{converged!}}.



\begin{sphinxadmonition}{warning}{Warning:}
\sphinxAtStartPar
If the SCF fails, the program prints \sphinxcode{\sphinxupquote{no convergence!}} (or, in older versions, \sphinxcode{\sphinxupquote{SCF is far from convergence}}) and
then \sphinxstylestrong{carries on}: it prints an energy, orbitals and properties, and exits normally. The numbers are not usable.
Always look for \sphinxcode{\sphinxupquote{converged!}} before trusting a result; {\hyperref[\detokenize{troubleshooting/common-problems:troubleshooting-common-problems}]{\sphinxcrossref{\DUrole{std}{\DUrole{std-ref}{Solving common problems}}}}} lists the fixes.
\end{sphinxadmonition}



\sphinxAtStartPar
The total energy follows, in hartree, on the line starting with \sphinxcode{\sphinxupquote{E  =}}, together with its components:

\begin{sphinxVerbatim}[commandchars=\\\{\}]
  E  =     \PYGZhy{}76.3653913723
  T  =      76.07611549  V  =    \PYGZhy{}199.28820003  J  =      46.93680394
  Ex =      \PYGZhy{}8.95038631  Ec =      \PYGZhy{}0.32925888

  dipole =   0.735449 au (  1.869339 D)
\end{sphinxVerbatim}



\sphinxAtStartPar
If your energy differs from this one in more than the last few digits, check the basis file and the version.

\sphinxAtStartPar
\sphinxcode{\sphinxupquote{E}} is the number you want. The components \sphinxcode{\sphinxupquote{T}}, \sphinxcode{\sphinxupquote{V}}, \sphinxcode{\sphinxupquote{J}}, \sphinxcode{\sphinxupquote{Ex}} and \sphinxcode{\sphinxupquote{Ec}} do not add up to \sphinxcode{\sphinxupquote{E}}, because the nuclear
repulsion energy is not printed on these lines. This is expected and is not an error.



\sphinxAtStartPar
Every run also prints its results as a restart block on lines prefixed with \sphinxcode{\sphinxupquote{eng\textgreater{}}}. For \sphinxcode{\sphinxupquote{task=energy}} it holds the
energy and the dipole moment; gradient and Hessian runs add more:

\begin{sphinxVerbatim}[commandchars=\\\{\}]
eng\PYGZgt{}\PYGZdl{}Energy
eng\PYGZgt{} E=\PYGZhy{}7.6365391372338e+01
eng\PYGZgt{} D=\PYGZhy{}2.96281302e\PYGZhy{}14  6.81315681e\PYGZhy{}13 \PYGZhy{}7.35448725e\PYGZhy{}01
eng\PYGZgt{}\PYGZdl{}end
\end{sphinxVerbatim}




\section{Optimizing a geometry and confirming a minimum}
\label{\detokenize{start/getting-started:optimizing-a-geometry-and-confirming-a-minimum}}
\sphinxAtStartPar
The standard beginner procedure, as taught on the KNC wiki, alternates between optimizations and Hessians:


\begin{enumerate}
\sphinxsetlistlabels{\arabic}{enumi}{enumii}{}{.}%
\item {} 
\sphinxAtStartPar
Optimize the geometry.

\item {} 
\sphinxAtStartPar
Check that the optimization converged. If it did not, restart from the last geometry.

\item {} 
\sphinxAtStartPar
Take the final geometry and compute a Hessian (force constants and vibrational frequencies) there.

\item {} 
\sphinxAtStartPar
If any frequency is imaginary, the structure is not a minimum: re\sphinxhyphen{}optimize with that Hessian in the input, then
compute a new Hessian.

\item {} 
\sphinxAtStartPar
When all frequencies are real, read the thermochemistry.

\end{enumerate}

\sphinxAtStartPar
The Hessian check matters most when the start geometry is symmetric. An experienced user reports that Priroda’s
optimizer rarely leaves imaginary frequencies, except when the starting structure is symmetric and the optimizer keeps
that symmetry; then the check is compulsory. 


\subsection{Step 1: optimize}
\label{\detokenize{start/getting-started:step-1-optimize}}
\sphinxAtStartPar
Save the following as \sphinxcode{\sphinxupquote{nh3.in}}:

\begin{sphinxVerbatim}[commandchars=\\\{\}]
\PYGZdl{}system mem=128 disk=1 \PYGZdl{}end
\PYGZdl{}control
 task=optimize
 theory=dft
 basis=/home/user/priroda/basis.in
\PYGZdl{}end
\PYGZdl{}optimize tol=1e\PYGZhy{}5 steps=100 \PYGZdl{}end
\PYGZdl{}molecule
 charge=0 mult=1
 z\PYGZhy{}matrix
 set=L1
   7
   1    1    1.02672223
   1    1    1.02672140    2    104.150240
   1    1    1.02671372    2    104.144601   \PYGZhy{}3    104.150290
\PYGZdl{}end
\end{sphinxVerbatim}



\sphinxAtStartPar
The third hydrogen is placed by one distance and two angles: the negative index \sphinxcode{\sphinxupquote{\sphinxhyphen{}3}} means “the angle to atom 3”
instead of a dihedral. 

\sphinxAtStartPar
{\hyperref[\detokenize{reference/optimize:optimize-tolerance}]{\sphinxcrossref{\DUrole{std}{\DUrole{std-ref}{tol}}}}} is the convergence threshold on the gradient. The displacement threshold is then set
to 100 times that value. The default is \sphinxcode{\sphinxupquote{1e\sphinxhyphen{}4}}; the tighter \sphinxcode{\sphinxupquote{1e\sphinxhyphen{}5}} used in most published examples is a better
starting point for a Hessian. {\hyperref[\detokenize{reference/optimize:optimize-steps}]{\sphinxcrossref{\DUrole{std}{\DUrole{std-ref}{steps}}}}} caps the number of cycles (128 by default in version 25;
older versions used 50). If an optimization needs far more than 50 steps, raising the limit rarely helps; supplying a
Hessian usually does. The output echoes all these settings under \sphinxcode{\sphinxupquote{Geometry optimization options:}}.



\sphinxAtStartPar
Run it with \sphinxcode{\sphinxupquote{p nh3.in nh3.out}}.


\subsection{Step 2: check convergence}
\label{\detokenize{start/getting-started:step-2-check-convergence}}
\sphinxAtStartPar
Each optimization cycle starts with a line such as \sphinxcode{\sphinxupquote{Geometry Optimization Step 1:}}, followed by the energy, the
largest gradient and the largest displacement. When both are below the thresholds, the program prints
\sphinxcode{\sphinxupquote{OPTIMIZATION CONVERGED in N steps!}}. Check for it:



\begin{sphinxVerbatim}[commandchars=\\\{\}]
grep\PYG{+w}{ }\PYG{l+s+s2}{\PYGZdq{}OPTIMIZATION CONVERGED\PYGZdq{}}\PYG{+w}{ }nh3.out
\end{sphinxVerbatim}

\sphinxAtStartPar
Right after the marker, the program prints the optimized geometry as a ready\sphinxhyphen{}made \sphinxcode{\sphinxupquote{\$molecule}} group, each line
prefixed with \sphinxcode{\sphinxupquote{MOL\textgreater{}}}, followed by the final energy. For \sphinxcode{\sphinxupquote{nh3.in}}, version 25 printed:

\begin{sphinxVerbatim}[commandchars=\\\{\}]
 OPTIMIZATION CONVERGED in 6 steps!
MOL\PYGZgt{}\PYGZdl{}molecule
MOL\PYGZgt{} z\PYGZhy{}matrix
MOL\PYGZgt{} set=L1
MOL\PYGZgt{}   7
MOL\PYGZgt{}   1    1    1.02753254
MOL\PYGZgt{}   1    1    1.02744669    2    104.120811
MOL\PYGZgt{}   1    1    1.02748748    2    104.112983   \PYGZhy{}3    104.135802
MOL\PYGZgt{}\PYGZdl{}end
MOL\PYGZgt{} Energy =    \PYGZhy{}56.4988651558
\end{sphinxVerbatim}



\sphinxAtStartPar
The block keeps the coordinate type of the input: a Z\sphinxhyphen{}matrix input gives a Z\sphinxhyphen{}matrix with the optimized values. Only
a charge or multiplicity that differs from the defaults (0 and 1) is written: \sphinxcode{\sphinxupquote{nh3.in}} gives no such line, and an
OH⁻ optimization gave \sphinxcode{\sphinxupquote{MOL\textgreater{} charge=\sphinxhyphen{}1}}. So the block can be pasted into the next input as it is.




\subsection{Step 3: take the final geometry}
\label{\detokenize{start/getting-started:step-3-take-the-final-geometry}}
\sphinxAtStartPar
Strip the \sphinxcode{\sphinxupquote{MOL\textgreater{}}} prefix to get a group you can paste into the next input. The following commands copy \sphinxcode{\sphinxupquote{nh3.in}} to
\sphinxcode{\sphinxupquote{nh3\_freq.in}}, switch the task to \sphinxcode{\sphinxupquote{hessian}}, drop the old \sphinxcode{\sphinxupquote{\$molecule}} group and everything after it, and append the
optimized geometry:

\begin{sphinxVerbatim}[commandchars=\\\{\}]
sed\PYG{+w}{ }\PYGZhy{}e\PYG{+w}{ }\PYG{l+s+s1}{\PYGZsq{}s/task=optimize/task=hessian/\PYGZsq{}}\PYG{+w}{ }\PYGZhy{}e\PYG{+w}{ }\PYG{l+s+s1}{\PYGZsq{}/\PYGZbs{}\PYGZdl{}molecule/,\PYGZdl{}d\PYGZsq{}}\PYG{+w}{ }nh3.in\PYG{+w}{ }\PYGZgt{}\PYG{+w}{ }nh3\PYGZus{}freq.in
sed\PYG{+w}{ }\PYGZhy{}n\PYG{+w}{ }\PYG{l+s+s1}{\PYGZsq{}s/\PYGZca{}MOL\PYGZgt{}//p\PYGZsq{}}\PYG{+w}{ }nh3.out\PYG{+w}{ }\PYGZgt{}\PYGZgt{}\PYG{+w}{ }nh3\PYGZus{}freq.in
\end{sphinxVerbatim}



\sphinxAtStartPar
The first command assumes that \sphinxcode{\sphinxupquote{\$molecule}} is the last group of the input, as in all examples here. The \sphinxcode{\sphinxupquote{Energy = ...}}
line that ends up after \sphinxcode{\sphinxupquote{\$end}} does no harm: text outside a group is ignored. A normal optimization prints the \sphinxcode{\sphinxupquote{MOL\textgreater{}}}
block only once. A scan or an IRC prints one per point, so there you would pick the block you need.



\begin{sphinxadmonition}{note}{Note:}
\sphinxAtStartPar
\sphinxstylestrong{If the optimization did not converge}, there are no \sphinxcode{\sphinxupquote{MOL\textgreater{}}} lines, and the program prints no message about it:
the output just ends after the last step. Do not restart from the original geometry: start from the last one
reached. Every step prints the geometry for the next step as a \sphinxcode{\sphinxupquote{\$molecule}} group on lines prefixed with lower\sphinxhyphen{}case
\sphinxcode{\sphinxupquote{mol\textgreater{}}}, so this keeps the last such block:

\begin{sphinxVerbatim}[commandchars=\\\{\}]
sed\PYG{+w}{ }\PYGZhy{}n\PYG{+w}{ }\PYG{l+s+s1}{\PYGZsq{}s/\PYGZca{}mol\PYGZgt{}//p\PYGZsq{}}\PYG{+w}{ }nh3.out\PYG{+w}{ }\PYG{p}{|}\PYG{+w}{ }awk\PYG{+w}{ }\PYG{l+s+s1}{\PYGZsq{}/\PYGZbs{}\PYGZdl{}molecule/\PYGZob{}buf=\PYGZdq{}\PYGZdq{}\PYGZcb{} \PYGZob{}buf=buf \PYGZdl{}0 \PYGZdq{}\PYGZbs{}n\PYGZdq{}\PYGZcb{} END\PYGZob{}printf \PYGZdq{}\PYGZpc{}s\PYGZdq{}, buf\PYGZcb{}\PYGZsq{}}
\end{sphinxVerbatim}

\sphinxAtStartPar
The official \sphinxcode{\sphinxupquote{gm}} script (below) falls back to the input echo (lines starting with \sphinxcode{\sphinxupquote{\textless{}\$molecule}}) when there are no
\sphinxcode{\sphinxupquote{MOL\textgreater{}}} lines.


\end{sphinxadmonition}


\subsubsection{The official helper scripts \sphinxstyleliteralintitle{\sphinxupquote{ge}} and \sphinxstyleliteralintitle{\sphinxupquote{gm}}}
\label{\detokenize{start/getting-started:the-official-helper-scripts-ge-and-gm}}
\sphinxAtStartPar
The download site ships two small bash scripts for exactly this loop. Both take a job name (with or without the \sphinxcode{\sphinxupquote{.in}}
or \sphinxcode{\sphinxupquote{.out}} extension) and work on \sphinxcode{\sphinxupquote{name.in}} and \sphinxcode{\sphinxupquote{name.out}} in place:


\begin{itemize}
\item {} 
\sphinxAtStartPar
\sphinxcode{\sphinxupquote{gm nh3}} replaces the \sphinxcode{\sphinxupquote{\$molecule}} group in \sphinxcode{\sphinxupquote{nh3.in}} with the optimized one from \sphinxcode{\sphinxupquote{nh3.out}} and keeps the old input as
\sphinxcode{\sphinxupquote{nh3.bak}}. Like the commands above, it requires \sphinxcode{\sphinxupquote{\$molecule}} to be the last group, because it replaces everything from
\sphinxcode{\sphinxupquote{\$molecule}} onwards. It also appends the final \sphinxcode{\sphinxupquote{Energy =}} line and, when non\sphinxhyphen{}zero, the dipole line.

\item {} 
\sphinxAtStartPar
\sphinxcode{\sphinxupquote{ge nh3}} appends the \sphinxcode{\sphinxupquote{eng\textgreater{}}} lines of \sphinxcode{\sphinxupquote{nh3.out}} (the \sphinxcode{\sphinxupquote{\$Energy}} group: energy, gradient and, after a Hessian run, the
Hessian) to the end of \sphinxcode{\sphinxupquote{nh3.in}}.

\end{itemize}

\sphinxAtStartPar
So a cycle with these scripts is: run, \sphinxcode{\sphinxupquote{gm}}, edit the task, run again. Note that \sphinxcode{\sphinxupquote{ge}} copies every \sphinxcode{\sphinxupquote{eng\textgreater{}}} line in the
output, so use it on a Hessian output, not on an optimization output that has one block per step.




\subsection{Step 4: compute the Hessian}
\label{\detokenize{start/getting-started:step-4-compute-the-hessian}}
\sphinxAtStartPar
After step 3, \sphinxcode{\sphinxupquote{nh3\_freq.in}} looks like this, with your optimized values in place of these coordinates. Two optional
lines are added here: \sphinxcode{\sphinxupquote{print=+molden}} (for viewing the vibrations, see {\hyperref[\detokenize{start/getting-started:qs-molden}]{\sphinxcrossref{\DUrole{std}{\DUrole{std-ref}{Viewing results in Molden}}}}}) and a \sphinxcode{\sphinxupquote{\$thermo}} group:

\begin{sphinxVerbatim}[commandchars=\\\{\}]
\PYGZdl{}system mem=128 disk=1 \PYGZdl{}end
\PYGZdl{}control
 task=hessian
 theory=dft
 basis=/home/user/priroda/basis.in
 print=+molden
\PYGZdl{}end
\PYGZdl{}optimize tol=1e\PYGZhy{}5 steps=100 \PYGZdl{}end
\PYGZdl{}thermo sigma=3 \PYGZdl{}end
\PYGZdl{}molecule
 charge=0 mult=1
 z\PYGZhy{}matrix
 set=L1
   7
   1    1    1.02672223
   1    1    1.02672140    2    104.150240
   1    1    1.02671372    2    104.144601   \PYGZhy{}3    104.150290
\PYGZdl{}end
\end{sphinxVerbatim}



\sphinxAtStartPar
The \sphinxcode{\sphinxupquote{\$optimize}} group is ignored by a Hessian run but does no harm, and it is useful again if you have to go back to
step 1. {\hyperref[\detokenize{reference/thermo:thermo-sigma}]{\sphinxcrossref{\DUrole{std}{\DUrole{std-ref}{sigma}}}}} is the rotational symmetry number used in the thermochemistry; for ammonia
(point group C3v) it is 3. \sphinxstylestrong{Priroda does not work it out itself}: without \sphinxcode{\sphinxupquote{sigma=}} it uses 1, and the entropy comes
out too large by R ln σ (2.18 cal/(mol·K) for ammonia). The default temperature is 298.15 K; other temperatures go in
{\hyperref[\detokenize{reference/thermo:thermo-temperature}]{\sphinxcrossref{\DUrole{std}{\DUrole{std-ref}{temperature}}}}}.



\sphinxAtStartPar
For DFT the Hessian is analytic; the output announces it with \sphinxcode{\sphinxupquote{ANALYTIC SECOND DERIVATIVES OF SCF ENERGY}}.




\subsection{Step 5: check the frequencies}
\label{\detokenize{start/getting-started:step-5-check-the-frequencies}}
\sphinxAtStartPar
For a minimum, \sphinxstylestrong{all frequencies must be real (positive)}. A single imaginary frequency marks a transition state;
several mean a higher\sphinxhyphen{}order saddle point. The frequencies are in the table after the line \sphinxcode{\sphinxupquote{Harmonic Frequencies}}.
For ammonia, version 25 printed:



\begin{sphinxVerbatim}[commandchars=\\\{\}]
 *\PYGZhy{}\PYGZhy{}\PYGZhy{}\PYGZhy{}\PYGZhy{}\PYGZhy{}*\PYGZhy{}\PYGZhy{}\PYGZhy{}\PYGZhy{}\PYGZhy{}\PYGZhy{}\PYGZhy{}\PYGZhy{}\PYGZhy{}*\PYGZhy{}\PYGZhy{}\PYGZhy{}\PYGZhy{}\PYGZhy{}\PYGZhy{}\PYGZhy{}\PYGZhy{}*\PYGZhy{}\PYGZhy{}\PYGZhy{}\PYGZhy{}\PYGZhy{}\PYGZhy{}\PYGZhy{}\PYGZhy{}\PYGZhy{}*\PYGZhy{}\PYGZhy{}\PYGZhy{}\PYGZhy{}\PYGZhy{}\PYGZhy{}\PYGZhy{}\PYGZhy{}\PYGZhy{}*\PYGZhy{}\PYGZhy{}\PYGZhy{}\PYGZhy{}\PYGZhy{}\PYGZhy{}\PYGZhy{}\PYGZhy{}*\PYGZhy{}\PYGZhy{}\PYGZhy{}\PYGZhy{}\PYGZhy{}\PYGZhy{}\PYGZhy{}\PYGZhy{}\PYGZhy{}\PYGZhy{}\PYGZhy{}*\PYGZhy{}\PYGZhy{}\PYGZhy{}\PYGZhy{}\PYGZhy{}\PYGZhy{}\PYGZhy{}\PYGZhy{}\PYGZhy{}\PYGZhy{}\PYGZhy{}*
 | Mode | Freq.   | Mass.  | IR Int. |Raman A. |Raman D.| ROA (90)  | ROA (180) |
 *\PYGZhy{}\PYGZhy{}\PYGZhy{}\PYGZhy{}\PYGZhy{}\PYGZhy{}*\PYGZhy{}\PYGZhy{}\PYGZhy{}\PYGZhy{}\PYGZhy{}\PYGZhy{}\PYGZhy{}\PYGZhy{}\PYGZhy{}*\PYGZhy{}\PYGZhy{}\PYGZhy{}\PYGZhy{}\PYGZhy{}\PYGZhy{}\PYGZhy{}\PYGZhy{}*\PYGZhy{}\PYGZhy{}\PYGZhy{}\PYGZhy{}\PYGZhy{}\PYGZhy{}\PYGZhy{}\PYGZhy{}\PYGZhy{}*\PYGZhy{}\PYGZhy{}\PYGZhy{}\PYGZhy{}\PYGZhy{}\PYGZhy{}\PYGZhy{}\PYGZhy{}\PYGZhy{}*\PYGZhy{}\PYGZhy{}\PYGZhy{}\PYGZhy{}\PYGZhy{}\PYGZhy{}\PYGZhy{}\PYGZhy{}*\PYGZhy{}\PYGZhy{}\PYGZhy{}\PYGZhy{}\PYGZhy{}\PYGZhy{}\PYGZhy{}\PYGZhy{}\PYGZhy{}\PYGZhy{}\PYGZhy{}*\PYGZhy{}\PYGZhy{}\PYGZhy{}\PYGZhy{}\PYGZhy{}\PYGZhy{}\PYGZhy{}\PYGZhy{}\PYGZhy{}\PYGZhy{}\PYGZhy{}*
 |    1 | 1095.36 |   1.18 | 111.670 |   0.000 | 0.0000 |   0.00000 |   0.00000 |
 |    2 | 1635.36 |   1.06 |  12.058 |   0.000 | 0.0000 |   0.00000 |   0.00000 |
 |    3 | 1635.49 |   1.06 |  12.054 |   0.000 | 0.0000 |   0.00000 |   0.00000 |
 |    4 | 3371.36 |   1.03 |   3.443 |   0.000 | 0.0000 |   0.00000 |   0.00000 |
 |    5 | 3488.82 |   1.09 |   0.493 |   0.000 | 0.0000 |   0.00000 |   0.00000 |
 |    6 | 3488.95 |   1.09 |   0.494 |   0.000 | 0.0000 |   0.00000 |   0.00000 |
 *\PYGZhy{}\PYGZhy{}\PYGZhy{}\PYGZhy{}\PYGZhy{}\PYGZhy{}*\PYGZhy{}\PYGZhy{}\PYGZhy{}\PYGZhy{}\PYGZhy{}\PYGZhy{}\PYGZhy{}\PYGZhy{}\PYGZhy{}*\PYGZhy{}\PYGZhy{}\PYGZhy{}\PYGZhy{}\PYGZhy{}\PYGZhy{}\PYGZhy{}\PYGZhy{}*\PYGZhy{}\PYGZhy{}\PYGZhy{}\PYGZhy{}\PYGZhy{}\PYGZhy{}\PYGZhy{}\PYGZhy{}\PYGZhy{}*\PYGZhy{}\PYGZhy{}\PYGZhy{}\PYGZhy{}\PYGZhy{}\PYGZhy{}\PYGZhy{}\PYGZhy{}\PYGZhy{}*\PYGZhy{}\PYGZhy{}\PYGZhy{}\PYGZhy{}\PYGZhy{}\PYGZhy{}\PYGZhy{}\PYGZhy{}*\PYGZhy{}\PYGZhy{}\PYGZhy{}\PYGZhy{}\PYGZhy{}\PYGZhy{}\PYGZhy{}\PYGZhy{}\PYGZhy{}\PYGZhy{}\PYGZhy{}*\PYGZhy{}\PYGZhy{}\PYGZhy{}\PYGZhy{}\PYGZhy{}\PYGZhy{}\PYGZhy{}\PYGZhy{}\PYGZhy{}\PYGZhy{}\PYGZhy{}*
\end{sphinxVerbatim}



\sphinxAtStartPar
The columns are the frequency in cm⁻\(\sp{\text{1}}\), the reduced mass in amu and the IR intensity in km/mol. The Raman columns
stay zero in a plain Hessian run; for Raman intensities, run \sphinxcode{\sphinxupquote{task=dipole}} first and append its \sphinxcode{\sphinxupquote{num\textgreater{}}} lines (see
{\hyperref[\detokenize{start/input-syntax::doc}]{\sphinxcrossref{\DUrole{doc}{Input syntax}}}}). A second table, \sphinxcode{\sphinxupquote{Normal Modes (mass\sphinxhyphen{}weighted) and Frequencies}}, follows with the mode vectors. To print the table:



\begin{sphinxVerbatim}[commandchars=\\\{\}]
grep\PYG{+w}{ }\PYGZhy{}A\PYG{+w}{ }\PYG{l+m}{12}\PYG{+w}{ }\PYG{l+s+s2}{\PYGZdq{}Harmonic Frequencies\PYGZdq{}}\PYG{+w}{ }nh3\PYGZus{}freq.out
\end{sphinxVerbatim}



\sphinxAtStartPar
An imaginary frequency is printed with a trailing \sphinxcode{\sphinxupquote{i}}. Planar ammonia, the transition state of the inversion,
gives one:

\begin{sphinxVerbatim}[commandchars=\\\{\}]
 |    1 |  845.22i|   1.21 | 357.748 |   0.000 | 0.0000 |   0.00000 |   0.00000 |
\end{sphinxVerbatim}



\sphinxAtStartPar
\sphinxcode{\sphinxupquote{grep "i|" nh3\_freq.out}} therefore finds the imaginary modes.


\subsection{Step 6: if a frequency is imaginary, re\sphinxhyphen{}optimize}
\label{\detokenize{start/getting-started:step-6-if-a-frequency-is-imaginary-re-optimize}}
\sphinxAtStartPar
Put the Hessian you just computed into the optimization input. It is printed as a \sphinxcode{\sphinxupquote{\$Energy}} group on \sphinxcode{\sphinxupquote{eng\textgreater{}}} lines.
Starting from a good Hessian, the optimizer can move off the saddle point. In a version 25 test that started
exactly on a saddle point (planar ammonia), it reached the pyramidal minimum in 19 steps; from a near\sphinxhyphen{}symmetric
structure it usually takes fewer.



\begin{sphinxVerbatim}[commandchars=\\\{\}]
sed\PYG{+w}{ }\PYGZhy{}e\PYG{+w}{ }\PYG{l+s+s1}{\PYGZsq{}s/task=hessian/task=optimize/\PYGZsq{}}\PYG{+w}{ }nh3\PYGZus{}freq.in\PYG{+w}{ }\PYGZgt{}\PYG{+w}{ }nh3\PYGZus{}opt2.in
sed\PYG{+w}{ }\PYGZhy{}n\PYG{+w}{ }\PYG{l+s+s1}{\PYGZsq{}s/\PYGZca{}eng\PYGZgt{}//p\PYGZsq{}}\PYG{+w}{ }nh3\PYGZus{}freq.out\PYG{+w}{ }\PYG{p}{|}\PYG{+w}{ }awk\PYG{+w}{ }\PYG{l+s+s1}{\PYGZsq{}/\PYGZbs{}\PYGZdl{}Energy/\PYGZob{}buf=\PYGZdq{}\PYGZdq{}\PYGZcb{} \PYGZob{}buf=buf \PYGZdl{}0 \PYGZdq{}\PYGZbs{}n\PYGZdq{}\PYGZcb{} END\PYGZob{}printf \PYGZdq{}\PYGZpc{}s\PYGZdq{}, buf\PYGZcb{}\PYGZsq{}}\PYG{+w}{ }\PYGZgt{}\PYGZgt{}\PYG{+w}{ }nh3\PYGZus{}opt2.in
\end{sphinxVerbatim}





\sphinxAtStartPar
The second command appends only the last \sphinxcode{\sphinxupquote{\$Energy}} group, after \sphinxcode{\sphinxupquote{\$molecule}}. Do not leave an older \sphinxcode{\sphinxupquote{\$Energy}} group in
the file. Then run \sphinxcode{\sphinxupquote{nh3\_opt2.in}} and go back to step 2. (\sphinxcode{\sphinxupquote{ge nh3\_freq}} does the appending in place, if you work with
the official scripts.)


\subsection{Step 7: read the thermochemistry}
\label{\detokenize{start/getting-started:step-7-read-the-thermochemistry}}
\sphinxAtStartPar
Once all frequencies are real, the end of the Hessian output holds the zero\sphinxhyphen{}point energy and a table of thermodynamic
functions. The KNC wiki points to the string \sphinxcode{\sphinxupquote{ZPVE}}. Version 25 prints the table at 298.150 K by default; this
older user output (molecule not stated) shows 298.000 K:



\begin{sphinxVerbatim}[commandchars=\\\{\}]
 E =  \PYGZhy{}196.328936    E0 =  \PYGZhy{}196.191669

 ZPVE = 0.137267 a.u. =   86.1366 kcal/mol

 \PYGZhy{}\PYGZhy{}\PYGZhy{}\PYGZhy{}\PYGZhy{}\PYGZhy{}\PYGZhy{}\PYGZhy{}\PYGZhy{}\PYGZhy{}\PYGZhy{}\PYGZhy{}\PYGZhy{}\PYGZhy{}\PYGZhy{}\PYGZhy{}\PYGZhy{}\PYGZhy{}\PYGZhy{}\PYGZhy{}\PYGZhy{}\PYGZhy{}\PYGZhy{}\PYGZhy{}\PYGZhy{}\PYGZhy{}\PYGZhy{}\PYGZhy{}\PYGZhy{}\PYGZhy{}\PYGZhy{}\PYGZhy{}\PYGZhy{}\PYGZhy{}\PYGZhy{}\PYGZhy{}\PYGZhy{}\PYGZhy{}\PYGZhy{}\PYGZhy{}\PYGZhy{}\PYGZhy{}\PYGZhy{}\PYGZhy{}\PYGZhy{}\PYGZhy{}\PYGZhy{}\PYGZhy{}\PYGZhy{}\PYGZhy{}\PYGZhy{}\PYGZhy{}\PYGZhy{}\PYGZhy{}\PYGZhy{}\PYGZhy{}\PYGZhy{}\PYGZhy{}\PYGZhy{}\PYGZhy{}\PYGZhy{}\PYGZhy{}\PYGZhy{}\PYGZhy{}\PYGZhy{}\PYGZhy{}\PYGZhy{}\PYGZhy{}\PYGZhy{}\PYGZhy{}
 T =  298.000 K      ln(Q)       S          U          H          G
                              cal/mol*K  kcal/mol   kcal/mol   kcal/mol
 \PYGZhy{}\PYGZhy{}\PYGZhy{}\PYGZhy{}\PYGZhy{}\PYGZhy{}\PYGZhy{}\PYGZhy{}\PYGZhy{}\PYGZhy{}\PYGZhy{}\PYGZhy{}\PYGZhy{}\PYGZhy{}\PYGZhy{}\PYGZhy{}\PYGZhy{}\PYGZhy{}\PYGZhy{}\PYGZhy{}\PYGZhy{}\PYGZhy{}\PYGZhy{}\PYGZhy{}\PYGZhy{}\PYGZhy{}\PYGZhy{}\PYGZhy{}\PYGZhy{}\PYGZhy{}\PYGZhy{}\PYGZhy{}\PYGZhy{}\PYGZhy{}\PYGZhy{}\PYGZhy{}\PYGZhy{}\PYGZhy{}\PYGZhy{}\PYGZhy{}\PYGZhy{}\PYGZhy{}\PYGZhy{}\PYGZhy{}\PYGZhy{}\PYGZhy{}\PYGZhy{}\PYGZhy{}\PYGZhy{}\PYGZhy{}\PYGZhy{}\PYGZhy{}\PYGZhy{}\PYGZhy{}\PYGZhy{}\PYGZhy{}\PYGZhy{}\PYGZhy{}\PYGZhy{}\PYGZhy{}\PYGZhy{}\PYGZhy{}\PYGZhy{}\PYGZhy{}\PYGZhy{}\PYGZhy{}\PYGZhy{}\PYGZhy{}\PYGZhy{}\PYGZhy{}
 translational      16.9514    38.6542     0.8883     1.4805   \PYGZhy{}10.0385
 rotational         11.1465    25.1313     0.8883     0.8883    \PYGZhy{}6.6008
 vibrational         0.6842     4.7490    87.1466    87.1466    85.7314
 total              28.7821    68.5345    88.9232    89.5154    69.0921
 \PYGZhy{}\PYGZhy{}\PYGZhy{}\PYGZhy{}\PYGZhy{}\PYGZhy{}\PYGZhy{}\PYGZhy{}\PYGZhy{}\PYGZhy{}\PYGZhy{}\PYGZhy{}\PYGZhy{}\PYGZhy{}\PYGZhy{}\PYGZhy{}\PYGZhy{}\PYGZhy{}\PYGZhy{}\PYGZhy{}\PYGZhy{}\PYGZhy{}\PYGZhy{}\PYGZhy{}\PYGZhy{}\PYGZhy{}\PYGZhy{}\PYGZhy{}\PYGZhy{}\PYGZhy{}\PYGZhy{}\PYGZhy{}\PYGZhy{}\PYGZhy{}\PYGZhy{}\PYGZhy{}\PYGZhy{}\PYGZhy{}\PYGZhy{}\PYGZhy{}\PYGZhy{}\PYGZhy{}\PYGZhy{}\PYGZhy{}\PYGZhy{}\PYGZhy{}\PYGZhy{}\PYGZhy{}\PYGZhy{}\PYGZhy{}\PYGZhy{}\PYGZhy{}\PYGZhy{}\PYGZhy{}\PYGZhy{}\PYGZhy{}\PYGZhy{}\PYGZhy{}\PYGZhy{}\PYGZhy{}\PYGZhy{}\PYGZhy{}\PYGZhy{}\PYGZhy{}\PYGZhy{}\PYGZhy{}\PYGZhy{}\PYGZhy{}\PYGZhy{}\PYGZhy{}
\end{sphinxVerbatim}



\sphinxAtStartPar
How to read it:
\begin{itemize}
\item {} 
\sphinxAtStartPar
\sphinxcode{\sphinxupquote{E}} is the electronic energy in hartree. \sphinxcode{\sphinxupquote{E0 = E + ZPVE}} is the energy corrected for the zero\sphinxhyphen{}point vibrational
energy (here −196.328936 + 0.137267 = −196.191669).

\item {} 
\sphinxAtStartPar
The \sphinxcode{\sphinxupquote{U}}, \sphinxcode{\sphinxupquote{H}} and \sphinxcode{\sphinxupquote{G}} columns are \sphinxstylestrong{thermal corrections} in kcal/mol, not total values. They already contain the
ZPVE. The program does not print the total enthalpy or free energy; you add them yourself.

\item {} 
\sphinxAtStartPar
Take the corrections from the \sphinxcode{\sphinxupquote{total}} row and add them to \sphinxstylestrong{E, not E0}. Adding them to \sphinxcode{\sphinxupquote{E0}} would count the
zero\sphinxhyphen{}point energy twice:

\sphinxAtStartPar
H = E + ddH,  G = E + ddG,  with ddH = \sphinxcode{\sphinxupquote{H}} (total) and ddG = \sphinxcode{\sphinxupquote{G}} (total).

\item {} 
\sphinxAtStartPar
Convert units before adding: 1 hartree = 627.51 kcal/mol. For the table above,
G = −196.328936 × 627.51 + 69.0921 = −123129.28 kcal/mol, or −196.328936 + 69.0921 / 627.51 = −196.218831 hartree.

\item {} 
\sphinxAtStartPar
The entropy \sphinxcode{\sphinxupquote{S}} (cal/(mol·K)) is taken from the \sphinxcode{\sphinxupquote{total}} row as it stands.

\end{itemize}



\begin{sphinxadmonition}{warning}{Warning:}
\sphinxAtStartPar
For a \sphinxstylestrong{single atom}, the table is wrong: the rotational row contains \sphinxcode{\sphinxupquote{\sphinxhyphen{}inf}}/\sphinxcode{\sphinxupquote{inf}} or a spurious RT term. Use only
the translational contribution for atoms. Older versions were also reported to put an extra RT into the enthalpy of
monoatomic species. 
\end{sphinxadmonition}

\sphinxAtStartPar
The temperature is printed in the \sphinxcode{\sphinxupquote{T =}} line. Version 25 uses 298.150 K; outputs posted by users of older versions
also show \sphinxcode{\sphinxupquote{298.000 K}}. Set {\hyperref[\detokenize{reference/thermo:thermo-temperature}]{\sphinxcrossref{\DUrole{std}{\DUrole{std-ref}{temperature}}}}} explicitly when you compare results across versions.




\subsection{The one\sphinxhyphen{}run alternative: \sphinxstyleliteralintitle{\sphinxupquote{task=optimize+hessian}}}
\label{\detokenize{start/getting-started:the-one-run-alternative-task-optimize-hessian}}
\sphinxAtStartPar
You can let one job optimize the structure and then compute the Hessian at the converged geometry. In \sphinxcode{\sphinxupquote{nh3.in}},
change the task line:

\begin{sphinxVerbatim}[commandchars=\\\{\}]
\PYGZdl{}control
 task=optimize+hessian
 theory=dft
 basis=/home/user/priroda/basis.in
 print=+molden
\PYGZdl{}end
\end{sphinxVerbatim}



\sphinxAtStartPar
This is the only combined task that works: two \sphinxcode{\sphinxupquote{task=}} lines do not chain jobs, and only the first one runs.
\sphinxcode{\sphinxupquote{optimize+hessian}} existed in Priroda 2, was missing in versions 6 to 8, and is available again from version 9.
If the resulting Hessian shows an imaginary frequency, you still need step 6.




\section{Reading the output}
\label{\detokenize{start/getting-started:reading-the-output}}\label{\detokenize{start/getting-started:qs-reading-output}}

\subsection{Where things are}
\label{\detokenize{start/getting-started:where-things-are}}
\sphinxAtStartPar
The output is plain text, and the results meant for reuse are on lines tagged with a prefix, so \sphinxcode{\sphinxupquote{grep}} finds them.




\begin{savenotes}\sphinxattablestart
\sphinxthistablewithglobalstyle
\centering
\begin{tabulary}{\linewidth}[t]{TT}
\sphinxtoprule
\sphinxstyletheadfamily 
\sphinxAtStartPar
Look for
&\sphinxstyletheadfamily 
\sphinxAtStartPar
What it is
\\
\sphinxmidrule
\sphinxtableatstartofbodyhook
\sphinxAtStartPar
the summary box (\sphinxcode{\sphinxupquote{Number of atoms}}, \sphinxcode{\sphinxupquote{Theoretical Method}}, …)
&
\sphinxAtStartPar
what the program understood from the input
\\
\sphinxhline
\sphinxAtStartPar
\sphinxcode{\sphinxupquote{converged!}} / \sphinxcode{\sphinxupquote{no convergence!}}
&
\sphinxAtStartPar
the end of each SCF
\\
\sphinxhline
\sphinxAtStartPar
\sphinxcode{\sphinxupquote{E  =}}
&
\sphinxAtStartPar
total energy in hartree, followed by its components (without nuclear repulsion)
\\
\sphinxhline
\sphinxAtStartPar
\sphinxcode{\sphinxupquote{eng\textgreater{}}}
&
\sphinxAtStartPar
the \sphinxcode{\sphinxupquote{\$Energy}} group: energy \sphinxcode{\sphinxupquote{E=}}, dipole \sphinxcode{\sphinxupquote{D=}}, gradient \sphinxcode{\sphinxupquote{G=}}, and the Hessian after \sphinxcode{\sphinxupquote{task=hessian}}
\\
\sphinxhline
\sphinxAtStartPar
\sphinxcode{\sphinxupquote{OPTIMIZATION CONVERGED in N steps!}}
&
\sphinxAtStartPar
a converged optimization
\\
\sphinxhline
\sphinxAtStartPar
\sphinxcode{\sphinxupquote{MOL\textgreater{}}}
&
\sphinxAtStartPar
the optimized geometry as a \sphinxcode{\sphinxupquote{\$molecule}} group, then \sphinxcode{\sphinxupquote{Energy =}}
\\
\sphinxhline
\sphinxAtStartPar
\sphinxcode{\sphinxupquote{Harmonic Frequencies}}
&
\sphinxAtStartPar
the frequency table (\sphinxcode{\sphinxupquote{Normal Modes}} follows with the mode vectors)
\\
\sphinxhline
\sphinxAtStartPar
\sphinxcode{\sphinxupquote{ZPVE}}, \sphinxcode{\sphinxupquote{T = ... K}}
&
\sphinxAtStartPar
zero\sphinxhyphen{}point energy and the thermochemistry table
\\
\sphinxhline
\sphinxAtStartPar
\sphinxcode{\sphinxupquote{Overlap populations:}}
&
\sphinxAtStartPar
Mulliken populations (\sphinxcode{\sphinxupquote{print=+bonds}})
\\
\sphinxhline
\sphinxAtStartPar
\sphinxcode{\sphinxupquote{Atomic partitions of the density:}}
&
\sphinxAtStartPar
Hirshfeld charges (\sphinxcode{\sphinxupquote{print=+charges}})
\\
\sphinxhline
\sphinxAtStartPar
\sphinxcode{\sphinxupquote{freq\textgreater{}}}
&
\sphinxAtStartPar
vibrations in Molden format (\sphinxcode{\sphinxupquote{print=+molden}}, Hessian runs)
\\
\sphinxhline
\sphinxAtStartPar
\sphinxcode{\sphinxupquote{mos\textgreater{}}}
&
\sphinxAtStartPar
orbitals in Molden format (\sphinxcode{\sphinxupquote{print=+vectors+molden}})
\\
\sphinxbottomrule
\end{tabulary}
\sphinxtableafterendhook\par
\sphinxattableend\end{savenotes}




\subsection{Charges and bond orders}
\label{\detokenize{start/getting-started:charges-and-bond-orders}}\label{\detokenize{start/getting-started:qs-charges}}
\sphinxAtStartPar
With \sphinxcode{\sphinxupquote{print=+bonds}}, SCF runs print Mulliken charges and spin populations under \sphinxcode{\sphinxupquote{Overlap populations:}}, followed by
a table of bond orders. With \sphinxcode{\sphinxupquote{print=+charges}}, DFT runs print Hirshfeld charges and spins under
\sphinxcode{\sphinxupquote{Atomic partitions of the density:}}. The flags can be combined, as in \sphinxcode{\sphinxupquote{water.in}}, which gives:

\begin{sphinxVerbatim}[commandchars=\\\{\}]
 Overlap populations:
  atom    charge   spin p.
 O1      \PYGZhy{}0.4167  \PYGZhy{}0.00000
 H1       0.2083  \PYGZhy{}0.00000
 H2       0.2083  \PYGZhy{}0.00000

   Bond orders:
  O1    |  H1   :  1.07  H2   :  1.07
  H1    |  O1   :  1.07
  H2    |  O1   :  1.07

 Atomic partitions of the density:
   atom    charge    spin
   1  O   \PYGZhy{}0.2905    0.0000
   2  H    0.1451    0.0000
   3  H    0.1451    0.0000
  total   10.0003    0.0000
\end{sphinxVerbatim}



\sphinxAtStartPar
The \sphinxcode{\sphinxupquote{total}} row of the Hirshfeld block sums the electron populations, not the charges.


\subsection{Viewing results in Molden}
\label{\detokenize{start/getting-started:viewing-results-in-molden}}\label{\detokenize{start/getting-started:qs-molden}}
\sphinxAtStartPar
The KNC wiki treats Molden as the reference viewer. Open files with the \sphinxcode{\sphinxupquote{READ}} button in the \sphinxcode{\sphinxupquote{Miscellaneous}} group.



\sphinxAtStartPar
\sphinxstylestrong{Vibrations.} Run the Hessian with {\hyperref[\detokenize{reference/control:control-print}]{\sphinxcrossref{\DUrole{std}{\DUrole{std-ref}{print=+molden}}}}} and cut out the \sphinxcode{\sphinxupquote{freq\textgreater{}}} lines:

\begin{sphinxVerbatim}[commandchars=\\\{\}]
sed\PYG{+w}{ }\PYGZhy{}n\PYG{+w}{ }\PYG{l+s+s1}{\PYGZsq{}s/\PYGZca{}freq\PYGZgt{}//p\PYGZsq{}}\PYG{+w}{ }nh3\PYGZus{}freq.out\PYG{+w}{ }\PYGZgt{}\PYG{+w}{ }nh3.freq
\end{sphinxVerbatim}

\sphinxAtStartPar
Open \sphinxcode{\sphinxupquote{nh3.freq}} in Molden to animate the normal modes.  

\sphinxAtStartPar
\sphinxstylestrong{Orbitals.} Add \sphinxcode{\sphinxupquote{print=+vectors+molden}} to \sphinxcode{\sphinxupquote{\$control}}. The orbitals then appear on \sphinxcode{\sphinxupquote{mos\textgreater{}}} lines:

\begin{sphinxVerbatim}[commandchars=\\\{\}]
grep\PYG{+w}{ }\PYG{l+s+s1}{\PYGZsq{}\PYGZca{}mos\PYGZgt{}\PYGZsq{}}\PYG{+w}{ }water.out\PYG{+w}{ }\PYG{p}{|}\PYG{+w}{ }sed\PYG{+w}{ }\PYG{l+s+s1}{\PYGZsq{}s/\PYGZca{}mos\PYGZgt{}//\PYGZsq{}}\PYG{+w}{ }\PYGZgt{}\PYG{+w}{ }water.mos
\end{sphinxVerbatim}



\sphinxAtStartPar
The \sphinxcode{\sphinxupquote{+molden}} part is essential: without it, there are no \sphinxcode{\sphinxupquote{mos\textgreater{}}} lines. The files are large, and Molden can be slow
with them. 

\sphinxAtStartPar
\sphinxstylestrong{Geometry.} If the optimized geometry is Cartesian, this one\sphinxhyphen{}liner turns the \sphinxcode{\sphinxupquote{MOL\textgreater{}}} block into an XYZ file that
Molden reads (the symbol list covers H to Ar; extend it for heavier elements):

\begin{sphinxVerbatim}[commandchars=\\\{\}]
sed\PYG{+w}{ }\PYGZhy{}n\PYG{+w}{ }\PYG{l+s+s1}{\PYGZsq{}s/\PYGZca{}MOL\PYGZgt{}//p\PYGZsq{}}\PYG{+w}{ }job.out\PYG{+w}{ }\PYG{p}{|}\PYG{+w}{ }awk\PYG{+w}{ }\PYG{l+s+s1}{\PYGZsq{}BEGIN\PYGZob{}split(\PYGZdq{}H He Li Be B C N O F Ne Na Mg Al Si P S Cl Ar\PYGZdq{},s,\PYGZdq{} \PYGZdq{})\PYGZcb{}}
\PYG{l+s+s1}{  NF==4 \PYGZam{}\PYGZam{} \PYGZdl{}1+0\PYGZgt{}0 \PYGZob{}n++; l[n]=s[\PYGZdl{}1]\PYGZdq{} \PYGZdq{}\PYGZdl{}2\PYGZdq{} \PYGZdq{}\PYGZdl{}3\PYGZdq{} \PYGZdq{}\PYGZdl{}4\PYGZcb{}}
\PYG{l+s+s1}{  END\PYGZob{}print n; print \PYGZdq{}optimized geometry\PYGZdq{}; for(i=1;i\PYGZlt{}=n;i++) print l[i]\PYGZcb{}\PYGZsq{}}\PYG{+w}{ }\PYGZgt{}\PYG{+w}{ }job.xyz
\end{sphinxVerbatim}





\sphinxAtStartPar
For a Z\sphinxhyphen{}matrix geometry such as \sphinxcode{\sphinxupquote{nh3.in}}, the \sphinxcode{\sphinxupquote{freq\textgreater{}}} file from the Hessian run is the easier way to look at the
structure.

\sphinxAtStartPar
\sphinxstylestrong{pri2mol.} The Perl script \sphinxcode{\sphinxupquote{pri2mol}} by the forum user amge, available from
\sphinxcode{\sphinxupquote{http://limor1.nioch.nsc.ru/quant/program/pri2mol/}}, converts a Priroda output into Molden and XYZ files: geometries,
vibrations, orbitals, NMR shifts and density critical points. Priroda’s own viewer, \sphinxcode{\sphinxupquote{xm}}, is on the download site.




\section{Where next}
\label{\detokenize{start/getting-started:where-next}}\begin{itemize}
\item {} 
\sphinxAtStartPar
{\hyperref[\detokenize{start/input-syntax::doc}]{\sphinxcrossref{\DUrole{doc}{Input syntax}}}} explains groups, keyword abbreviation, arrays and restart blocks.

\item {} 
\sphinxAtStartPar
{\hyperref[\detokenize{start/running::doc}]{\sphinxcrossref{\DUrole{doc}{Installing and running}}}} covers memory, scratch disks and parallel runs.

\item {} 
\sphinxAtStartPar
The {\hyperref[\detokenize{reference/index::doc}]{\sphinxcrossref{\DUrole{doc}{input reference}}}} describes every known keyword, group by group.

\item {} 
\sphinxAtStartPar
{\hyperref[\detokenize{troubleshooting/common-problems:troubleshooting-common-problems}]{\sphinxcrossref{\DUrole{std}{\DUrole{std-ref}{Solving common problems}}}}} helps when a run fails.

\end{itemize}

\sphinxstepscope


\chapter{Input syntax}
\label{\detokenize{start/input-syntax:input-syntax}}\label{\detokenize{start/input-syntax::doc}}
\sphinxAtStartPar
\sphinxcode{\sphinxupquote{input.txt}} describes the input as “a kind of NAMELIST free format”. It is close to the GAMESS\sphinxhyphen{}US format: keywords are
grouped, and each group starts with \sphinxcode{\sphinxupquote{\$name}} and ends with \sphinxcode{\sphinxupquote{\$end}}. This chapter covers the rules that apply to all
groups. The individual keywords are in the {\hyperref[\detokenize{reference/index::doc}]{\sphinxcrossref{\DUrole{doc}{input reference}}}}.



\sphinxAtStartPar
A complete input for an NH3 optimization, adapted from the KNC wiki FAQ:

\begin{sphinxVerbatim}[commandchars=\\\{\}]
\PYGZdl{}system mem=512 disk=\PYGZhy{}512 path=/tmp \PYGZdl{}end
\PYGZdl{}control
 task=optimize
 theory=dft
 basis=/home/user/priroda/basis.in
\PYGZdl{}end
\PYGZdl{}optimize tol=1e\PYGZhy{}5 steps=100 \PYGZdl{}end
\PYGZdl{}molecule
 charge=0 mult=1
 z\PYGZhy{}matrix
 set=L1
   7
   1    1    1.02672223
   1    1    1.02672140    2    104.150240
   1    1    1.02671372    2    104.144601   \PYGZhy{}3    104.150290
\PYGZdl{}end
\end{sphinxVerbatim}




\section{Groups}
\label{\detokenize{start/input-syntax:groups}}
\sphinxAtStartPar
A group starts with \sphinxcode{\sphinxupquote{\$}} followed by the group name and ends with \sphinxcode{\sphinxupquote{\$end}}. It may sit on one line or span several:

\begin{sphinxVerbatim}[commandchars=\\\{\}]
\PYGZdl{}scf conv=1e\PYGZhy{}7 \PYGZdl{}end

\PYGZdl{}scf
 convergence=1e\PYGZhy{}7
 procedure=bfgs
\PYGZdl{}end
\end{sphinxVerbatim}

\sphinxAtStartPar
Inside a group, \sphinxcode{\sphinxupquote{keyword=value}} pairs are separated by spaces or line breaks, and several pairs can share a line
(\sphinxcode{\sphinxupquote{theory=qm\_n3 basis=qm.in}}). Blank lines inside a group are allowed.



\sphinxAtStartPar
The order of the groups does not matter, and neither does the order of keywords within a group. The exceptions are
\sphinxcode{\sphinxupquote{\$molecule}} and \sphinxcode{\sphinxupquote{\$Energy}}, whose contents are ordered data.



\sphinxAtStartPar
Groups you do not need can be left out; most keywords have sensible defaults. \sphinxcode{\sphinxupquote{input.txt}} advises changing only what
the problem requires.



\sphinxAtStartPar
Two groups can also share a line: \sphinxcode{\sphinxupquote{\$scf conv=1e\sphinxhyphen{}7 \$end \$grid acc=1e\sphinxhyphen{}9 \$end}} sets both values in version 25. Every
input in the sources starts each group on a new line, which is easier to read.




\section{Keywords}
\label{\detokenize{start/input-syntax:keywords}}

\subsection{Case}
\label{\detokenize{start/input-syntax:case}}
\sphinxAtStartPar
Group names, keyword names and keyword values are case\sphinxhyphen{}insensitive: \sphinxcode{\sphinxupquote{\$System}} and \sphinxcode{\sphinxupquote{\$system}}, \sphinxcode{\sphinxupquote{theory=riDFT}} and
\sphinxcode{\sphinxupquote{theory=ridft}}, \sphinxcode{\sphinxupquote{functional=pbe}} and \sphinxcode{\sphinxupquote{functional=PBE}} are the same. File names are the exception, because Linux file
names are case\sphinxhyphen{}sensitive. This applies to \sphinxcode{\sphinxupquote{basis=}}, \sphinxcode{\sphinxupquote{save=}} and \sphinxcode{\sphinxupquote{path=}}.




\subsection{Abbreviations}
\label{\detokenize{start/input-syntax:abbreviations}}
\sphinxAtStartPar
Within a group, a keyword can be shortened to a prefix. The example in \sphinxcode{\sphinxupquote{input.txt}} is \sphinxcode{\sphinxupquote{mem}} for \sphinxcode{\sphinxupquote{memory}}. Inputs in the
sources use, among others:




\begin{savenotes}\sphinxattablestart
\sphinxthistablewithglobalstyle
\centering
\begin{tabulary}{\linewidth}[t]{TT}
\sphinxtoprule
\sphinxstyletheadfamily 
\sphinxAtStartPar
Short form
&\sphinxstyletheadfamily 
\sphinxAtStartPar
Keyword
\\
\sphinxmidrule
\sphinxtableatstartofbodyhook
\sphinxAtStartPar
\sphinxcode{\sphinxupquote{mem=}}
&
\sphinxAtStartPar
{\hyperref[\detokenize{reference/system:system-memory}]{\sphinxcrossref{\DUrole{std}{\DUrole{std-ref}{\$system memory}}}}}
\\
\sphinxhline
\sphinxAtStartPar
\sphinxcode{\sphinxupquote{conv=}}, \sphinxcode{\sphinxupquote{iter=}}, \sphinxcode{\sphinxupquote{proc=}}
&
\sphinxAtStartPar
{\hyperref[\detokenize{reference/scf:scf-convergence}]{\sphinxcrossref{\DUrole{std}{\DUrole{std-ref}{\$scf convergence}}}}}, {\hyperref[\detokenize{reference/scf:scf-iterations}]{\sphinxcrossref{\DUrole{std}{\DUrole{std-ref}{iterations}}}}}, {\hyperref[\detokenize{reference/scf:scf-procedure}]{\sphinxcrossref{\DUrole{std}{\DUrole{std-ref}{procedure}}}}}
\\
\sphinxhline
\sphinxAtStartPar
\sphinxcode{\sphinxupquote{func=}}
&
\sphinxAtStartPar
{\hyperref[\detokenize{reference/dft:dft-functional}]{\sphinxcrossref{\DUrole{std}{\DUrole{std-ref}{\$dft functional}}}}}
\\
\sphinxhline
\sphinxAtStartPar
\sphinxcode{\sphinxupquote{acc=}}, \sphinxcode{\sphinxupquote{accur=}}
&
\sphinxAtStartPar
{\hyperref[\detokenize{reference/grid:grid-accuracy}]{\sphinxcrossref{\DUrole{std}{\DUrole{std-ref}{\$grid accuracy}}}}}
\\
\sphinxhline
\sphinxAtStartPar
\sphinxcode{\sphinxupquote{tol=}}, \sphinxcode{\sphinxupquote{cart=}}, \sphinxcode{\sphinxupquote{val=}}
&
\sphinxAtStartPar
{\hyperref[\detokenize{reference/optimize:optimize-tolerance}]{\sphinxcrossref{\DUrole{std}{\DUrole{std-ref}{\$optimize tolerance}}}}}, {\hyperref[\detokenize{reference/optimize:optimize-cartesian}]{\sphinxcrossref{\DUrole{std}{\DUrole{std-ref}{cartesian}}}}}, {\hyperref[\detokenize{reference/optimize:optimize-value}]{\sphinxcrossref{\DUrole{std}{\DUrole{std-ref}{value}}}}}
\\
\sphinxhline
\sphinxAtStartPar
\sphinxcode{\sphinxupquote{disp=}}
&
\sphinxAtStartPar
{\hyperref[\detokenize{reference/d2edr2:d2edr2-displace}]{\sphinxcrossref{\DUrole{std}{\DUrole{std-ref}{\$d2edr2 displace}}}}}
\\
\sphinxhline
\sphinxAtStartPar
\sphinxcode{\sphinxupquote{nuc=}}
&
\sphinxAtStartPar
{\hyperref[\detokenize{reference/control:control-nucleus}]{\sphinxcrossref{\DUrole{std}{\DUrole{std-ref}{\$control nucleus}}}}}
\\
\sphinxhline
\sphinxAtStartPar
\sphinxcode{\sphinxupquote{rout=}}
&
\sphinxAtStartPar
{\hyperref[\detokenize{reference/integral:integral-routine}]{\sphinxcrossref{\DUrole{std}{\DUrole{std-ref}{\$integral routine}}}}}
\\
\sphinxbottomrule
\end{tabulary}
\sphinxtableafterendhook\par
\sphinxattableend\end{savenotes}



\sphinxAtStartPar
Keywords belong to their group. The same name can mean different things in different groups: \sphinxcode{\sphinxupquote{steps}} in
{\hyperref[\detokenize{reference/optimize:optimize-steps}]{\sphinxcrossref{\DUrole{std}{\DUrole{std-ref}{\$optimize}}}}} is the maximum number of optimization cycles, while \sphinxcode{\sphinxupquote{steps}} in
{\hyperref[\detokenize{reference/d2edr2:d2edr2-steps}]{\sphinxcrossref{\DUrole{std}{\DUrole{std-ref}{\$d2edr2}}}}} selects displacement steps of a numerical Hessian.

\begin{sphinxadmonition}{warning}{Warning:}
\sphinxAtStartPar
A prefix must stay unambiguous within its group. How Priroda resolves a prefix that matches two keywords is not
documented. In version 6, when \sphinxcode{\sphinxupquote{\$optimize}} contained both \sphinxcode{\sphinxupquote{step=}} and \sphinxcode{\sphinxupquote{steps=}}, the maximum number of steps was taken
from \sphinxcode{\sphinxupquote{step}}. In inputs you keep or share, prefer the full keyword names; they are also easier to search for.
\end{sphinxadmonition}



\sphinxAtStartPar
A misspelled or unknown keyword name stops the run right after the basis\sphinxhyphen{}set echo, with a message naming the group
and the keyword:

\begin{sphinxVerbatim}[commandchars=\\\{\}]
*** error in \PYGZsq{}\PYGZdl{}scf\PYGZsq{} : unknown name nosuchkeyword
\end{sphinxVerbatim}



\sphinxAtStartPar
Only groups that the task reads are checked. In version 25, an unknown keyword in \sphinxcode{\sphinxupquote{\$optimize}} stopped an
optimization (\sphinxcode{\sphinxupquote{unknown name coordinated}}) but was ignored in an energy run, which does not read \sphinxcode{\sphinxupquote{\$optimize}}. A
clean energy run therefore does not prove that the rest of the input is correct.



\sphinxAtStartPar
The same message can appear for a keyword you spelled correctly. In version 25,
\sphinxcode{\sphinxupquote{\$scf convergence=1e\sphinxhyphen{}7 iterations=250 \$end}} stops with \sphinxcode{\sphinxupquote{unknown name terations}}: the value list of \sphinxcode{\sphinxupquote{convergence}}
takes the first letter of the next keyword. Give \sphinxcode{\sphinxupquote{convergence}} both of its values, or put it last in the group;
see {\hyperref[\detokenize{reference/scf:scf-convergence}]{\sphinxcrossref{\DUrole{std}{\DUrole{std-ref}{\$scf convergence}}}}}.




\section{Values}
\label{\detokenize{start/input-syntax:values}}\begin{description}
\sphinxlineitem{Scalars}
\sphinxAtStartPar
Integers (\sphinxcode{\sphinxupquote{steps=100}}), reals in plain or exponent notation (\sphinxcode{\sphinxupquote{trust=0.2}}, \sphinxcode{\sphinxupquote{tol=1e\sphinxhyphen{}5}}) and words
(\sphinxcode{\sphinxupquote{task=optimize}}). A sign is allowed: \sphinxcode{\sphinxupquote{charge=+1}}.



\sphinxlineitem{Lists}
\sphinxAtStartPar
Values separated by commas, for keywords that take several numbers: \sphinxcode{\sphinxupquote{tol=1e\sphinxhyphen{}5,1e\sphinxhyphen{}4}}, \sphinxcode{\sphinxupquote{mem=2048,256}},
\sphinxcode{\sphinxupquote{temperature=203, 223, 248, 298}}. Spaces after the commas are allowed, and some lists are also written with spaces
between groups of values (\sphinxcode{\sphinxupquote{fix=1,1,2,0,0 3,2,1,4,3}}). If you give fewer values than the keyword holds, the rest
keep their defaults: \sphinxcode{\sphinxupquote{\$scf conv=1e\sphinxhyphen{}10}} sets only the SCF threshold and leaves the CP\sphinxhyphen{}SCF threshold at \sphinxcode{\sphinxupquote{1e\sphinxhyphen{}3}}.



\sphinxlineitem{Indexed arrays}
\sphinxAtStartPar
\sphinxcode{\sphinxupquote{name(i)=a,b,…}} assigns values to an array starting at element \sphinxstyleemphasis{i} and leaves the other elements unchanged.
\sphinxcode{\sphinxupquote{\$guess na(85)=86,85 \$end}} sets element 85 of \sphinxcode{\sphinxupquote{na}} to 86 and element 86 to 85, which swaps two orbitals. The same
form is used for {\hyperref[\detokenize{reference/optimize:optimize-value}]{\sphinxcrossref{\DUrole{std}{\DUrole{std-ref}{\$optimize value}}}}} (\sphinxcode{\sphinxupquote{val(5)=2.0,2.1}}),
{\hyperref[\detokenize{reference/nmr:nmr-standard}]{\sphinxcrossref{\DUrole{std}{\DUrole{std-ref}{\$nmr standard}}}}} (\sphinxcode{\sphinxupquote{standard(6)=181.9921}}) and {\hyperref[\detokenize{reference/atoms:atoms-core}]{\sphinxcrossref{\DUrole{std}{\DUrole{std-ref}{\$atoms core}}}}}.



\sphinxlineitem{Flag strings}
\sphinxAtStartPar
{\hyperref[\detokenize{reference/control:control-print}]{\sphinxcrossref{\DUrole{std}{\DUrole{std-ref}{\$control print}}}}} takes a string of options, each switched on with \sphinxcode{\sphinxupquote{+}} or off with \sphinxcode{\sphinxupquote{\sphinxhyphen{}}}:
\sphinxcode{\sphinxupquote{print=+vectors+molden}}, \sphinxcode{\sphinxupquote{print=\sphinxhyphen{}internals+molden}}. The leading \sphinxcode{\sphinxupquote{+}} may be omitted (\sphinxcode{\sphinxupquote{print=charges+bonds}}).



\sphinxlineitem{Quoted strings}
\sphinxAtStartPar
A file name or path may be put in double quotes, and must be if it contains special characters:
\sphinxcode{\sphinxupquote{basis="3z.bas"}}, \sphinxcode{\sphinxupquote{path="/home/user/priroda/scratch"}}. Use the straight ASCII double quote \sphinxcode{\sphinxupquote{"}}. Typographic quotes
break the name, and two single quotes do not count as a double quote: \sphinxcode{\sphinxupquote{basis=\textquotesingle{}\textquotesingle{}3z.bas\textquotesingle{}\textquotesingle{}}} fails with
\sphinxcode{\sphinxupquote{error opening file \textquotesingle{}\textquotesingle{}3z.bas\textquotesingle{}\textquotesingle{}}}.



\end{description}


\section{The \sphinxstyleliteralintitle{\sphinxupquote{\$molecule}} group}
\label{\detokenize{start/input-syntax:the-molecule-group}}
\sphinxAtStartPar
{\hyperref[\detokenize{reference/molecule:group-molecule}]{\sphinxcrossref{\DUrole{std}{\DUrole{std-ref}{\$molecule}}}}} is read as ordered data rather than as keyword pairs:
\begin{enumerate}
\sphinxsetlistlabels{\arabic}{enumi}{enumii}{}{.}%
\item {} 
\sphinxAtStartPar
\sphinxcode{\sphinxupquote{charge=}} and \sphinxcode{\sphinxupquote{mult=}} (spin multiplicity 2S+1), both optional, with defaults 0 and 1;

\item {} 
\sphinxAtStartPar
one of the words \sphinxcode{\sphinxupquote{cartesian}} or \sphinxcode{\sphinxupquote{z\sphinxhyphen{}matrix}};

\item {} 
\sphinxAtStartPar
one line per atom, optionally preceded by \sphinxcode{\sphinxupquote{set=}} lines that choose the basis set for the atoms that follow.

\end{enumerate}

\sphinxAtStartPar
The first items may share the opening line:

\begin{sphinxVerbatim}[commandchars=\\\{\}]
\PYGZdl{}molecule charge=0 mult=1 cartesian
 set=L1
  8    0.000000    0.000000    0.117300
  1    0.000000    0.757200   \PYGZhy{}0.469200
  1    0.000000   \PYGZhy{}0.757200   \PYGZhy{}0.469200
\PYGZdl{}end
\end{sphinxVerbatim}



\sphinxAtStartPar
Each atom line begins with the atomic number written as an integer. Cartesian coordinates are in Å unless \sphinxcode{\sphinxupquote{unit=b}} is
given. The Z\sphinxhyphen{}matrix forms, mixed Cartesian and Z\sphinxhyphen{}matrix input, isotope masses and per\sphinxhyphen{}atom basis sets are described
in the {\hyperref[\detokenize{reference/molecule:group-molecule}]{\sphinxcrossref{\DUrole{std}{\DUrole{std-ref}{\$molecule}}}}} reference.




\subsection{Several \sphinxstyleliteralintitle{\sphinxupquote{\$molecule}} blocks}
\label{\detokenize{start/input-syntax:several-molecule-blocks}}
\sphinxAtStartPar
Some tasks read more than one \sphinxcode{\sphinxupquote{\$molecule}} group from one input:
\begin{itemize}
\item {} 
\sphinxAtStartPar
\sphinxcode{\sphinxupquote{task=transit}}, \sphinxcode{\sphinxupquote{task=saddle}} and \sphinxcode{\sphinxupquote{task=path}} take two blocks, the start and the end structure. The atoms must be
numbered the same way in both.

\item {} 
\sphinxAtStartPar
\sphinxcode{\sphinxupquote{task=uniq}} takes any number of blocks and filters out duplicate structures; the official example has 157.

\end{itemize}

\sphinxAtStartPar
These tasks are undocumented; see {\hyperref[\detokenize{reference/control:control-task}]{\sphinxcrossref{\DUrole{std}{\DUrole{std-ref}{\$control task}}}}} and {\hyperref[\detokenize{reference/uniq:group-uniq}]{\sphinxcrossref{\DUrole{std}{\DUrole{std-ref}{\$uniq}}}}}.




\section{Text outside groups}
\label{\detokenize{start/input-syntax:text-outside-groups}}
\sphinxAtStartPar
Everything outside a \sphinxcode{\sphinxupquote{\$name … \$end}} group is ignored. There is no documented comment character, but this rule serves
the same purpose, and the official example inputs use it to switch input off:


\begin{itemize}
\item {} 
\sphinxAtStartPar
\sphinxstylestrong{Delete the \sphinxcode{\sphinxupquote{\$}}} to disable a whole group. \sphinxcode{\sphinxupquote{integral dir=1 \$end}} in the official 2023\textendash{}2025 example inputs is plain text,
not an \sphinxcode{\sphinxupquote{\$integral}} group.

\item {} 
\sphinxAtStartPar
\sphinxstylestrong{Move keywords after \sphinxcode{\sphinxupquote{\$end}}} to disable them. In the example below, \sphinxcode{\sphinxupquote{v=}} and \sphinxcode{\sphinxupquote{h=}} are outside \sphinxcode{\sphinxupquote{\$dynamic}} and
have no effect.

\item {} 
\sphinxAtStartPar
\sphinxstylestrong{Keep notes after \sphinxcode{\sphinxupquote{\$end}}.} Helper scripts append result lines such as \sphinxcode{\sphinxupquote{Energy = …}} after the \sphinxcode{\sphinxupquote{\$molecule}} group,
where they do no harm.

\end{itemize}

\begin{sphinxVerbatim}[commandchars=\\\{\}]
integral dir=1 \PYGZdl{}end
\PYGZdl{}dynamic
 e=0.001
 d=0.03
 n=1000,100
\PYGZdl{}end
 v=500.0
 h=0.1
\PYGZdl{}molecule
 ...
\PYGZdl{}end
 Energy =   \PYGZhy{}783.83738426
\end{sphinxVerbatim}



\sphinxAtStartPar
In version 25 a line of text before the first group and a disabled group with its leftover \sphinxcode{\sphinxupquote{\$end}}
(\sphinxcode{\sphinxupquote{integral dir=1 \$end}}) are both ignored: the result is identical to the same input without them. Text inside a
group is parsed as keywords, so do not put notes there.



\sphinxAtStartPar
A group opened with \sphinxcode{\sphinxupquote{\&}} instead of \sphinxcode{\sphinxupquote{\$}} is ignored like any other text: in version 25, \sphinxcode{\sphinxupquote{\&dft func=BLYP \$end}} left the
functional at PBE, with no message. Always use \sphinxcode{\sphinxupquote{\$}}.




\section{Appending results of earlier runs}
\label{\detokenize{start/input-syntax:appending-results-of-earlier-runs}}
\sphinxAtStartPar
Priroda writes results that can be reused as input into the main output file. Each such line carries a prefix, for
example \sphinxcode{\sphinxupquote{eng\textgreater{}}} or \sphinxcode{\sphinxupquote{num\textgreater{}}}. Strip the prefix and append the lines to the next input. The program has no separate
restart file for these data; SCF orbitals are the exception, saved with {\hyperref[\detokenize{reference/control:control-save}]{\sphinxcrossref{\DUrole{std}{\DUrole{std-ref}{\$control save}}}}}.




\subsection{\sphinxstyleliteralintitle{\sphinxupquote{\$Energy}} from \sphinxstyleliteralintitle{\sphinxupquote{eng\textgreater{}}} lines}
\label{\detokenize{start/input-syntax:energy-from-eng-lines}}
\sphinxAtStartPar
A run prints its energy, dipole, gradient and, after \sphinxcode{\sphinxupquote{task=hessian}}, the Hessian as an
{\hyperref[\detokenize{reference/Energy:group-energy}]{\sphinxcrossref{\DUrole{std}{\DUrole{std-ref}{\$Energy}}}}} group on lines prefixed with \sphinxcode{\sphinxupquote{eng\textgreater{}}}:

\begin{sphinxVerbatim}[commandchars=\\\{\}]
eng\PYGZgt{}\PYGZdl{}Energy
eng\PYGZgt{} E=\PYGZhy{}2.894457077327e+04
eng\PYGZgt{} D= 0.00000000e+00  0.00000000e+00  0.00000000e+00
eng\PYGZgt{} G= 3.53877831e\PYGZhy{}09  1.65504544e\PYGZhy{}08 \PYGZhy{}3.64771621e\PYGZhy{}07
eng\PYGZgt{}   \PYGZhy{}1.31128901e\PYGZhy{}07 \PYGZhy{}6.73350427e\PYGZhy{}09 \PYGZhy{}2.19780956e\PYGZhy{}06
eng\PYGZgt{}    1.27590122e\PYGZhy{}07 \PYGZhy{}9.81695012e\PYGZhy{}09  2.56258122e\PYGZhy{}06
eng\PYGZgt{}\PYGZdl{}end
\end{sphinxVerbatim}



\sphinxAtStartPar
To reuse it, remove the prefix and append the block to the end of the next input, after \sphinxcode{\sphinxupquote{\$molecule}}:

\begin{sphinxVerbatim}[commandchars=\\\{\}]
sed \PYGZhy{}n \PYGZsq{}s/\PYGZca{}eng\PYGZgt{}//p\PYGZsq{} mol.out \PYGZgt{}\PYGZgt{} mol.in
\end{sphinxVerbatim}

\sphinxAtStartPar
The \sphinxcode{\sphinxupquote{ge}} script on the official site does the same. It also skips the first block when the output contains a
\sphinxcode{\sphinxupquote{NUMERICAL AND ANALYTICAL GRADIENT}} section.



\begin{sphinxadmonition}{warning}{Warning:}
\sphinxAtStartPar
An optimization prints an \sphinxcode{\sphinxupquote{eng\textgreater{}}} block at every step, and the command above appends all of them. Append only the
block you need, normally the one from a \sphinxcode{\sphinxupquote{task=hessian}} run. If the input holds several \sphinxcode{\sphinxupquote{\$Energy}} groups, the
program uses the \sphinxstylestrong{last} one, even when an earlier one matches the geometry.
\end{sphinxadmonition}



\sphinxAtStartPar
A stored Hessian is required by {\hyperref[\detokenize{reference/optimize:optimize-saddle}]{\sphinxcrossref{\DUrole{std}{\DUrole{std-ref}{\$optimize saddle=1}}}}} and \sphinxcode{\sphinxupquote{task=irc}}. It also helps difficult
minimizations, and it lets you redo the vibrational analysis or the thermochemistry without recomputing.




\subsection{\sphinxstyleliteralintitle{\sphinxupquote{num\textgreater{}}} lines}
\label{\detokenize{start/input-syntax:num-lines}}
\sphinxAtStartPar
Partial results of a numerical calculation are printed on lines prefixed with \sphinxcode{\sphinxupquote{num\textgreater{}}}. There are two uses:
\begin{itemize}
\item {} 
\sphinxAtStartPar
A numerical Hessian split over several runs with {\hyperref[\detokenize{reference/d2edr2:d2edr2-steps}]{\sphinxcrossref{\DUrole{std}{\DUrole{std-ref}{\$d2edr2 steps}}}}}. Collect the \sphinxcode{\sphinxupquote{num\textgreater{}}} lines of
all partial outputs into one input, add \sphinxcode{\sphinxupquote{\$d2edr2 steps=\sphinxhyphen{}1 \$end}}, and run once more:

\begin{sphinxVerbatim}[commandchars=\\\{\}]
grep \PYGZhy{}h \PYGZdq{}\PYGZca{}num\PYGZgt{}\PYGZdq{} mol.out1 mol.out2 mol.out3 mol.out4 | sed \PYGZsq{}s/\PYGZca{}num\PYGZgt{}//\PYGZsq{} \PYGZgt{}\PYGZgt{} mol.in
\end{sphinxVerbatim}

\sphinxAtStartPar
The lines hold numbered \sphinxcode{\sphinxupquote{\$Energy000}}, \sphinxcode{\sphinxupquote{\$Energy001}}, … groups, one per displaced geometry. The combining run does
no SCF of its own. The command given in \sphinxcode{\sphinxupquote{input.txt}}, \sphinxcode{\sphinxupquote{sed "s/\textbackslash{}.*num\textgreater{}//g"}}, leaves the file names that \sphinxcode{\sphinxupquote{grep}} adds
for several files at the start of each line; use \sphinxcode{\sphinxupquote{grep \sphinxhyphen{}h}} as above.

\item {} 
\sphinxAtStartPar
Raman intensities. A \sphinxcode{\sphinxupquote{task=dipole}} run writes \sphinxcode{\sphinxupquote{num\textgreater{}}} lines; append them to the input of a \sphinxcode{\sphinxupquote{task=hessian}} run, and
the \sphinxcode{\sphinxupquote{Raman A.}} and \sphinxcode{\sphinxupquote{Raman D.}} columns of the frequency table are filled in.

\end{itemize}




\subsection{Optimized geometry from \sphinxstyleliteralintitle{\sphinxupquote{MOL\textgreater{}}} lines}
\label{\detokenize{start/input-syntax:optimized-geometry-from-mol-lines}}
\sphinxAtStartPar
After \sphinxcode{\sphinxupquote{OPTIMIZATION CONVERGED}}, the final geometry is printed as a ready\sphinxhyphen{}made \sphinxcode{\sphinxupquote{\$molecule}} group on lines prefixed with
\sphinxcode{\sphinxupquote{MOL\textgreater{}}}. Strip the prefix to continue from that geometry. The \sphinxcode{\sphinxupquote{gm}} script on the official site replaces the \sphinxcode{\sphinxupquote{\$molecule}}
group of the input with it. Because \sphinxcode{\sphinxupquote{gm}} replaces everything from \sphinxcode{\sphinxupquote{\$molecule}} to the end of the file, keep
\sphinxcode{\sphinxupquote{\$molecule}} as the last group of your inputs, followed only by appended result blocks.




\section{Common mistakes}
\label{\detokenize{start/input-syntax:common-mistakes}}\begin{description}
\sphinxlineitem{Non\sphinxhyphen{}integer atomic numbers}
\sphinxAtStartPar
\sphinxcode{\sphinxupquote{25.0}} or \sphinxcode{\sphinxupquote{8.0}} instead of \sphinxcode{\sphinxupquote{25}} or \sphinxcode{\sphinxupquote{8}} in \sphinxcode{\sphinxupquote{\$molecule}}. The atom lines are read as a stream of numbers, so the
rest of the group shifts. In version 25 this gives no error at all: water written with \sphinxcode{\sphinxupquote{8.0}} and \sphinxcode{\sphinxupquote{1.0}} ran as an
oxygen atom plus three empty centres with Z=0, and printed a meaningless energy. Older versions were reported to
stop with \sphinxcode{\sphinxupquote{illegal charge/multiplicity}}. Check the \sphinxcode{\sphinxupquote{formula:}} line near the top of the output. Element symbols
(\sphinxcode{\sphinxupquote{O}}) are not accepted either.



\sphinxlineitem{Typographic quotes}
\sphinxAtStartPar
Word processors and some web pages turn \sphinxcode{\sphinxupquote{"}} into curly or angle quotes, which break file names in \sphinxcode{\sphinxupquote{basis=}} and
\sphinxcode{\sphinxupquote{path=}}. Retype them as straight quotes.



\sphinxlineitem{Two \sphinxcode{\sphinxupquote{task=}} lines}
\sphinxAtStartPar
A second \sphinxcode{\sphinxupquote{task=}} line does not start a second job; only the first task runs. \sphinxcode{\sphinxupquote{task=optimize+hessian}} is the only
combination the program understands. Run other sequences as separate jobs.



\sphinxlineitem{Windows line endings (CRLF)}
\sphinxAtStartPar
The run stops in \sphinxcode{\sphinxupquote{\$molecule}} with \sphinxcode{\sphinxupquote{molecule input: \textquotesingle{}\textquotesingle{}}} followed by a garbled \sphinxcode{\sphinxupquote{\textquotesingle{} ?matrix}} or \sphinxcode{\sphinxupquote{?rtesian}}. Convert the
file with \sphinxcode{\sphinxupquote{dos2unix}}; see {\hyperref[\detokenize{start/running::doc}]{\sphinxcrossref{\DUrole{doc}{Installing and running}}}}.



\sphinxlineitem{Misspelled task or theory name}
\sphinxAtStartPar
Gives no error. An unknown \sphinxcode{\sphinxupquote{task=}} value sets up the run and stops quietly after the SCF and grid settings; an
unknown \sphinxcode{\sphinxupquote{theory=}} value runs Hartree\textendash{}Fock, with the misspelled name in the summary box. Check the
\sphinxcode{\sphinxupquote{Theoretical Method}} line and that the output goes on to an SCF.



\sphinxlineitem{Hybrid functional without \sphinxcode{\sphinxupquote{theory=ridft}}}
\sphinxAtStartPar
Older versions ran without the exact exchange. With \sphinxcode{\sphinxupquote{theory=dft}}, version 25 stops after the basis\sphinxhyphen{}set echo with
no message; with \sphinxcode{\sphinxupquote{theory=tddft}} it still runs without the exact exchange and gives wrong energies. See
{\hyperref[\detokenize{reference/control:control-theory}]{\sphinxcrossref{\DUrole{std}{\DUrole{std-ref}{\$control theory}}}}}.



\sphinxlineitem{\sphinxcode{\sphinxupquote{basis4.in}} without \sphinxcode{\sphinxupquote{four=1}}}
\sphinxAtStartPar
The relativistic basis file needs \sphinxcode{\sphinxupquote{four=1}} in \sphinxcode{\sphinxupquote{\$control}}.



\end{description}

\sphinxAtStartPar
More input problems, with their exact error messages, are listed in {\hyperref[\detokenize{troubleshooting/common-problems:troubleshooting-input}]{\sphinxcrossref{\DUrole{std}{\DUrole{std-ref}{The input is not read, or is read wrongly}}}}}.

\sphinxstepscope




\chapter{Input reference}
\label{\detokenize{reference/index:input-reference}}\label{\detokenize{reference/index::doc}}

\begin{savenotes}\sphinxattablestart
\sphinxthistablewithglobalstyle
\centering
\begin{tabulary}{\linewidth}[t]{TT}
\sphinxtoprule
\sphinxstyletheadfamily 
\sphinxAtStartPar
Group
&\sphinxstyletheadfamily 
\sphinxAtStartPar
Purpose
\\
\sphinxmidrule
\sphinxtableatstartofbodyhook
\sphinxAtStartPar
{\hyperref[\detokenize{reference/system:group-system}]{\sphinxcrossref{\DUrole{std}{\DUrole{std-ref}{\$system}}}}}
&
\sphinxAtStartPar
Memory, disk and scratch\sphinxhyphen{}directory limits.
\\
\sphinxhline
\sphinxAtStartPar
{\hyperref[\detokenize{reference/control:group-control}]{\sphinxcrossref{\DUrole{std}{\DUrole{std-ref}{\$control}}}}}
&
\sphinxAtStartPar
Job type, level of theory, basis\sphinxhyphen{}set file, Hamiltonian and printing.
\\
\sphinxhline
\sphinxAtStartPar
{\hyperref[\detokenize{reference/molecule:group-molecule}]{\sphinxcrossref{\DUrole{std}{\DUrole{std-ref}{\$molecule}}}}}
&
\sphinxAtStartPar
Charge, multiplicity, geometry (Cartesian and/or Z\sphinxhyphen{}matrix) and per\sphinxhyphen{}atom basis sets.
\\
\sphinxhline
\sphinxAtStartPar
{\hyperref[\detokenize{reference/guess:group-guess}]{\sphinxcrossref{\DUrole{std}{\DUrole{std-ref}{\$guess}}}}}
&
\sphinxAtStartPar
Initial guess for the SCF vectors.
\\
\sphinxhline
\sphinxAtStartPar
{\hyperref[\detokenize{reference/scf:group-scf}]{\sphinxcrossref{\DUrole{std}{\DUrole{std-ref}{\$scf}}}}}
&
\sphinxAtStartPar
Spin restriction, convergence, iteration limits and SCF solver.
\\
\sphinxhline
\sphinxAtStartPar
{\hyperref[\detokenize{reference/dft:group-dft}]{\sphinxcrossref{\DUrole{std}{\DUrole{std-ref}{\$dft}}}}}
&
\sphinxAtStartPar
Exchange\textendash{}correlation functional.
\\
\sphinxhline
\sphinxAtStartPar
{\hyperref[\detokenize{reference/grid:group-grid}]{\sphinxcrossref{\DUrole{std}{\DUrole{std-ref}{\$grid}}}}}
&
\sphinxAtStartPar
Adaptive integration grid for DFT.
\\
\sphinxhline
\sphinxAtStartPar
{\hyperref[\detokenize{reference/integral:group-integral}]{\sphinxcrossref{\DUrole{std}{\DUrole{std-ref}{\$integral}}}}}
&
\sphinxAtStartPar
Integral evaluation for RI methods.
\\
\sphinxhline
\sphinxAtStartPar
{\hyperref[\detokenize{reference/atoms:group-atoms}]{\sphinxcrossref{\DUrole{std}{\DUrole{std-ref}{\$atoms}}}}}
&
\sphinxAtStartPar
Per\sphinxhyphen{}element core definitions for correlated (HF\sphinxhyphen{}based) methods.
\\
\sphinxhline
\sphinxAtStartPar
{\hyperref[\detokenize{reference/optimize:group-optimize}]{\sphinxcrossref{\DUrole{std}{\DUrole{std-ref}{\$optimize}}}}}
&
\sphinxAtStartPar
Geometry optimization, saddle\sphinxhyphen{}point search, IRC and relaxed scans.
\\
\sphinxhline
\sphinxAtStartPar
{\hyperref[\detokenize{reference/d2edr2:group-d2edr2}]{\sphinxcrossref{\DUrole{std}{\DUrole{std-ref}{\$d2edr2}}}}}
&
\sphinxAtStartPar
Numerical Hessian by finite differences of the gradient.
\\
\sphinxhline
\sphinxAtStartPar
{\hyperref[\detokenize{reference/thermo:group-thermo}]{\sphinxcrossref{\DUrole{std}{\DUrole{std-ref}{\$thermo}}}}}
&
\sphinxAtStartPar
Thermochemistry after a Hessian calculation.
\\
\sphinxhline
\sphinxAtStartPar
{\hyperref[\detokenize{reference/nmr:group-nmr}]{\sphinxcrossref{\DUrole{std}{\DUrole{std-ref}{\$nmr}}}}}
&
\sphinxAtStartPar
Reference shieldings and extra points for NMR shielding calculations.
\\
\sphinxhline
\sphinxAtStartPar
{\hyperref[\detokenize{reference/Energy:group-energy}]{\sphinxcrossref{\DUrole{std}{\DUrole{std-ref}{\$Energy}}}}}
&
\sphinxAtStartPar
Results of an earlier run (energy, gradient, Hessian, dipole data), read back as input.
\\
\sphinxhline
\sphinxAtStartPar
{\hyperref[\detokenize{reference/dynamic:group-dynamic}]{\sphinxcrossref{\DUrole{std}{\DUrole{std-ref}{\$dynamic}}}}}
&
\sphinxAtStartPar
Molecular dynamics with a thermostat.
\\
\sphinxhline
\sphinxAtStartPar
{\hyperref[\detokenize{reference/uniq:group-uniq}]{\sphinxcrossref{\DUrole{std}{\DUrole{std-ref}{\$uniq}}}}}
&
\sphinxAtStartPar
Duplicate\sphinxhyphen{}structure filtering over several \sphinxcode{\sphinxupquote{\$molecule}} blocks.
\\
\sphinxbottomrule
\end{tabulary}
\sphinxtableafterendhook\par
\sphinxattableend\end{savenotes}

\sphinxstepscope




\section{\sphinxstyleliteralintitle{\sphinxupquote{\$system}}}
\label{\detokenize{reference/system:system}}\label{\detokenize{reference/system:group-system}}\label{\detokenize{reference/system::doc}}
\index{\$system@\spxentry{\$system}}\ignorespaces 
\sphinxAtStartPar
Memory, disk and scratch\sphinxhyphen{}directory limits.

\sphinxAtStartPar
All limits apply \sphinxstylestrong{per MPI process}. For example, \sphinxcode{\sphinxupquote{disk=10}} on four processes can use 40 GB. In a parallel run,
most of the memory and all of the disk storage are used by the master process (rank 0).


\subsection{\sphinxstyleliteralintitle{\sphinxupquote{memory}}}
\label{\detokenize{reference/system:memory}}\label{\detokenize{reference/system:system-memory}}
\index{memory (\$system)@\spxentry{memory}\spxextra{\$system}}\ignorespaces 
\sphinxAtStartPar
\sphinxstylestrong{Type:} list\\
\sphinxstylestrong{Default:} \sphinxcode{\sphinxupquote{256}}\\
\sphinxstylestrong{Units:} MB\\
\sphinxstylestrong{Status:} current · verified on 25 (2025.02.28)\\
\sphinxstylestrong{Sources:} input.txt; 01 \S{}4.1 (two values); 06 \S{}3 (estimate); 03 \S{}4, 04 \#196 (per process); 08 t02

\sphinxAtStartPar
Default: echoed by version 25 when \sphinxcode{\sphinxupquote{\$system}} is omitted.

\sphinxAtStartPar
Memory limit in megabytes. A second value sets the limit for the other MPI ranks: \sphinxcode{\sphinxupquote{mem=2048,256}} gives rank 0
2048 MB and every other rank 256 MB. A single value applies to all ranks. The output echoes the result as
\sphinxcode{\sphinxupquote{Memory: \textless{}mem\textgreater{}m + \textless{}|disk|\textgreater{}m (\textless{}mem2\textgreater{}m + \textless{}|disk2|\textgreater{}m)}}.

\begin{sphinxadmonition}{note}{Note:}
\sphinxAtStartPar
For DFT, the v5 documentation estimates the memory as about 8·N\(\sp{\text{2}}\)·3 bytes for energy and gradient, and twice that
for an analytic Hessian, where N is the basis dimension. \sphinxcode{\sphinxupquote{Dft.rtf}} gives an old default of 16 MB.
Without a \sphinxcode{\sphinxupquote{\$system}} group, version 25 echoes \sphinxcode{\sphinxupquote{Memory: 256m + 512m (256m + 512m)}}: 256 MB of memory plus a
512 MB RAM disk per process, read from the echo format.
\end{sphinxadmonition}


\subsection{\sphinxstyleliteralintitle{\sphinxupquote{disk}}}
\label{\detokenize{reference/system:disk}}\label{\detokenize{reference/system:system-disk}}
\index{disk (\$system)@\spxentry{disk}\spxextra{\$system}}\ignorespaces 
\sphinxAtStartPar
\sphinxstylestrong{Type:} list\\
\sphinxstylestrong{Default:} \sphinxcode{\sphinxupquote{\sphinxhyphen{}512}}\\
\sphinxstylestrong{Status:} current · not verified\\
\sphinxstylestrong{Sources:} input.txt; Dft.rtf (one file per GB); 01 \S{}4.1; 04 \S{}0, \#468; 03 \S{}17; 08 t02

\sphinxAtStartPar
Default: a 512 MB RAM disk, read from the version 25 echo \sphinxcode{\sphinxupquote{Memory: 256m + 512m}} when \sphinxcode{\sphinxupquote{\$system}} is omitted.

\sphinxAtStartPar
Positive values set the disk space in gigabytes; one scratch file is opened per gigabyte, which sidesteps
per\sphinxhyphen{}file size limits. Negative values set the amount of memory, in megabytes, used to emulate disk
(a RAM disk). A second value applies to the other MPI ranks, as for \sphinxcode{\sphinxupquote{memory}}; the official example inputs typically give
them \sphinxcode{\sphinxupquote{\sphinxhyphen{}1}}.

\begin{sphinxadmonition}{note}{Note:}
\sphinxAtStartPar
With a RAM disk, \sphinxcode{\sphinxupquote{memory}} plus |\sphinxcode{\sphinxupquote{disk}}| (times the number of processes) must fit in physical RAM, or the run
fails with \sphinxcode{\sphinxupquote{malloc(...) failed}}. Example for a 24 GB node: \sphinxcode{\sphinxupquote{memory=9000 disk=\sphinxhyphen{}14000}}.
Keeping scratch in RAM gave a 6× speed\sphinxhyphen{}up for RI\sphinxhyphen{}MP2 on one cluster.
\end{sphinxadmonition}

\sphinxAtStartPar
\sphinxstylestrong{Example:}

\begin{sphinxVerbatim}[commandchars=\\\{\}]
\PYGZdl{}system mem=2048,256 disk=\PYGZhy{}256,\PYGZhy{}1 \PYGZdl{}end
\end{sphinxVerbatim}


\subsection{\sphinxstyleliteralintitle{\sphinxupquote{path}}}
\label{\detokenize{reference/system:path}}\label{\detokenize{reference/system:system-path}}
\index{path (\$system)@\spxentry{path}\spxextra{\$system}}\ignorespaces 
\sphinxAtStartPar
\sphinxstylestrong{Type:} string\\
\sphinxstylestrong{Default:} \sphinxstyleemphasis{none stated}\\
\sphinxstylestrong{Status:} current · verified on 25 (2025.02.28)\\
\sphinxstylestrong{Sources:} input.txt; 03 \S{}2; 04 \#465; 08 t01, t64

\sphinxAtStartPar
Default: the \sphinxcode{\sphinxupquote{TMPDIR}} environment variable, otherwise \sphinxcode{\sphinxupquote{/tmp}} (version 25). \sphinxcode{\sphinxupquote{input.txt}} says the current working directory.

\sphinxAtStartPar
Directory for temporary disk files. For parallel runs you can give a colon\sphinxhyphen{}separated list with one path per
process: \sphinxcode{\sphinxupquote{path=/tmpa:/tmpb:/tmpc:/tmpd}} makes process 1 write to \sphinxcode{\sphinxupquote{/tmpa}}, process 2 to \sphinxcode{\sphinxupquote{/tmpb}}, and so on.
Those directories should be on different physical drives. The value may be double\sphinxhyphen{}quoted.

\begin{sphinxadmonition}{note}{Note:}
\sphinxAtStartPar
For all methods except GGA DFT, performance degrades “terribly” if more than one process reads and writes
large files on the same physical drive. Never put temporary files on NFS. On Windows (Cygwin), which has no
\sphinxcode{\sphinxupquote{/tmp}}, use \sphinxcode{\sphinxupquote{path=.}}. Version 25, run without \sphinxcode{\sphinxupquote{path=}}, used \sphinxcode{\sphinxupquote{TMPDIR}} when it was set and \sphinxcode{\sphinxupquote{/tmp}} when it was
not (echo \sphinxcode{\sphinxupquote{Disk:   1g \textquotesingle{}/tmp\textquotesingle{}}}).
\end{sphinxadmonition}

\sphinxAtStartPar
\sphinxstylestrong{Example:}

\begin{sphinxVerbatim}[commandchars=\\\{\}]
\PYGZdl{}system path=/tmpa:/tmpb:/tmpc:/tmpd \PYGZdl{}end
\end{sphinxVerbatim}

\sphinxstepscope




\section{\sphinxstyleliteralintitle{\sphinxupquote{\$control}}}
\label{\detokenize{reference/control:control}}\label{\detokenize{reference/control:group-control}}\label{\detokenize{reference/control::doc}}
\index{\$control@\spxentry{\$control}}\ignorespaces 
\sphinxAtStartPar
Job type, level of theory, basis\sphinxhyphen{}set file, Hamiltonian and printing.


\subsection{\sphinxstyleliteralintitle{\sphinxupquote{task}}}
\label{\detokenize{reference/control:task}}\label{\detokenize{reference/control:control-task}}
\index{task (\$control)@\spxentry{task}\spxextra{\$control}}\ignorespaces 
\sphinxAtStartPar
\sphinxstylestrong{Type:} enum\\
\sphinxstylestrong{Default:} \sphinxcode{\sphinxupquote{gradient}}\\
\sphinxstylestrong{Status:} current · verified on 25 (2025.02.28)\\
\sphinxstylestrong{Sources:} input.txt; Dft.rtf (default gradient, check); 02 \S{}3.2 (pri: default gradient); 03 \S{}5; 04 \#69\sphinxhyphen{}\#73, \#145, \#478, \#497; 01 \S{}4.2; 08 t02, t18, t26, t34, t36, t37; 09; 08 t50, t57a

\sphinxAtStartPar
Job type.


\begin{savenotes}\sphinxattablestart
\sphinxthistablewithglobalstyle
\sphinxthistablewithnovlinesstyle
\centering
\begin{tabulary}{\linewidth}[t]{\Y{0.3}\Y{0.7}}
\sphinxtoprule
\sphinxstyletheadfamily 
\sphinxAtStartPar
Value
&\sphinxstyletheadfamily 
\sphinxAtStartPar
Meaning
\\
\sphinxmidrule
\sphinxtableatstartofbodyhook
\sphinxAtStartPar
\sphinxcode{\sphinxupquote{energy}}
&
\sphinxAtStartPar
Energy at the given geometry.
\\
\sphinxhline
\sphinxAtStartPar
\sphinxcode{\sphinxupquote{gradient}}
&
\sphinxAtStartPar
Energy and nuclear gradient. Costs little more than the energy. (\sphinxstylestrong{default})
\\
\sphinxhline
\sphinxAtStartPar
\sphinxcode{\sphinxupquote{hessian}}
&
\sphinxAtStartPar
Second derivatives (force constants) followed by vibrational analysis and thermochemistry.
\\
\sphinxhline
\sphinxAtStartPar
\sphinxcode{\sphinxupquote{optimize}}
&
\sphinxAtStartPar
Geometry optimization, including saddle points with \sphinxcode{\sphinxupquote{\$optimize saddle=1}}.
\\
\sphinxhline
\sphinxAtStartPar
\sphinxcode{\sphinxupquote{optimize+hessian}}
&
\sphinxAtStartPar
Optimization followed by a Hessian at the final geometry. The only combined task that works. Present in v2 and v9 and later, absent in v6.
\\
\sphinxhline
\sphinxAtStartPar
\sphinxcode{\sphinxupquote{irc}}
&
\sphinxAtStartPar
Intrinsic reaction coordinate. Needs a Hessian at the TS in the input (\sphinxcode{\sphinxupquote{\$Energy}}). Each point is printed on \sphinxcode{\sphinxupquote{IRC\textgreater{}}} lines with its distance and energy.
\\
\sphinxhline
\sphinxAtStartPar
\sphinxcode{\sphinxupquote{scan}}
&
\sphinxAtStartPar
Relaxed potential\sphinxhyphen{}energy\sphinxhyphen{}surface scan over the coordinates in \sphinxcode{\sphinxupquote{\$optimize fix}}. The structures and energies are summarized on \sphinxcode{\sphinxupquote{SCAN\textgreater{}}} lines at the end.
\\
\sphinxhline
\sphinxAtStartPar
\sphinxcode{\sphinxupquote{nmr}}
&
\sphinxAtStartPar
NMR shielding tensors. Works only with \sphinxcode{\sphinxupquote{theory=dft}}: with \sphinxcode{\sphinxupquote{ridft}}, version 25 silently runs only the ground\sphinxhyphen{}state SCF.
\\
\sphinxhline
\sphinxAtStartPar
\sphinxcode{\sphinxupquote{dipole}}
&
\sphinxAtStartPar
Numerical dipole polarizabilities and their derivatives with respect to nuclear coordinates (used for Raman intensities), from SCF runs in finite electric fields. The result is a \sphinxcode{\sphinxupquote{num\textgreater{}}} block.
\\
\sphinxhline
\sphinxAtStartPar
\sphinxcode{\sphinxupquote{check}}
&
\sphinxAtStartPar
Not a task name in the version 25 program. Like any unrecognized task, \sphinxcode{\sphinxupquote{task=check}} sets up the run (geometry, basis sets, SCF and grid settings) and stops before the SCF without a message, which makes it a usable input check. \sphinxstyleemphasis{(undocumented)}
\\
\sphinxhline
\sphinxAtStartPar
\sphinxcode{\sphinxupquote{extremal}}
&
\sphinxAtStartPar
Gradient\sphinxhyphen{}extremal path (v9 banner “Calculation of Gradient Extremal Path”); slow. \sphinxstyleemphasis{(undocumented)}
\\
\sphinxhline
\sphinxAtStartPar
\sphinxcode{\sphinxupquote{transit}}
&
\sphinxAtStartPar
Synchronous transit in internal coordinates. Needs two \sphinxcode{\sphinxupquote{\$molecule}} blocks (start and end). Reported working in v8. \sphinxstyleemphasis{(undocumented)}
\\
\sphinxhline
\sphinxAtStartPar
\sphinxcode{\sphinxupquote{saddle}}
&
\sphinxAtStartPar
Approximate saddle\sphinxhyphen{}point search from two \sphinxcode{\sphinxupquote{\$molecule}} blocks (start and end). \sphinxstyleemphasis{(undocumented)}
\\
\sphinxhline
\sphinxAtStartPar
\sphinxcode{\sphinxupquote{path}}
&
\sphinxAtStartPar
Minimum\sphinxhyphen{}energy\sphinxhyphen{}path optimization from two \sphinxcode{\sphinxupquote{\$molecule}} blocks (start and end). \sphinxstyleemphasis{(undocumented)}
\\
\sphinxhline
\sphinxAtStartPar
\sphinxcode{\sphinxupquote{conical}}
&
\sphinxAtStartPar
Conical intersection of states 0 and 1 (v9 banner). Untested. \sphinxstyleemphasis{(undocumented)}
\\
\sphinxhline
\sphinxAtStartPar
\sphinxcode{\sphinxupquote{anharmonic}}
&
\sphinxAtStartPar
Anharmonic corrections to vibrational energies; slow. \sphinxstyleemphasis{(undocumented)}
\\
\sphinxhline
\sphinxAtStartPar
\sphinxcode{\sphinxupquote{opt\_n1}}
&
\sphinxAtStartPar
Optimizer for very large systems (used for 8060 atoms with \sphinxcode{\sphinxupquote{theory=qm\_n2}}). Pairs with \sphinxcode{\sphinxupquote{\$optimize mod=1}}. \sphinxstyleemphasis{(undocumented)}
\\
\sphinxhline
\sphinxAtStartPar
\sphinxcode{\sphinxupquote{thermostat}}
&
\sphinxAtStartPar
Molecular dynamics with a thermostat; see \sphinxcode{\sphinxupquote{\$dynamic}}. \sphinxstyleemphasis{(undocumented)}
\\
\sphinxhline
\sphinxAtStartPar
\sphinxcode{\sphinxupquote{uniq}}
&
\sphinxAtStartPar
Filter duplicate structures among several \sphinxcode{\sphinxupquote{\$molecule}} blocks; see \sphinxcode{\sphinxupquote{\$uniq}}. \sphinxstyleemphasis{(undocumented)}
\\
\sphinxbottomrule
\end{tabulary}
\sphinxtableafterendhook\par
\sphinxattableend\end{savenotes}

\begin{sphinxadmonition}{note}{Note:}
\sphinxAtStartPar
Two \sphinxcode{\sphinxupquote{task=}} lines do not chain jobs: only the first one runs. Apart from \sphinxcode{\sphinxupquote{optimize+hessian}}, sequences such
as optimize\(\rightarrow\)NMR need separate runs or an external driver script.
\sphinxstylestrong{A misspelled task name gives no error}: version 25 sets up the run and stops quietly, as for \sphinxcode{\sphinxupquote{check}}.
\sphinxcode{\sphinxupquote{task=critical}} does not work; density critical points come from \sphinxcode{\sphinxupquote{print=+density}}.
\sphinxcode{\sphinxupquote{task=QM}} appears only in the header of the shipped \sphinxcode{\sphinxupquote{qm.in}} parameter file and is not a user task.
The binary contains further task words (\sphinxcode{\sphinxupquote{IRC4}}, \sphinxcode{\sphinxupquote{PATHX}}, \sphinxcode{\sphinxupquote{PATHR}}, \sphinxcode{\sphinxupquote{STERIC}}, \sphinxcode{\sphinxupquote{POLARIZABILITY}}, \sphinxcode{\sphinxupquote{MMDYNAMIC}},
and others) that are untested or need data that is not distributed.
\end{sphinxadmonition}

\sphinxAtStartPar
\sphinxstylestrong{See also:} {\hyperref[\detokenize{reference/optimize:optimize-saddle}]{\sphinxcrossref{\DUrole{std}{\DUrole{std-ref}{saddle (\$optimize)}}}}}, {\hyperref[\detokenize{reference/optimize:optimize-points}]{\sphinxcrossref{\DUrole{std}{\DUrole{std-ref}{points (\$optimize)}}}}}, {\hyperref[\detokenize{reference/control:control-print}]{\sphinxcrossref{\DUrole{std}{\DUrole{std-ref}{print (\$control)}}}}}


\subsection{\sphinxstyleliteralintitle{\sphinxupquote{theory}}}
\label{\detokenize{reference/control:theory}}\label{\detokenize{reference/control:control-theory}}
\index{theory (\$control)@\spxentry{theory}\spxextra{\$control}}\ignorespaces 
\sphinxAtStartPar
\sphinxstylestrong{Type:} enum\\
\sphinxstylestrong{Default:} \sphinxcode{\sphinxupquote{dft}}\\
\sphinxstylestrong{Status:} current · verified on 25 (2025.02.28)\\
\sphinxstylestrong{Sources:} input.txt; Dft.rtf (default dft); 06 \S{}3 (v5 ab initio list); 03 \S{}5 (Evgeniy 2011 test); 04 \#96, \#99, \#307\sphinxhyphen{}\#310; 01 \S{}3.1, \S{}4.2; 08 t02, t08, t17, t31\sphinxhyphen{}t33, t38, t39, t46; 08 t50b

\sphinxAtStartPar
Theoretical method.


\begin{savenotes}\sphinxattablestart
\sphinxthistablewithglobalstyle
\sphinxthistablewithnovlinesstyle
\centering
\begin{tabulary}{\linewidth}[t]{\Y{0.3}\Y{0.7}}
\sphinxtoprule
\sphinxstyletheadfamily 
\sphinxAtStartPar
Value
&\sphinxstyletheadfamily 
\sphinxAtStartPar
Meaning
\\
\sphinxmidrule
\sphinxtableatstartofbodyhook
\sphinxAtStartPar
\sphinxcode{\sphinxupquote{dft}}
&
\sphinxAtStartPar
GGA DFT with density fitting. Hybrid functionals are \sphinxstylestrong{not} applied here: see notes. (\sphinxstylestrong{default})
\\
\sphinxhline
\sphinxAtStartPar
\sphinxcode{\sphinxupquote{ridft}}
&
\sphinxAtStartPar
Hybrid DFT with exact exchange via the resolution\sphinxhyphen{}of\sphinxhyphen{}identity (RI) approximation. Required for hybrid functionals; the default functional is then \sphinxcode{\sphinxupquote{PBE1}}. Hessians are numerical.
\\
\sphinxhline
\sphinxAtStartPar
\sphinxcode{\sphinxupquote{DFT\_n4}}
&
\sphinxAtStartPar
Newer DFT variant, used in the official example inputs with \sphinxcode{\sphinxupquote{func=LPBEVV}} (2025). Runs in version 25 (summary box: \sphinxcode{\sphinxupquote{DFT\_N4}}); what it changes is not documented. \sphinxstyleemphasis{(undocumented)}
\\
\sphinxhline
\sphinxAtStartPar
\sphinxcode{\sphinxupquote{tddft}}
&
\sphinxAtStartPar
GGA time\sphinxhyphen{}dependent DFT (no hybrid TDDFT).
\\
\sphinxhline
\sphinxAtStartPar
\sphinxcode{\sphinxupquote{rimp2}}
&
\sphinxAtStartPar
RI\sphinxhyphen{}MP2. Hessians are numerical.
\\
\sphinxhline
\sphinxAtStartPar
\sphinxcode{\sphinxupquote{hf}}
&
\sphinxAtStartPar
Hartree\textendash{}Fock.
\\
\sphinxhline
\sphinxAtStartPar
\sphinxcode{\sphinxupquote{rihf}}
&
\sphinxAtStartPar
RI Hartree\textendash{}Fock.
\\
\sphinxhline
\sphinxAtStartPar
\sphinxcode{\sphinxupquote{cis}}
&
\sphinxAtStartPar
CIS (RHF/UHF). Prints a table of excited states; \sphinxcode{\sphinxupquote{E =}} is the total energy of the first excited state. Slow and disk\sphinxhyphen{}hungry.
\\
\sphinxhline
\sphinxAtStartPar
\sphinxcode{\sphinxupquote{mp2}}
&
\sphinxAtStartPar
Conventional MP2 (RHF, UHF, ROHF). Very disk\sphinxhyphen{}hungry.
\\
\sphinxhline
\sphinxAtStartPar
\sphinxcode{\sphinxupquote{mp3}}
&
\sphinxAtStartPar
MP3.
\\
\sphinxhline
\sphinxAtStartPar
\sphinxcode{\sphinxupquote{mp4dq}}
&
\sphinxAtStartPar
MP4(DQ).
\\
\sphinxhline
\sphinxAtStartPar
\sphinxcode{\sphinxupquote{mp4sdq}}
&
\sphinxAtStartPar
MP4(SDQ).
\\
\sphinxhline
\sphinxAtStartPar
\sphinxcode{\sphinxupquote{mp4sdtq}}
&
\sphinxAtStartPar
MP4(SDTQ) (RHF).
\\
\sphinxhline
\sphinxAtStartPar
\sphinxcode{\sphinxupquote{ccs}}
&
\sphinxAtStartPar
CCS. With an HF reference it returns the HF energy.
\\
\sphinxhline
\sphinxAtStartPar
\sphinxcode{\sphinxupquote{ccd}}
&
\sphinxAtStartPar
CCD.
\\
\sphinxhline
\sphinxAtStartPar
\sphinxcode{\sphinxupquote{ccsd}}
&
\sphinxAtStartPar
CCSD (RHF/UHF).
\\
\sphinxhline
\sphinxAtStartPar
\sphinxcode{\sphinxupquote{ccsd4t}}
&
\sphinxAtStartPar
CCSD(T).
\\
\sphinxhline
\sphinxAtStartPar
\sphinxcode{\sphinxupquote{ccsdsd}}
&
\sphinxAtStartPar
Linear\sphinxhyphen{}response (EOM) CCSD for excited states (RHF). \sphinxstyleemphasis{(undocumented)}
\\
\sphinxhline
\sphinxAtStartPar
\sphinxcode{\sphinxupquote{qm\_n3}}
&
\sphinxAtStartPar
Laikov’s semiempirical QM model. Needs \sphinxcode{\sphinxupquote{basis=qm.in}} and no \sphinxcode{\sphinxupquote{set=}}; parameters exist for H, Li\textendash{}F, Na\textendash{}Cl, Sc\textendash{}Cu. Runs in version 25, with a short output of its own (no summary box, no \sphinxcode{\sphinxupquote{E  =}} line). \sphinxstyleemphasis{(undocumented)}
\\
\sphinxhline
\sphinxAtStartPar
\sphinxcode{\sphinxupquote{qm\_n2}}
&
\sphinxAtStartPar
Earlier variant of the QM model, used for very large systems with \sphinxcode{\sphinxupquote{task=opt\_n1}}. \sphinxstyleemphasis{(undocumented)}
\\
\sphinxbottomrule
\end{tabulary}
\sphinxtableafterendhook\par
\sphinxattableend\end{savenotes}

\begin{sphinxadmonition}{note}{Note:}
\sphinxAtStartPar
\sphinxstylestrong{A misspelled theory name gives no error}: version 25 prints it in the summary box and runs Hartree\textendash{}Fock.

\sphinxAtStartPar
\sphinxstylestrong{Hybrid trap.} Hybrid functionals (\sphinxcode{\sphinxupquote{PBE1}}, \sphinxcode{\sphinxupquote{B3LYP}}, …) take effect only with \sphinxcode{\sphinxupquote{theory=ridft}}. With \sphinxcode{\sphinxupquote{theory=dft}}
or \sphinxcode{\sphinxupquote{tddft}}, older versions silently drop the exact exchange. In version 25, \sphinxcode{\sphinxupquote{theory=dft}} stops right after
printing the basis sets, with no error message in the output and exit status 1. \sphinxcode{\sphinxupquote{theory=tddft}} still runs:
it reports \sphinxcode{\sphinxupquote{PBE1}} in the summary box but leaves out the exact exchange, so the energies are wrong (water:
\sphinxcode{\sphinxupquote{Ex}} \sphinxhyphen{}6.61 instead of about \sphinxhyphen{}8.95 hartree, total energy about 2.2 hartree too high), and the run ends
normally.

\sphinxAtStartPar
The ab initio list comes from the v6 documentation and a 2011 user test, which was energy\sphinxhyphen{}only. In version 25,
water energies run to completion with \sphinxcode{\sphinxupquote{hf}}, \sphinxcode{\sphinxupquote{rihf}}, \sphinxcode{\sphinxupquote{cis}}, \sphinxcode{\sphinxupquote{mp2}}, \sphinxcode{\sphinxupquote{mp3}}, \sphinxcode{\sphinxupquote{mp4dq}}, \sphinxcode{\sphinxupquote{mp4sdq}}, \sphinxcode{\sphinxupquote{mp4sdtq}}, \sphinxcode{\sphinxupquote{ccs}},
\sphinxcode{\sphinxupquote{ccd}}, \sphinxcode{\sphinxupquote{ccsd}} and \sphinxcode{\sphinxupquote{ccsd4t}}; an \sphinxcode{\sphinxupquote{hf}} optimization also worked, so HF gradients exist. For \sphinxcode{\sphinxupquote{ridft}} and
\sphinxcode{\sphinxupquote{rimp2}}, \sphinxcode{\sphinxupquote{task=hessian}} works by finite differences of the gradient (\sphinxcode{\sphinxupquote{NUMERICAL COMPUTATION OF SECOND DERIVATIVES}}); for \sphinxcode{\sphinxupquote{dft}} it is analytic. These
were reported \sphinxstylestrong{not} working in v6: \sphinxcode{\sphinxupquote{mp4}}, \sphinxcode{\sphinxupquote{rimp4}}, \sphinxcode{\sphinxupquote{rimp4sdtq}}, \sphinxcode{\sphinxupquote{rimp4sdq}}, \sphinxcode{\sphinxupquote{riccsd}}, \sphinxcode{\sphinxupquote{riccsd4t}}. They
either crash or return the HF energy. The binary contains many more method words (CASSCF, MRCISD,
CCSDT, …); none has been reported working.
\end{sphinxadmonition}

\sphinxAtStartPar
\sphinxstylestrong{See also:} {\hyperref[\detokenize{reference/dft:dft-functional}]{\sphinxcrossref{\DUrole{std}{\DUrole{std-ref}{functional (\$dft)}}}}}


\subsection{\sphinxstyleliteralintitle{\sphinxupquote{state}}}
\label{\detokenize{reference/control:state}}\label{\detokenize{reference/control:control-state}}
\index{state (\$control)@\spxentry{state}\spxextra{\$control}}\ignorespaces 
\sphinxAtStartPar
\sphinxstylestrong{Type:} list\\
\sphinxstylestrong{Default:} \sphinxstyleemphasis{none stated}\\
\sphinxstylestrong{Status:} current · verified on 25 (2025.02.28)\\
\sphinxstylestrong{Sources:} input.txt; 08 t33

\sphinxAtStartPar
For excited\sphinxhyphen{}state methods. The first value is the excited state n whose properties are computed; the second
is the number of states N included in the iterative solution of the equations (\sphinxcode{\sphinxupquote{state=n,N}}). In version 25,
\sphinxcode{\sphinxupquote{theory=tddft}} with \sphinxcode{\sphinxupquote{state=1,5}} printed all five states in an \sphinxcode{\sphinxupquote{Electronic Spectrum:}} table (excitation energy,
wavelength, transition moment, oscillator strength), and \sphinxcode{\sphinxupquote{E  =}} was the total energy of state 1.


\subsection{\sphinxstyleliteralintitle{\sphinxupquote{basis}}}
\label{\detokenize{reference/control:basis}}\label{\detokenize{reference/control:control-basis}}
\index{basis (\$control)@\spxentry{basis}\spxextra{\$control}}\ignorespaces 
\sphinxAtStartPar
\sphinxstylestrong{Type:} string\\
\sphinxstylestrong{Default:} \sphinxstyleemphasis{none stated}\\
\sphinxstylestrong{Status:} current · verified on 25 (2025.02.28)\\
\sphinxstylestrong{Sources:} input.txt; Dft.rtf; 01 \S{}4.11, \S{}6; 03 \S{}5; 04 \S{}0; 08 t03

\sphinxAtStartPar
Default: basis\sphinxhyphen{}set data are read from the input file itself.

\sphinxAtStartPar
Name of the file containing the basis\sphinxhyphen{}set data (and ECPs, or QM parameters for \sphinxcode{\sphinxupquote{qm.in}}). The name may be
double\sphinxhyphen{}quoted, which is needed when the path contains special characters. Use straight quotes: typographic
quotes break parsing.

\begin{sphinxadmonition}{note}{Note:}
\sphinxAtStartPar
Shipped files include \sphinxcode{\sphinxupquote{basis.in}} (Λ/L family, default set \sphinxcode{\sphinxupquote{L1}}), \sphinxcode{\sphinxupquote{basis4.in}} (scalar\sphinxhyphen{}relativistic L family;
needs \sphinxcode{\sphinxupquote{four=1}}), \sphinxcode{\sphinxupquote{3z.bas}} (DFT only, Z ≤ 54), \sphinxcode{\sphinxupquote{sbk.bas}} (SBK ECPs, up to Z = 86), \sphinxcode{\sphinxupquote{b2*.in}}, \sphinxcode{\sphinxupquote{b2u*.in}},
\sphinxcode{\sphinxupquote{b4s*.in}}, and \sphinxcode{\sphinxupquote{qm.in}}. In version 25 a relative name such as \sphinxcode{\sphinxupquote{basis=basis.in}} works when the file is in the
directory the program is started from. An absolute path is still the safest choice for parallel and batch
runs.
\end{sphinxadmonition}

\sphinxAtStartPar
\sphinxstylestrong{See also:} {\hyperref[\detokenize{reference/molecule:molecule-set}]{\sphinxcrossref{\DUrole{std}{\DUrole{std-ref}{set (\$molecule)}}}}}, {\hyperref[\detokenize{reference/control:control-four}]{\sphinxcrossref{\DUrole{std}{\DUrole{std-ref}{four (\$control)}}}}}


\subsection{\sphinxstyleliteralintitle{\sphinxupquote{four}}}
\label{\detokenize{reference/control:four}}\label{\detokenize{reference/control:control-four}}
\index{four (\$control)@\spxentry{four}\spxextra{\$control}}\ignorespaces 
\sphinxAtStartPar
\sphinxstylestrong{Type:} switch\\
\sphinxstylestrong{Default:} \sphinxcode{\sphinxupquote{0}}\\
\sphinxstylestrong{Status:} current · verified on 25 (2025.02.28)\\
\sphinxstylestrong{Sources:} input.txt; 01 \S{}3.3; 04 \#207\sphinxhyphen{}\#208, \#336; 08 t35, t49

\sphinxAtStartPar
Four\sphinxhyphen{}component option.


\begin{savenotes}\sphinxattablestart
\sphinxthistablewithglobalstyle
\sphinxthistablewithnovlinesstyle
\centering
\begin{tabulary}{\linewidth}[t]{\Y{0.3}\Y{0.7}}
\sphinxtoprule
\sphinxstyletheadfamily 
\sphinxAtStartPar
Value
&\sphinxstyletheadfamily 
\sphinxAtStartPar
Meaning
\\
\sphinxmidrule
\sphinxtableatstartofbodyhook
\sphinxAtStartPar
\sphinxcode{\sphinxupquote{0}}
&
\sphinxAtStartPar
Non\sphinxhyphen{}relativistic theory. (\sphinxstylestrong{default})
\\
\sphinxhline
\sphinxAtStartPar
\sphinxcode{\sphinxupquote{1}}
&
\sphinxAtStartPar
Scalar\sphinxhyphen{}relativistic theory (K. G. Dyall, J. Chem. Phys. 100 (1994) 2118\textendash{}2127). The output reports “Hamiltonian: scalar\sphinxhyphen{}relativistic”.
\\
\sphinxbottomrule
\end{tabulary}
\sphinxtableafterendhook\par
\sphinxattableend\end{savenotes}

\begin{sphinxadmonition}{note}{Note:}
\sphinxAtStartPar
Required with the relativistic basis file \sphinxcode{\sphinxupquote{basis4.in}}. Without it, version 25 stops right after the \sphinxcode{\sphinxupquote{Basis set input:}} line with \sphinxcode{\sphinxupquote{1?}}.
\end{sphinxadmonition}

\sphinxAtStartPar
\sphinxstylestrong{See also:} {\hyperref[\detokenize{reference/control:control-nucleus}]{\sphinxcrossref{\DUrole{std}{\DUrole{std-ref}{nucleus (\$control)}}}}}, {\hyperref[\detokenize{reference/scf:scf-d1small}]{\sphinxcrossref{\DUrole{std}{\DUrole{std-ref}{d1small (\$scf)}}}}}


\subsection{\sphinxstyleliteralintitle{\sphinxupquote{nucleus}}}
\label{\detokenize{reference/control:nucleus}}\label{\detokenize{reference/control:control-nucleus}}
\index{nucleus (\$control)@\spxentry{nucleus}\spxextra{\$control}}\ignorespaces 
\sphinxAtStartPar
\sphinxstylestrong{Type:} enum\\
\sphinxstylestrong{Default:} \sphinxstyleemphasis{none stated}\\
\sphinxstylestrong{Status:} current · verified on 25 (2025.02.28)\\
\sphinxstylestrong{Sources:} input.txt; 01 \S{}4.2; 06 \S{}1; 08 t35

\sphinxAtStartPar
Default: \sphinxcode{\sphinxupquote{finite}} for four\sphinxhyphen{}component theory, \sphinxcode{\sphinxupquote{point}} for non\sphinxhyphen{}relativistic theory.

\sphinxAtStartPar
Nuclear model. The official example inputs write it as \sphinxcode{\sphinxupquote{nuc=finite}}.


\begin{savenotes}\sphinxattablestart
\sphinxthistablewithglobalstyle
\sphinxthistablewithnovlinesstyle
\centering
\begin{tabulary}{\linewidth}[t]{\Y{0.3}\Y{0.7}}
\sphinxtoprule
\sphinxstyletheadfamily 
\sphinxAtStartPar
Value
&\sphinxstyletheadfamily 
\sphinxAtStartPar
Meaning
\\
\sphinxmidrule
\sphinxtableatstartofbodyhook
\sphinxAtStartPar
\sphinxcode{\sphinxupquote{finite}}
&
\sphinxAtStartPar
Gaussian nuclear model (L. Visscher, K. G. Dyall, At. Data Nucl. Data Tables 67 (1997) 207\textendash{}224).
\\
\sphinxhline
\sphinxAtStartPar
\sphinxcode{\sphinxupquote{point}}
&
\sphinxAtStartPar
Point\sphinxhyphen{}nucleus model.
\\
\sphinxbottomrule
\end{tabulary}
\sphinxtableafterendhook\par
\sphinxattableend\end{savenotes}

\begin{sphinxadmonition}{note}{Note:}
\sphinxAtStartPar
A basis optimized for a point nucleus must not be used with the finite\sphinxhyphen{}nucleus model, and vice versa. With
\sphinxcode{\sphinxupquote{four=1}} and no \sphinxcode{\sphinxupquote{nucleus=}}, version 25 prints \sphinxcode{\sphinxupquote{Finite nucleus model (1):}} with the Gaussian parameters of each
element.
\end{sphinxadmonition}

\sphinxAtStartPar
\sphinxstylestrong{See also:} {\hyperref[\detokenize{reference/control:control-four}]{\sphinxcrossref{\DUrole{std}{\DUrole{std-ref}{four (\$control)}}}}}


\subsection{\sphinxstyleliteralintitle{\sphinxupquote{print}}}
\label{\detokenize{reference/control:print}}\label{\detokenize{reference/control:control-print}}
\index{print (\$control)@\spxentry{print}\spxextra{\$control}}\ignorespaces 
\sphinxAtStartPar
\sphinxstylestrong{Type:} flags\\
\sphinxstylestrong{Default:} \sphinxstyleemphasis{none stated}\\
\sphinxstylestrong{Status:} current · verified on 25 (2025.02.28)\\
\sphinxstylestrong{Sources:} input.txt; Dft.rtf; 06 \S{}3; 04 \#376, \#460, \#498\sphinxhyphen{}\#502, \#509; 08 t01, t45

\sphinxAtStartPar
Printing options. Flags are appended one after another; \sphinxcode{\sphinxupquote{+}} switches an option on and \sphinxcode{\sphinxupquote{\sphinxhyphen{}}} switches it off.
The leading \sphinxcode{\sphinxupquote{+}} can be omitted (\sphinxcode{\sphinxupquote{print=charges+bonds}}).


\begin{savenotes}\sphinxattablestart
\sphinxthistablewithglobalstyle
\sphinxthistablewithnovlinesstyle
\centering
\begin{tabulary}{\linewidth}[t]{\Y{0.3}\Y{0.7}}
\sphinxtoprule
\sphinxstyletheadfamily 
\sphinxAtStartPar
Value
&\sphinxstyletheadfamily 
\sphinxAtStartPar
Meaning
\\
\sphinxmidrule
\sphinxtableatstartofbodyhook
\sphinxAtStartPar
\sphinxcode{\sphinxupquote{+charges}}
&
\sphinxAtStartPar
Hirshfeld charges (DFT methods), printed as “Atomic partitions of the density”.
\\
\sphinxhline
\sphinxAtStartPar
\sphinxcode{\sphinxupquote{+bonds}}
&
\sphinxAtStartPar
Mulliken populations and bond orders (SCF only), printed as “Overlap populations”.
\\
\sphinxhline
\sphinxAtStartPar
\sphinxcode{\sphinxupquote{+esr}}
&
\sphinxAtStartPar
ESR parameters.
\\
\sphinxhline
\sphinxAtStartPar
\sphinxcode{\sphinxupquote{+vectors}}
&
\sphinxAtStartPar
SCF vectors.
\\
\sphinxhline
\sphinxAtStartPar
\sphinxcode{\sphinxupquote{+molden}}
&
\sphinxAtStartPar
MOLDEN data: vibrations on \sphinxcode{\sphinxupquote{freq\textgreater{}}} lines for \sphinxcode{\sphinxupquote{task=hessian}}; with \sphinxcode{\sphinxupquote{+vectors+molden}} also orbitals on \sphinxcode{\sphinxupquote{mos\textgreater{}}} lines (\sphinxcode{\sphinxupquote{{[}Molden Format{]}}}, \sphinxcode{\sphinxupquote{{[}Atoms{]} AU}}, \sphinxcode{\sphinxupquote{{[}5D{]}}}, \sphinxcode{\sphinxupquote{{[}GTO{]}}}, \sphinxcode{\sphinxupquote{{[}MO{]}}} sections).
\\
\sphinxhline
\sphinxAtStartPar
\sphinxcode{\sphinxupquote{+density}}
&
\sphinxAtStartPar
Critical points of the electron density (QTAIM). Combine with \sphinxcode{\sphinxupquote{+charges}} to label them. \sphinxstyleemphasis{(undocumented)}
\\
\sphinxhline
\sphinxAtStartPar
\sphinxcode{\sphinxupquote{\sphinxhyphen{}internals}}
&
\sphinxAtStartPar
Suppress the printout of internal coordinates. \sphinxstyleemphasis{(undocumented)}
\\
\sphinxbottomrule
\end{tabulary}
\sphinxtableafterendhook\par
\sphinxattableend\end{savenotes}

\sphinxAtStartPar
\sphinxstylestrong{Example:}

\begin{sphinxVerbatim}[commandchars=\\\{\}]
\PYGZdl{}control print=+vectors+molden \PYGZdl{}end
\end{sphinxVerbatim}


\subsection{\sphinxstyleliteralintitle{\sphinxupquote{derivative}}}
\label{\detokenize{reference/control:derivative}}\label{\detokenize{reference/control:control-derivative}}
\index{derivative (\$control)@\spxentry{derivative}\spxextra{\$control}}\ignorespaces 
\sphinxAtStartPar
\sphinxstylestrong{Type:} enum\\
\sphinxstylestrong{Default:} \sphinxstyleemphasis{none stated}\\
\sphinxstylestrong{Status:} undocumented · verified on 25 (2025.02.28)\\
\sphinxstylestrong{Sources:} 09; 08 t53

\sphinxAtStartPar
Default: analytic second derivatives where the method has them (GGA DFT), numerical otherwise.

\sphinxAtStartPar
How \sphinxcode{\sphinxupquote{task=hessian}} computes the second derivatives. Not in any written source; found in the program and
tested in version 25, where \sphinxcode{\sphinxupquote{derivative=numeric}} with \sphinxcode{\sphinxupquote{theory=dft}} gave the same water frequencies as the
analytic Hessian to within 0.2 cm⁻\(\sp{\text{1}}\).


\begin{savenotes}\sphinxattablestart
\sphinxthistablewithglobalstyle
\sphinxthistablewithnovlinesstyle
\centering
\begin{tabulary}{\linewidth}[t]{\Y{0.3}\Y{0.7}}
\sphinxtoprule
\sphinxstyletheadfamily 
\sphinxAtStartPar
Value
&\sphinxstyletheadfamily 
\sphinxAtStartPar
Meaning
\\
\sphinxmidrule
\sphinxtableatstartofbodyhook
\sphinxAtStartPar
\sphinxcode{\sphinxupquote{analytic}}
&
\sphinxAtStartPar
Analytic second derivatives.
\\
\sphinxhline
\sphinxAtStartPar
\sphinxcode{\sphinxupquote{numeric}}
&
\sphinxAtStartPar
Second derivatives by finite differences of the gradient (\sphinxcode{\sphinxupquote{NUMERICAL COMPUTATION OF SECOND DERIVATIVES}}).
\\
\sphinxbottomrule
\end{tabulary}
\sphinxtableafterendhook\par
\sphinxattableend\end{savenotes}

\sphinxAtStartPar
\sphinxstylestrong{See also:} {\hyperref[\detokenize{reference/d2edr2:d2edr2-steps}]{\sphinxcrossref{\DUrole{std}{\DUrole{std-ref}{steps (\$d2edr2)}}}}}


\subsection{\sphinxstyleliteralintitle{\sphinxupquote{save}}}
\label{\detokenize{reference/control:save}}\label{\detokenize{reference/control:control-save}}
\index{save (\$control)@\spxentry{save}\spxextra{\$control}}\ignorespaces 
\sphinxAtStartPar
\sphinxstylestrong{Type:} string\\
\sphinxstylestrong{Default:} \sphinxstyleemphasis{none stated}\\
\sphinxstylestrong{Status:} current · verified on 25 (2025.02.28)\\
\sphinxstylestrong{Sources:} input.txt; Dft.rtf; 04 \#306; 08 t27

\sphinxAtStartPar
File name for saving SCF vectors. With \sphinxcode{\sphinxupquote{\$guess read=1}} the vectors are read back from the same file.
This is the basis of the SCF restart technique. In version 25, a run with \sphinxcode{\sphinxupquote{\$guess read=1}} prints
\sphinxcode{\sphinxupquote{Initial Guess (\sphinxhyphen{}1)}} and \sphinxcode{\sphinxupquote{vectors from \textquotesingle{}\textless{}file\textgreater{}\textquotesingle{}}}; the normal guess is printed as \sphinxcode{\sphinxupquote{Initial Guess (2)}}.

\sphinxAtStartPar
\sphinxstylestrong{See also:} {\hyperref[\detokenize{reference/guess:guess-read}]{\sphinxcrossref{\DUrole{std}{\DUrole{std-ref}{read (\$guess)}}}}}

\sphinxstepscope




\section{\sphinxstyleliteralintitle{\sphinxupquote{\$molecule}}}
\label{\detokenize{reference/molecule:molecule}}\label{\detokenize{reference/molecule:group-molecule}}\label{\detokenize{reference/molecule::doc}}
\index{\$molecule@\spxentry{\$molecule}}\ignorespaces 
\sphinxAtStartPar
Charge, multiplicity, geometry (Cartesian and/or Z\sphinxhyphen{}matrix) and per\sphinxhyphen{}atom basis sets.

\sphinxAtStartPar
Unlike the other groups, \sphinxcode{\sphinxupquote{\$molecule}} holds \sphinxstylestrong{ordered data}. It starts with \sphinxcode{\sphinxupquote{charge=}} and \sphinxcode{\sphinxupquote{mult=}}, then one of the
keywords \sphinxcode{\sphinxupquote{cartesian}} or \sphinxcode{\sphinxupquote{z\sphinxhyphen{}matrix}}, followed by one line per atom.

\sphinxAtStartPar
Cartesian lines: \sphinxcode{\sphinxupquote{Z x y z}}. \sphinxstylestrong{Z is the atomic number written as an integer}: \sphinxcode{\sphinxupquote{8}}, not \sphinxcode{\sphinxupquote{8.0}} and not \sphinxcode{\sphinxupquote{O}}. The
atom lines are read as a stream of numbers, so a value such as \sphinxcode{\sphinxupquote{8.0}} shifts everything after it. In version 25
this gives \sphinxstylestrong{no error}: the run continues with a wrong molecule (extra atoms with Z=0) and a meaningless
energy. Check the \sphinxcode{\sphinxupquote{formula:}} and atom\sphinxhyphen{}count lines at the top of the output. The coordinates are in Å unless
\sphinxcode{\sphinxupquote{unit=b}} is given.

\sphinxAtStartPar
Z\sphinxhyphen{}matrix lines, \sphinxcode{\sphinxupquote{Z j r k a l b}}, can combine distances, angles and torsions in four ways:


\begin{savenotes}\sphinxattablestart
\sphinxthistablewithglobalstyle
\centering
\begin{tabulary}{\linewidth}[t]{TT}
\sphinxtoprule
\sphinxstyletheadfamily 
\sphinxAtStartPar
Form
&\sphinxstyletheadfamily 
\sphinxAtStartPar
Meaning
\\
\sphinxmidrule
\sphinxtableatstartofbodyhook
\sphinxAtStartPar
\sphinxcode{\sphinxupquote{q\_i j r\_ij  k a\_ijk  l b\_ijkl}}
&
\sphinxAtStartPar
distance, angle, torsion
\\
\sphinxhline
\sphinxAtStartPar
\sphinxcode{\sphinxupquote{q\_i j r\_ij  k a\_ijk \sphinxhyphen{}l a\_ijl}}
&
\sphinxAtStartPar
distance and two angles
\\
\sphinxhline
\sphinxAtStartPar
\sphinxcode{\sphinxupquote{q\_i j r\_ij \sphinxhyphen{}k r\_ik   l b\_ijkl}}
&
\sphinxAtStartPar
two distances and torsion
\\
\sphinxhline
\sphinxAtStartPar
\sphinxcode{\sphinxupquote{q\_i j r\_ij \sphinxhyphen{}k r\_ik  \sphinxhyphen{}l r\_il}}
&
\sphinxAtStartPar
three distances
\\
\sphinxbottomrule
\end{tabulary}
\sphinxtableafterendhook\par
\sphinxattableend\end{savenotes}

\sphinxAtStartPar
Distances are in Å, and angles and torsions are in degrees. In a pure Z\sphinxhyphen{}matrix the first atom gives only its
atomic number, the second only a distance, and the third two internal coordinates. The two formats can be mixed:
put the Cartesian atoms first, then the remaining atoms as a Z\sphinxhyphen{}matrix.

\sphinxAtStartPar
The program does not use or impose point\sphinxhyphen{}group symmetry unless the newer \sphinxcode{\sphinxupquote{sym=}} option is given. Dummy atoms
are not supported, so counterpoise (BSSE) calculations cannot be done; NICS probe points go in \sphinxcode{\sphinxupquote{\$nmr points=}}
instead. Some tasks (\sphinxcode{\sphinxupquote{transit}}, \sphinxcode{\sphinxupquote{saddle}}, \sphinxcode{\sphinxupquote{path}}, \sphinxcode{\sphinxupquote{uniq}}) read \sphinxstylestrong{several} \sphinxcode{\sphinxupquote{\$molecule}} blocks from one input.
Text after \sphinxcode{\sphinxupquote{\$end}}, such as the \sphinxcode{\sphinxupquote{Energy = …}} lines that helper scripts append, is ignored. The line
\sphinxcode{\sphinxupquote{molecule input: \textquotesingle{}\textquotesingle{}}} near the top of every output is normal and not an error.


\subsection{\sphinxstyleliteralintitle{\sphinxupquote{charge}}}
\label{\detokenize{reference/molecule:charge}}\label{\detokenize{reference/molecule:molecule-charge}}
\index{charge (\$molecule)@\spxentry{charge}\spxextra{\$molecule}}\ignorespaces 
\sphinxAtStartPar
\sphinxstylestrong{Type:} integer\\
\sphinxstylestrong{Default:} \sphinxcode{\sphinxupquote{0}}\\
\sphinxstylestrong{Status:} current · not verified\\
\sphinxstylestrong{Sources:} input.txt; Dft.rtf (default)

\sphinxAtStartPar
Molecular charge.


\subsection{\sphinxstyleliteralintitle{\sphinxupquote{mult}}}
\label{\detokenize{reference/molecule:mult}}\label{\detokenize{reference/molecule:molecule-mult}}
\index{mult (\$molecule)@\spxentry{mult}\spxextra{\$molecule}}\ignorespaces 
\sphinxAtStartPar
\sphinxstylestrong{Type:} integer\\
\sphinxstylestrong{Default:} \sphinxcode{\sphinxupquote{1}}\\
\sphinxstylestrong{Status:} current · verified on 25 (2025.02.28)\\
\sphinxstylestrong{Sources:} input.txt; Dft.rtf (default); 08 t10

\sphinxAtStartPar
Spin multiplicity, 2S+1.


\subsection{\sphinxstyleliteralintitle{\sphinxupquote{mass}}}
\label{\detokenize{reference/molecule:mass}}\label{\detokenize{reference/molecule:molecule-mass}}
\index{mass (\$molecule)@\spxentry{mass}\spxextra{\$molecule}}\ignorespaces 
\sphinxAtStartPar
\sphinxstylestrong{Type:} real\\
\sphinxstylestrong{Default:} \sphinxstyleemphasis{none stated}\\
\sphinxstylestrong{Units:} amu\\
\sphinxstylestrong{Status:} current · not verified\\
\sphinxstylestrong{Sources:} input.txt

\sphinxAtStartPar
Optional atomic mass written after an atom’s coordinates, for isotopes.


\subsection{\sphinxstyleliteralintitle{\sphinxupquote{set}}}
\label{\detokenize{reference/molecule:set}}\label{\detokenize{reference/molecule:molecule-set}}
\index{set (\$molecule)@\spxentry{set}\spxextra{\$molecule}}\ignorespaces 
\sphinxAtStartPar
\sphinxstylestrong{Type:} string\\
\sphinxstylestrong{Default:} \sphinxstyleemphasis{none stated}\\
\sphinxstylestrong{Status:} current · not verified\\
\sphinxstylestrong{Sources:} input.txt; 03 \S{}5 (default set); 01 \S{}4.11; 06 \S{}1

\sphinxAtStartPar
Default: the set named by \sphinxcode{\sphinxupquote{default=}} in the basis\sphinxhyphen{}file header (\sphinxcode{\sphinxupquote{L1}} in \sphinxcode{\sphinxupquote{basis.in}}).

\sphinxAtStartPar
Basis set for the atoms that follow. The name must match one in the basis\sphinxhyphen{}set file. It can appear several
times in the group. The orbital basis is paired with its own auxiliary (fitting) basis, so you cannot
combine arbitrary orbital and auxiliary sets.

\sphinxAtStartPar
\sphinxstylestrong{Example:}

\begin{sphinxVerbatim}[commandchars=\\\{\}]
\PYGZdl{}molecule
 z\PYGZhy{}matrix
 set=B2
   8
 set=B1
   1    1    0.96252725
   1    1    0.96252725    2    102.943872
\PYGZdl{}end
\end{sphinxVerbatim}

\sphinxAtStartPar
\sphinxstylestrong{See also:} {\hyperref[\detokenize{reference/control:control-basis}]{\sphinxcrossref{\DUrole{std}{\DUrole{std-ref}{basis (\$control)}}}}}


\subsection{\sphinxstyleliteralintitle{\sphinxupquote{unit}}}
\label{\detokenize{reference/molecule:unit}}\label{\detokenize{reference/molecule:molecule-unit}}
\index{unit (\$molecule)@\spxentry{unit}\spxextra{\$molecule}}\ignorespaces 
\sphinxAtStartPar
\sphinxstylestrong{Type:} string\\
\sphinxstylestrong{Default:} \sphinxstyleemphasis{none stated}\\
\sphinxstylestrong{Status:} undocumented · verified on 25 (2025.02.28)\\
\sphinxstylestrong{Sources:} 01 \S{}4.7; 08 t47

\sphinxAtStartPar
Length unit of the coordinates. The default unit is Å. In version 25, water given in bohr with \sphinxcode{\sphinxupquote{unit=b}} gave the same geometry and energy as in Å.


\begin{savenotes}\sphinxattablestart
\sphinxthistablewithglobalstyle
\sphinxthistablewithnovlinesstyle
\centering
\begin{tabulary}{\linewidth}[t]{\Y{0.3}\Y{0.7}}
\sphinxtoprule
\sphinxstyletheadfamily 
\sphinxAtStartPar
Value
&\sphinxstyletheadfamily 
\sphinxAtStartPar
Meaning
\\
\sphinxmidrule
\sphinxtableatstartofbodyhook
\sphinxAtStartPar
\sphinxcode{\sphinxupquote{b}}
&
\sphinxAtStartPar
Coordinates in bohr.
\\
\sphinxbottomrule
\end{tabulary}
\sphinxtableafterendhook\par
\sphinxattableend\end{savenotes}


\subsection{\sphinxstyleliteralintitle{\sphinxupquote{sym}}}
\label{\detokenize{reference/molecule:sym}}\label{\detokenize{reference/molecule:molecule-sym}}
\index{sym (\$molecule)@\spxentry{sym}\spxextra{\$molecule}}\ignorespaces 
\sphinxAtStartPar
\sphinxstylestrong{Type:} list\\
\sphinxstylestrong{Default:} \sphinxstyleemphasis{none stated}\\
\sphinxstylestrong{Status:} undocumented · verified on 25 (2025.02.28)\\
\sphinxstylestrong{Sources:} 01 \S{}4.7; 04 \#481; 08 t15; 08 t30a

\sphinxAtStartPar
Used as \sphinxcode{\sphinxupquote{sym=1,0.01}} in the official 2023\textendash{}2025 example inputs, apparently to enable symmetry with tolerance 0.01.
Newer versions report the detected point group (\sphinxcode{\sphinxupquote{symmetry: D6h}}). In version 25, \sphinxcode{\sphinxupquote{sym=1,0.01}} on a slightly
distorted NH3 printed \sphinxcode{\sphinxupquote{symmetry: C3v    6, unique:    2}} and echoed the geometry symmetrized. On its own it
does not keep the symmetry during an optimization; add \sphinxcode{\sphinxupquote{\$optimize sym=1}} for that.

\sphinxAtStartPar
\sphinxstylestrong{See also:} {\hyperref[\detokenize{reference/optimize:optimize-symmetry}]{\sphinxcrossref{\DUrole{std}{\DUrole{std-ref}{symmetry (\$optimize)}}}}}

\sphinxstepscope




\section{\sphinxstyleliteralintitle{\sphinxupquote{\$guess}}}
\label{\detokenize{reference/guess:guess}}\label{\detokenize{reference/guess:group-guess}}\label{\detokenize{reference/guess::doc}}
\index{\$guess@\spxentry{\$guess}}\ignorespaces 
\sphinxAtStartPar
Initial guess for the SCF vectors.


\subsection{\sphinxstyleliteralintitle{\sphinxupquote{read}}}
\label{\detokenize{reference/guess:read}}\label{\detokenize{reference/guess:guess-read}}
\index{read (\$guess)@\spxentry{read}\spxextra{\$guess}}\ignorespaces 
\sphinxAtStartPar
\sphinxstylestrong{Type:} switch\\
\sphinxstylestrong{Default:} \sphinxstyleemphasis{none stated}\\
\sphinxstylestrong{Status:} current · verified on 25 (2025.02.28)\\
\sphinxstylestrong{Sources:} input.txt; 04 \#306; 05 \S{}1; 08 t27

\sphinxAtStartPar
Source of the starting SCF vectors.


\begin{savenotes}\sphinxattablestart
\sphinxthistablewithglobalstyle
\sphinxthistablewithnovlinesstyle
\centering
\begin{tabulary}{\linewidth}[t]{\Y{0.3}\Y{0.7}}
\sphinxtoprule
\sphinxstyletheadfamily 
\sphinxAtStartPar
Value
&\sphinxstyletheadfamily 
\sphinxAtStartPar
Meaning
\\
\sphinxmidrule
\sphinxtableatstartofbodyhook
\sphinxAtStartPar
\sphinxcode{\sphinxupquote{1}}
&
\sphinxAtStartPar
Read SCF vectors from the file named by \sphinxcode{\sphinxupquote{\$control save=}}.
\\
\sphinxhline
\sphinxAtStartPar
\sphinxcode{\sphinxupquote{0}}
&
\sphinxAtStartPar
Extended Hückel guess.
\\
\sphinxbottomrule
\end{tabulary}
\sphinxtableafterendhook\par
\sphinxattableend\end{savenotes}

\begin{sphinxadmonition}{note}{Note:}
\sphinxAtStartPar
Reading saved vectors is the standard way to restart a stalled SCF, and to reuse converged orbitals in a
follow\sphinxhyphen{}up job such as a Hessian after an optimization.
\end{sphinxadmonition}

\sphinxAtStartPar
\sphinxstylestrong{See also:} {\hyperref[\detokenize{reference/control:control-save}]{\sphinxcrossref{\DUrole{std}{\DUrole{std-ref}{save (\$control)}}}}}, {\hyperref[\detokenize{reference/scf:scf-iterations}]{\sphinxcrossref{\DUrole{std}{\DUrole{std-ref}{iterations (\$scf)}}}}}


\subsection{\sphinxstyleliteralintitle{\sphinxupquote{na}}}
\label{\detokenize{reference/guess:na}}\label{\detokenize{reference/guess:guess-na}}
\index{na (\$guess)@\spxentry{na}\spxextra{\$guess}}\ignorespaces 
\sphinxAtStartPar
\sphinxstylestrong{Type:} indexed\sphinxhyphen{}array\\
\sphinxstylestrong{Default:} \sphinxstyleemphasis{none stated}\\
\sphinxstylestrong{Status:} current · not verified\\
\sphinxstylestrong{Sources:} input.txt

\sphinxAtStartPar
Default: \sphinxcode{\sphinxupquote{na(1)=1,2,3,…,N}} (the natural order).

\sphinxAtStartPar
Reordering of the starting vectors for alpha spin.

\sphinxAtStartPar
\sphinxstylestrong{Example:}

\begin{sphinxVerbatim}[commandchars=\\\{\}]
\PYGZdl{}guess na(25)=26,25 \PYGZdl{}end    \PYGZsh{} swap the 25th and 26th vectors
\end{sphinxVerbatim}

\sphinxAtStartPar
\sphinxstylestrong{See also:} {\hyperref[\detokenize{reference/guess:guess-nb}]{\sphinxcrossref{\DUrole{std}{\DUrole{std-ref}{nb (\$guess)}}}}}


\subsection{\sphinxstyleliteralintitle{\sphinxupquote{nb}}}
\label{\detokenize{reference/guess:nb}}\label{\detokenize{reference/guess:guess-nb}}
\index{nb (\$guess)@\spxentry{nb}\spxextra{\$guess}}\ignorespaces 
\sphinxAtStartPar
\sphinxstylestrong{Type:} indexed\sphinxhyphen{}array\\
\sphinxstylestrong{Default:} \sphinxstyleemphasis{none stated}\\
\sphinxstylestrong{Status:} current · not verified\\
\sphinxstylestrong{Sources:} input.txt

\sphinxAtStartPar
Default: \sphinxcode{\sphinxupquote{nb(1)=1,2,3,…,N}}.

\sphinxAtStartPar
Reordering of the starting vectors for beta spin.

\sphinxAtStartPar
\sphinxstylestrong{See also:} {\hyperref[\detokenize{reference/guess:guess-na}]{\sphinxcrossref{\DUrole{std}{\DUrole{std-ref}{na (\$guess)}}}}}


\subsection{\sphinxstyleliteralintitle{\sphinxupquote{mix}}}
\label{\detokenize{reference/guess:mix}}\label{\detokenize{reference/guess:guess-mix}}
\index{mix (\$guess)@\spxentry{mix}\spxextra{\$guess}}\ignorespaces 
\sphinxAtStartPar
\sphinxstylestrong{Type:} switch\\
\sphinxstylestrong{Default:} \sphinxcode{\sphinxupquote{0}}\\
\sphinxstylestrong{Status:} undocumented · not verified\\
\sphinxstylestrong{Sources:} input.txt (annotation corrupted); 03 \#0, \S{}7; 06 \S{}3

\sphinxAtStartPar
Breaks the equivalence of the alpha and beta electrons in the highest occupied orbital. It is meant for
open\sphinxhyphen{}shell singlets (for example singlet biradicals) and is only sensible with \sphinxcode{\sphinxupquote{\$scf restrict=0}}.


\begin{savenotes}\sphinxattablestart
\sphinxthistablewithglobalstyle
\sphinxthistablewithnovlinesstyle
\centering
\begin{tabulary}{\linewidth}[t]{\Y{0.3}\Y{0.7}}
\sphinxtoprule
\sphinxstyletheadfamily 
\sphinxAtStartPar
Value
&\sphinxstyletheadfamily 
\sphinxAtStartPar
Meaning
\\
\sphinxmidrule
\sphinxtableatstartofbodyhook
\sphinxAtStartPar
\sphinxcode{\sphinxupquote{0}}
&
\sphinxAtStartPar
No mixing. (\sphinxstylestrong{default})
\\
\sphinxhline
\sphinxAtStartPar
\sphinxcode{\sphinxupquote{1}}
&
\sphinxAtStartPar
Mix the HOMO and LUMO to break alpha/beta equivalence.
\\
\sphinxbottomrule
\end{tabulary}
\sphinxtableafterendhook\par
\sphinxattableend\end{savenotes}

\begin{sphinxadmonition}{note}{Note:}
\sphinxAtStartPar
Users report that with non\sphinxhyphen{}RI DFT, \sphinxcode{\sphinxupquote{mix=1}} can converge back to the restricted solution (same energy as
RKS). RI\sphinxhyphen{}DFT, or a hybrid, broke the symmetry in the reported cases.
\end{sphinxadmonition}

\sphinxAtStartPar
\sphinxstylestrong{See also:} {\hyperref[\detokenize{reference/scf:scf-restrict}]{\sphinxcrossref{\DUrole{std}{\DUrole{std-ref}{restrict (\$scf)}}}}}

\sphinxstepscope




\section{\sphinxstyleliteralintitle{\sphinxupquote{\$scf}}}
\label{\detokenize{reference/scf:scf}}\label{\detokenize{reference/scf:group-scf}}\label{\detokenize{reference/scf::doc}}
\index{\$scf@\spxentry{\$scf}}\ignorespaces 
\sphinxAtStartPar
Spin restriction, convergence, iteration limits and SCF solver.


\subsection{\sphinxstyleliteralintitle{\sphinxupquote{restrict}}}
\label{\detokenize{reference/scf:restrict}}\label{\detokenize{reference/scf:scf-restrict}}
\index{restrict (\$scf)@\spxentry{restrict}\spxextra{\$scf}}\ignorespaces 
\sphinxAtStartPar
\sphinxstylestrong{Type:} switch\\
\sphinxstylestrong{Default:} \sphinxstyleemphasis{none stated}\\
\sphinxstylestrong{Status:} current · verified on 25 (2025.02.28)\\
\sphinxstylestrong{Sources:} input.txt; 08 t48

\sphinxAtStartPar
Default: \sphinxcode{\sphinxupquote{1}} for a closed\sphinxhyphen{}shell singlet, otherwise \sphinxcode{\sphinxupquote{0}}. Version 25 ran O2 with \sphinxcode{\sphinxupquote{mult=3}} (an even number of electrons) unrestricted without being told.

\sphinxAtStartPar
Spin restriction.


\begin{savenotes}\sphinxattablestart
\sphinxthistablewithglobalstyle
\sphinxthistablewithnovlinesstyle
\centering
\begin{tabulary}{\linewidth}[t]{\Y{0.3}\Y{0.7}}
\sphinxtoprule
\sphinxstyletheadfamily 
\sphinxAtStartPar
Value
&\sphinxstyletheadfamily 
\sphinxAtStartPar
Meaning
\\
\sphinxmidrule
\sphinxtableatstartofbodyhook
\sphinxAtStartPar
\sphinxcode{\sphinxupquote{0}}
&
\sphinxAtStartPar
Spin\sphinxhyphen{}unrestricted.
\\
\sphinxhline
\sphinxAtStartPar
\sphinxcode{\sphinxupquote{1}}
&
\sphinxAtStartPar
Spin\sphinxhyphen{}restricted (available only for a few methods).
\\
\sphinxbottomrule
\end{tabulary}
\sphinxtableafterendhook\par
\sphinxattableend\end{savenotes}

\sphinxAtStartPar
\sphinxstylestrong{See also:} {\hyperref[\detokenize{reference/guess:guess-mix}]{\sphinxcrossref{\DUrole{std}{\DUrole{std-ref}{mix (\$guess)}}}}}


\subsection{\sphinxstyleliteralintitle{\sphinxupquote{convergence}}}
\label{\detokenize{reference/scf:convergence}}\label{\detokenize{reference/scf:scf-convergence}}
\index{convergence (\$scf)@\spxentry{convergence}\spxextra{\$scf}}\ignorespaces 
\sphinxAtStartPar
\sphinxstylestrong{Type:} list\\
\sphinxstylestrong{Default:} \sphinxcode{\sphinxupquote{1e\sphinxhyphen{}6,1e\sphinxhyphen{}3}}\\
\sphinxstylestrong{Status:} current · verified on 25 (2025.02.28)\\
\sphinxstylestrong{Sources:} input.txt; 01 \S{}4.5 (output echo); 05 \S{}2; 04 \#322\sphinxhyphen{}\#324; 08 t01, t07a, t07b; 08 t22a\sphinxhyphen{}c

\sphinxAtStartPar
Default: the first value is for SCF and the second for CP\sphinxhyphen{}SCF.

\sphinxAtStartPar
SCF and CP\sphinxhyphen{}SCF convergence thresholds. The defaults are recommended. To save time in difficult cases you may
loosen the SCF value to \sphinxcode{\sphinxupquote{1e\sphinxhyphen{}5}} for preliminary geometry optimizations.

\begin{sphinxadmonition}{note}{Note:}
\sphinxAtStartPar
\sphinxstylestrong{Give both values, or put \sphinxcode{\sphinxupquote{convergence}} last.} In version 25, a one\sphinxhyphen{}value \sphinxcode{\sphinxupquote{convergence=1e\sphinxhyphen{}7}} followed by
\sphinxcode{\sphinxupquote{iterations=250}} stops the run with \sphinxcode{\sphinxupquote{*** error in \textquotesingle{}\$scf\textquotesingle{} : unknown name terations}}, also when the two
keywords are on separate lines. \sphinxcode{\sphinxupquote{convergence=1e\sphinxhyphen{}7,1e\sphinxhyphen{}3 iterations=250}} works, and so does
\sphinxcode{\sphinxupquote{iterations=250 convergence=1e\sphinxhyphen{}7}}; \sphinxcode{\sphinxupquote{convergence=1e\sphinxhyphen{}7 procedure=bfgs iterations=250}} also worked. Older
versions gave the same error as \sphinxcode{\sphinxupquote{unknown name erations}}. Check the \sphinxcode{\sphinxupquote{scf: conv = …; iter=…}} echo line in the
output.
\end{sphinxadmonition}

\sphinxAtStartPar
\sphinxstylestrong{See also:} {\hyperref[\detokenize{reference/scf:scf-iterations}]{\sphinxcrossref{\DUrole{std}{\DUrole{std-ref}{iterations (\$scf)}}}}}, {\hyperref[\detokenize{reference/scf:scf-procedure}]{\sphinxcrossref{\DUrole{std}{\DUrole{std-ref}{procedure (\$scf)}}}}}


\subsection{\sphinxstyleliteralintitle{\sphinxupquote{iterations}}}
\label{\detokenize{reference/scf:iterations}}\label{\detokenize{reference/scf:scf-iterations}}
\index{iterations (\$scf)@\spxentry{iterations}\spxextra{\$scf}}\ignorespaces 
\sphinxAtStartPar
\sphinxstylestrong{Type:} list\\
\sphinxstylestrong{Default:} \sphinxcode{\sphinxupquote{32,16}}\\
\sphinxstylestrong{Status:} current · verified on 25 (2025.02.28)\\
\sphinxstylestrong{Sources:} input.txt; 01 \S{}4.5; 04 \#319; 05 \S{}5; Dft.rtf; 02 \S{}3.2 (pri help); 04 \#306; 08 t01

\sphinxAtStartPar
Default: Version 25 echoes \sphinxcode{\sphinxupquote{iter=32,16}} for DFT. v6, v13 and v15 echoed \sphinxcode{\sphinxupquote{16,16}}, a HF run for CCSD showed \sphinxcode{\sphinxupquote{48,16}}, and
\sphinxcode{\sphinxupquote{Dft.rtf}} gives 30 (older version). The default may depend on the method and version.

\sphinxAtStartPar
Maximum number of iterations and of microiterations (\sphinxcode{\sphinxupquote{iterations=it,micro}}).

\begin{sphinxadmonition}{note}{Note:}
\sphinxAtStartPar
Do not raise this limit hoping for convergence: understand why SCF is not converging. A recommended
technique for difficult cases is to cap the run at \sphinxcode{\sphinxupquote{iterations=10}}, save the vectors (\sphinxcode{\sphinxupquote{\$control save=}}) and
restart with \sphinxcode{\sphinxupquote{\$guess read=1}}, repeating as needed.
\end{sphinxadmonition}

\sphinxAtStartPar
\sphinxstylestrong{See also:} {\hyperref[\detokenize{reference/scf:scf-convergence}]{\sphinxcrossref{\DUrole{std}{\DUrole{std-ref}{convergence (\$scf)}}}}}, {\hyperref[\detokenize{reference/guess:guess-read}]{\sphinxcrossref{\DUrole{std}{\DUrole{std-ref}{read (\$guess)}}}}}


\subsection{\sphinxstyleliteralintitle{\sphinxupquote{procedure}}}
\label{\detokenize{reference/scf:procedure}}\label{\detokenize{reference/scf:scf-procedure}}
\index{procedure (\$scf)@\spxentry{procedure}\spxextra{\$scf}}\ignorespaces 
\sphinxAtStartPar
\sphinxstylestrong{Type:} enum\\
\sphinxstylestrong{Default:} \sphinxcode{\sphinxupquote{nr}}\\
\sphinxstylestrong{Status:} current · verified on 25 (2025.02.28)\\
\sphinxstylestrong{Sources:} input.txt; Dft.rtf; 04 \#320; 08 t07b

\sphinxAtStartPar
Default: \sphinxcode{\sphinxupquote{nr}} for GGA DFT.

\sphinxAtStartPar
Algorithm for solving the SCF equations. If one procedure fails, try the other.


\begin{savenotes}\sphinxattablestart
\sphinxthistablewithglobalstyle
\sphinxthistablewithnovlinesstyle
\centering
\begin{tabulary}{\linewidth}[t]{\Y{0.3}\Y{0.7}}
\sphinxtoprule
\sphinxstyletheadfamily 
\sphinxAtStartPar
Value
&\sphinxstyletheadfamily 
\sphinxAtStartPar
Meaning
\\
\sphinxmidrule
\sphinxtableatstartofbodyhook
\sphinxAtStartPar
\sphinxcode{\sphinxupquote{nr}}
&
\sphinxAtStartPar
Newton\textendash{}Raphson quadratically convergent. Can recover even from a wrong starting electronic configuration but may blow up in hard cases. (\sphinxstylestrong{default})
\\
\sphinxhline
\sphinxAtStartPar
\sphinxcode{\sphinxupquote{bfgs}}
&
\sphinxAtStartPar
BFGS Hessian update. Less prone to blowing up but usually slower and sometimes converges poorly.
\\
\sphinxbottomrule
\end{tabulary}
\sphinxtableafterendhook\par
\sphinxattableend\end{savenotes}

\sphinxAtStartPar
\sphinxstylestrong{Example:}

\begin{sphinxVerbatim}[commandchars=\\\{\}]
\PYGZdl{}scf convergence=1e\PYGZhy{}7 procedure=bfgs iterations=250 \PYGZdl{}end
\PYGZdl{}grid accuracy=1e\PYGZhy{}9 \PYGZdl{}end
\end{sphinxVerbatim}


\subsection{\sphinxstyleliteralintitle{\sphinxupquote{d1small}}}
\label{\detokenize{reference/scf:d1small}}\label{\detokenize{reference/scf:scf-d1small}}
\index{d1small (\$scf)@\spxentry{d1small}\spxextra{\$scf}}\ignorespaces 
\sphinxAtStartPar
\sphinxstylestrong{Type:} switch\\
\sphinxstylestrong{Default:} \sphinxcode{\sphinxupquote{0}}\\
\sphinxstylestrong{Status:} current · not verified\\
\sphinxstylestrong{Sources:} input.txt

\sphinxAtStartPar
Small\sphinxhyphen{}component density in quadratically convergent SCF.


\begin{savenotes}\sphinxattablestart
\sphinxthistablewithglobalstyle
\sphinxthistablewithnovlinesstyle
\centering
\begin{tabulary}{\linewidth}[t]{\Y{0.3}\Y{0.7}}
\sphinxtoprule
\sphinxstyletheadfamily 
\sphinxAtStartPar
Value
&\sphinxstyletheadfamily 
\sphinxAtStartPar
Meaning
\\
\sphinxmidrule
\sphinxtableatstartofbodyhook
\sphinxAtStartPar
\sphinxcode{\sphinxupquote{1}}
&
\sphinxAtStartPar
Include the small\sphinxhyphen{}component density in the linear equations of quadratically convergent SCF (scalar\sphinxhyphen{}relativistic pure DFT).
\\
\sphinxhline
\sphinxAtStartPar
\sphinxcode{\sphinxupquote{0}}
&
\sphinxAtStartPar
Discard these contributions. This is faster, but SCF may fail to converge for molecules containing the heaviest atoms. (\sphinxstylestrong{default})
\\
\sphinxbottomrule
\end{tabulary}
\sphinxtableafterendhook\par
\sphinxattableend\end{savenotes}

\sphinxAtStartPar
\sphinxstylestrong{See also:} {\hyperref[\detokenize{reference/control:control-four}]{\sphinxcrossref{\DUrole{std}{\DUrole{std-ref}{four (\$control)}}}}}


\subsection{\sphinxstyleliteralintitle{\sphinxupquote{core}}}
\label{\detokenize{reference/scf:core}}\label{\detokenize{reference/scf:scf-core}}
\index{core (\$scf)@\spxentry{core}\spxextra{\$scf}}\ignorespaces 
\sphinxAtStartPar
\sphinxstylestrong{Type:} integer\\
\sphinxstylestrong{Default:} \sphinxstyleemphasis{none stated}\\
\sphinxstylestrong{Status:} current · not verified\\
\sphinxstylestrong{Sources:} input.txt

\sphinxAtStartPar
Default: computed from the atom table; see \sphinxcode{\sphinxupquote{\$atoms core}}.

\sphinxAtStartPar
Number of core electrons for HF\sphinxhyphen{}based methods.

\sphinxAtStartPar
\sphinxstylestrong{See also:} {\hyperref[\detokenize{reference/atoms:atoms-core}]{\sphinxcrossref{\DUrole{std}{\DUrole{std-ref}{core (\$atoms)}}}}}

\sphinxstepscope




\section{\sphinxstyleliteralintitle{\sphinxupquote{\$dft}}}
\label{\detokenize{reference/dft:dft}}\label{\detokenize{reference/dft:group-dft}}\label{\detokenize{reference/dft::doc}}
\index{\$dft@\spxentry{\$dft}}\ignorespaces 
\sphinxAtStartPar
Exchange\textendash{}correlation functional.


\subsection{\sphinxstyleliteralintitle{\sphinxupquote{functional}}}
\label{\detokenize{reference/dft:functional}}\label{\detokenize{reference/dft:dft-functional}}
\index{functional (\$dft)@\spxentry{functional}\spxextra{\$dft}}\ignorespaces 
\sphinxAtStartPar
\sphinxstylestrong{Type:} enum\\
\sphinxstylestrong{Default:} \sphinxcode{\sphinxupquote{PBE}}\\
\sphinxstylestrong{Status:} current · verified on 25 (2025.02.28)\\
\sphinxstylestrong{Sources:} input.txt; 01 \S{}3.2; 03 \S{}5 (strings); 04 \S{}0; 08 t02; 08 t31, t38

\sphinxAtStartPar
Default: \sphinxcode{\sphinxupquote{PBE}} (the default GGA); \sphinxcode{\sphinxupquote{PBE1}} for \sphinxcode{\sphinxupquote{theory=ridft}}.

\sphinxAtStartPar
Exchange\textendash{}correlation approximation for DFT methods.


\begin{savenotes}\sphinxattablestart
\sphinxthistablewithglobalstyle
\sphinxthistablewithnovlinesstyle
\centering
\begin{tabulary}{\linewidth}[t]{\Y{0.3}\Y{0.7}}
\sphinxtoprule
\sphinxstyletheadfamily 
\sphinxAtStartPar
Value
&\sphinxstyletheadfamily 
\sphinxAtStartPar
Meaning
\\
\sphinxmidrule
\sphinxtableatstartofbodyhook
\sphinxAtStartPar
\sphinxcode{\sphinxupquote{PBE}}
&
\sphinxAtStartPar
Perdew\textendash{}Burke\textendash{}Ernzerhof, Phys. Rev. Lett. 77 (1996) 3865. (\sphinxstylestrong{default})
\\
\sphinxhline
\sphinxAtStartPar
\sphinxcode{\sphinxupquote{mPBE}}
&
\sphinxAtStartPar
Adamo\textendash{}Barone modification of PBE, J. Chem. Phys. 116 (2002) 5933.
\\
\sphinxhline
\sphinxAtStartPar
\sphinxcode{\sphinxupquote{BLYP}}
&
\sphinxAtStartPar
Becke\textendash{}Lee\textendash{}Yang\textendash{}Parr, Phys. Rev. A 38 (1988) 3098; Phys. Rev. B 37 (1988) 785.
\\
\sphinxhline
\sphinxAtStartPar
\sphinxcode{\sphinxupquote{OLYP}}
&
\sphinxAtStartPar
Handy\textendash{}Cohen exchange + LYP correlation, J. Chem. Phys. 117 (2002) 1441.
\\
\sphinxhline
\sphinxAtStartPar
\sphinxcode{\sphinxupquote{PBE1}}
&
\sphinxAtStartPar
PBE hybrid with 1/4 exact exchange, J. Chem. Phys. 110 (1999) 6158. Default for \sphinxcode{\sphinxupquote{theory=ridft}}.
\\
\sphinxhline
\sphinxAtStartPar
\sphinxcode{\sphinxupquote{B3LYP\_PW}}
&
\sphinxAtStartPar
Becke three\sphinxhyphen{}parameter hybrid with PW local correlation (see notes).
\\
\sphinxhline
\sphinxAtStartPar
\sphinxcode{\sphinxupquote{B3LYP}}
&
\sphinxAtStartPar
Becke three\sphinxhyphen{}parameter hybrid (see notes).
\\
\sphinxhline
\sphinxAtStartPar
\sphinxcode{\sphinxupquote{LPBEVV}}
&
\sphinxAtStartPar
Newer functional used in the official example inputs with \sphinxcode{\sphinxupquote{theory=DFT\_n4}} (2025). Runs in version 25; its definition is not documented. \sphinxstyleemphasis{(undocumented)}
\\
\sphinxhline
\sphinxAtStartPar
\sphinxcode{\sphinxupquote{LPBEVVV}}
&
\sphinxAtStartPar
Newer functional used in the official example inputs with \sphinxcode{\sphinxupquote{theory=ridft}} (2023\textendash{}2025). Runs in version 25; its definition is not documented. \sphinxstyleemphasis{(undocumented)}
\\
\sphinxbottomrule
\end{tabulary}
\sphinxtableafterendhook\par
\sphinxattableend\end{savenotes}

\begin{sphinxadmonition}{note}{Note:}
\sphinxAtStartPar
The B3LYP form in Priroda is
Exc = 0.2·Ex(HF) + 0.8·(0.9·Ex(B) + 0.1·Ex(S)) + 0.81·Ec(LYP) + 0.19·Ec(LSD).
Ec(LSD) is the local correlation of either Phys. Rev. B 45 (1992) 13244 (PW) or Can. J. Phys. 58 (1980) 1200 (VWN).
\sphinxcode{\sphinxupquote{input.txt}} points out that some commercial programs use VWN, although Becke himself used PW
(J. Chem. Phys. 98 (1993) 5648).

\sphinxAtStartPar
\sphinxstylestrong{Hybrids (\sphinxcode{\sphinxupquote{PBE1}}, \sphinxcode{\sphinxupquote{B3LYP}}, \sphinxcode{\sphinxupquote{B3LYP\_PW}}) require \sphinxcode{\sphinxupquote{theory=ridft}}.} With \sphinxcode{\sphinxupquote{theory=dft}} or \sphinxcode{\sphinxupquote{tddft}} they either
run silently without exact exchange or abort; see \sphinxcode{\sphinxupquote{\$control theory}}.

\sphinxAtStartPar
No empirical dispersion correction (D3 or similar) is available. The v6 binary contains further functional
names (\sphinxcode{\sphinxupquote{VWN5}}, \sphinxcode{\sphinxupquote{SLYP0}}, \sphinxcode{\sphinxupquote{BPBE}}, \sphinxcode{\sphinxupquote{HFLSD}}, \sphinxcode{\sphinxupquote{HFPBE}}, \sphinxcode{\sphinxupquote{B2PBE}}, \sphinxcode{\sphinxupquote{B3PBE}}, \sphinxcode{\sphinxupquote{HFLYP}}, \sphinxcode{\sphinxupquote{B1LYP}}, \sphinxcode{\sphinxupquote{B2LYP}}), none of
them documented or tested.
\end{sphinxadmonition}

\sphinxAtStartPar
\sphinxstylestrong{See also:} {\hyperref[\detokenize{reference/control:control-theory}]{\sphinxcrossref{\DUrole{std}{\DUrole{std-ref}{theory (\$control)}}}}}

\sphinxstepscope




\section{\sphinxstyleliteralintitle{\sphinxupquote{\$grid}}}
\label{\detokenize{reference/grid:grid}}\label{\detokenize{reference/grid:group-grid}}\label{\detokenize{reference/grid::doc}}
\index{\$grid@\spxentry{\$grid}}\ignorespaces 
\sphinxAtStartPar
Adaptive integration grid for DFT.


\subsection{\sphinxstyleliteralintitle{\sphinxupquote{accuracy}}}
\label{\detokenize{reference/grid:accuracy}}\label{\detokenize{reference/grid:grid-accuracy}}
\index{accuracy (\$grid)@\spxentry{accuracy}\spxextra{\$grid}}\ignorespaces 
\sphinxAtStartPar
\sphinxstylestrong{Type:} list\\
\sphinxstylestrong{Default:} \sphinxcode{\sphinxupquote{1e\sphinxhyphen{}8,1e\sphinxhyphen{}5}}\\
\sphinxstylestrong{Status:} current · verified on 25 (2025.02.28)\\
\sphinxstylestrong{Sources:} input.txt; 01 \S{}4.4 (output echo); Dft.rtf; 08 t01, t04

\sphinxAtStartPar
Default: \sphinxcode{\sphinxupquote{1e\sphinxhyphen{}8}} for the final grid and \sphinxcode{\sphinxupquote{1e\sphinxhyphen{}5}} for the coarse grid used in the first SCF iterations, as echoed by v13, v15
and v25 (\sphinxcode{\sphinxupquote{grid: accuracy=1.0e\sphinxhyphen{}08,1.0e\sphinxhyphen{}05}}). \sphinxcode{\sphinxupquote{Dft.rtf}} gives \sphinxcode{\sphinxupquote{1e\sphinxhyphen{}7}} (older version).

\sphinxAtStartPar
Target accuracy of the exchange\textendash{}correlation energy per atom for the adaptively generated grid (DFT
methods). The default, 1e\sphinxhyphen{}8, is a fine grid. Use 1e\sphinxhyphen{}7 to speed up geometry optimizations with
\sphinxcode{\sphinxupquote{tolerance\textless{}=1e\sphinxhyphen{}4}}, and 1e\sphinxhyphen{}9 for very accurate calculations. The grid grows during the SCF; the output
reports each size on a line such as \sphinxcode{\sphinxupquote{grid:    34000   53}} (number of points) in the SCF table.

\begin{sphinxadmonition}{note}{Note:}
\sphinxAtStartPar
\sphinxcode{\sphinxupquote{Dft.rtf}} suggests 1e\sphinxhyphen{}6 for rough optimizations (\sphinxcode{\sphinxupquote{tol=1e\sphinxhyphen{}3}}) and 1e\sphinxhyphen{}8 for floppy molecules optimized to
\sphinxcode{\sphinxupquote{tol=1e\sphinxhyphen{}5}}. It says that asking for more than 1e\sphinxhyphen{}8 “makes no deep sense”, because the integration formulas
cannot reach unlimited accuracy. This conflicts with \sphinxcode{\sphinxupquote{input.txt}}’s 1e\sphinxhyphen{}9 for very accurate work.
\end{sphinxadmonition}

\sphinxAtStartPar
\sphinxstylestrong{See also:} {\hyperref[\detokenize{reference/optimize:optimize-tolerance}]{\sphinxcrossref{\DUrole{std}{\DUrole{std-ref}{tolerance (\$optimize)}}}}}

\sphinxstepscope




\section{\sphinxstyleliteralintitle{\sphinxupquote{\$integral}}}
\label{\detokenize{reference/integral:integral}}\label{\detokenize{reference/integral:group-integral}}\label{\detokenize{reference/integral::doc}}
\index{\$integral@\spxentry{\$integral}}\ignorespaces 
\sphinxAtStartPar
Integral evaluation for RI methods.


\subsection{\sphinxstyleliteralintitle{\sphinxupquote{direct}}}
\label{\detokenize{reference/integral:direct}}\label{\detokenize{reference/integral:integral-direct}}
\index{direct (\$integral)@\spxentry{direct}\spxextra{\$integral}}\ignorespaces 
\sphinxAtStartPar
\sphinxstylestrong{Type:} switch\\
\sphinxstylestrong{Default:} \sphinxcode{\sphinxupquote{0}}\\
\sphinxstylestrong{Status:} undocumented · not verified\\
\sphinxstylestrong{Sources:} input.txt (annotation corrupted); 03 \#0, \S{}17; 05 \S{}2

\sphinxAtStartPar
Selects direct (recomputed) Coulomb integrals for RI methods.


\begin{savenotes}\sphinxattablestart
\sphinxthistablewithglobalstyle
\sphinxthistablewithnovlinesstyle
\centering
\begin{tabulary}{\linewidth}[t]{\Y{0.3}\Y{0.7}}
\sphinxtoprule
\sphinxstyletheadfamily 
\sphinxAtStartPar
Value
&\sphinxstyletheadfamily 
\sphinxAtStartPar
Meaning
\\
\sphinxmidrule
\sphinxtableatstartofbodyhook
\sphinxAtStartPar
\sphinxcode{\sphinxupquote{0}}
&
\sphinxAtStartPar
Store integrals. (\sphinxstylestrong{default})
\\
\sphinxhline
\sphinxAtStartPar
\sphinxcode{\sphinxupquote{1}}
&
\sphinxAtStartPar
Recompute Coulomb integrals (direct mode) in \sphinxcode{\sphinxupquote{theory=rimp2}} and \sphinxcode{\sphinxupquote{theory=ridft}}.
\\
\sphinxbottomrule
\end{tabulary}
\sphinxtableafterendhook\par
\sphinxattableend\end{savenotes}

\begin{sphinxadmonition}{note}{Note:}
\sphinxAtStartPar
Direct mode roughly halves disk use, but it is slower. It is also the reported workaround when a small RI\sphinxhyphen{}MP2
job claims not to fit in memory, or crashes with \sphinxcode{\sphinxupquote{write VMNN(…)?}}.
\end{sphinxadmonition}

\sphinxAtStartPar
\sphinxstylestrong{Example:}

\begin{sphinxVerbatim}[commandchars=\\\{\}]
\PYGZdl{}integral direct=1 rout=2 \PYGZdl{}end
\end{sphinxVerbatim}


\subsection{\sphinxstyleliteralintitle{\sphinxupquote{routine}}}
\label{\detokenize{reference/integral:routine}}\label{\detokenize{reference/integral:integral-routine}}
\index{routine (\$integral)@\spxentry{routine}\spxextra{\$integral}}\ignorespaces 
\sphinxAtStartPar
\sphinxstylestrong{Type:} integer\\
\sphinxstylestrong{Default:} \sphinxcode{\sphinxupquote{0}}\\
\sphinxstylestrong{Status:} undocumented · not verified\\
\sphinxstylestrong{Sources:} input.txt (annotation corrupted); 03 \#0

\sphinxAtStartPar
Algorithm for the two\sphinxhyphen{}electron integrals in correlated methods. \sphinxcode{\sphinxupquote{input.txt}}’s example uses \sphinxcode{\sphinxupquote{rout=2}}, but no
source explains the value 2.


\begin{savenotes}\sphinxattablestart
\sphinxthistablewithglobalstyle
\sphinxthistablewithnovlinesstyle
\centering
\begin{tabulary}{\linewidth}[t]{\Y{0.3}\Y{0.7}}
\sphinxtoprule
\sphinxstyletheadfamily 
\sphinxAtStartPar
Value
&\sphinxstyletheadfamily 
\sphinxAtStartPar
Meaning
\\
\sphinxmidrule
\sphinxtableatstartofbodyhook
\sphinxAtStartPar
\sphinxcode{\sphinxupquote{0}}
&
\sphinxAtStartPar
Default algorithm. (\sphinxstylestrong{default})
\\
\sphinxhline
\sphinxAtStartPar
\sphinxcode{\sphinxupquote{1}}
&
\sphinxAtStartPar
Slightly slower but parallelizes better over many processes.
\\
\sphinxbottomrule
\end{tabulary}
\sphinxtableafterendhook\par
\sphinxattableend\end{savenotes}

\sphinxstepscope




\section{\sphinxstyleliteralintitle{\sphinxupquote{\$atoms}}}
\label{\detokenize{reference/atoms:atoms}}\label{\detokenize{reference/atoms:group-atoms}}\label{\detokenize{reference/atoms::doc}}
\index{\$atoms@\spxentry{\$atoms}}\ignorespaces 
\sphinxAtStartPar
Per\sphinxhyphen{}element core definitions for correlated (HF\sphinxhyphen{}based) methods.


\subsection{\sphinxstyleliteralintitle{\sphinxupquote{core}}}
\label{\detokenize{reference/atoms:core}}\label{\detokenize{reference/atoms:atoms-core}}
\index{core (\$atoms)@\spxentry{core}\spxextra{\$atoms}}\ignorespaces 
\sphinxAtStartPar
\sphinxstylestrong{Type:} indexed\sphinxhyphen{}array\\
\sphinxstylestrong{Default:} \sphinxstyleemphasis{none stated}\\
\sphinxstylestrong{Status:} current · verified on 25 (2025.02.28)\\
\sphinxstylestrong{Sources:} input.txt; 08 t16

\sphinxAtStartPar
For HF\sphinxhyphen{}based methods, the number of core shells for L=0,1,2,3, given for each atomic number. The values are
used only to compute the total number of core electrons in the molecule.

\begin{sphinxadmonition}{note}{Note:}
\sphinxAtStartPar
The index is 10·Z: \sphinxcode{\sphinxupquote{core(280)}} is nickel, \sphinxcode{\sphinxupquote{core(80)}} oxygen. In version 25, \sphinxcode{\sphinxupquote{core(80)=0,0,0,0}} removed the
oxygen 1s from the frozen core of an RI\sphinxhyphen{}MP2 run. The output reports the frozen core as \sphinxcode{\sphinxupquote{atom   8 core =  1}}
and \sphinxcode{\sphinxupquote{Number of core electrons 1}}; the number there counts core orbitals, not electrons.
\end{sphinxadmonition}

\sphinxAtStartPar
\sphinxstylestrong{Example:}

\begin{sphinxVerbatim}[commandchars=\\\{\}]
\PYGZdl{}atoms core(280)=2,1,0,0 \PYGZdl{}end
\PYGZsh{} 1s, 2s, 2p (10 electrons) of nickel are core
\end{sphinxVerbatim}

\sphinxAtStartPar
\sphinxstylestrong{See also:} {\hyperref[\detokenize{reference/scf:scf-core}]{\sphinxcrossref{\DUrole{std}{\DUrole{std-ref}{core (\$scf)}}}}}, {\hyperref[\detokenize{reference/atoms:atoms-mcore}]{\sphinxcrossref{\DUrole{std}{\DUrole{std-ref}{mcore (\$atoms)}}}}}


\subsection{\sphinxstyleliteralintitle{\sphinxupquote{mcore}}}
\label{\detokenize{reference/atoms:mcore}}\label{\detokenize{reference/atoms:atoms-mcore}}
\index{mcore (\$atoms)@\spxentry{mcore}\spxextra{\$atoms}}\ignorespaces 
\sphinxAtStartPar
\sphinxstylestrong{Type:} integer\\
\sphinxstylestrong{Default:} \sphinxstyleemphasis{none stated}\\
\sphinxstylestrong{Status:} current · not verified\\
\sphinxstylestrong{Sources:} input.txt

\sphinxAtStartPar
A simple way to add outer\sphinxhyphen{}core shells to the correlation treatment: \sphinxcode{\sphinxupquote{mcore=M}} adds all shells whose principal
quantum number is lower than that of the outermost valence shell by less than M.

\sphinxAtStartPar
\sphinxstylestrong{See also:} {\hyperref[\detokenize{reference/atoms:atoms-core}]{\sphinxcrossref{\DUrole{std}{\DUrole{std-ref}{core (\$atoms)}}}}}

\sphinxstepscope




\section{\sphinxstyleliteralintitle{\sphinxupquote{\$optimize}}}
\label{\detokenize{reference/optimize:optimize}}\label{\detokenize{reference/optimize:group-optimize}}\label{\detokenize{reference/optimize::doc}}
\index{\$optimize@\spxentry{\$optimize}}\ignorespaces 
\sphinxAtStartPar
Geometry optimization, saddle\sphinxhyphen{}point search, IRC and relaxed scans.

\sphinxAtStartPar
\sphinxstylestrong{Relevant when:} \sphinxcode{\sphinxupquote{task = optimize | optimize+hessian | irc | scan | opt\_n1}}


\subsection{\sphinxstyleliteralintitle{\sphinxupquote{saddle}}}
\label{\detokenize{reference/optimize:saddle}}\label{\detokenize{reference/optimize:optimize-saddle}}
\index{saddle (\$optimize)@\spxentry{saddle}\spxextra{\$optimize}}\ignorespaces 
\sphinxAtStartPar
\sphinxstylestrong{Type:} switch\\
\sphinxstylestrong{Default:} \sphinxcode{\sphinxupquote{0}}\\
\sphinxstylestrong{Status:} current · verified on 25 (2025.02.28)\\
\sphinxstylestrong{Sources:} input.txt; Dft.rtf; 04 \#216; 03 \S{}9; 08 t25

\sphinxAtStartPar
Type of stationary point sought.


\begin{savenotes}\sphinxattablestart
\sphinxthistablewithglobalstyle
\sphinxthistablewithnovlinesstyle
\centering
\begin{tabulary}{\linewidth}[t]{\Y{0.3}\Y{0.7}}
\sphinxtoprule
\sphinxstyletheadfamily 
\sphinxAtStartPar
Value
&\sphinxstyletheadfamily 
\sphinxAtStartPar
Meaning
\\
\sphinxmidrule
\sphinxtableatstartofbodyhook
\sphinxAtStartPar
\sphinxcode{\sphinxupquote{0}}
&
\sphinxAtStartPar
Optimize to a minimum. (\sphinxstylestrong{default})
\\
\sphinxhline
\sphinxAtStartPar
\sphinxcode{\sphinxupquote{1}}
&
\sphinxAtStartPar
Optimize to a saddle point.
\\
\sphinxbottomrule
\end{tabulary}
\sphinxtableafterendhook\par
\sphinxattableend\end{savenotes}

\begin{sphinxadmonition}{note}{Note:}
\sphinxAtStartPar
A saddle\sphinxhyphen{}point search needs an accurate Hessian in the input: compute one with \sphinxcode{\sphinxupquote{task=hessian}} and append its
\sphinxcode{\sphinxupquote{\$Energy}} block (the \sphinxcode{\sphinxupquote{eng\textgreater{}}} lines). A usual route is scan \(\rightarrow\) Hessian at the scan maximum \(\rightarrow\) \sphinxcode{\sphinxupquote{saddle=1}} \(\rightarrow\)
\sphinxcode{\sphinxupquote{task=hessian}} to confirm one imaginary frequency \(\rightarrow\) IRC. In version 25 the output confirms the search with
\sphinxcode{\sphinxupquote{Saddle point search: follow = 1}} and, when a Hessian was read from \sphinxcode{\sphinxupquote{\$Energy}}, \sphinxcode{\sphinxupquote{Lowest eigenvalues:}} with one
negative value.
\end{sphinxadmonition}

\sphinxAtStartPar
\sphinxstylestrong{See also:} {\hyperref[\detokenize{reference/optimize:optimize-follow}]{\sphinxcrossref{\DUrole{std}{\DUrole{std-ref}{follow (\$optimize)}}}}}, {\hyperref[\detokenize{reference/Energy:energy-h}]{\sphinxcrossref{\DUrole{std}{\DUrole{std-ref}{H (\$Energy)}}}}}


\subsection{\sphinxstyleliteralintitle{\sphinxupquote{tolerance}}}
\label{\detokenize{reference/optimize:tolerance}}\label{\detokenize{reference/optimize:optimize-tolerance}}
\index{tolerance (\$optimize)@\spxentry{tolerance}\spxextra{\$optimize}}\ignorespaces 
\sphinxAtStartPar
\sphinxstylestrong{Type:} list\\
\sphinxstylestrong{Default:} \sphinxcode{\sphinxupquote{0.0001,0.01}}\\
\sphinxstylestrong{Units:} au\\
\sphinxstylestrong{Status:} current · verified on 25 (2025.02.28)\\
\sphinxstylestrong{Sources:} input.txt; 01 \S{}4.6 (output echo); 08 t11, t12

\sphinxAtStartPar
Default: if only the first value is given, the second is 100 times larger (the recommended setting).

\sphinxAtStartPar
Convergence tolerance on the gradient (first value) and on the displacement (second value).

\sphinxAtStartPar
\sphinxstylestrong{See also:} {\hyperref[\detokenize{reference/grid:grid-accuracy}]{\sphinxcrossref{\DUrole{std}{\DUrole{std-ref}{accuracy (\$grid)}}}}}


\subsection{\sphinxstyleliteralintitle{\sphinxupquote{trust}}}
\label{\detokenize{reference/optimize:trust}}\label{\detokenize{reference/optimize:optimize-trust}}
\index{trust (\$optimize)@\spxentry{trust}\spxextra{\$optimize}}\ignorespaces 
\sphinxAtStartPar
\sphinxstylestrong{Type:} list\\
\sphinxstylestrong{Default:} \sphinxcode{\sphinxupquote{0.125,0.015625}}\\
\sphinxstylestrong{Status:} current · verified on 25 (2025.02.28)\\
\sphinxstylestrong{Sources:} input.txt; 01 \S{}4.6; 04 \#319; Dft.rtf; 08 t11, t12

\sphinxAtStartPar
Default: echoed by version 25 as \sphinxcode{\sphinxupquote{Trust radius = 0.12500000 (max), 0.01562500 (min)}}, both in internal coordinates and
with symmetry (\sphinxcode{\sphinxupquote{sym=1}}). v6, v13 and v15 echoed \sphinxcode{\sphinxupquote{0.2 (max), 0.0125 (min)}} in internal coordinates, and
0.1 / 0.00625 with \sphinxcode{\sphinxupquote{cartesian=1}}. If only the maximum is given, the minimum is 1/16 of it and the initial value 1/4.
\sphinxcode{\sphinxupquote{Dft.rtf}} gives 0.5 (older version).

\sphinxAtStartPar
Trust radius (the largest allowed step), given as maximum, minimum and initial values.


\subsection{\sphinxstyleliteralintitle{\sphinxupquote{follow}}}
\label{\detokenize{reference/optimize:follow}}\label{\detokenize{reference/optimize:optimize-follow}}
\index{follow (\$optimize)@\spxentry{follow}\spxextra{\$optimize}}\ignorespaces 
\sphinxAtStartPar
\sphinxstylestrong{Type:} integer\\
\sphinxstylestrong{Default:} \sphinxcode{\sphinxupquote{1}}\\
\sphinxstylestrong{Status:} current · verified on 25 (2025.02.28)\\
\sphinxstylestrong{Sources:} input.txt; Dft.rtf; 08 t25

\sphinxAtStartPar
Default: \sphinxcode{\sphinxupquote{1}} according to \sphinxcode{\sphinxupquote{Dft.rtf}}; version 25 echoes \sphinxcode{\sphinxupquote{Saddle point search: follow = 1}}.

\sphinxAtStartPar
Follow the n\sphinxhyphen{}th Hessian eigenvector in a saddle\sphinxhyphen{}point search.

\sphinxAtStartPar
\sphinxstylestrong{See also:} {\hyperref[\detokenize{reference/optimize:optimize-saddle}]{\sphinxcrossref{\DUrole{std}{\DUrole{std-ref}{saddle (\$optimize)}}}}}


\subsection{\sphinxstyleliteralintitle{\sphinxupquote{steps}}}
\label{\detokenize{reference/optimize:steps}}\label{\detokenize{reference/optimize:optimize-steps}}
\index{steps (\$optimize)@\spxentry{steps}\spxextra{\$optimize}}\ignorespaces 
\sphinxAtStartPar
\sphinxstylestrong{Type:} integer\\
\sphinxstylestrong{Default:} \sphinxcode{\sphinxupquote{128}}\\
\sphinxstylestrong{Status:} current · verified on 25 (2025.02.28)\\
\sphinxstylestrong{Sources:} input.txt; Dft.rtf; 03 \S{}8 (step/steps bug); 08 t14, t15

\sphinxAtStartPar
Maximum number of optimization cycles.

\begin{sphinxadmonition}{note}{Note:}
\sphinxAtStartPar
Version 25 echoes \sphinxcode{\sphinxupquote{Maximum number of steps = 128}} when \sphinxcode{\sphinxupquote{steps}} is not given; older sources give 50. When the
limit is reached the program prints no message: the output just ends, and the last geometry appears only in
\sphinxcode{\sphinxupquote{mol\textgreater{}}} lines, not \sphinxcode{\sphinxupquote{MOL\textgreater{}}}. If an optimization has not converged within about 50 steps, something is probably
wrong; raising \sphinxcode{\sphinxupquote{steps}} rarely helps. Computing a Hessian first and supplying it often does. In v6, if both \sphinxcode{\sphinxupquote{steps}} and \sphinxcode{\sphinxupquote{step}} are
given, the value of \sphinxcode{\sphinxupquote{step}} is used; give only one of them.
\end{sphinxadmonition}

\sphinxAtStartPar
\sphinxstylestrong{See also:} {\hyperref[\detokenize{reference/optimize:optimize-tolerance}]{\sphinxcrossref{\DUrole{std}{\DUrole{std-ref}{tolerance (\$optimize)}}}}}, {\hyperref[\detokenize{reference/d2edr2:d2edr2-steps}]{\sphinxcrossref{\DUrole{std}{\DUrole{std-ref}{steps (\$d2edr2)}}}}}


\subsection{\sphinxstyleliteralintitle{\sphinxupquote{cartesian}}}
\label{\detokenize{reference/optimize:cartesian}}\label{\detokenize{reference/optimize:optimize-cartesian}}
\index{cartesian (\$optimize)@\spxentry{cartesian}\spxextra{\$optimize}}\ignorespaces 
\sphinxAtStartPar
\sphinxstylestrong{Type:} switch\\
\sphinxstylestrong{Default:} \sphinxcode{\sphinxupquote{0}}\\
\sphinxstylestrong{Status:} current · not verified\\
\sphinxstylestrong{Sources:} input.txt; 01 \S{}4.6; 04 \#518

\sphinxAtStartPar
Coordinate system for the optimization. The output reports \sphinxcode{\sphinxupquote{Coordinates: internal}} or \sphinxcode{\sphinxupquote{cartesian}}.


\begin{savenotes}\sphinxattablestart
\sphinxthistablewithglobalstyle
\sphinxthistablewithnovlinesstyle
\centering
\begin{tabulary}{\linewidth}[t]{\Y{0.3}\Y{0.7}}
\sphinxtoprule
\sphinxstyletheadfamily 
\sphinxAtStartPar
Value
&\sphinxstyletheadfamily 
\sphinxAtStartPar
Meaning
\\
\sphinxmidrule
\sphinxtableatstartofbodyhook
\sphinxAtStartPar
\sphinxcode{\sphinxupquote{0}}
&
\sphinxAtStartPar
Redundant internal coordinates. (\sphinxstylestrong{default})
\\
\sphinxhline
\sphinxAtStartPar
\sphinxcode{\sphinxupquote{1}}
&
\sphinxAtStartPar
Cartesian coordinates (last resort).
\\
\sphinxbottomrule
\end{tabulary}
\sphinxtableafterendhook\par
\sphinxattableend\end{savenotes}

\begin{sphinxadmonition}{note}{Note:}
\sphinxAtStartPar
\sphinxcode{\sphinxupquote{cartesian=1}} is the fix when the program crashes with \sphinxcode{\sphinxupquote{process 0 got signal 6}} after a \sphinxcode{\sphinxupquote{too few}} message in
the internal\sphinxhyphen{}coordinate table.
\end{sphinxadmonition}


\subsection{\sphinxstyleliteralintitle{\sphinxupquote{back}}}
\label{\detokenize{reference/optimize:back}}\label{\detokenize{reference/optimize:optimize-back}}
\index{back (\$optimize)@\spxentry{back}\spxextra{\$optimize}}\ignorespaces 
\sphinxAtStartPar
\sphinxstylestrong{Type:} switch\\
\sphinxstylestrong{Default:} \sphinxcode{\sphinxupquote{0}}\\
\sphinxstylestrong{Status:} current · verified on 25 (2025.02.28)\\
\sphinxstylestrong{Sources:} input.txt; 02 \S{}3.2; 04 \S{}0; 08 t26, t42

\sphinxAtStartPar
IRC direction. A full IRC takes two runs, one per direction: version 25 with the default \sphinxcode{\sphinxupquote{back=0}}
followed one direction only (\sphinxcode{\sphinxupquote{n = 1, back = 0, down = 1}} in the output), and \sphinxcode{\sphinxupquote{back=1}} the other one
(negative \sphinxcode{\sphinxupquote{distance}} values on the \sphinxcode{\sphinxupquote{IRC\textgreater{}}} lines). Re\sphinxhyphen{}optimize the endpoints afterwards, because
an IRC does not end at a converged minimum.


\begin{savenotes}\sphinxattablestart
\sphinxthistablewithglobalstyle
\sphinxthistablewithnovlinesstyle
\centering
\begin{tabulary}{\linewidth}[t]{\Y{0.3}\Y{0.7}}
\sphinxtoprule
\sphinxstyletheadfamily 
\sphinxAtStartPar
Value
&\sphinxstyletheadfamily 
\sphinxAtStartPar
Meaning
\\
\sphinxmidrule
\sphinxtableatstartofbodyhook
\sphinxAtStartPar
\sphinxcode{\sphinxupquote{0}}
&
\sphinxAtStartPar
Trace the IRC path forwards. (\sphinxstylestrong{default})
\\
\sphinxhline
\sphinxAtStartPar
\sphinxcode{\sphinxupquote{1}}
&
\sphinxAtStartPar
Trace the IRC path backwards.
\\
\sphinxbottomrule
\end{tabulary}
\sphinxtableafterendhook\par
\sphinxattableend\end{savenotes}

\sphinxAtStartPar
\sphinxstylestrong{See also:} {\hyperref[\detokenize{reference/optimize:optimize-points}]{\sphinxcrossref{\DUrole{std}{\DUrole{std-ref}{points (\$optimize)}}}}}, {\hyperref[\detokenize{reference/optimize:optimize-radius}]{\sphinxcrossref{\DUrole{std}{\DUrole{std-ref}{radius (\$optimize)}}}}}


\subsection{\sphinxstyleliteralintitle{\sphinxupquote{points}}}
\label{\detokenize{reference/optimize:points}}\label{\detokenize{reference/optimize:optimize-points}}
\index{points (\$optimize)@\spxentry{points}\spxextra{\$optimize}}\ignorespaces 
\sphinxAtStartPar
\sphinxstylestrong{Type:} integer\\
\sphinxstylestrong{Default:} \sphinxstyleemphasis{none stated}\\
\sphinxstylestrong{Status:} current · verified on 25 (2025.02.28)\\
\sphinxstylestrong{Sources:} input.txt; 08 t26, t37

\sphinxAtStartPar
Number of points to optimize in an IRC run, or number of constrained geometry optimizations in a scan run.
In a scan the points include both ends: \sphinxcode{\sphinxupquote{points=4}} gives s = 0, 1/3, 2/3 and 1. An IRC can stop before
\sphinxcode{\sphinxupquote{points}} is reached when it arrives at a minimum.

\sphinxAtStartPar
\sphinxstylestrong{See also:} {\hyperref[\detokenize{reference/optimize:optimize-value}]{\sphinxcrossref{\DUrole{std}{\DUrole{std-ref}{value (\$optimize)}}}}}


\subsection{\sphinxstyleliteralintitle{\sphinxupquote{radius}}}
\label{\detokenize{reference/optimize:radius}}\label{\detokenize{reference/optimize:optimize-radius}}
\index{radius (\$optimize)@\spxentry{radius}\spxextra{\$optimize}}\ignorespaces 
\sphinxAtStartPar
\sphinxstylestrong{Type:} real\\
\sphinxstylestrong{Default:} \sphinxcode{\sphinxupquote{0.1}}\\
\sphinxstylestrong{Status:} current · verified on 25 (2025.02.28)\\
\sphinxstylestrong{Sources:} input.txt; Dft.rtf; 04 \#298; 08 t26

\sphinxAtStartPar
Default: \sphinxcode{\sphinxupquote{0.1}} as echoed by version 25 (\sphinxcode{\sphinxupquote{r = 0.1000000000}}); \sphinxcode{\sphinxupquote{Dft.rtf}} gives 0.2 (older version).

\sphinxAtStartPar
Radius, in mass\sphinxhyphen{}weighted coordinates, for IRC path optimization. Adjacent IRC points end up roughly twice
this distance apart. The IRC coordinate is reported in bohr·amu\textasciicircum{}1/2.

\begin{sphinxadmonition}{note}{Note:}
\sphinxAtStartPar
\sphinxcode{\sphinxupquote{Vicinity of a Stationary Point(?) \sphinxhyphen{}\sphinxhyphen{}\sphinxhyphen{} STOP.}} ends an IRC that has reached the region of a minimum; in
version 25 it was the normal end of a test IRC. Users have also seen it early, when the energy did not
decrease on a step; changing the step size is the reported fix there.
\end{sphinxadmonition}


\subsection{\sphinxstyleliteralintitle{\sphinxupquote{coordinates}}}
\label{\detokenize{reference/optimize:coordinates}}\label{\detokenize{reference/optimize:optimize-coordinates}}
\index{coordinates (\$optimize)@\spxentry{coordinates}\spxextra{\$optimize}}\ignorespaces 
\sphinxAtStartPar
\sphinxstylestrong{Type:} list\\
\sphinxstylestrong{Default:} \sphinxstyleemphasis{none stated}\\
\sphinxstylestrong{Status:} current · verified on 25 (2025.02.28)\\
\sphinxstylestrong{Sources:} input.txt; 03 \S{}5 (type 4); 08 t40b

\sphinxAtStartPar
Default: generated automatically.

\sphinxAtStartPar
List of redundant internal coordinates to use, each given as five integers.


\begin{savenotes}\sphinxattablestart
\sphinxthistablewithglobalstyle
\sphinxthistablewithnovlinesstyle
\centering
\begin{tabulary}{\linewidth}[t]{\Y{0.3}\Y{0.7}}
\sphinxtoprule
\sphinxstyletheadfamily 
\sphinxAtStartPar
Value
&\sphinxstyletheadfamily 
\sphinxAtStartPar
Meaning
\\
\sphinxmidrule
\sphinxtableatstartofbodyhook
\sphinxAtStartPar
\sphinxcode{\sphinxupquote{1,i,j,0,0}}
&
\sphinxAtStartPar
Distance i\textendash{}j.
\\
\sphinxhline
\sphinxAtStartPar
\sphinxcode{\sphinxupquote{2,i,j,k,0}}
&
\sphinxAtStartPar
Angle i\textendash{}j\textendash{}k.
\\
\sphinxhline
\sphinxAtStartPar
\sphinxcode{\sphinxupquote{3,i,j,k,l}}
&
\sphinxAtStartPar
Torsion i\textendash{}j\textendash{}k\textendash{}l.
\\
\sphinxhline
\sphinxAtStartPar
\sphinxcode{\sphinxupquote{4,i,j,k,l}}
&
\sphinxAtStartPar
Special coordinate for linear fragments. A user describes it as the angle between line l\textendash{}k and the normal to plane k\textendash{}j\textendash{}i.
\\
\sphinxhline
\sphinxAtStartPar
\sphinxcode{\sphinxupquote{5,i,j,k,l}}
&
\sphinxAtStartPar
Special coordinate for linear fragments (meaning not described).
\\
\sphinxbottomrule
\end{tabulary}
\sphinxtableafterendhook\par
\sphinxattableend\end{savenotes}

\begin{sphinxadmonition}{note}{Note:}
\sphinxAtStartPar
\sphinxcode{\sphinxupquote{input.txt}} spells this keyword \sphinxcode{\sphinxupquote{coordinated=}}, which is a typo: version 25 stops with \sphinxcode{\sphinxupquote{*** error in \textquotesingle{}\$optimize\textquotesingle{} : unknown name coordinated}}. Use \sphinxcode{\sphinxupquote{coordinates=}}.
\end{sphinxadmonition}

\sphinxAtStartPar
\sphinxstylestrong{See also:} {\hyperref[\detokenize{reference/optimize:optimize-fix}]{\sphinxcrossref{\DUrole{std}{\DUrole{std-ref}{fix (\$optimize)}}}}}


\subsection{\sphinxstyleliteralintitle{\sphinxupquote{fix}}}
\label{\detokenize{reference/optimize:fix}}\label{\detokenize{reference/optimize:optimize-fix}}
\index{fix (\$optimize)@\spxentry{fix}\spxextra{\$optimize}}\ignorespaces 
\sphinxAtStartPar
\sphinxstylestrong{Type:} list\\
\sphinxstylestrong{Default:} \sphinxstyleemphasis{none stated}\\
\sphinxstylestrong{Status:} current · verified on 25 (2025.02.28)\\
\sphinxstylestrong{Sources:} input.txt; 02 \S{}3.2; 03 \S{}5 (wish list); 04 \#372\sphinxhyphen{}\#373; 08 t37

\sphinxAtStartPar
Internal coordinates to keep fixed during the optimization, in the same five\sphinxhyphen{}integer format as
\sphinxcode{\sphinxupquote{coordinates}}. Several entries are separated by commas or spaces.

\begin{sphinxadmonition}{note}{Note:}
\sphinxAtStartPar
Atoms cannot be frozen in Cartesian space; only internal coordinates can be fixed. Two coordinates also cannot
be constrained to stay equal.
\end{sphinxadmonition}

\sphinxAtStartPar
\sphinxstylestrong{See also:} {\hyperref[\detokenize{reference/optimize:optimize-value}]{\sphinxcrossref{\DUrole{std}{\DUrole{std-ref}{value (\$optimize)}}}}}, {\hyperref[\detokenize{reference/optimize:optimize-coordinates}]{\sphinxcrossref{\DUrole{std}{\DUrole{std-ref}{coordinates (\$optimize)}}}}}


\subsection{\sphinxstyleliteralintitle{\sphinxupquote{value}}}
\label{\detokenize{reference/optimize:value}}\label{\detokenize{reference/optimize:optimize-value}}
\index{value (\$optimize)@\spxentry{value}\spxextra{\$optimize}}\ignorespaces 
\sphinxAtStartPar
\sphinxstylestrong{Type:} indexed\sphinxhyphen{}array\\
\sphinxstylestrong{Default:} \sphinxstyleemphasis{none stated}\\
\sphinxstylestrong{Units:} Å and degrees\\
\sphinxstylestrong{Status:} current · verified on 25 (2025.02.28)\\
\sphinxstylestrong{Sources:} input.txt; 03 \S{}5, \S{}8; 04 \#278\sphinxhyphen{}\#279, \#373; 08 t37

\sphinxAtStartPar
Default: computed from the input geometry.

\sphinxAtStartPar
Values of the fixed internal coordinates. For \sphinxcode{\sphinxupquote{task=scan}}, give all the initial values followed by all the
final values. \sphinxcode{\sphinxupquote{val(i)=}} fills the list starting at position i, so values left unset keep their input\sphinxhyphen{}geometry
values.

\begin{sphinxadmonition}{note}{Note:}
\sphinxAtStartPar
\sphinxstylestrong{Priroda’s dihedral angles have the opposite sign to the convention used by Molden, MOPAC and GAMESS.} Flip
the sign of any torsion value taken from another program, or the scan starts from a mangled geometry.
\end{sphinxadmonition}

\sphinxAtStartPar
\sphinxstylestrong{Example:}

\begin{sphinxVerbatim}[commandchars=\\\{\}]
\PYGZdl{}optimize
 fix   =1,7,9,0,0, 1,8,10,0,0, 3,8,7,9,0, 3,10,9,7,0
 val(1)=1.6,       1.7
 val(5)=2.0,       2.1
 points=5
\PYGZdl{}end
\PYGZsh{} 5 optimizations: r(7\PYGZhy{}9) = 1.6 … 2.0, r(8\PYGZhy{}10) = 1.7 … 2.1;
\PYGZsh{} the last two fixed coordinates keep their input values
\end{sphinxVerbatim}

\sphinxAtStartPar
\sphinxstylestrong{See also:} {\hyperref[\detokenize{reference/optimize:optimize-fix}]{\sphinxcrossref{\DUrole{std}{\DUrole{std-ref}{fix (\$optimize)}}}}}, {\hyperref[\detokenize{reference/optimize:optimize-points}]{\sphinxcrossref{\DUrole{std}{\DUrole{std-ref}{points (\$optimize)}}}}}


\subsection{\sphinxstyleliteralintitle{\sphinxupquote{gradient}}}
\label{\detokenize{reference/optimize:gradient}}\label{\detokenize{reference/optimize:optimize-gradient}}
\index{gradient (\$optimize)@\spxentry{gradient}\spxextra{\$optimize}}\ignorespaces 
\sphinxAtStartPar
\sphinxstylestrong{Type:} switch\\
\sphinxstylestrong{Default:} \sphinxstyleemphasis{none stated}\\
\sphinxstylestrong{Status:} current · not verified\\
\sphinxstylestrong{Sources:} input.txt

\sphinxAtStartPar
Set \sphinxcode{\sphinxupquote{gradient=0}} when the starting second derivatives were computed with a different method or basis set.
This excludes the gradient terms from the internal/Cartesian transformation at the first geometry.
Otherwise keep the default.


\subsection{\sphinxstyleliteralintitle{\sphinxupquote{print}}}
\label{\detokenize{reference/optimize:print}}\label{\detokenize{reference/optimize:optimize-print}}
\index{print (\$optimize)@\spxentry{print}\spxextra{\$optimize}}\ignorespaces 
\sphinxAtStartPar
\sphinxstylestrong{Type:} switch\\
\sphinxstylestrong{Default:} \sphinxstyleemphasis{none stated}\\
\sphinxstylestrong{Status:} current · not verified\\
\sphinxstylestrong{Sources:} input.txt

\sphinxAtStartPar
Amount of optimization output.


\begin{savenotes}\sphinxattablestart
\sphinxthistablewithglobalstyle
\sphinxthistablewithnovlinesstyle
\centering
\begin{tabulary}{\linewidth}[t]{\Y{0.3}\Y{0.7}}
\sphinxtoprule
\sphinxstyletheadfamily 
\sphinxAtStartPar
Value
&\sphinxstyletheadfamily 
\sphinxAtStartPar
Meaning
\\
\sphinxmidrule
\sphinxtableatstartofbodyhook
\sphinxAtStartPar
\sphinxcode{\sphinxupquote{0}}
&
\sphinxAtStartPar
Minimal.
\\
\sphinxhline
\sphinxAtStartPar
\sphinxcode{\sphinxupquote{1}}
&
\sphinxAtStartPar
Detailed
\\
\sphinxbottomrule
\end{tabulary}
\sphinxtableafterendhook\par
\sphinxattableend\end{savenotes}


\subsection{\sphinxstyleliteralintitle{\sphinxupquote{symmetry}}}
\label{\detokenize{reference/optimize:symmetry}}\label{\detokenize{reference/optimize:optimize-symmetry}}
\index{symmetry (\$optimize)@\spxentry{symmetry}\spxextra{\$optimize}}\ignorespaces 
\sphinxAtStartPar
\sphinxstylestrong{Type:} switch\\
\sphinxstylestrong{Default:} \sphinxstyleemphasis{none stated}\\
\sphinxstylestrong{Status:} undocumented · verified on 25 (2025.02.28)\\
\sphinxstylestrong{Sources:} 01 \S{}4.6; 04 \#481; 09; 08 t15; 08 t30b

\sphinxAtStartPar
The full name of the \sphinxcode{\sphinxupquote{sym=1}} used in the official 2023\textendash{}2025 example inputs (any unique prefix is
accepted). In version 25, together with \sphinxcode{\sphinxupquote{\$molecule sym=1,0.01}},
the optimizer switched to Cartesian coordinates (\sphinxcode{\sphinxupquote{Coordinates: cartesian}}), printed \sphinxcode{\sphinxupquote{symmetry: C3v (6)}} at
every step, and returned an exactly symmetric structure followed by \sphinxcode{\sphinxupquote{MOL\textgreater{} symmetry: C3v}}. Without
\sphinxcode{\sphinxupquote{\$molecule sym=}}, \sphinxcode{\sphinxupquote{\$optimize sym=1}} only switches to Cartesian coordinates and reports \sphinxcode{\sphinxupquote{symmetry: C1 (1)}},
so give both options.

\sphinxAtStartPar
\sphinxstylestrong{See also:} {\hyperref[\detokenize{reference/molecule:molecule-sym}]{\sphinxcrossref{\DUrole{std}{\DUrole{std-ref}{sym (\$molecule)}}}}}


\subsection{\sphinxstyleliteralintitle{\sphinxupquote{model}}}
\label{\detokenize{reference/optimize:model}}\label{\detokenize{reference/optimize:optimize-model}}
\index{model (\$optimize)@\spxentry{model}\spxextra{\$optimize}}\ignorespaces 
\sphinxAtStartPar
\sphinxstylestrong{Type:} integer\\
\sphinxstylestrong{Default:} \sphinxstyleemphasis{none stated}\\
\sphinxstylestrong{Status:} undocumented · not verified\\
\sphinxstylestrong{Sources:} 01 \S{}4.6; 09

\sphinxAtStartPar
Used as \sphinxcode{\sphinxupquote{mod=1}} with \sphinxcode{\sphinxupquote{task=opt\_n1}} (very large systems). In version 25, \sphinxcode{\sphinxupquote{model}} is read only by the separate
optimizer of \sphinxcode{\sphinxupquote{task=opt\_n1}}, together with \sphinxcode{\sphinxupquote{function}}, \sphinxcode{\sphinxupquote{h0}}, \sphinxcode{\sphinxupquote{delta}}, \sphinxcode{\sphinxupquote{fragment}}, \sphinxcode{\sphinxupquote{steps}}, \sphinxcode{\sphinxupquote{trust}},
\sphinxcode{\sphinxupquote{tolerance}} and \sphinxcode{\sphinxupquote{print}}. Meaning not documented.

\sphinxstepscope




\section{\sphinxstyleliteralintitle{\sphinxupquote{\$d2edr2}}}
\label{\detokenize{reference/d2edr2:d2edr2}}\label{\detokenize{reference/d2edr2:group-d2edr2}}\label{\detokenize{reference/d2edr2::doc}}
\index{\$d2edr2@\spxentry{\$d2edr2}}\ignorespaces 
\sphinxAtStartPar
Numerical Hessian by finite differences of the gradient.


\subsection{\sphinxstyleliteralintitle{\sphinxupquote{length}}}
\label{\detokenize{reference/d2edr2:length}}\label{\detokenize{reference/d2edr2:d2edr2-length}}
\index{length (\$d2edr2)@\spxentry{length}\spxextra{\$d2edr2}}\ignorespaces 
\sphinxAtStartPar
\sphinxstylestrong{Type:} real\\
\sphinxstylestrong{Default:} \sphinxstyleemphasis{none stated}\\
\sphinxstylestrong{Status:} current · not verified\\
\sphinxstylestrong{Sources:} input.txt

\sphinxAtStartPar
Step length for numerical differentiation of the potential\sphinxhyphen{}energy surface.


\subsection{\sphinxstyleliteralintitle{\sphinxupquote{displace}}}
\label{\detokenize{reference/d2edr2:displace}}\label{\detokenize{reference/d2edr2:d2edr2-displace}}
\index{displace (\$d2edr2)@\spxentry{displace}\spxextra{\$d2edr2}}\ignorespaces 
\sphinxAtStartPar
\sphinxstylestrong{Type:} switch\\
\sphinxstylestrong{Default:} \sphinxcode{\sphinxupquote{2}}\\
\sphinxstylestrong{Status:} current · verified on 25 (2025.02.28)\\
\sphinxstylestrong{Sources:} input.txt; 08 t55

\sphinxAtStartPar
Number of displacements per coordinate. More accurate numerical Hessians may also need tighter SCF convergence.


\begin{savenotes}\sphinxattablestart
\sphinxthistablewithglobalstyle
\sphinxthistablewithnovlinesstyle
\centering
\begin{tabulary}{\linewidth}[t]{\Y{0.3}\Y{0.7}}
\sphinxtoprule
\sphinxstyletheadfamily 
\sphinxAtStartPar
Value
&\sphinxstyletheadfamily 
\sphinxAtStartPar
Meaning
\\
\sphinxmidrule
\sphinxtableatstartofbodyhook
\sphinxAtStartPar
\sphinxcode{\sphinxupquote{1}}
&
\sphinxAtStartPar
One displacement per coordinate.
\\
\sphinxhline
\sphinxAtStartPar
\sphinxcode{\sphinxupquote{2}}
&
\sphinxAtStartPar
Two displacements per coordinate. More accurate but twice as slow. (\sphinxstylestrong{default})
\\
\sphinxbottomrule
\end{tabulary}
\sphinxtableafterendhook\par
\sphinxattableend\end{savenotes}

\sphinxAtStartPar
\sphinxstylestrong{See also:} {\hyperref[\detokenize{reference/scf:scf-convergence}]{\sphinxcrossref{\DUrole{std}{\DUrole{std-ref}{convergence (\$scf)}}}}}


\subsection{\sphinxstyleliteralintitle{\sphinxupquote{steps}}}
\label{\detokenize{reference/d2edr2:steps}}\label{\detokenize{reference/d2edr2:d2edr2-steps}}
\index{steps (\$d2edr2)@\spxentry{steps}\spxextra{\$d2edr2}}\ignorespaces 
\sphinxAtStartPar
\sphinxstylestrong{Type:} list\\
\sphinxstylestrong{Default:} \sphinxcode{\sphinxupquote{0,\sphinxhyphen{}1}}\\
\sphinxstylestrong{Status:} current · verified on 25 (2025.02.28)\\
\sphinxstylestrong{Sources:} input.txt; 08 t55

\sphinxAtStartPar
Default: \sphinxcode{\sphinxupquote{0,\sphinxhyphen{}1}} runs every step from 0 to 3N (or 6N) in a single run.

\sphinxAtStartPar
The first and last displacement steps to do in this run. The reference geometry is step 0, and the displaced
geometries are 1 … 3N (one displacement) or 1 … 6N (two displacements). This lets a numerical Hessian be split
across separate runs and then combined.

\begin{sphinxadmonition}{note}{Note:}
\sphinxAtStartPar
Not the same keyword as \sphinxcode{\sphinxupquote{\$optimize steps}}. Tested in version 25 with water (\sphinxcode{\sphinxupquote{ridft}}, \sphinxcode{\sphinxupquote{displace=1}}, steps 0\textendash{}9):
two partial runs (\sphinxcode{\sphinxupquote{steps=0,4}}, \sphinxcode{\sphinxupquote{steps=5,9}}) and a combining run with \sphinxcode{\sphinxupquote{steps=\sphinxhyphen{}1}} gave the same frequencies as a
one\sphinxhyphen{}run Hessian to within 0.4 cm⁻\(\sp{\text{1}}\). The combining run does no SCF. The collecting command in \sphinxcode{\sphinxupquote{input.txt}}
(\sphinxcode{\sphinxupquote{grep "num\textgreater{}" … | sed "s/\textbackslash{}.*num\textgreater{}//g"}}) leaves the file names that \sphinxcode{\sphinxupquote{grep}} adds at the start of each line;
the example uses \sphinxcode{\sphinxupquote{grep \sphinxhyphen{}h}} instead.
\end{sphinxadmonition}

\sphinxAtStartPar
\sphinxstylestrong{Example:}

\begin{sphinxVerbatim}[commandchars=\\\{\}]
\PYGZsh{} 12\PYGZhy{}atom molecule, four inputs run in parallel on different machines:
\PYGZdl{}d2edr2 disp=1 steps= 0, 9 \PYGZdl{}end    \PYGZsh{} mol.in1
\PYGZdl{}d2edr2 disp=1 steps=10,18 \PYGZdl{}end    \PYGZsh{} mol.in2
\PYGZdl{}d2edr2 disp=1 steps=19,27 \PYGZdl{}end    \PYGZsh{} mol.in3
\PYGZdl{}d2edr2 disp=1 steps=28,36 \PYGZdl{}end    \PYGZsh{} mol.in4
\PYGZsh{} collect the partial results into one input:
grep \PYGZhy{}h \PYGZdq{}\PYGZca{}num\PYGZgt{}\PYGZdq{} mol.out1 mol.out2 mol.out3 mol.out4 | sed \PYGZsq{}s/\PYGZca{}num\PYGZgt{}//\PYGZsq{} \PYGZgt{}\PYGZgt{} mol.in
\PYGZsh{} add the line below and run once more to process them:
\PYGZdl{}d2edr2 steps=\PYGZhy{}1 \PYGZdl{}end
\end{sphinxVerbatim}

\sphinxstepscope




\section{\sphinxstyleliteralintitle{\sphinxupquote{\$thermo}}}
\label{\detokenize{reference/thermo:thermo}}\label{\detokenize{reference/thermo:group-thermo}}\label{\detokenize{reference/thermo::doc}}
\index{\$thermo@\spxentry{\$thermo}}\ignorespaces 
\sphinxAtStartPar
Thermochemistry after a Hessian calculation.


\subsection{\sphinxstyleliteralintitle{\sphinxupquote{sigma}}}
\label{\detokenize{reference/thermo:sigma}}\label{\detokenize{reference/thermo:thermo-sigma}}
\index{sigma (\$thermo)@\spxentry{sigma}\spxextra{\$thermo}}\ignorespaces 
\sphinxAtStartPar
\sphinxstylestrong{Type:} integer\\
\sphinxstylestrong{Default:} \sphinxcode{\sphinxupquote{1}}\\
\sphinxstylestrong{Status:} current · verified on 25 (2025.02.28)\\
\sphinxstylestrong{Sources:} input.txt; 08 t13; 08 t44

\sphinxAtStartPar
Default: 1 (version 25 does not detect the symmetry number).

\sphinxAtStartPar
Rotational symmetry number for thermochemistry. Version 25 does not detect it and uses 1 unless \sphinxcode{\sphinxupquote{sigma=}} is
given, which overestimates the rotational entropy by R ln σ (2.18 cal/(mol·K) for NH3, σ = 3). With
\sphinxcode{\sphinxupquote{sigma=3}}, NH3 gave S = 46.03 cal/(mol·K) at 298.15 K.


\subsection{\sphinxstyleliteralintitle{\sphinxupquote{temperature}}}
\label{\detokenize{reference/thermo:temperature}}\label{\detokenize{reference/thermo:thermo-temperature}}
\index{temperature (\$thermo)@\spxentry{temperature}\spxextra{\$thermo}}\ignorespaces 
\sphinxAtStartPar
\sphinxstylestrong{Type:} list\\
\sphinxstylestrong{Default:} \sphinxcode{\sphinxupquote{298.15}}\\
\sphinxstylestrong{Units:} K\\
\sphinxstylestrong{Status:} current · verified on 25 (2025.02.28)\\
\sphinxstylestrong{Sources:} input.txt; Dft.rtf; 08 t13; 08 t44

\sphinxAtStartPar
Default: \sphinxcode{\sphinxupquote{298.15}}: version 25 prints \sphinxcode{\sphinxupquote{T =  298.150 K}}, as \sphinxcode{\sphinxupquote{Dft.rtf}} states. Older posted outputs also show \sphinxcode{\sphinxupquote{298.000 K}}
(2011, 2021), so set \sphinxcode{\sphinxupquote{temperature=}} explicitly when comparing with other versions.

\sphinxAtStartPar
Temperatures at which thermochemistry is printed, up to 50 values. Each temperature gets its own table:
\sphinxcode{\sphinxupquote{temperature=200,298.15,400}} printed three.


\subsection{\sphinxstyleliteralintitle{\sphinxupquote{mass}}}
\label{\detokenize{reference/thermo:mass}}\label{\detokenize{reference/thermo:thermo-mass}}
\index{mass (\$thermo)@\spxentry{mass}\spxextra{\$thermo}}\ignorespaces 
\sphinxAtStartPar
\sphinxstylestrong{Type:} list\\
\sphinxstylestrong{Default:} \sphinxstyleemphasis{none stated}\\
\sphinxstylestrong{Units:} amu\\
\sphinxstylestrong{Status:} undocumented · not verified\\
\sphinxstylestrong{Sources:} Dft.rtf

\sphinxAtStartPar
Atomic masses for vibrational analysis and thermochemistry. Use it to redo the analysis for other isotopes
from a stored Hessian (\sphinxcode{\sphinxupquote{\$Energy}}) without recomputing it.

\sphinxAtStartPar
\sphinxstylestrong{See also:} {\hyperref[\detokenize{reference/molecule:molecule-mass}]{\sphinxcrossref{\DUrole{std}{\DUrole{std-ref}{mass (\$molecule)}}}}}, {\hyperref[\detokenize{reference/Energy:energy-h}]{\sphinxcrossref{\DUrole{std}{\DUrole{std-ref}{H (\$Energy)}}}}}

\sphinxstepscope




\section{\sphinxstyleliteralintitle{\sphinxupquote{\$nmr}}}
\label{\detokenize{reference/nmr:nmr}}\label{\detokenize{reference/nmr:group-nmr}}\label{\detokenize{reference/nmr::doc}}
\index{\$nmr@\spxentry{\$nmr}}\ignorespaces 
\sphinxAtStartPar
Reference shieldings and extra points for NMR shielding calculations.

\sphinxAtStartPar
\sphinxstylestrong{Relevant when:} \sphinxcode{\sphinxupquote{task = nmr (theory = dft only)}}


\subsection{\sphinxstyleliteralintitle{\sphinxupquote{standard}}}
\label{\detokenize{reference/nmr:standard}}\label{\detokenize{reference/nmr:nmr-standard}}
\index{standard (\$nmr)@\spxentry{standard}\spxextra{\$nmr}}\ignorespaces 
\sphinxAtStartPar
\sphinxstylestrong{Type:} indexed\sphinxhyphen{}array\\
\sphinxstylestrong{Default:} \sphinxstyleemphasis{none stated}\\
\sphinxstylestrong{Status:} current · verified on 25 (2025.02.28)\\
\sphinxstylestrong{Sources:} input.txt; 04 \#69; 02 \S{}5.7; 08 t34; 08 t57b

\sphinxAtStartPar
Reference shielding constants, indexed by atomic number, for computing chemical shifts. Values are needed
only for the elements present in the molecule.

\begin{sphinxadmonition}{note}{Note:}
\sphinxAtStartPar
A reference shielding is only valid for the method, basis set and grid it was computed with. Users quote
different TMS values (e.g. \sphinxcode{\sphinxupquote{standard(1)=31.3867 standard(6)=181.9781}}) for other settings. Compute TMS with
exactly the settings of the production run. Without \sphinxcode{\sphinxupquote{standard}}, version 25 prints \sphinxcode{\sphinxupquote{shift=}} as minus the
isotropic shielding; with it, \sphinxcode{\sphinxupquote{shift = standard − isotropic shielding}}.
\end{sphinxadmonition}

\sphinxAtStartPar
\sphinxstylestrong{Example:}

\begin{sphinxVerbatim}[commandchars=\\\{\}]
\PYGZdl{}nmr
 standard(1)=31.3915
 standard(6)=181.9921
\PYGZdl{}end
\PYGZsh{} computed for Si(CH3)4; usable for a series of hydrocarbons
\end{sphinxVerbatim}


\subsection{\sphinxstyleliteralintitle{\sphinxupquote{points}}}
\label{\detokenize{reference/nmr:points}}\label{\detokenize{reference/nmr:nmr-points}}
\index{points (\$nmr)@\spxentry{points}\spxextra{\$nmr}}\ignorespaces 
\sphinxAtStartPar
\sphinxstylestrong{Type:} list\\
\sphinxstylestrong{Default:} \sphinxstyleemphasis{none stated}\\
\sphinxstylestrong{Units:} Å\\
\sphinxstylestrong{Status:} current · not verified\\
\sphinxstylestrong{Sources:} input.txt; 02 \S{}3.2; 03 \S{}12

\sphinxAtStartPar
Extra points in space, given as x,y,z triples, at which magnetic shielding tensors are computed in
addition to the nuclear centres. This is the way to compute NICS: NICS is the negative of the isotropic
shielding printed for the point.

\sphinxAtStartPar
\sphinxstylestrong{Example:}

\begin{sphinxVerbatim}[commandchars=\\\{\}]
\PYGZdl{}nmr
 points=
      0.00000000    0.00000000    0.00000000
      0.00000000    0.00000000    1.00000000
\PYGZdl{}end
\end{sphinxVerbatim}

\sphinxstepscope




\section{\sphinxstyleliteralintitle{\sphinxupquote{\$Energy}}}
\label{\detokenize{reference/Energy:energy}}\label{\detokenize{reference/Energy:group-energy}}\label{\detokenize{reference/Energy::doc}}
\index{\$Energy@\spxentry{\$Energy}}\ignorespaces 
\sphinxAtStartPar
Results of an earlier run (energy, gradient, Hessian, dipole data), read back as input.

\sphinxAtStartPar
The program prints this group on lines prefixed with \sphinxcode{\sphinxupquote{eng\textgreater{}}} (an energy/gradient/Hessian run), or \sphinxcode{\sphinxupquote{num\textgreater{}}} (the
pieces of a distributed numerical Hessian; see \sphinxcode{\sphinxupquote{\$d2edr2 steps}}). Strip the prefix and append the group to the
next input. The shipped helper script \sphinxcode{\sphinxupquote{ge}} does this.

\sphinxAtStartPar
A stored Hessian is required for \sphinxcode{\sphinxupquote{\$optimize saddle=1}} and \sphinxcode{\sphinxupquote{task=irc}}. It also speeds up difficult minimizations,
and it lets you redo vibrational analysis or thermochemistry (for example for other isotopes, \sphinxcode{\sphinxupquote{\$thermo mass=}})
without recomputing.

\begin{sphinxVerbatim}[commandchars=\\\{\}]
eng\PYGZgt{}\PYGZdl{}Energy
eng\PYGZgt{} E=\PYGZhy{}2.284057276940e+03
eng\PYGZgt{} D= 4.22290114e\PYGZhy{}08 \PYGZhy{}1.82061636e\PYGZhy{}08  2.32661117e\PYGZhy{}08
eng\PYGZgt{} G= 5.13271929e\PYGZhy{}06 \PYGZhy{}5.52234429e\PYGZhy{}06 \PYGZhy{}6.15997899e\PYGZhy{}06
eng\PYGZgt{}    ...
eng\PYGZgt{}\PYGZdl{}end
\end{sphinxVerbatim}


\subsection{\sphinxstyleliteralintitle{\sphinxupquote{E}}}
\label{\detokenize{reference/Energy:e}}\label{\detokenize{reference/Energy:energy-e}}
\index{E (\$Energy)@\spxentry{E}\spxextra{\$Energy}}\ignorespaces 
\sphinxAtStartPar
\sphinxstylestrong{Type:} real\\
\sphinxstylestrong{Default:} \sphinxstyleemphasis{none stated}\\
\sphinxstylestrong{Units:} hartree\\
\sphinxstylestrong{Status:} undocumented · not verified\\
\sphinxstylestrong{Sources:} 01 \S{}4.10; 02 \S{}3.2

\sphinxAtStartPar
Total energy.


\subsection{\sphinxstyleliteralintitle{\sphinxupquote{G}}}
\label{\detokenize{reference/Energy:g}}\label{\detokenize{reference/Energy:energy-g}}
\index{G (\$Energy)@\spxentry{G}\spxextra{\$Energy}}\ignorespaces 
\sphinxAtStartPar
\sphinxstylestrong{Type:} list\\
\sphinxstylestrong{Default:} \sphinxstyleemphasis{none stated}\\
\sphinxstylestrong{Units:} au\\
\sphinxstylestrong{Status:} undocumented · not verified\\
\sphinxstylestrong{Sources:} 01 \S{}4.10; 02 \S{}3.2

\sphinxAtStartPar
Cartesian gradient (3N values).


\subsection{\sphinxstyleliteralintitle{\sphinxupquote{H}}}
\label{\detokenize{reference/Energy:h}}\label{\detokenize{reference/Energy:energy-h}}
\index{H (\$Energy)@\spxentry{H}\spxextra{\$Energy}}\ignorespaces 
\sphinxAtStartPar
\sphinxstylestrong{Type:} list\\
\sphinxstylestrong{Default:} \sphinxstyleemphasis{none stated}\\
\sphinxstylestrong{Units:} au\\
\sphinxstylestrong{Status:} undocumented · verified on 25 (2025.02.28)\\
\sphinxstylestrong{Sources:} 02 \S{}3.2; 08 t24, t25, t41, t41b

\sphinxAtStartPar
Hessian as a lower triangle, 3N(3N+1)/2 values. Written by \sphinxcode{\sphinxupquote{task=hessian}} on \sphinxcode{\sphinxupquote{eng\textgreater{}}} lines; in version 25 a
\sphinxcode{\sphinxupquote{saddle=1}} optimization and an IRC read it back from an appended \sphinxcode{\sphinxupquote{\$Energy}} group. If the input holds several
\sphinxcode{\sphinxupquote{\$Energy}} groups, the last one is used.

\sphinxAtStartPar
\sphinxstylestrong{See also:} {\hyperref[\detokenize{reference/optimize:optimize-saddle}]{\sphinxcrossref{\DUrole{std}{\DUrole{std-ref}{saddle (\$optimize)}}}}}


\subsection{\sphinxstyleliteralintitle{\sphinxupquote{D}}}
\label{\detokenize{reference/Energy:d}}\label{\detokenize{reference/Energy:energy-d}}
\index{D (\$Energy)@\spxentry{D}\spxextra{\$Energy}}\ignorespaces 
\sphinxAtStartPar
\sphinxstylestrong{Type:} list\\
\sphinxstylestrong{Default:} \sphinxstyleemphasis{none stated}\\
\sphinxstylestrong{Units:} au\\
\sphinxstylestrong{Status:} undocumented · not verified\\
\sphinxstylestrong{Sources:} 01 \S{}4.10; 02 \S{}3.2 (inferred)

\sphinxAtStartPar
Dipole moment.


\subsection{\sphinxstyleliteralintitle{\sphinxupquote{Dx}}}
\label{\detokenize{reference/Energy:dx}}\label{\detokenize{reference/Energy:energy-dx}}
\index{Dx (\$Energy)@\spxentry{Dx}\spxextra{\$Energy}}\ignorespaces 
\sphinxAtStartPar
\sphinxstylestrong{Type:} list\\
\sphinxstylestrong{Default:} \sphinxstyleemphasis{none stated}\\
\sphinxstylestrong{Status:} undocumented · not verified\\
\sphinxstylestrong{Sources:} 02 \S{}3.2 (inferred)

\sphinxAtStartPar
Dipole derivatives with respect to nuclear coordinates (inferred from use: IR/Raman intensities).


\subsection{\sphinxstyleliteralintitle{\sphinxupquote{A}}}
\label{\detokenize{reference/Energy:a}}\label{\detokenize{reference/Energy:energy-a}}
\index{A (\$Energy)@\spxentry{A}\spxextra{\$Energy}}\ignorespaces 
\sphinxAtStartPar
\sphinxstylestrong{Type:} list\\
\sphinxstylestrong{Default:} \sphinxstyleemphasis{none stated}\\
\sphinxstylestrong{Status:} undocumented · not verified\\
\sphinxstylestrong{Sources:} 02 \S{}3.2 (inferred); 04 \#381

\sphinxAtStartPar
Polarizability, and its derivatives after \sphinxcode{\sphinxupquote{task=dipole}} (inferred; used for Raman intensities).

\sphinxstepscope




\section{\sphinxstyleliteralintitle{\sphinxupquote{\$dynamic}}}
\label{\detokenize{reference/dynamic:dynamic}}\label{\detokenize{reference/dynamic:group-dynamic}}\label{\detokenize{reference/dynamic::doc}}
\index{\$dynamic@\spxentry{\$dynamic}}\ignorespaces 
\sphinxAtStartPar
Molecular dynamics with a thermostat.

\sphinxAtStartPar
\sphinxstylestrong{Relevant when:} \sphinxcode{\sphinxupquote{task = thermostat}}

\sphinxAtStartPar
Known only from the official example inputs and a third\sphinxhyphen{}party input generator. Every meaning below is
inferred, except \sphinxcode{\sphinxupquote{file}}.

\begin{sphinxVerbatim}[commandchars=\\\{\}]
\PYGZdl{}dynamic
 e=0.001
 d=0.03
 n=1000,100
 m=10
 i=1
 file=H40C20/H40C20\PYGZus{}\PYGZpc{}06d.qm.in
\PYGZdl{}end
\end{sphinxVerbatim}


\subsection{\sphinxstyleliteralintitle{\sphinxupquote{e}}}
\label{\detokenize{reference/dynamic:e}}\label{\detokenize{reference/dynamic:dynamic-e}}
\index{e (\$dynamic)@\spxentry{e}\spxextra{\$dynamic}}\ignorespaces 
\sphinxAtStartPar
\sphinxstylestrong{Type:} real\\
\sphinxstylestrong{Default:} \sphinxstyleemphasis{none stated}\\
\sphinxstylestrong{Units:} hartree\\
\sphinxstylestrong{Status:} undocumented · not verified\\
\sphinxstylestrong{Sources:} 01 \S{}4.8; 02 \S{}3.2

\sphinxAtStartPar
Thermostat temperature as kT (inferred: a generator script sets e = T/315775).


\subsection{\sphinxstyleliteralintitle{\sphinxupquote{d}}}
\label{\detokenize{reference/dynamic:d}}\label{\detokenize{reference/dynamic:dynamic-d}}
\index{d (\$dynamic)@\spxentry{d}\spxextra{\$dynamic}}\ignorespaces 
\sphinxAtStartPar
\sphinxstylestrong{Type:} real\\
\sphinxstylestrong{Default:} \sphinxstyleemphasis{none stated}\\
\sphinxstylestrong{Status:} undocumented · not verified\\
\sphinxstylestrong{Sources:} 01 \S{}4.8; 02 \S{}3.2

\sphinxAtStartPar
Time step (inferred).


\subsection{\sphinxstyleliteralintitle{\sphinxupquote{n}}}
\label{\detokenize{reference/dynamic:n}}\label{\detokenize{reference/dynamic:dynamic-n}}
\index{n (\$dynamic)@\spxentry{n}\spxextra{\$dynamic}}\ignorespaces 
\sphinxAtStartPar
\sphinxstylestrong{Type:} list\\
\sphinxstylestrong{Default:} \sphinxstyleemphasis{none stated}\\
\sphinxstylestrong{Status:} undocumented · not verified\\
\sphinxstylestrong{Sources:} 01 \S{}4.8; 02 \S{}3.2

\sphinxAtStartPar
Number of steps and print interval (inferred).


\subsection{\sphinxstyleliteralintitle{\sphinxupquote{m}}}
\label{\detokenize{reference/dynamic:m}}\label{\detokenize{reference/dynamic:dynamic-m}}
\index{m (\$dynamic)@\spxentry{m}\spxextra{\$dynamic}}\ignorespaces 
\sphinxAtStartPar
\sphinxstylestrong{Type:} integer\\
\sphinxstylestrong{Default:} \sphinxstyleemphasis{none stated}\\
\sphinxstylestrong{Status:} undocumented · not verified\\
\sphinxstylestrong{Sources:} 01 \S{}4.8; 02 \S{}3.2

\sphinxAtStartPar
Meaning not known.


\subsection{\sphinxstyleliteralintitle{\sphinxupquote{i}}}
\label{\detokenize{reference/dynamic:i}}\label{\detokenize{reference/dynamic:dynamic-i}}
\index{i (\$dynamic)@\spxentry{i}\spxextra{\$dynamic}}\ignorespaces 
\sphinxAtStartPar
\sphinxstylestrong{Type:} integer\\
\sphinxstylestrong{Default:} \sphinxstyleemphasis{none stated}\\
\sphinxstylestrong{Status:} undocumented · not verified\\
\sphinxstylestrong{Sources:} 01 \S{}4.8; 02 \S{}3.2

\sphinxAtStartPar
Run index or random seed (inferred).


\subsection{\sphinxstyleliteralintitle{\sphinxupquote{v}}}
\label{\detokenize{reference/dynamic:v}}\label{\detokenize{reference/dynamic:dynamic-v}}
\index{v (\$dynamic)@\spxentry{v}\spxextra{\$dynamic}}\ignorespaces 
\sphinxAtStartPar
\sphinxstylestrong{Type:} real\\
\sphinxstylestrong{Default:} \sphinxstyleemphasis{none stated}\\
\sphinxstylestrong{Units:} bohr\(\sp{\text{3}}\) per atom\\
\sphinxstylestrong{Status:} undocumented · not verified\\
\sphinxstylestrong{Sources:} 01 \S{}4.8; 02 \S{}3.2

\sphinxAtStartPar
Volume per atom of a spherical confining cavity, e.g. for simulating a liquid (inferred).


\subsection{\sphinxstyleliteralintitle{\sphinxupquote{h}}}
\label{\detokenize{reference/dynamic:h}}\label{\detokenize{reference/dynamic:dynamic-h}}
\index{h (\$dynamic)@\spxentry{h}\spxextra{\$dynamic}}\ignorespaces 
\sphinxAtStartPar
\sphinxstylestrong{Type:} real\\
\sphinxstylestrong{Default:} \sphinxstyleemphasis{none stated}\\
\sphinxstylestrong{Status:} undocumented · not verified\\
\sphinxstylestrong{Sources:} 01 \S{}4.8; 02 \S{}3.2

\sphinxAtStartPar
Stiffness of the cavity wall (inferred).


\subsection{\sphinxstyleliteralintitle{\sphinxupquote{file}}}
\label{\detokenize{reference/dynamic:file}}\label{\detokenize{reference/dynamic:dynamic-file}}
\index{file (\$dynamic)@\spxentry{file}\spxextra{\$dynamic}}\ignorespaces 
\sphinxAtStartPar
\sphinxstylestrong{Type:} string\\
\sphinxstylestrong{Default:} \sphinxstyleemphasis{none stated}\\
\sphinxstylestrong{Status:} undocumented · not verified\\
\sphinxstylestrong{Sources:} 01 \S{}4.8

\sphinxAtStartPar
printf\sphinxhyphen{}style pattern for snapshot files, which are written as ready\sphinxhyphen{}to\sphinxhyphen{}run inputs; \sphinxcode{\sphinxupquote{\%06d}} is the frame number.

\sphinxstepscope




\section{\sphinxstyleliteralintitle{\sphinxupquote{\$uniq}}}
\label{\detokenize{reference/uniq:uniq}}\label{\detokenize{reference/uniq:group-uniq}}\label{\detokenize{reference/uniq::doc}}
\index{\$uniq@\spxentry{\$uniq}}\ignorespaces 
\sphinxAtStartPar
Duplicate\sphinxhyphen{}structure filtering over several \sphinxcode{\sphinxupquote{\$molecule}} blocks.

\sphinxAtStartPar
\sphinxstylestrong{Relevant when:} \sphinxcode{\sphinxupquote{task = uniq}}


\subsection{\sphinxstyleliteralintitle{\sphinxupquote{distance}}}
\label{\detokenize{reference/uniq:distance}}\label{\detokenize{reference/uniq:uniq-distance}}
\index{distance (\$uniq)@\spxentry{distance}\spxextra{\$uniq}}\ignorespaces 
\sphinxAtStartPar
\sphinxstylestrong{Type:} real\\
\sphinxstylestrong{Default:} \sphinxstyleemphasis{none stated}\\
\sphinxstylestrong{Status:} undocumented · not verified\\
\sphinxstylestrong{Sources:} 01 \S{}4.9; 09

\sphinxAtStartPar
Threshold for treating two structures as identical (inferred from the official example \sphinxcode{\sphinxupquote{dist=0.01}}, run on 157
MD snapshots, each in its own \sphinxcode{\sphinxupquote{\$molecule}} block). \sphinxcode{\sphinxupquote{dist}} is a prefix of the full name \sphinxcode{\sphinxupquote{distance}}; the group
also reads \sphinxcode{\sphinxupquote{file}}.

\sphinxstepscope


\chapter{Solving common problems}
\label{\detokenize{troubleshooting/common-problems:solving-common-problems}}\label{\detokenize{troubleshooting/common-problems:troubleshooting-common-problems}}\label{\detokenize{troubleshooting/common-problems::doc}}
\sphinxAtStartPar
This chapter is arranged by symptom. Each entry names the likely cause and the fix that users have reported.

\sphinxAtStartPar
Almost everything here comes from user forums, not from the developer. Each fix says how well it is supported:
“confirmed by several users” means more than one person reported that the fix worked, and “single report” means one
user posted it and nobody confirmed it. Many reports date from Priroda 6 to 10 (2009\textendash{}2015). Where a problem may
have been fixed in later builds, the entry says so.

\begin{sphinxadmonition}{note}{Note:}
\sphinxAtStartPar
Most fixes below have not been checked against a current Priroda binary. Where a test run on version 25
(2025.02.28) confirmed or contradicted a report, the entry says so.
\end{sphinxadmonition}


\section{The input is not read, or is read wrongly}
\label{\detokenize{troubleshooting/common-problems:the-input-is-not-read-or-is-read-wrongly}}\label{\detokenize{troubleshooting/common-problems:troubleshooting-input}}
\sphinxAtStartPar
When Priroda cannot parse the input it usually stops early, often with no explanation. Before anything else, open
the output and check what it echoed back: the \sphinxcode{\sphinxupquote{Memory:}} and \sphinxcode{\sphinxupquote{Disk:}} lines, the \sphinxcode{\sphinxupquote{formula:}} line, the summary box
(number of atoms, charge, multiplicity, method, functional), the \sphinxcode{\sphinxupquote{scf:}} line and the \sphinxcode{\sphinxupquote{grid:}} line. If a value you set
is not shown there, Priroda did not read it.




\subsection{Output stops after \sphinxstyleliteralintitle{\sphinxupquote{molecule input: \textquotesingle{}\textquotesingle{}}}}
\label{\detokenize{troubleshooting/common-problems:output-stops-after-molecule-input}}
\sphinxAtStartPar
The line \sphinxcode{\sphinxupquote{molecule input: \textquotesingle{}\textquotesingle{}}} itself is printed in every run. It is a problem only when the output ends there or
continues with garbled text.



\sphinxAtStartPar
\sphinxstylestrong{Cause.} In almost every reported case the input file had Windows (CRLF) line endings. The tail of the output then
shows a mangled keyword followed by \sphinxcode{\sphinxupquote{?}}, for example \sphinxcode{\sphinxupquote{\textquotesingle{} ?matrix}}, \sphinxcode{\sphinxupquote{?rtesian}} or \sphinxcode{\sphinxupquote{\textquotesingle{}z\sphinxhyphen{}matrix\textquotesingle{} ?}}. Inputs saved as
Unicode, or edited in Word, fail the same way. Word’s “save as text” still writes CRLF line endings.



\sphinxAtStartPar
\sphinxstylestrong{Fix.} Save the input as plain ASCII text with Unix (LF) line endings, for example with \sphinxcode{\sphinxupquote{dos2unix input.in}}.
Notepad++ and Far Manager (Shift+F2 in the editor) can also convert line endings. Use straight quotes only. Several
users confirmed this fix. Running \sphinxcode{\sphinxupquote{dos2unix}} on a file that already has LF endings does no harm.




\subsection{The molecule is wrong: odd formula, extra atoms}
\label{\detokenize{troubleshooting/common-problems:the-molecule-is-wrong-odd-formula-extra-atoms}}
\sphinxAtStartPar
Atomic numbers written as real numbers (\sphinxcode{\sphinxupquote{8.0}} instead of \sphinxcode{\sphinxupquote{8}}) in \sphinxcode{\sphinxupquote{\$molecule}}. The atom lines are read as a stream of
numbers, so everything after the first \sphinxcode{\sphinxupquote{.0}} shifts. In version 25 there is \sphinxstylestrong{no error}: water written with \sphinxcode{\sphinxupquote{8.0}} and
\sphinxcode{\sphinxupquote{1.0}} ran as four atoms, an oxygen and three empty centres with Z=0, with the formula \sphinxcode{\sphinxupquote{\_\_3O}} and a meaningless
energy. An older version was reported to stop with \sphinxcode{\sphinxupquote{illegal charge/multiplicity}} instead. Write atomic numbers as
integers, and check the \sphinxcode{\sphinxupquote{formula:}} line and the atom count at the top of every output.



\sphinxAtStartPar
Avoid Cyrillic letters and spaces in the computer name, directory names and file names. One user’s problem with this
output went away after he had dealt with both line endings and such names.




\subsection{\sphinxstyleliteralintitle{\sphinxupquote{task=scan}} runs an optimization, or nothing works on one machine}
\label{\detokenize{troubleshooting/common-problems:task-scan-runs-an-optimization-or-nothing-works-on-one-machine}}
\sphinxAtStartPar
An input that runs on one computer but not on another, or a \sphinxcode{\sphinxupquote{task=scan}} that silently runs a plain optimization, has
been traced to the same encoding problem: an input edited in Word or saved as Unicode. Fix it as above.



\sphinxAtStartPar
Two \sphinxcode{\sphinxupquote{task=}} lines in one input do not chain jobs. Only the first one runs, and no error is printed. The only combined
task that works is \sphinxcode{\sphinxupquote{task=optimize+hessian}}; see {\hyperref[\detokenize{reference/control:control-task}]{\sphinxcrossref{\DUrole{std}{\DUrole{std-ref}{task}}}}}. Run other sequences, such as optimization
followed by NMR, as separate jobs.




\subsection{The basis file is not found}
\label{\detokenize{troubleshooting/common-problems:the-basis-file-is-not-found}}\begin{itemize}
\item {} 
\sphinxAtStartPar
Quote file names with straight double quotes: \sphinxcode{\sphinxupquote{basis="3z.bas"}}. Two single quotes (\sphinxcode{\sphinxupquote{basis=\textquotesingle{}\textquotesingle{}3z.bas\textquotesingle{}\textquotesingle{}}}) give
\sphinxcode{\sphinxupquote{error opening file \textquotesingle{}\textquotesingle{}3z.bas\textquotesingle{}\textquotesingle{}}}, and typographic quotes copied from a web page or word processor fail too.



\item {} 
\sphinxAtStartPar
Give an absolute path to the basis file if Priroda is not started from the directory that contains it. One user’s
parallel run aborted with \sphinxcode{\sphinxupquote{MPI\_Abort}} until he wrote the full basis path in the input. In version 25 a relative
name worked with the file in the current directory, in a serial and in a 2\sphinxhyphen{}process run.



\item {} 
\sphinxAtStartPar
The relativistic basis file \sphinxcode{\sphinxupquote{basis4.in}} needs {\hyperref[\detokenize{reference/control:control-four}]{\sphinxcrossref{\DUrole{std}{\DUrole{std-ref}{four=1}}}}} in {\hyperref[\detokenize{reference/control:group-control}]{\sphinxcrossref{\DUrole{std}{\DUrole{std-ref}{\$control}}}}}.
Without it the job stops right after the \sphinxcode{\sphinxupquote{Basis set input:}} line, with \sphinxcode{\sphinxupquote{1?}} as the only message (version 25).
The user who first reported this confirmed the fix.



\item {} 
\sphinxAtStartPar
The name given in {\hyperref[\detokenize{reference/molecule:molecule-set}]{\sphinxcrossref{\DUrole{std}{\DUrole{std-ref}{set}}}}} must exist in the basis file.



\end{itemize}


\subsection{The input or output file cannot be opened}
\label{\detokenize{troubleshooting/common-problems:the-input-or-output-file-cannot-be-opened}}
\sphinxAtStartPar
On Linux the usual cause is a wrong path. Check it with \sphinxcode{\sphinxupquote{ls \sphinxhyphen{}l}} before blaming Priroda. In one reported case the
directory in the command line simply did not exist. Work in a directory you own and can write to (for example
\sphinxcode{\sphinxupquote{\textasciitilde{}/priroda}}), \sphinxcode{\sphinxupquote{cd}} into it, and run the binary with an explicit path, since the current directory is normally not in
\sphinxcode{\sphinxupquote{PATH}}:



\begin{sphinxVerbatim}[commandchars=\\\{\}]
cd \PYGZti{}/priroda
./p job.in job.out
\end{sphinxVerbatim}

\sphinxAtStartPar
File names are case\sphinxhyphen{}sensitive, unlike keywords. One user of an older version saw the extra complaint
\sphinxcode{\sphinxupquote{0: output file \textquotesingle{}tasK.out\textquotesingle{}?}} when the file name contained capital letters. Version 25 wrote an output file with
capital letters in its name without complaint.



\sphinxAtStartPar
On Windows, \sphinxcode{\sphinxupquote{cd}} into the folder with the input before running \sphinxcode{\sphinxupquote{p.exe input.in output.out}}. One user still got
\sphinxcode{\sphinxupquote{cannot open input file}} after doing that, and the thread never found the cause.




\subsection{Output stops at \sphinxstyleliteralintitle{\sphinxupquote{0: opening file /tmp/...}} (Windows)}
\label{\detokenize{troubleshooting/common-problems:output-stops-at-0-opening-file-tmp-windows}}
\sphinxAtStartPar
Old Windows builds put scratch files in \sphinxcode{\sphinxupquote{/tmp}} when nothing else is set (the output echoes \sphinxcode{\sphinxupquote{Disk:   1 GB \textquotesingle{}/tmp\textquotesingle{}}}),
and Windows has no \sphinxcode{\sphinxupquote{/tmp}}. Set the scratch directory
explicitly, for example \sphinxcode{\sphinxupquote{\$system path=. \$end}}, or set the \sphinxcode{\sphinxupquote{TMPDIR}} environment variable before starting the program.
This is the oldest known problem in both forums and the fix is well established. See {\hyperref[\detokenize{reference/system:system-path}]{\sphinxcrossref{\DUrole{std}{\DUrole{std-ref}{path}}}}}.




\subsection{\sphinxstyleliteralintitle{\sphinxupquote{\$scf}} settings are ignored or rejected}
\label{\detokenize{troubleshooting/common-problems:scf-settings-are-ignored-or-rejected}}
\sphinxAtStartPar
The parser has a problem when {\hyperref[\detokenize{reference/scf:scf-convergence}]{\sphinxcrossref{\DUrole{std}{\DUrole{std-ref}{convergence}}}}} is directly followed by another keyword in
{\hyperref[\detokenize{reference/scf:group-scf}]{\sphinxcrossref{\DUrole{std}{\DUrole{std-ref}{\$scf}}}}}. In version 25,

\begin{sphinxVerbatim}[commandchars=\\\{\}]
\PYGZdl{}scf convergence=1e\PYGZhy{}7 iterations=250 \PYGZdl{}end
\end{sphinxVerbatim}

\sphinxAtStartPar
stops the run after the basis\sphinxhyphen{}set echo with \sphinxcode{\sphinxupquote{*** error in \textquotesingle{}\$scf\textquotesingle{} : unknown name terations}}: the value list of
\sphinxcode{\sphinxupquote{convergence}} takes the first letter of the next keyword. With {\hyperref[\detokenize{reference/scf:scf-procedure}]{\sphinxcrossref{\DUrole{std}{\DUrole{std-ref}{procedure=bfgs}}}}} between them,
\sphinxcode{\sphinxupquote{\$scf convergence=1e\sphinxhyphen{}7 procedure=bfgs iterations=250 \$end}}, both values are read. Forum users of 2015 builds saw the
same error as \sphinxcode{\sphinxupquote{unknown name erations}}, and one found that the pair worked only with \sphinxcode{\sphinxupquote{procedure=bfgs}}, which this
explains.



\sphinxAtStartPar
Further tests on version 25 narrowed it down. The problem is a \sphinxcode{\sphinxupquote{convergence=}} with \sphinxstylestrong{one} value followed by another
keyword, also on the next line. Both of these work:

\begin{sphinxVerbatim}[commandchars=\\\{\}]
\PYGZdl{}scf convergence=1e\PYGZhy{}7,1e\PYGZhy{}3 iterations=250 \PYGZdl{}end
\PYGZdl{}scf iterations=250 convergence=1e\PYGZhy{}7 \PYGZdl{}end
\end{sphinxVerbatim}

\sphinxAtStartPar
So give both values, or put \sphinxcode{\sphinxupquote{convergence}} last, and confirm the values in the \sphinxcode{\sphinxupquote{scf:}} echo line of the output, for
example \sphinxcode{\sphinxupquote{scf: conv = 1.0e\sphinxhyphen{}07, 1.0e\sphinxhyphen{}03; iter=250,16}}.




\section{SCF does not converge or diverges}
\label{\detokenize{troubleshooting/common-problems:scf-does-not-converge-or-diverges}}\label{\detokenize{troubleshooting/common-problems:troubleshooting-scf}}

\subsection{Recognizing the problem}
\label{\detokenize{troubleshooting/common-problems:recognizing-the-problem}}
\sphinxAtStartPar
Users report that Priroda’s SCF usually converges well, better than in many other programs. Trouble mostly comes
with metals in low oxidation states, where the electronic structure has some multiconfigurational character, with
very large basis sets, and with diffuse functions.



\sphinxAtStartPar
Signs of trouble in the SCF table:
\begin{itemize}
\item {} 
\sphinxAtStartPar
\sphinxcode{\sphinxupquote{hmin}} values that keep growing in magnitude, or a line such as \sphinxcode{\sphinxupquote{h = \sphinxhyphen{}25.15!}};

\item {} 
\sphinxAtStartPar
\sphinxcode{\sphinxupquote{lowest eigenvalue = ...!!!}} lines with large negative numbers;

\item {} 
\sphinxAtStartPar
an energy that rises from one macro\sphinxhyphen{}iteration to the next;

\item {} 
\sphinxAtStartPar
\sphinxcode{\sphinxupquote{no convergence!}}, in older versions followed by \sphinxcode{\sphinxupquote{SCF is far from convergence}}.

\end{itemize}



\begin{sphinxadmonition}{warning}{Warning:}
\sphinxAtStartPar
In version 25, a single\sphinxhyphen{}point run whose SCF hits the iteration limit prints \sphinxcode{\sphinxupquote{no convergence!}} and then carries on:
it prints an energy, orbitals and properties and exits normally. An OH radical with default settings did this after
the default 32 iterations, still creeping towards convergence. Always check for \sphinxcode{\sphinxupquote{converged!}}.
\end{sphinxadmonition}



\sphinxAtStartPar
The DFT SCF runs on a coarse grid first and then switches to the final grid; the switch is marked by a new
\sphinxcode{\sphinxupquote{grid:}} line with the number of points. A short jump in the gradient at that point is normal. Steady growth after it is not (see
{\hyperref[\detokenize{troubleshooting/common-problems:troubleshooting-fine-grid}]{\sphinxcrossref{\DUrole{std}{\DUrole{std-ref}{Divergence after the switch to the fine grid (open\sphinxhyphen{}shell transition\sphinxhyphen{}metal complex)}}}}}).




\subsection{Check the input first}
\label{\detokenize{troubleshooting/common-problems:check-the-input-first}}
\sphinxAtStartPar
In several reported failures the SCF was not the real problem. Before changing SCF settings, check:
\begin{itemize}
\item {} 
\sphinxAtStartPar
\sphinxstylestrong{Multiplicity.} {\hyperref[\detokenize{reference/molecule:molecule-mult}]{\sphinxcrossref{\DUrole{std}{\DUrole{std-ref}{mult}}}}} must match the electron count. One failing input gave \sphinxcode{\sphinxupquote{mult=0}} for
an odd number of electrons, and the summary box showed \sphinxcode{\sphinxupquote{Spin multiplicity 0}}.

\item {} 
\sphinxAtStartPar
\sphinxstylestrong{Theory and functional.} A hybrid functional with \sphinxcode{\sphinxupquote{theory=dft}} is wrong; see
{\hyperref[\detokenize{troubleshooting/common-problems:troubleshooting-hybrid}]{\sphinxcrossref{\DUrole{std}{\DUrole{std-ref}{Hybrid functional silently ignored}}}}}.

\item {} 
\sphinxAtStartPar
\sphinxstylestrong{Basis.} Non\sphinxhyphen{}native basis sets and diffuse functions converge badly. Priroda’s auxiliary (fitting) bases are not
designed for foreign orbital bases. Use a native basis, or start from a geometry optimized in one.

\item {} 
\sphinxAtStartPar
\sphinxstylestrong{Geometry.} A very poor starting geometry can by itself prevent SCF convergence. Pre\sphinxhyphen{}optimize with a cheap pure
functional (for example PBE with \sphinxcode{\sphinxupquote{theory=dft}}) before switching to the target method.

\end{itemize}




\subsection{Restart from saved vectors}
\label{\detokenize{troubleshooting/common-problems:restart-from-saved-vectors}}
\sphinxAtStartPar
This is the technique most often recommended by experienced users, on both forums, and the limor1 \sphinxcode{\sphinxupquote{pri}} driver
script automates it:


\begin{enumerate}
\sphinxsetlistlabels{\arabic}{enumi}{enumii}{}{.}%
\item {} 
\sphinxAtStartPar
Cap the number of SCF iterations at a moderate value, for example \sphinxcode{\sphinxupquote{iterations=10}} in {\hyperref[\detokenize{reference/scf:group-scf}]{\sphinxcrossref{\DUrole{std}{\DUrole{std-ref}{\$scf}}}}}.
Choose the cap so that \sphinxcode{\sphinxupquote{hmin}} does not run away.

\item {} 
\sphinxAtStartPar
Have Priroda write its vectors with {\hyperref[\detokenize{reference/control:control-save}]{\sphinxcrossref{\DUrole{std}{\DUrole{std-ref}{save}}}}} in {\hyperref[\detokenize{reference/control:group-control}]{\sphinxcrossref{\DUrole{std}{\DUrole{std-ref}{\$control}}}}}.

\item {} 
\sphinxAtStartPar
When the run stops with \sphinxcode{\sphinxupquote{SCF is far from convergence}}, add \sphinxcode{\sphinxupquote{\$guess read=1 \$end}} (see {\hyperref[\detokenize{reference/guess:guess-read}]{\sphinxcrossref{\DUrole{std}{\DUrole{std-ref}{read}}}}})
and rerun with the latest geometry.

\item {} 
\sphinxAtStartPar
Repeat until the SCF converges.

\end{enumerate}

\begin{sphinxVerbatim}[commandchars=\\\{\}]
\PYGZdl{}control
 task=optimize
 theory=dft
 basis=\PYGZdq{}basis.in\PYGZdq{}
 save=job.vec
\PYGZdl{}end
\PYGZdl{}scf iterations=10 \PYGZdl{}end
\PYGZdl{}guess read=1 \PYGZdl{}end
\end{sphinxVerbatim}



\sphinxAtStartPar
In version 25 the output of a restarted run shows \sphinxcode{\sphinxupquote{Initial Guess (\sphinxhyphen{}1)}} and \sphinxcode{\sphinxupquote{vectors from \textquotesingle{}job.vec\textquotesingle{}}}; a normal start
shows \sphinxcode{\sphinxupquote{Initial Guess (2)}}.



\sphinxAtStartPar
Leave out the \sphinxcode{\sphinxupquote{\$guess}} line in the first run. Vectors from a converged run of an easier case, or of the same molecule
at another geometry, also work as a starting point.




\subsection{Switch the SCF procedure and tighten the grid}
\label{\detokenize{troubleshooting/common-problems:switch-the-scf-procedure-and-tighten-the-grid}}
\sphinxAtStartPar
The default {\hyperref[\detokenize{reference/scf:scf-procedure}]{\sphinxcrossref{\DUrole{std}{\DUrole{std-ref}{procedure}}}}} is Newton\textendash{}Raphson (\sphinxcode{\sphinxupquote{nr}}). If it blows up, try \sphinxcode{\sphinxupquote{bfgs}}, and the other way
round. One experienced user recommends this setting for difficult cases, “which is nearly all of them”:



\begin{sphinxVerbatim}[commandchars=\\\{\}]
\PYGZdl{}scf
 convergence=1e\PYGZhy{}7
 procedure=bfgs
 iterations=250
\PYGZdl{}end
\PYGZdl{}grid
 accuracy=1e\PYGZhy{}9
\PYGZdl{}end
\end{sphinxVerbatim}



\sphinxAtStartPar
He notes that the default 16 iterations of the versions of that time are too few for many jobs; version 25
defaults to 32. If you tighten {\hyperref[\detokenize{reference/scf:scf-convergence}]{\sphinxcrossref{\DUrole{std}{\DUrole{std-ref}{convergence}}}}} to \sphinxcode{\sphinxupquote{1e\sphinxhyphen{}8}}, make the grid finer as well, \sphinxcode{\sphinxupquote{1e\sphinxhyphen{}9}} or \sphinxcode{\sphinxupquote{1e\sphinxhyphen{}10}} (see {\hyperref[\detokenize{reference/grid:grid-accuracy}]{\sphinxcrossref{\DUrole{std}{\DUrole{std-ref}{accuracy}}}}}).



\begin{sphinxadmonition}{note}{Note:}
\sphinxAtStartPar
The user who asked never reported whether this fixed the failing job. \sphinxcode{\sphinxupquote{bfgs}} is not a cure\sphinxhyphen{}all: in the case in the
next section it may have been part of the problem.
\end{sphinxadmonition}




\subsection{Divergence after the switch to the fine grid (open\sphinxhyphen{}shell transition\sphinxhyphen{}metal complex)}
\label{\detokenize{troubleshooting/common-problems:divergence-after-the-switch-to-the-fine-grid-open-shell-transition-metal-complex}}\label{\detokenize{troubleshooting/common-problems:troubleshooting-fine-grid}}
\sphinxAtStartPar
\sphinxstylestrong{Symptom.} A doublet Cu(II) complex (OLYP/L2, \sphinxcode{\sphinxupquote{restrict=0}}, \sphinxcode{\sphinxupquote{procedure=BFGS}}, \sphinxcode{\sphinxupquote{\$grid acc=1e\sphinxhyphen{}12}}, \sphinxcode{\sphinxupquote{task=hessian}})
converged on the coarse grid. After the switch to the final grid it oscillated for about 30 iterations and then
diverged, with the total energy reaching +29 hartree. BLYP behaved the same way.



\sphinxAtStartPar
\sphinxstylestrong{Fixes.} Two experienced users solved it in two ways, which they considered equivalent:
\begin{itemize}
\item {} 
\sphinxAtStartPar
Run \sphinxcode{\sphinxupquote{task=optimize+hessian}} instead of a single\sphinxhyphen{}point Hessian, and drop \sphinxcode{\sphinxupquote{procedure=BFGS}} and \sphinxcode{\sphinxupquote{iterations}}:

\begin{sphinxVerbatim}[commandchars=\\\{\}]
\PYGZdl{}control task=optimize+hessian \PYGZdl{}end
\PYGZdl{}scf restrict=0 convergence=1e\PYGZhy{}10 \PYGZdl{}end
\end{sphinxVerbatim}

\sphinxAtStartPar
Each optimization step then starts from the previous wavefunction.  

\item {} 
\sphinxAtStartPar
Save the wavefunction at the start of the optimization with {\hyperref[\detokenize{reference/control:control-save}]{\sphinxcrossref{\DUrole{std}{\DUrole{std-ref}{save}}}}} and read it into the Hessian
job with \sphinxcode{\sphinxupquote{\$guess read=1}}. For an already optimized structure, the guess built from scratch turned out to be poor.



\end{itemize}


\subsection{Open\sphinxhyphen{}shell singlets and broken symmetry}
\label{\detokenize{troubleshooting/common-problems:open-shell-singlets-and-broken-symmetry}}
\sphinxAtStartPar
For an open\sphinxhyphen{}shell (broken\sphinxhyphen{}symmetry) singlet, use an unrestricted calculation and mix the HOMO and LUMO of the guess:

\begin{sphinxVerbatim}[commandchars=\\\{\}]
\PYGZdl{}scf restrict=0 \PYGZdl{}end
\PYGZdl{}guess mix=1 \PYGZdl{}end
\end{sphinxVerbatim}



\sphinxAtStartPar
See {\hyperref[\detokenize{reference/scf:scf-restrict}]{\sphinxcrossref{\DUrole{std}{\DUrole{std-ref}{restrict}}}}} and {\hyperref[\detokenize{reference/guess:guess-mix}]{\sphinxcrossref{\DUrole{std}{\DUrole{std-ref}{mix}}}}}. With non\sphinxhyphen{}RI pure DFT (\sphinxcode{\sphinxupquote{theory=dft}}, PBE or BLYP), one
user found that \sphinxcode{\sphinxupquote{restrict=0 mix=1}} converged back to the restricted solution, with identical energies and identical
alpha and beta orbital energies. RI\sphinxhyphen{}DFT broke the symmetry in his case. Another user notes that pure GGA functionals
often do not break symmetry, and suggests a hybrid (with \sphinxcode{\sphinxupquote{theory=ridft}}) or a Hartree\textendash{}Fock calculation first.



\sphinxAtStartPar
In a version 25 test on singlet O2, \sphinxcode{\sphinxupquote{restrict=0 mix=1}} converged back to the restricted solution with both
\sphinxcode{\sphinxupquote{theory=dft}} (PBE) and \sphinxcode{\sphinxupquote{theory=ridft}} (PBE1): alpha and beta orbital energies were identical, and the PBE energy
equalled the closed\sphinxhyphen{}shell one. So \sphinxcode{\sphinxupquote{mix=1}} alone does not guarantee a broken\sphinxhyphen{}symmetry solution; compare the alpha and
beta orbital energies in the output.




\subsection{Heavy elements}
\label{\detokenize{troubleshooting/common-problems:heavy-elements}}
\sphinxAtStartPar
For scalar\sphinxhyphen{}relativistic pure DFT with the heaviest atoms, the v5 documentation suggests
{\hyperref[\detokenize{reference/scf:scf-d1small}]{\sphinxcrossref{\DUrole{std}{\DUrole{std-ref}{d1small=1}}}}}; without it, the quadratically convergent SCF may fail.



\sphinxAtStartPar
If an orbital ordering problem is suspected, {\hyperref[\detokenize{reference/guess:guess-na}]{\sphinxcrossref{\DUrole{std}{\DUrole{std-ref}{na}}}}} and {\hyperref[\detokenize{reference/guess:guess-nb}]{\sphinxcrossref{\DUrole{std}{\DUrole{std-ref}{nb}}}}} reorder the starting
vectors. 


\subsection{Cases without a known fix}
\label{\detokenize{troubleshooting/common-problems:cases-without-a-known-fix}}\begin{itemize}
\item {} 
\sphinxAtStartPar
An HF reference for CCSD diverged for a small closed\sphinxhyphen{}shell molecule (7 atoms, \sphinxcode{\sphinxupquote{set=L1}}, \sphinxcode{\sphinxupquote{restrict=0}}). The
\sphinxcode{\sphinxupquote{change}} column stuck at \sphinxcode{\sphinxupquote{0.50000000}}. Nobody replied. Running a closed\sphinxhyphen{}shell molecule with \sphinxcode{\sphinxupquote{restrict=0}} was
unusual and may be worth changing.



\item {} 
\sphinxAtStartPar
An HfCl2 molecule (PBE, \sphinxcode{\sphinxupquote{four=1}}, \sphinxcode{\sphinxupquote{basis4.in}}, \sphinxcode{\sphinxupquote{set=L3}}, \sphinxcode{\sphinxupquote{accuracy=1e\sphinxhyphen{}10}}) would not converge. Nobody replied.



\end{itemize}


\section{Geometry optimization problems}
\label{\detokenize{troubleshooting/common-problems:geometry-optimization-problems}}\label{\detokenize{troubleshooting/common-problems:troubleshooting-optimization}}

\subsection{The optimization does not converge within \sphinxstyleliteralintitle{\sphinxupquote{steps}}}
\label{\detokenize{troubleshooting/common-problems:the-optimization-does-not-converge-within-steps}}
\sphinxAtStartPar
The default limit is 128 steps in version 25 and 50 in older versions (see {\hyperref[\detokenize{reference/optimize:optimize-steps}]{\sphinxcrossref{\DUrole{std}{\DUrole{std-ref}{steps}}}}}). When the
limit is reached, the output just ends: there is no message and no \sphinxcode{\sphinxupquote{MOL\textgreater{}}} block. The documentation’s advice is that
an optimization still running after about 50 steps usually has some other problem, and raising \sphinxcode{\sphinxupquote{steps}} rarely helps.
What users do instead:


\begin{itemize}
\item {} 
\sphinxAtStartPar
\sphinxstylestrong{Restart from the last geometry.} The limor1 \sphinxcode{\sphinxupquote{pri}} script restarts Priroda from the last geometry whenever an
optimization runs out of steps. This helps when the optimizer cycles around one geometry even after 300 steps.



\item {} 
\sphinxAtStartPar
\sphinxstylestrong{Supply a Hessian.} Run \sphinxcode{\sphinxupquote{task=hessian}}, append its \sphinxcode{\sphinxupquote{\$Energy}} block (the \sphinxcode{\sphinxupquote{eng\textgreater{}}} lines) to the input, and optimize
again. See {\hyperref[\detokenize{reference/Energy:group-energy}]{\sphinxcrossref{\DUrole{std}{\DUrole{std-ref}{\$Energy}}}}}. 

\item {} 
\sphinxAtStartPar
\sphinxstylestrong{Use Cartesian coordinates.} {\hyperref[\detokenize{reference/optimize:optimize-cartesian}]{\sphinxcrossref{\DUrole{std}{\DUrole{std-ref}{cartesian=1}}}}} is described in the documentation as a
last resort, but users find it more robust for automated runs and for floppy systems such as explicit solvation
shells. 

\end{itemize}

\sphinxAtStartPar
Users also say that Priroda’s optimizer is weaker than Gaussian’s, while its SCF is stronger.




\subsection{Crash with \sphinxstyleliteralintitle{\sphinxupquote{process 0 got signal 6}} and \sphinxstyleliteralintitle{\sphinxupquote{too few}}}
\label{\detokenize{troubleshooting/common-problems:crash-with-process-0-got-signal-6-and-too-few}}
\sphinxAtStartPar
\sphinxstylestrong{Symptom.} The optimization dies with \sphinxcode{\sphinxupquote{process 0 got signal 6}}. In the output of the last step, the
\sphinxcode{\sphinxupquote{INTERNAL COORDINATES}} table ends with \sphinxcode{\sphinxupquote{too few}}.

\sphinxAtStartPar
\sphinxstylestrong{Fix.} Optimize in Cartesian coordinates:

\begin{sphinxVerbatim}[commandchars=\\\{\}]
\PYGZdl{}optimize cartesian=1 \PYGZdl{}end
\end{sphinxVerbatim}

\sphinxAtStartPar
This was the answer from one experienced user; the asker did not report back.




\subsection{The optimization “converges” in one step}
\label{\detokenize{troubleshooting/common-problems:the-optimization-converges-in-one-step}}
\sphinxAtStartPar
A run that reports convergence after a single step, with an energy that is not a number, did not optimize anything:
the SCF failed on the first step. See {\hyperref[\detokenize{troubleshooting/common-problems:troubleshooting-nan}]{\sphinxcrossref{\DUrole{std}{\DUrole{std-ref}{NaN energy, or an optimization that stops after one step}}}}}.




\subsection{A symmetric structure stays symmetric}
\label{\detokenize{troubleshooting/common-problems:a-symmetric-structure-stays-symmetric}}
\sphinxAtStartPar
Without the symmetry options, Priroda works in C1, but a symmetric starting geometry often stays symmetric through
an optimization that converges normally, and the Hessian then shows a saddle point. If the optimization did not break the symmetry, always
compute the Hessian. If it has an imaginary frequency, paste its \sphinxcode{\sphinxupquote{\$Energy}} block into the input, displace the
geometry along that mode if needed, and optimize again.



\sphinxAtStartPar
Version 25 can also use symmetry on purpose: \sphinxcode{\sphinxupquote{\$molecule sym=1,0.01}} detects the point group and symmetrizes the
input geometry, and \sphinxcode{\sphinxupquote{\$optimize sym=1}} keeps that symmetry during the optimization. Both options are needed. See
{\hyperref[\detokenize{reference/molecule:molecule-sym}]{\sphinxcrossref{\DUrole{std}{\DUrole{std-ref}{\$molecule sym}}}}} and {\hyperref[\detokenize{reference/optimize:optimize-symmetry}]{\sphinxcrossref{\DUrole{std}{\DUrole{std-ref}{\$optimize symmetry}}}}}.




\subsection{A scan starts from a mangled geometry}
\label{\detokenize{troubleshooting/common-problems:a-scan-starts-from-a-mangled-geometry}}
\sphinxAtStartPar
\sphinxstylestrong{Cause.} Priroda’s dihedral angles have the opposite sign to the convention used by Molden, MOPAC, GAMESS and most
other programs. A dihedral {\hyperref[\detokenize{reference/optimize:optimize-value}]{\sphinxcrossref{\DUrole{std}{\DUrole{std-ref}{value}}}}} copied from another program makes the first scan point a
scrambled structure that has nothing to do with the input geometry.

\sphinxAtStartPar
\sphinxstylestrong{Fix.} Flip the sign of every dihedral value taken from another program. The user who hit this confirmed that
flipping the signs fixed it, and the sign difference is reported independently on the other forum.



\sphinxAtStartPar
A scan that requested 10 points but computed only 7 (and 11 of 20) was reported once and never explained.




\subsection{IRC gives only two points, or no curve}
\label{\detokenize{troubleshooting/common-problems:irc-gives-only-two-points-or-no-curve}}
\sphinxAtStartPar
An IRC must start from a properly optimized transition state with an accurate Hessian in the input. Find the TS
first (\sphinxcode{\sphinxupquote{task=optimize}} with {\hyperref[\detokenize{reference/optimize:optimize-saddle}]{\sphinxcrossref{\DUrole{std}{\DUrole{std-ref}{saddle=1}}}}}), compute the Hessian there, check that it has exactly
one imaginary frequency and that this mode is the reaction you want, and only then run \sphinxcode{\sphinxupquote{task=irc}}. An IRC started
from an unoptimized geometry gives two points with {\hyperref[\detokenize{reference/optimize:optimize-back}]{\sphinxcrossref{\DUrole{std}{\DUrole{std-ref}{back=0}}}}} and no curve with \sphinxcode{\sphinxupquote{back=1}}.




\subsection{IRC stops with \sphinxstyleliteralintitle{\sphinxupquote{Vicinity of a Stationary Point(?) \sphinxhyphen{}\sphinxhyphen{}\sphinxhyphen{} STOP.}}}
\label{\detokenize{troubleshooting/common-problems:irc-stops-with-vicinity-of-a-stationary-point-stop}}
\sphinxAtStartPar
This message is also the normal end of an IRC that has reached a minimum: a version 25 test IRC of the ammonia
inversion ended with it after 7 points, at the minimum. It is a problem only when it comes early.



\sphinxAtStartPar
On the first IRC step this means that the energy did not go down on that step. The reported fix is to change the
step size; see {\hyperref[\detokenize{reference/optimize:optimize-radius}]{\sphinxcrossref{\DUrole{std}{\DUrole{std-ref}{radius}}}}}:
\begin{itemize}
\item {} 
\sphinxAtStartPar
on a very flat barrier top, increase the step so that it reaches the slope;

\item {} 
\sphinxAtStartPar
on a very narrow barrier, such as a proton transfer, decrease it, because a step suited to heavy atoms throws the
proton too far.

\end{itemize}

\sphinxAtStartPar
This advice comes from one experienced user.  


\subsection{IRC ends near a minimum but not in it}
\label{\detokenize{troubleshooting/common-problems:irc-ends-near-a-minimum-but-not-in-it}}
\sphinxAtStartPar
This is normal: always re\sphinxhyphen{}optimize both IRC endpoints. An IRC can also stop on a flat shelf where the gradient is
small, or at another TS. Then compute a Hessian at the stopping point and run a new IRC from there; this sometimes
has to be repeated three or four times. Priroda has no option to recompute the Hessian every n steps.




\section{Wrong or suspicious results}
\label{\detokenize{troubleshooting/common-problems:wrong-or-suspicious-results}}\label{\detokenize{troubleshooting/common-problems:troubleshooting-results}}

\subsection{Hybrid functional silently ignored}
\label{\detokenize{troubleshooting/common-problems:hybrid-functional-silently-ignored}}\label{\detokenize{troubleshooting/common-problems:troubleshooting-hybrid}}
\sphinxAtStartPar
\sphinxstylestrong{Cause.} Hybrid functionals (B3LYP, PBE0/PBE1, …) are applied only with {\hyperref[\detokenize{reference/control:control-theory}]{\sphinxcrossref{\DUrole{std}{\DUrole{std-ref}{theory=ridft}}}}}.
With \sphinxcode{\sphinxupquote{theory=dft}} or \sphinxcode{\sphinxupquote{theory=tddft}}, older versions (reported for v6 to v8) ran without any error message but used
zero exact exchange, so the result was a pure\sphinxhyphen{}GGA\sphinxhyphen{}like calculation labelled as a hybrid. The developer was told about
it and, according to a user who corresponded with him, did not intend to change it.



\sphinxAtStartPar
Newer builds (reported in 2014) refuse with \sphinxcode{\sphinxupquote{theory=dft}}: the job prints \sphinxcode{\sphinxupquote{Basis set input:}} twice and stops with
\sphinxcode{\sphinxupquote{application called MPI\_Abort(MPI\_COMM\_WORLD, 1) \sphinxhyphen{} process 0}} (on stderr), before the summary box. Several users
reproduced this, and version 25 does the same.



\begin{sphinxadmonition}{warning}{Warning:}
\sphinxAtStartPar
\sphinxstylestrong{With \sphinxcode{\sphinxupquote{theory=tddft}} the trap is still open in version 25.} A water run with \sphinxcode{\sphinxupquote{functional=PBE1}} showed \sphinxcode{\sphinxupquote{PBE1}} in the
summary box, finished normally and printed excited states, but the exchange energy was −6.61 instead of about
−8.95 hartree and the total energy was about 2.2 hartree too high: the exact\sphinxhyphen{}exchange part was simply left out. Any
TDDFT result with a hybrid functional is wrong.
\end{sphinxadmonition}



\sphinxAtStartPar
\sphinxstylestrong{Fix.} Use \sphinxcode{\sphinxupquote{theory=ridft}} for every hybrid ground\sphinxhyphen{}state calculation, and a pure GGA functional (PBE, the default)
for TDDFT. There is no hybrid TDDFT in Priroda.




\subsection{NaN energy, or an optimization that stops after one step}
\label{\detokenize{troubleshooting/common-problems:nan-energy-or-an-optimization-that-stops-after-one-step}}\label{\detokenize{troubleshooting/common-problems:troubleshooting-nan}}
\sphinxAtStartPar
\sphinxstylestrong{Cause.} The SCF energy on the first step is not a number, so the optimizer stops at once and reports convergence.
The reported causes are all basis problems:
\begin{itemize}
\item {} 
\sphinxAtStartPar
\sphinxstylestrong{No basis for an element.} A molecule containing Bi was run with \sphinxcode{\sphinxupquote{3z.bas}}, which covers only Z ≤ 54. Use
\sphinxcode{\sphinxupquote{sbk.bas}} (ECPs on every element) or \sphinxcode{\sphinxupquote{basis4.in}} with {\hyperref[\detokenize{reference/control:control-four}]{\sphinxcrossref{\DUrole{std}{\DUrole{std-ref}{four=1}}}}}.



\item {} 
\sphinxAtStartPar
\sphinxstylestrong{ECP not read.} In one case nothing was computed and there was no error message. The output showed zeros for one
element in the basis table, and no sign that the ECP had been read. The cause was a defect in the user’s basis file;
\sphinxcode{\sphinxupquote{20+sbk.bas}} worked. 

\item {} 
\sphinxAtStartPar
A NaN at the first SCF iteration for {[}PtCl6{]}2\sphinxhyphen{} with the St\sphinxhyphen{}rsc\sphinxhyphen{}97 ECP basis went away after switching to the
\sphinxcode{\sphinxupquote{sbk(u)}} basis. Why that basis failed is not known. 

\end{itemize}

\sphinxAtStartPar
\sphinxstylestrong{What to check.} In the output, the ECP should be read before the basis, and every atom in the basis table should
have a non\sphinxhyphen{}zero number of basis functions.




\subsection{TDDFT oscillator strengths are all zero}
\label{\detokenize{troubleshooting/common-problems:tddft-oscillator-strengths-are-all-zero}}
\sphinxAtStartPar
A closed\sphinxhyphen{}shell TDDFT run on version 25 printed non\sphinxhyphen{}zero oscillator strengths in its \sphinxcode{\sphinxupquote{Electronic Spectrum:}} table
with the default settings. 

\sphinxAtStartPar
The reported cause of all\sphinxhyphen{}zero values: with \sphinxcode{\sphinxupquote{restrict=0}}, singlet and triplet excitations get mixed. Use \sphinxcode{\sphinxupquote{restrict=1}} and compute one multiplicity per run.
The user who reported the problem confirmed the fix.




\subsection{The dipole moment is zero}
\label{\detokenize{troubleshooting/common-problems:the-dipole-moment-is-zero}}
\sphinxAtStartPar
Version 6 printed \sphinxcode{\sphinxupquote{dipole =   0.000000 au (  0.000000 D)}} for every molecule when relativistic basis sets were used,
even with \sphinxcode{\sphinxupquote{task=dipole}}. Version 8 and later print a value.




\subsection{Energies differ between Priroda versions}
\label{\detokenize{troubleshooting/common-problems:energies-differ-between-priroda-versions}}
\sphinxAtStartPar
Newer builds are faster but give slightly different total energies for the same input. Do not compare close
energies computed with different versions; within one version the comparison is fine.




\subsection{The printed energy components do not add up to E}
\label{\detokenize{troubleshooting/common-problems:the-printed-energy-components-do-not-add-up-to-e}}
\sphinxAtStartPar
The line \sphinxcode{\sphinxupquote{T = ... V = ... J = ... Ex = ... Ec = ...}} does not include the nuclear repulsion energy, so the
components do not sum to \sphinxcode{\sphinxupquote{E}}. This is not an error: E = T + V + J + Ex + Ec + E(nn). The answer came from the developer,
relayed by a user.




\subsection{Thermochemistry of single atoms}
\label{\detokenize{troubleshooting/common-problems:thermochemistry-of-single-atoms}}
\sphinxAtStartPar
For a monoatomic species the enthalpy has an extra RT, and the rotational entropy is printed as \sphinxcode{\sphinxupquote{\sphinxhyphen{}inf}}. Correct these
by hand and keep only the translational terms.




\subsection{Results depend on the number of processes}
\label{\detokenize{troubleshooting/common-problems:results-depend-on-the-number-of-processes}}
\sphinxAtStartPar
In one unstable TS search started from a poor guess, the optimization followed different paths with different
\sphinxcode{\sphinxupquote{\sphinxhyphen{}np}}. Normally the results differ only in insignificant digits.




\section{Memory, disk and scratch errors}
\label{\detokenize{troubleshooting/common-problems:memory-disk-and-scratch-errors}}\label{\detokenize{troubleshooting/common-problems:troubleshooting-memory}}
\sphinxAtStartPar
All {\hyperref[\detokenize{reference/system:group-system}]{\sphinxcrossref{\DUrole{std}{\DUrole{std-ref}{\$system}}}}} limits apply per MPI process. With \sphinxcode{\sphinxupquote{disk=10}} on four processes, up to 40 GB can be
used. 


\subsection{\sphinxstyleliteralintitle{\sphinxupquote{need … bytes.}}}
\label{\detokenize{troubleshooting/common-problems:need-bytes}}
\sphinxAtStartPar
\sphinxstylestrong{Cause.} The {\hyperref[\detokenize{reference/system:system-memory}]{\sphinxcrossref{\DUrole{std}{\DUrole{std-ref}{memory}}}}} limit in \sphinxcode{\sphinxupquote{\$system}} is too small for the job. Version 25 stops right after
the basis\sphinxhyphen{}set echo with a line such as \sphinxcode{\sphinxupquote{need 12544 + 7905408 = 7917952 bytes.}} (here for water with \sphinxcode{\sphinxupquote{mem=2}}); the
last number is what the job needs at that point. Raise \sphinxcode{\sphinxupquote{memory}} above it.




\subsection{\sphinxstyleliteralintitle{\sphinxupquote{malloc(...) failed}}}
\label{\detokenize{troubleshooting/common-problems:malloc-failed}}
\sphinxAtStartPar
\sphinxstylestrong{Cause.} Priroda could not allocate the memory requested in {\hyperref[\detokenize{reference/system:system-memory}]{\sphinxcrossref{\DUrole{std}{\DUrole{std-ref}{memory}}}}}, plus any RAM disk
requested with a negative {\hyperref[\detokenize{reference/system:system-disk}]{\sphinxcrossref{\DUrole{std}{\DUrole{std-ref}{disk}}}}}. The number in the message is the request in bytes.



\sphinxAtStartPar
\sphinxstylestrong{Fix.} Ask for less than the node really has free:
\begin{itemize}
\item {} 
\sphinxAtStartPar
With a RAM disk, \sphinxcode{\sphinxupquote{memory}} plus |\sphinxcode{\sphinxupquote{disk}}| (both in MB), times the number of processes, must fit in physical RAM.
An example from a 24 GB node is \sphinxcode{\sphinxupquote{memory=9000 disk=\sphinxhyphen{}14000}}. Several users confirmed this rule.



\item {} 
\sphinxAtStartPar
The 32\sphinxhyphen{}bit Windows (Cygwin) builds cannot allocate more than about 1 GB. 512 MB was enough for most DFT jobs.
Several users confirmed this limit. The 64\sphinxhyphen{}bit Windows build with \sphinxcode{\sphinxupquote{cygwin1.dll}} is reported to have no such limit.



\end{itemize}

\sphinxAtStartPar
A message beginning \sphinxcode{\sphinxupquote{We need 1.5 GB of memory...}} was reported once. One reply guessed that it is about RAM
(\sphinxcode{\sphinxupquote{memory=}}), not disk space. It may be an older form of the \sphinxcode{\sphinxupquote{need … bytes.}} message above.




\subsection{\sphinxstyleliteralintitle{\sphinxupquote{out of disk space}}}
\label{\detokenize{troubleshooting/common-problems:out-of-disk-space}}
\sphinxAtStartPar
\sphinxstylestrong{Symptom.} \sphinxcode{\sphinxupquote{0: out of disk space (VNNNN), \textless{}N\textgreater{}KB}}, or a similar message naming another array (version 25:
\sphinxcode{\sphinxupquote{0: out of disk space (grid\_W), 8213824}} for a 1 MB RAM disk), such as
\sphinxcode{\sphinxupquote{(ci.Rnvvo)}}.

\sphinxAtStartPar
\sphinxstylestrong{Fix.}
\begin{itemize}
\item {} 
\sphinxAtStartPar
Raise {\hyperref[\detokenize{reference/system:system-disk}]{\sphinxcrossref{\DUrole{std}{\DUrole{std-ref}{disk}}}}} above the amount reported. Remember that the value is per process. One user’s
correlated Hessian that ran out of disk on 16 processes then ran with \sphinxcode{\sphinxupquote{disk=40}} on one fast drive.



\item {} 
\sphinxAtStartPar
Conventional MP2 is very disk\sphinxhyphen{}hungry. Use \sphinxcode{\sphinxupquote{theory=rimp2}} instead.



\item {} 
\sphinxAtStartPar
Check that \sphinxcode{\sphinxupquote{cartesian}} or \sphinxcode{\sphinxupquote{z\sphinxhyphen{}matrix}} in \sphinxcode{\sphinxupquote{\$molecule}} matches the coordinates you actually gave; wrong coordinates
can make the job far bigger than intended. 

\end{itemize}

\sphinxAtStartPar
Each scratch file is created for 1 GB per process and filled as needed. Jobs needing more than 20\textendash{}45 GB of scratch
were, in one experienced user’s words, not worth their cost.



\sphinxAtStartPar
A failure to read a scratch file during MP2 (\sphinxcode{\sphinxupquote{0: read file ./tmp/...}}) was given the same fix.




\subsection{A small RI\sphinxhyphen{}MP2 job does not fit in memory}
\label{\detokenize{troubleshooting/common-problems:a-small-ri-mp2-job-does-not-fit-in-memory}}
\sphinxAtStartPar
A 7\sphinxhyphen{}atom RI\sphinxhyphen{}MP2 job with \sphinxcode{\sphinxupquote{four=1}}, \sphinxcode{\sphinxupquote{basis4.in}} and \sphinxcode{\sphinxupquote{set=L22}} claimed not to fit in memory however much was given,
while smaller and larger basis sets ran. A second user reproduced it and got \sphinxcode{\sphinxupquote{0: write VMNN(11753704:14042536)?}}.
Both found that direct integrals avoid the problem:

\begin{sphinxVerbatim}[commandchars=\\\{\}]
\PYGZdl{}integral direct=1 \PYGZdl{}end
\end{sphinxVerbatim}

\sphinxAtStartPar
See {\hyperref[\detokenize{reference/integral:integral-direct}]{\sphinxcrossref{\DUrole{std}{\DUrole{std-ref}{direct}}}}}. 


\subsection{Where scratch files go}
\label{\detokenize{troubleshooting/common-problems:where-scratch-files-go}}
\sphinxAtStartPar
Priroda writes scratch files to the directory given by {\hyperref[\detokenize{reference/system:system-path}]{\sphinxcrossref{\DUrole{std}{\DUrole{std-ref}{path}}}}}; if \sphinxcode{\sphinxupquote{path}} is not set, to
\sphinxcode{\sphinxupquote{\$TMPDIR}}; otherwise, in version 25, to \sphinxcode{\sphinxupquote{/tmp}} (a 2020 user report said the directory of the input file). On a
cluster, \sphinxcode{\sphinxupquote{/tmp}} is often too small or quota\sphinxhyphen{}limited for large scratch files. Use the cluster’s local scratch area, or
a RAM disk with a negative \sphinxcode{\sphinxupquote{disk}} value. Never put scratch on NFS.



\sphinxAtStartPar
A job that stops while writing a scratch file, with a message like
\sphinxcode{\sphinxupquote{0: write file /tmp/tmp9633.11 (leAX): fd=21 ...}}, was reported once during an analytic Hessian and not answered.
A full scratch directory or a file\sphinxhyphen{}size limit is a plausible cause, but this is a guess; check the free space and the
\sphinxcode{\sphinxupquote{disk=}} setting. 

\sphinxAtStartPar
If Priroda does not get enough memory or scratch, it stops with an error; it does not fall back to a slower
algorithm. 


\subsection{Slow runs: disk speed}
\label{\detokenize{troubleshooting/common-problems:slow-runs-disk-speed}}
\sphinxAtStartPar
Priroda is very demanding of disk speed. For all methods except GGA DFT, the documentation warns that performance
degrades “terribly” if more than one process reads and writes large files on the same physical drive. Options, in order of
preference:
\begin{itemize}
\item {} 
\sphinxAtStartPar
keep scratch in RAM with a negative {\hyperref[\detokenize{reference/system:system-disk}]{\sphinxcrossref{\DUrole{std}{\DUrole{std-ref}{disk}}}}} value (one cluster reported a 6× speed\sphinxhyphen{}up for
RI\sphinxhyphen{}MP2);

\item {} 
\sphinxAtStartPar
give each process its own physical drive with a colon\sphinxhyphen{}separated {\hyperref[\detokenize{reference/system:system-path}]{\sphinxcrossref{\DUrole{std}{\DUrole{std-ref}{path}}}}} list;

\item {} 
\sphinxAtStartPar
if the job is disk\sphinxhyphen{}bound, switch to direct integrals with {\hyperref[\detokenize{reference/integral:integral-direct}]{\sphinxcrossref{\DUrole{std}{\DUrole{std-ref}{direct=1}}}}}. This costs CPU time but
can be much faster overall: one RI\sphinxhyphen{}MP2 optimization went from 2 to 20 steps per night.

\end{itemize}



\sphinxAtStartPar
One user with 16 processes ignored the warning and put all scratch on one fast disk.




\section{Parallel and MPI problems}
\label{\detokenize{troubleshooting/common-problems:parallel-and-mpi-problems}}\label{\detokenize{troubleshooting/common-problems:troubleshooting-parallel}}

\subsection{A parallel RI job hangs with no message}
\label{\detokenize{troubleshooting/common-problems:a-parallel-ri-job-hangs-with-no-message}}
\sphinxAtStartPar
\sphinxstylestrong{Cause.} RI methods, RI\sphinxhyphen{}MP2 and RI\sphinxhyphen{}DFT included, hang forever if any process gets no work. This happens when the job
is too small for the number of processes. Two experienced users confirmed it and called it a bug.

\sphinxAtStartPar
\sphinxstylestrong{How to recognize it.} After \sphinxcode{\sphinxupquote{integral estimates}}, the output lists the blocks per node. Look for a line with
\sphinxcode{\sphinxupquote{0 blocks}}:

\begin{sphinxVerbatim}[commandchars=\\\{\}]
    2892490 blocks,    103625770 integrals on node  0
    3681560 blocks,     91256706 integrals on node  1
\end{sphinxVerbatim}

\sphinxAtStartPar
\sphinxstylestrong{Fix.} Run small RI jobs on fewer processes, or on one core.



\sphinxAtStartPar
Older reports describe similar hangs: small jobs on many processors hang (v6/v7); MP2 on a 102\sphinxhyphen{}atom complex died with
signal 8 on more than about 24 cores (v7, MVAPICH); a v6 Hessian on 4 cores could hang at \sphinxcode{\sphinxupquote{Solving for responses ...}}.
The v6 advice was at most 2 cores for a Hessian and at most 4 for an optimization. Another user ran DFT, including
Hessians, on 24\textendash{}32 cores without trouble, so this advice is disputed. Parallel runs work properly only from version
16 on.




\subsection{Little or no speed\sphinxhyphen{}up from more cores}
\label{\detokenize{troubleshooting/common-problems:little-or-no-speed-up-from-more-cores}}\begin{itemize}
\item {} 
\sphinxAtStartPar
DFT scales well within one node but gains little beyond 6\textendash{}8 cores: one user measured at most 10\% more from 32
cores than from 6\textendash{}8. Another measured 101, 63 and 37 s for the same optimization on 8, 12 and 24 cores.



\item {} 
\sphinxAtStartPar
Running across several nodes gives no gain. Several experienced users report this.



\item {} 
\sphinxAtStartPar
In a DFT run most of the memory and all the disk I/O stay on rank 0, so spreading a DFT job over nodes does not
spread its memory. RI\sphinxhyphen{}MP2 distributes memory better.



\item {} 
\sphinxAtStartPar
For throughput, several single\sphinxhyphen{}core jobs at once can beat one multi\sphinxhyphen{}core job when the jobs are not disk\sphinxhyphen{}bound.



\item {} 
\sphinxAtStartPar
For RI\sphinxhyphen{}MP2, the v5 documentation suggests a number of processes that divides the number of occupied orbitals.



\end{itemize}


\subsection{Starting a parallel run on Linux}
\label{\detokenize{troubleshooting/common-problems:starting-a-parallel-run-on-linux}}
\sphinxAtStartPar
Version 25 must be started through an MPI launcher, such as the \sphinxcode{\sphinxupquote{mpiexec}} from the official site:
\sphinxcode{\sphinxupquote{./mpiexec \sphinxhyphen{}np 4 ./p job.in job.out}}. \sphinxcode{\sphinxupquote{./p \sphinxhyphen{}np 4 job.in job.out}}, which worked in versions 6 and 10, now stops with
\sphinxcode{\sphinxupquote{input file \textquotesingle{}\sphinxhyphen{}np\textquotesingle{}?.}}. With version 10, starting through the bundled \sphinxcode{\sphinxupquote{mpiexec}} once failed with
\sphinxcode{\sphinxupquote{Error in system call mpiexec: No such file or directory}} or \sphinxcode{\sphinxupquote{mpiexec could not exec p}}. Use absolute paths for the
binary, the input, the output and the basis file.



\sphinxAtStartPar
The message \sphinxcode{\sphinxupquote{OOPS: semop lock failed}} (with \sphinxcode{\sphinxupquote{process N received signal 2}} or \sphinxcode{\sphinxupquote{6}}) has two reported contexts:
\begin{itemize}
\item {} 
\sphinxAtStartPar
as a secondary message after another fatal error, such as running out of disk. Fix the first error;



\item {} 
\sphinxAtStartPar
in v6 shared\sphinxhyphen{}memory runs on 7 cores. Fewer cores avoided it; cleaning semaphores with \sphinxcode{\sphinxupquote{ipcrm}} did not help.



\end{itemize}

\sphinxAtStartPar
It also appeared for one user who ran as a normal user while the same input worked as root. That case was never
resolved. 

\sphinxAtStartPar
A run that stops producing output right after the header, while the process keeps using CPU, has been reported and
called common. The only advice was to look for stray files named after the input or output in the input, output,
program and home directories, and then to rule out causes one by one. 


\subsection{Windows}
\label{\detokenize{troubleshooting/common-problems:windows}}
\sphinxAtStartPar
Old Windows builds were serial only. For a parallel 64\sphinxhyphen{}bit run (version 16 or later), users report this setup:


\begin{enumerate}
\sphinxsetlistlabels{\arabic}{enumi}{enumii}{}{.}%
\item {} 
\sphinxAtStartPar
From a 64\sphinxhyphen{}bit Cygwin installation with OpenMPI, copy \sphinxcode{\sphinxupquote{cygwin1.dll}}, \sphinxcode{\sphinxupquote{cyggcc\_s\sphinxhyphen{}seh\sphinxhyphen{}1.dll}}, \sphinxcode{\sphinxupquote{cygmpi\sphinxhyphen{}12.dll}},
\sphinxcode{\sphinxupquote{cygopen\sphinxhyphen{}pal\sphinxhyphen{}13.dll}}, \sphinxcode{\sphinxupquote{cygopen\sphinxhyphen{}rte\sphinxhyphen{}12.dll}} and \sphinxcode{\sphinxupquote{orterun.exe}} next to the Priroda executable. Take
\sphinxcode{\sphinxupquote{cygwin1.dll}} from the same Cygwin installation, because the other DLLs work only with it. These are the names
from a 2017 Cygwin. A 2018 report says newer OpenMPI packages ship \sphinxcode{\sphinxupquote{cygopen\sphinxhyphen{}pal\sphinxhyphen{}40.dll}} and
\sphinxcode{\sphinxupquote{cygopen\sphinxhyphen{}rte\sphinxhyphen{}40.dll}} instead; take those.

\item {} 
\sphinxAtStartPar
To run outside the Cygwin terminal, create a \sphinxcode{\sphinxupquote{dev}} directory one level above the directory you start Priroda from,
with subdirectories \sphinxcode{\sphinxupquote{shm}} and \sphinxcode{\sphinxupquote{mqueue}}. For Priroda in \sphinxcode{\sphinxupquote{C:\textbackslash{}Priroda 16}} that means \sphinxcode{\sphinxupquote{C:\textbackslash{}dev\textbackslash{}shm}} and
\sphinxcode{\sphinxupquote{C:\textbackslash{}dev\textbackslash{}mqueue}}. Without them Cygwin prints \sphinxcode{\sphinxupquote{Warning: \textquotesingle{}/dev/shm\textquotesingle{} does not exists or is not a directory.}}

\item {} 
\sphinxAtStartPar
Set \sphinxcode{\sphinxupquote{TMPDIR}} and start the run from \sphinxcode{\sphinxupquote{cmd.exe}}:

\begin{sphinxVerbatim}[commandchars=\\\{\}]
cd /d \PYGZdq{}C:\PYGZbs{}Priroda 16\PYGZdq{}
set TMPDIR=\PYGZdq{}\PYGZpc{}TEMP\PYGZpc{}\PYGZdq{}
\PYGZdq{}orterun.exe\PYGZdq{} \PYGZhy{}np 2 ./p16\PYGZus{}64\PYGZus{}parallel.exe \PYGZdq{}C:\PYGZbs{}Work\PYGZbs{}file.in\PYGZdq{} \PYGZdq{}C:\PYGZbs{}Work\PYGZbs{}file.out\PYGZdq{}
\end{sphinxVerbatim}

\item {} 
\sphinxAtStartPar
If the program is in a system directory such as \sphinxcode{\sphinxupquote{C:\textbackslash{}Program Files}}, run it as administrator.

\end{enumerate}



\sphinxAtStartPar
Network problems reported with this setup:
\begin{itemize}
\item {} 
\sphinxAtStartPar
On a PC connected by WiFi that also has a LAN adapter, OpenMPI failed with
\sphinxcode{\sphinxupquote{tcp\_peer\_send\_blocking: send() to socket 12 failed}}. Disabling the LAN adapter fixed it (single report).



\item {} 
\sphinxAtStartPar
With the Ethernet cable unplugged, only some cores were used. Plugging it back in fixed it (single report).



\end{itemize}

\sphinxAtStartPar
A 2021 attempt to follow this recipe on Windows 10 failed with an error that was posted only as a screenshot; no fix
was posted. One user runs Priroda under WSL on Windows instead; no details of that setup were posted.



\begin{sphinxadmonition}{warning}{Warning:}
\sphinxAtStartPar
Not every build on the official site is reliable. A note there once said that the \sphinxcode{\sphinxupquote{w64/*.exe}} 64\sphinxhyphen{}bit Cygwin builds
had a DFT/ECP bug. If a Windows build misbehaves, ask the developer which build you should be running.
\end{sphinxadmonition}



\sphinxstepscope


\chapter{Keywords in the program}
\label{\detokenize{binary-keywords:keywords-in-the-program}}\label{\detokenize{binary-keywords:binary-keywords}}\label{\detokenize{binary-keywords::doc}}
\sphinxAtStartPar
This page lists every keyword name that the Priroda 25 (2025.02.28) executable reads, found by examining the
program file itself. Names that this manual documents are links into the {\hyperref[\detokenize{reference/index::doc}]{\sphinxcrossref{\DUrole{doc}{input reference}}}};
the others are only known to exist. Their meaning is unknown, so do not use them in production work.



\sphinxAtStartPar
A name listed here is accepted by the program; a name that is not listed stops the run with
\sphinxcode{\sphinxupquote{*** error in \textquotesingle{}\$group\textquotesingle{} : unknown name …}}, as long as the task reads that group. Any unique prefix of a name is
accepted, so \sphinxcode{\sphinxupquote{sym=1}} in \sphinxcode{\sphinxupquote{\$optimize}} means \sphinxcode{\sphinxupquote{symmetry=1}} and \sphinxcode{\sphinxupquote{dist=}} in \sphinxcode{\sphinxupquote{\$uniq}} means \sphinxcode{\sphinxupquote{distance=}}.




\section{Groups for everyday work}
\label{\detokenize{binary-keywords:groups-for-everyday-work}}

\begin{savenotes}\sphinxattablestart
\sphinxthistablewithglobalstyle
\centering
\begin{tabulary}{\linewidth}[t]{TT}
\sphinxtoprule
\sphinxstyletheadfamily 
\sphinxAtStartPar
Group
&\sphinxstyletheadfamily 
\sphinxAtStartPar
Keywords
\\
\sphinxmidrule
\sphinxtableatstartofbodyhook
\sphinxAtStartPar
{\hyperref[\detokenize{reference/system:group-system}]{\sphinxcrossref{\DUrole{std}{\DUrole{std-ref}{\$system}}}}}
&
\sphinxAtStartPar
{\hyperref[\detokenize{reference/system:system-memory}]{\sphinxcrossref{\DUrole{std}{\DUrole{std-ref}{memory}}}}}, {\hyperref[\detokenize{reference/system:system-path}]{\sphinxcrossref{\DUrole{std}{\DUrole{std-ref}{path}}}}}, {\hyperref[\detokenize{reference/system:system-disk}]{\sphinxcrossref{\DUrole{std}{\DUrole{std-ref}{disk}}}}}
\\
\sphinxhline
\sphinxAtStartPar
{\hyperref[\detokenize{reference/control:group-control}]{\sphinxcrossref{\DUrole{std}{\DUrole{std-ref}{\$control}}}}}
&
\sphinxAtStartPar
{\hyperref[\detokenize{reference/control:control-task}]{\sphinxcrossref{\DUrole{std}{\DUrole{std-ref}{task}}}}}, {\hyperref[\detokenize{reference/control:control-theory}]{\sphinxcrossref{\DUrole{std}{\DUrole{std-ref}{theory}}}}}, {\hyperref[\detokenize{reference/control:control-four}]{\sphinxcrossref{\DUrole{std}{\DUrole{std-ref}{four}}}}}, \sphinxcode{\sphinxupquote{a137}}, {\hyperref[\detokenize{reference/control:control-state}]{\sphinxcrossref{\DUrole{std}{\DUrole{std-ref}{state}}}}}, \sphinxcode{\sphinxupquote{weight}}, {\hyperref[\detokenize{reference/control:control-basis}]{\sphinxcrossref{\DUrole{std}{\DUrole{std-ref}{basis}}}}}, {\hyperref[\detokenize{reference/control:control-nucleus}]{\sphinxcrossref{\DUrole{std}{\DUrole{std-ref}{nucleus}}}}}, \sphinxcode{\sphinxupquote{field}}, \sphinxcode{\sphinxupquote{derivative}}, \sphinxcode{\sphinxupquote{fit}}, \sphinxcode{\sphinxupquote{wavelength}}, {\hyperref[\detokenize{reference/control:control-save}]{\sphinxcrossref{\DUrole{std}{\DUrole{std-ref}{save}}}}}, {\hyperref[\detokenize{reference/control:control-print}]{\sphinxcrossref{\DUrole{std}{\DUrole{std-ref}{print}}}}}
\\
\sphinxhline
\sphinxAtStartPar
{\hyperref[\detokenize{reference/scf:group-scf}]{\sphinxcrossref{\DUrole{std}{\DUrole{std-ref}{\$scf}}}}}
&
\sphinxAtStartPar
{\hyperref[\detokenize{reference/scf:scf-restrict}]{\sphinxcrossref{\DUrole{std}{\DUrole{std-ref}{restrict}}}}}, {\hyperref[\detokenize{reference/scf:scf-convergence}]{\sphinxcrossref{\DUrole{std}{\DUrole{std-ref}{convergence}}}}}, {\hyperref[\detokenize{reference/scf:scf-iterations}]{\sphinxcrossref{\DUrole{std}{\DUrole{std-ref}{iterations}}}}}, {\hyperref[\detokenize{reference/scf:scf-procedure}]{\sphinxcrossref{\DUrole{std}{\DUrole{std-ref}{procedure}}}}}, {\hyperref[\detokenize{reference/scf:scf-d1small}]{\sphinxcrossref{\DUrole{std}{\DUrole{std-ref}{d1small}}}}}, {\hyperref[\detokenize{reference/scf:scf-core}]{\sphinxcrossref{\DUrole{std}{\DUrole{std-ref}{core}}}}}, \sphinxcode{\sphinxupquote{occupation}}, \sphinxcode{\sphinxupquote{fraction}}, \sphinxcode{\sphinxupquote{temperature}}, \sphinxcode{\sphinxupquote{ttype}}, \sphinxcode{\sphinxupquote{diagm}}, \sphinxcode{\sphinxupquote{xmax}}, \sphinxcode{\sphinxupquote{n1}}, \sphinxcode{\sphinxupquote{fitxc}}, \sphinxcode{\sphinxupquote{eemin}}, \sphinxcode{\sphinxupquote{mab}}, \sphinxcode{\sphinxupquote{aw}}, \sphinxcode{\sphinxupquote{wv}}
\\
\sphinxhline
\sphinxAtStartPar
{\hyperref[\detokenize{reference/dft:group-dft}]{\sphinxcrossref{\DUrole{std}{\DUrole{std-ref}{\$dft}}}}}
&
\sphinxAtStartPar
{\hyperref[\detokenize{reference/dft:dft-functional}]{\sphinxcrossref{\DUrole{std}{\DUrole{std-ref}{functional}}}}}, \sphinxcode{\sphinxupquote{zero}}, \sphinxcode{\sphinxupquote{n}}, \sphinxcode{\sphinxupquote{p}}
\\
\sphinxhline
\sphinxAtStartPar
{\hyperref[\detokenize{reference/grid:group-grid}]{\sphinxcrossref{\DUrole{std}{\DUrole{std-ref}{\$grid}}}}}
&
\sphinxAtStartPar
{\hyperref[\detokenize{reference/grid:grid-accuracy}]{\sphinxcrossref{\DUrole{std}{\DUrole{std-ref}{accuracy}}}}}, \sphinxcode{\sphinxupquote{type}}, \sphinxcode{\sphinxupquote{cutoff}}, \sphinxcode{\sphinxupquote{nradial}}, \sphinxcode{\sphinxupquote{nadd}}, \sphinxcode{\sphinxupquote{use43}}
\\
\sphinxhline
\sphinxAtStartPar
{\hyperref[\detokenize{reference/integral:group-integral}]{\sphinxcrossref{\DUrole{std}{\DUrole{std-ref}{\$integral}}}}}
&
\sphinxAtStartPar
{\hyperref[\detokenize{reference/integral:integral-direct}]{\sphinxcrossref{\DUrole{std}{\DUrole{std-ref}{direct}}}}}, {\hyperref[\detokenize{reference/integral:integral-routine}]{\sphinxcrossref{\DUrole{std}{\DUrole{std-ref}{routine}}}}}, \sphinxcode{\sphinxupquote{storage}}, \sphinxcode{\sphinxupquote{invert}}, \sphinxcode{\sphinxupquote{cutij}}, \sphinxcode{\sphinxupquote{vcut}}
\\
\sphinxhline
\sphinxAtStartPar
{\hyperref[\detokenize{reference/atoms:group-atoms}]{\sphinxcrossref{\DUrole{std}{\DUrole{std-ref}{\$atoms}}}}}
&
\sphinxAtStartPar
{\hyperref[\detokenize{reference/atoms:atoms-core}]{\sphinxcrossref{\DUrole{std}{\DUrole{std-ref}{core}}}}}, {\hyperref[\detokenize{reference/atoms:atoms-mcore}]{\sphinxcrossref{\DUrole{std}{\DUrole{std-ref}{mcore}}}}}, \sphinxcode{\sphinxupquote{nc}}, \sphinxcode{\sphinxupquote{no}}, \sphinxcode{\sphinxupquote{active}}, \sphinxcode{\sphinxupquote{mocc}}, \sphinxcode{\sphinxupquote{occupation}}, \sphinxcode{\sphinxupquote{diagonalize}}
\\
\sphinxhline
\sphinxAtStartPar
{\hyperref[\detokenize{reference/guess:group-guess}]{\sphinxcrossref{\DUrole{std}{\DUrole{std-ref}{\$guess}}}}}
&
\sphinxAtStartPar
{\hyperref[\detokenize{reference/guess:guess-read}]{\sphinxcrossref{\DUrole{std}{\DUrole{std-ref}{read}}}}}, {\hyperref[\detokenize{reference/guess:guess-mix}]{\sphinxcrossref{\DUrole{std}{\DUrole{std-ref}{mix}}}}}, {\hyperref[\detokenize{reference/guess:guess-na}]{\sphinxcrossref{\DUrole{std}{\DUrole{std-ref}{na}}}}}, {\hyperref[\detokenize{reference/guess:guess-nb}]{\sphinxcrossref{\DUrole{std}{\DUrole{std-ref}{nb}}}}}, \sphinxcode{\sphinxupquote{n}}, \sphinxcode{\sphinxupquote{a}}, \sphinxcode{\sphinxupquote{type}}, \sphinxcode{\sphinxupquote{gradient}}, \sphinxcode{\sphinxupquote{localize}}
\\
\sphinxhline
\sphinxAtStartPar
{\hyperref[\detokenize{reference/optimize:group-optimize}]{\sphinxcrossref{\DUrole{std}{\DUrole{std-ref}{\$optimize}}}}}
&
\sphinxAtStartPar
{\hyperref[\detokenize{reference/optimize:optimize-saddle}]{\sphinxcrossref{\DUrole{std}{\DUrole{std-ref}{saddle}}}}}, {\hyperref[\detokenize{reference/optimize:optimize-tolerance}]{\sphinxcrossref{\DUrole{std}{\DUrole{std-ref}{tolerance}}}}}, {\hyperref[\detokenize{reference/optimize:optimize-trust}]{\sphinxcrossref{\DUrole{std}{\DUrole{std-ref}{trust}}}}}, {\hyperref[\detokenize{reference/optimize:optimize-steps}]{\sphinxcrossref{\DUrole{std}{\DUrole{std-ref}{steps}}}}}, {\hyperref[\detokenize{reference/optimize:optimize-follow}]{\sphinxcrossref{\DUrole{std}{\DUrole{std-ref}{follow}}}}}, {\hyperref[\detokenize{reference/optimize:optimize-points}]{\sphinxcrossref{\DUrole{std}{\DUrole{std-ref}{points}}}}}, {\hyperref[\detokenize{reference/optimize:optimize-back}]{\sphinxcrossref{\DUrole{std}{\DUrole{std-ref}{back}}}}}, \sphinxcode{\sphinxupquote{down}}, {\hyperref[\detokenize{reference/optimize:optimize-radius}]{\sphinxcrossref{\DUrole{std}{\DUrole{std-ref}{radius}}}}}, {\hyperref[\detokenize{reference/optimize:optimize-fix}]{\sphinxcrossref{\DUrole{std}{\DUrole{std-ref}{fix}}}}}, {\hyperref[\detokenize{reference/optimize:optimize-value}]{\sphinxcrossref{\DUrole{std}{\DUrole{std-ref}{value}}}}}, {\hyperref[\detokenize{reference/optimize:optimize-coordinates}]{\sphinxcrossref{\DUrole{std}{\DUrole{std-ref}{coordinates}}}}}, {\hyperref[\detokenize{reference/optimize:optimize-cartesian}]{\sphinxcrossref{\DUrole{std}{\DUrole{std-ref}{cartesian}}}}}, {\hyperref[\detokenize{reference/optimize:optimize-symmetry}]{\sphinxcrossref{\DUrole{std}{\DUrole{std-ref}{symmetry}}}}}, {\hyperref[\detokenize{reference/optimize:optimize-gradient}]{\sphinxcrossref{\DUrole{std}{\DUrole{std-ref}{gradient}}}}}, {\hyperref[\detokenize{reference/optimize:optimize-print}]{\sphinxcrossref{\DUrole{std}{\DUrole{std-ref}{print}}}}}, \sphinxcode{\sphinxupquote{exact}}, \sphinxcode{\sphinxupquote{bend}}, \sphinxcode{\sphinxupquote{local}}, \sphinxcode{\sphinxupquote{circle}}, \sphinxcode{\sphinxupquote{weight}}, \sphinxcode{\sphinxupquote{massweight}}, \sphinxcode{\sphinxupquote{minimize}}, \sphinxcode{\sphinxupquote{maximize}}, \sphinxcode{\sphinxupquote{cubicfit}}, \sphinxcode{\sphinxupquote{energy}}, \sphinxcode{\sphinxupquote{de0}}, \sphinxcode{\sphinxupquote{fragment}}, \sphinxcode{\sphinxupquote{noise}}, \sphinxcode{\sphinxupquote{state}}
\\
\sphinxhline
\sphinxAtStartPar
{\hyperref[\detokenize{reference/d2edr2:group-d2edr2}]{\sphinxcrossref{\DUrole{std}{\DUrole{std-ref}{\$d2edr2}}}}}
&
\sphinxAtStartPar
{\hyperref[\detokenize{reference/d2edr2:d2edr2-length}]{\sphinxcrossref{\DUrole{std}{\DUrole{std-ref}{length}}}}}, {\hyperref[\detokenize{reference/d2edr2:d2edr2-steps}]{\sphinxcrossref{\DUrole{std}{\DUrole{std-ref}{steps}}}}}, {\hyperref[\detokenize{reference/d2edr2:d2edr2-displace}]{\sphinxcrossref{\DUrole{std}{\DUrole{std-ref}{displace}}}}}
\\
\sphinxhline
\sphinxAtStartPar
{\hyperref[\detokenize{reference/thermo:group-thermo}]{\sphinxcrossref{\DUrole{std}{\DUrole{std-ref}{\$thermo}}}}}
&
\sphinxAtStartPar
{\hyperref[\detokenize{reference/thermo:thermo-sigma}]{\sphinxcrossref{\DUrole{std}{\DUrole{std-ref}{sigma}}}}}, {\hyperref[\detokenize{reference/thermo:thermo-temperature}]{\sphinxcrossref{\DUrole{std}{\DUrole{std-ref}{temperature}}}}}, {\hyperref[\detokenize{reference/thermo:thermo-mass}]{\sphinxcrossref{\DUrole{std}{\DUrole{std-ref}{mass}}}}}, \sphinxcode{\sphinxupquote{project}}, \sphinxcode{\sphinxupquote{localize}}, \sphinxcode{\sphinxupquote{diagonalize}}
\\
\sphinxhline
\sphinxAtStartPar
{\hyperref[\detokenize{reference/nmr:group-nmr}]{\sphinxcrossref{\DUrole{std}{\DUrole{std-ref}{\$nmr}}}}}
&
\sphinxAtStartPar
{\hyperref[\detokenize{reference/nmr:nmr-standard}]{\sphinxcrossref{\DUrole{std}{\DUrole{std-ref}{standard}}}}}, {\hyperref[\detokenize{reference/nmr:nmr-points}]{\sphinxcrossref{\DUrole{std}{\DUrole{std-ref}{points}}}}}
\\
\sphinxhline
\sphinxAtStartPar
{\hyperref[\detokenize{reference/Energy:group-energy}]{\sphinxcrossref{\DUrole{std}{\DUrole{std-ref}{\$Energy}}}}}
&
\sphinxAtStartPar
{\hyperref[\detokenize{reference/Energy:energy-e}]{\sphinxcrossref{\DUrole{std}{\DUrole{std-ref}{e}}}}}, {\hyperref[\detokenize{reference/Energy:energy-g}]{\sphinxcrossref{\DUrole{std}{\DUrole{std-ref}{g}}}}}, {\hyperref[\detokenize{reference/Energy:energy-h}]{\sphinxcrossref{\DUrole{std}{\DUrole{std-ref}{h}}}}}, {\hyperref[\detokenize{reference/Energy:energy-d}]{\sphinxcrossref{\DUrole{std}{\DUrole{std-ref}{d}}}}}, {\hyperref[\detokenize{reference/Energy:energy-a}]{\sphinxcrossref{\DUrole{std}{\DUrole{std-ref}{a}}}}}, {\hyperref[\detokenize{reference/Energy:energy-dx}]{\sphinxcrossref{\DUrole{std}{\DUrole{std-ref}{dx}}}}}, \sphinxcode{\sphinxupquote{ax}}, \sphinxcode{\sphinxupquote{r}}
\\
\sphinxhline
\sphinxAtStartPar
{\hyperref[\detokenize{reference/uniq:group-uniq}]{\sphinxcrossref{\DUrole{std}{\DUrole{std-ref}{\$uniq}}}}}
&
\sphinxAtStartPar
{\hyperref[\detokenize{reference/uniq:uniq-distance}]{\sphinxcrossref{\DUrole{std}{\DUrole{std-ref}{distance}}}}}, \sphinxcode{\sphinxupquote{file}}
\\
\sphinxhline
\sphinxAtStartPar
{\hyperref[\detokenize{reference/dynamic:group-dynamic}]{\sphinxcrossref{\DUrole{std}{\DUrole{std-ref}{\$dynamic}}}}}
&
\sphinxAtStartPar
\sphinxcode{\sphinxupquote{t}}, {\hyperref[\detokenize{reference/dynamic:dynamic-e}]{\sphinxcrossref{\DUrole{std}{\DUrole{std-ref}{e}}}}}, {\hyperref[\detokenize{reference/dynamic:dynamic-n}]{\sphinxcrossref{\DUrole{std}{\DUrole{std-ref}{n}}}}}, {\hyperref[\detokenize{reference/dynamic:dynamic-m}]{\sphinxcrossref{\DUrole{std}{\DUrole{std-ref}{m}}}}}, \sphinxcode{\sphinxupquote{a}}, \sphinxcode{\sphinxupquote{b}}, \sphinxcode{\sphinxupquote{c}}, {\hyperref[\detokenize{reference/dynamic:dynamic-i}]{\sphinxcrossref{\DUrole{std}{\DUrole{std-ref}{i}}}}}, \sphinxcode{\sphinxupquote{o}}, \sphinxcode{\sphinxupquote{z}}, \sphinxcode{\sphinxupquote{y}}, \sphinxcode{\sphinxupquote{dx}}, \sphinxcode{\sphinxupquote{kd}}, \sphinxcode{\sphinxupquote{lw}}, \sphinxcode{\sphinxupquote{hb}}, \sphinxcode{\sphinxupquote{rb}}, \sphinxcode{\sphinxupquote{pb}}, \sphinxcode{\sphinxupquote{v1}}, \sphinxcode{\sphinxupquote{frag}}, {\hyperref[\detokenize{reference/dynamic:dynamic-file}]{\sphinxcrossref{\DUrole{std}{\DUrole{std-ref}{file}}}}}, \sphinxcode{\sphinxupquote{fprint}}
\\
\sphinxbottomrule
\end{tabulary}
\sphinxtableafterendhook\par
\sphinxattableend\end{savenotes}

\sphinxAtStartPar
\sphinxcode{\sphinxupquote{\$molecule}} is read differently from the other groups, so its keywords (\sphinxcode{\sphinxupquote{charge}}, \sphinxcode{\sphinxupquote{mult}}, \sphinxcode{\sphinxupquote{cartesian}}, \sphinxcode{\sphinxupquote{z\sphinxhyphen{}matrix}},
\sphinxcode{\sphinxupquote{set}}, \sphinxcode{\sphinxupquote{unit}}, \sphinxcode{\sphinxupquote{sym}}) could not be listed from the program; see {\hyperref[\detokenize{reference/molecule:group-molecule}]{\sphinxcrossref{\DUrole{std}{\DUrole{std-ref}{\$molecule}}}}}.
\sphinxcode{\sphinxupquote{\$energy}} is also read as \sphinxcode{\sphinxupquote{\$energy001}}, \sphinxcode{\sphinxupquote{\$energy002}} and so on, for tasks with several structures.




\section{Groups of undocumented tasks}
\label{\detokenize{binary-keywords:groups-of-undocumented-tasks}}
\sphinxAtStartPar
These groups belong to tasks and methods that no source describes. They are listed only so that a name met in an
example input can be recognized.


\begin{savenotes}\sphinxattablestart
\sphinxthistablewithglobalstyle
\centering
\begin{tabulary}{\linewidth}[t]{TTT}
\sphinxtoprule
\sphinxstyletheadfamily 
\sphinxAtStartPar
Group
&\sphinxstyletheadfamily 
\sphinxAtStartPar
Keywords
&\sphinxstyletheadfamily 
\sphinxAtStartPar
Task or method
\\
\sphinxmidrule
\sphinxtableatstartofbodyhook
\sphinxAtStartPar
\sphinxcode{\sphinxupquote{\$mcscf}}
&
\sphinxAtStartPar
\sphinxcode{\sphinxupquote{core}}, \sphinxcode{\sphinxupquote{occupied}}, \sphinxcode{\sphinxupquote{virtual}}, \sphinxcode{\sphinxupquote{excitation}}, \sphinxcode{\sphinxupquote{program}}, \sphinxcode{\sphinxupquote{convergence}}, \sphinxcode{\sphinxupquote{iterations}}, \sphinxcode{\sphinxupquote{diagonalize}}
&
\sphinxAtStartPar
MCSCF methods
\\
\sphinxhline
\sphinxAtStartPar
\sphinxcode{\sphinxupquote{\$ci}}
&
\sphinxAtStartPar
\sphinxcode{\sphinxupquote{iterations}}, \sphinxcode{\sphinxupquote{ndiis}}, \sphinxcode{\sphinxupquote{ecut}}, \sphinxcode{\sphinxupquote{v4}}
&
\sphinxAtStartPar
CI methods
\\
\sphinxhline
\sphinxAtStartPar
\sphinxcode{\sphinxupquote{\$espmap}}
&
\sphinxAtStartPar
\sphinxcode{\sphinxupquote{d}}, \sphinxcode{\sphinxupquote{n}}
&
\sphinxAtStartPar
electrostatic\sphinxhyphen{}potential map
\\
\sphinxhline
\sphinxAtStartPar
\sphinxcode{\sphinxupquote{\$optimize}} (second reader)
&
\sphinxAtStartPar
\sphinxcode{\sphinxupquote{function}}, \sphinxcode{\sphinxupquote{model}}, \sphinxcode{\sphinxupquote{h0}}, \sphinxcode{\sphinxupquote{delta}}, \sphinxcode{\sphinxupquote{fragment}}, \sphinxcode{\sphinxupquote{steps}}, \sphinxcode{\sphinxupquote{trust}}, \sphinxcode{\sphinxupquote{tolerance}}, \sphinxcode{\sphinxupquote{print}}, \sphinxcode{\sphinxupquote{fix}}, \sphinxcode{\sphinxupquote{value}}
&
\sphinxAtStartPar
\sphinxcode{\sphinxupquote{task=opt\_n1}}
\\
\sphinxhline
\sphinxAtStartPar
\sphinxcode{\sphinxupquote{\$path}}
&
\sphinxAtStartPar
\sphinxcode{\sphinxupquote{distance}}, \sphinxcode{\sphinxupquote{points}}, \sphinxcode{\sphinxupquote{spring}}, \sphinxcode{\sphinxupquote{steps}}, \sphinxcode{\sphinxupquote{trust}}, \sphinxcode{\sphinxupquote{tolerance}}, \sphinxcode{\sphinxupquote{tension}}
&
\sphinxAtStartPar
\sphinxcode{\sphinxupquote{task=path}}, \sphinxcode{\sphinxupquote{path0}}
\\
\sphinxhline
\sphinxAtStartPar
\sphinxcode{\sphinxupquote{\$chain}}
&
\sphinxAtStartPar
\sphinxcode{\sphinxupquote{points}}
&
\sphinxAtStartPar
\sphinxcode{\sphinxupquote{task=chain}}
\\
\sphinxhline
\sphinxAtStartPar
\sphinxcode{\sphinxupquote{\$walk}}
&
\sphinxAtStartPar
\sphinxcode{\sphinxupquote{n}}, \sphinxcode{\sphinxupquote{m}}, \sphinxcode{\sphinxupquote{e}}, \sphinxcode{\sphinxupquote{i}}, \sphinxcode{\sphinxupquote{b}}, \sphinxcode{\sphinxupquote{dx}}, \sphinxcode{\sphinxupquote{lw}}, \sphinxcode{\sphinxupquote{hb}}, \sphinxcode{\sphinxupquote{v1}}, \sphinxcode{\sphinxupquote{fp}}
&
\sphinxAtStartPar
\sphinxcode{\sphinxupquote{task=walk}}
\\
\sphinxhline
\sphinxAtStartPar
\sphinxcode{\sphinxupquote{\$orient}}, \sphinxcode{\sphinxupquote{\$mol1}}, \sphinxcode{\sphinxupquote{\$mol2}}
&
\sphinxAtStartPar
orient: \sphinxcode{\sphinxupquote{r}}, \sphinxcode{\sphinxupquote{lz}}, \sphinxcode{\sphinxupquote{nr}}, \sphinxcode{\sphinxupquote{e}}, \sphinxcode{\sphinxupquote{d0}}, \sphinxcode{\sphinxupquote{amax}}, \sphinxcode{\sphinxupquote{set}}, \sphinxcode{\sphinxupquote{list}}, \sphinxcode{\sphinxupquote{path}}, \sphinxcode{\sphinxupquote{file}}, \sphinxcode{\sphinxupquote{text}}, \sphinxcode{\sphinxupquote{print}}; mol1, mol2: \sphinxcode{\sphinxupquote{e}}, \sphinxcode{\sphinxupquote{z}}, \sphinxcode{\sphinxupquote{x}}, \sphinxcode{\sphinxupquote{q}}, \sphinxcode{\sphinxupquote{m}}, \sphinxcode{\sphinxupquote{n}}
&
\sphinxAtStartPar
\sphinxcode{\sphinxupquote{task=orient}}
\\
\sphinxhline
\sphinxAtStartPar
\sphinxcode{\sphinxupquote{\$press}}
&
\sphinxAtStartPar
\sphinxcode{\sphinxupquote{nadd}}, \sphinxcode{\sphinxupquote{radius}}, \sphinxcode{\sphinxupquote{steps}}, \sphinxcode{\sphinxupquote{trust}}, \sphinxcode{\sphinxupquote{tolerance}}, \sphinxcode{\sphinxupquote{sign}}, \sphinxcode{\sphinxupquote{energy}}, \sphinxcode{\sphinxupquote{vector}}
&
\sphinxAtStartPar
\sphinxcode{\sphinxupquote{task=press}}
\\
\sphinxhline
\sphinxAtStartPar
\sphinxcode{\sphinxupquote{\$coord}}
&
\sphinxAtStartPar
\sphinxcode{\sphinxupquote{index}}
&
\sphinxAtStartPar
\sphinxcode{\sphinxupquote{task=forceconst}}
\\
\sphinxhline
\sphinxAtStartPar
\sphinxcode{\sphinxupquote{\$steric}}
&
\sphinxAtStartPar
\sphinxcode{\sphinxupquote{g}}, \sphinxcode{\sphinxupquote{v}}, \sphinxcode{\sphinxupquote{w}}, \sphinxcode{\sphinxupquote{n}}, \sphinxcode{\sphinxupquote{r}}, \sphinxcode{\sphinxupquote{e}}, \sphinxcode{\sphinxupquote{i}}
&
\sphinxAtStartPar
\sphinxcode{\sphinxupquote{task=steric}}
\\
\sphinxhline
\sphinxAtStartPar
\sphinxcode{\sphinxupquote{\$dimer}}
&
\sphinxAtStartPar
\sphinxcode{\sphinxupquote{job}}, \sphinxcode{\sphinxupquote{order}}, \sphinxcode{\sphinxupquote{distance}}, \sphinxcode{\sphinxupquote{radius}}, \sphinxcode{\sphinxupquote{step}}, \sphinxcode{\sphinxupquote{trust}}, \sphinxcode{\sphinxupquote{gradient}}, \sphinxcode{\sphinxupquote{energy}}, \sphinxcode{\sphinxupquote{iterations}}, \sphinxcode{\sphinxupquote{tolerance}}, \sphinxcode{\sphinxupquote{delta}}, \sphinxcode{\sphinxupquote{analytic}}, \sphinxcode{\sphinxupquote{file}}, \sphinxcode{\sphinxupquote{print}}
&
\sphinxAtStartPar
\sphinxcode{\sphinxupquote{task=dimer}}
\\
\sphinxhline
\sphinxAtStartPar
\sphinxcode{\sphinxupquote{\$anchor}}, \sphinxcode{\sphinxupquote{\$confine}}
&
\sphinxAtStartPar
anchor: \sphinxcode{\sphinxupquote{u}}, \sphinxcode{\sphinxupquote{k}} and two more; confine: names not separated
&
\sphinxAtStartPar
restraints in molecular dynamics
\\
\sphinxhline
\sphinxAtStartPar
\sphinxcode{\sphinxupquote{\$surface}}
&
\sphinxAtStartPar
\sphinxcode{\sphinxupquote{vx}}, \sphinxcode{\sphinxupquote{x}}, \sphinxcode{\sphinxupquote{nx}}, \sphinxcode{\sphinxupquote{mv}}, \sphinxcode{\sphinxupquote{r}}, \sphinxcode{\sphinxupquote{h}}, \sphinxcode{\sphinxupquote{vo}}, \sphinxcode{\sphinxupquote{nv}}
&
\sphinxAtStartPar
\sphinxcode{\sphinxupquote{task=anharmonic}}
\\
\sphinxhline
\sphinxAtStartPar
\sphinxcode{\sphinxupquote{\$fmatrix}}, \sphinxcode{\sphinxupquote{\$ce}}, \sphinxcode{\sphinxupquote{\$ca}}, \sphinxcode{\sphinxupquote{\$cb}}
&
\sphinxAtStartPar
fmatrix: \sphinxcode{\sphinxupquote{n}}, \sphinxcode{\sphinxupquote{f}}, \sphinxcode{\sphinxupquote{s}}, \sphinxcode{\sphinxupquote{fa}}, \sphinxcode{\sphinxupquote{fb}}; ce: \sphinxcode{\sphinxupquote{n}}, \sphinxcode{\sphinxupquote{c}}, \sphinxcode{\sphinxupquote{e}}, \sphinxcode{\sphinxupquote{ca}}, \sphinxcode{\sphinxupquote{cb}}, \sphinxcode{\sphinxupquote{ea}}, \sphinxcode{\sphinxupquote{eb}}; ca, cb: \sphinxcode{\sphinxupquote{n}}, \sphinxcode{\sphinxupquote{c}}, \sphinxcode{\sphinxupquote{e}}
&
\sphinxAtStartPar
starting orbitals given in the input
\\
\sphinxhline
\sphinxAtStartPar
\sphinxcode{\sphinxupquote{\$scf}} (QM model)
&
\sphinxAtStartPar
\sphinxcode{\sphinxupquote{restrict}}, \sphinxcode{\sphinxupquote{convergence}}, \sphinxcode{\sphinxupquote{iterations}}, \sphinxcode{\sphinxupquote{direct}}, \sphinxcode{\sphinxupquote{print}}, \sphinxcode{\sphinxupquote{mc}}, \sphinxcode{\sphinxupquote{amix}}, \sphinxcode{\sphinxupquote{ijmix}}, \sphinxcode{\sphinxupquote{cut}}, \sphinxcode{\sphinxupquote{nf}}, \sphinxcode{\sphinxupquote{nmax}}, \sphinxcode{\sphinxupquote{nb}}, \sphinxcode{\sphinxupquote{bond}}, \sphinxcode{\sphinxupquote{penalty}}, \sphinxcode{\sphinxupquote{distance}}, \sphinxcode{\sphinxupquote{option}}
&
\sphinxAtStartPar
\sphinxcode{\sphinxupquote{theory=qm\_n*}}
\\
\sphinxhline
\sphinxAtStartPar
\sphinxcode{\sphinxupquote{\$isogrid}}
&
\sphinxAtStartPar
\sphinxcode{\sphinxupquote{d}}, \sphinxcode{\sphinxupquote{s}}, \sphinxcode{\sphinxupquote{c}}, \sphinxcode{\sphinxupquote{r}}, \sphinxcode{\sphinxupquote{m}}, \sphinxcode{\sphinxupquote{n}}, \sphinxcode{\sphinxupquote{e}}, \sphinxcode{\sphinxupquote{p}}, \sphinxcode{\sphinxupquote{iter}}
&
\sphinxAtStartPar
\sphinxcode{\sphinxupquote{task=isogrid}}
\\
\sphinxbottomrule
\end{tabulary}
\sphinxtableafterendhook\par
\sphinxattableend\end{savenotes}

\sphinxAtStartPar
Further groups (\sphinxcode{\sphinxupquote{\$molbf}}, \sphinxcode{\sphinxupquote{\$moldf}}, \sphinxcode{\sphinxupquote{\$mbf}}, \sphinxcode{\sphinxupquote{\$mdf}}, \sphinxcode{\sphinxupquote{\$esp}}, \sphinxcode{\sphinxupquote{\$bond}}, \sphinxcode{\sphinxupquote{\$e2r6}}, \sphinxcode{\sphinxupquote{\$geometry}}, \sphinxcode{\sphinxupquote{\$dmtype}}, \sphinxcode{\sphinxupquote{\$mm}}, \sphinxcode{\sphinxupquote{\$order}},
\sphinxcode{\sphinxupquote{\$bf}}, \sphinxcode{\sphinxupquote{\$c6p}}, \sphinxcode{\sphinxupquote{\$hferror}}, \sphinxcode{\sphinxupquote{\$mpitest}}, \sphinxcode{\sphinxupquote{\$multiply}}, \sphinxcode{\sphinxupquote{\$gtest}}) serve the parametrization of the model Hamiltonians or
the developer’s own tests.




\section{Task names}
\label{\detokenize{binary-keywords:task-names}}
\sphinxAtStartPar
The program recognizes these values of {\hyperref[\detokenize{reference/control:control-task}]{\sphinxcrossref{\DUrole{std}{\DUrole{std-ref}{task}}}}}. Only the first group is documented or tested.
\begin{itemize}
\item {} 
\sphinxAtStartPar
Documented: \sphinxcode{\sphinxupquote{energy}}, \sphinxcode{\sphinxupquote{gradient}}, \sphinxcode{\sphinxupquote{hessian}}, \sphinxcode{\sphinxupquote{optimize}}, \sphinxcode{\sphinxupquote{optimize+hessian}}, \sphinxcode{\sphinxupquote{irc}}, \sphinxcode{\sphinxupquote{scan}}, \sphinxcode{\sphinxupquote{nmr}}, \sphinxcode{\sphinxupquote{dipole}},
\sphinxcode{\sphinxupquote{anharmonic}}, \sphinxcode{\sphinxupquote{thermostat}}, \sphinxcode{\sphinxupquote{uniq}}, \sphinxcode{\sphinxupquote{opt\_n1}}, \sphinxcode{\sphinxupquote{saddle}}, \sphinxcode{\sphinxupquote{transit}}, \sphinxcode{\sphinxupquote{extremal}}, \sphinxcode{\sphinxupquote{conical}}, \sphinxcode{\sphinxupquote{path}}.

\item {} 
\sphinxAtStartPar
Not documented: \sphinxcode{\sphinxupquote{irc4}}, \sphinxcode{\sphinxupquote{dipole1}}, \sphinxcode{\sphinxupquote{dipole1x}}, \sphinxcode{\sphinxupquote{dipole2}}, \sphinxcode{\sphinxupquote{polarizability}}, \sphinxcode{\sphinxupquote{numeric}}, \sphinxcode{\sphinxupquote{loop}}, \sphinxcode{\sphinxupquote{stretch}},
\sphinxcode{\sphinxupquote{extremal2}}, \sphinxcode{\sphinxupquote{path0}}, \sphinxcode{\sphinxupquote{pathr}}, \sphinxcode{\sphinxupquote{pathx}}, \sphinxcode{\sphinxupquote{chain}}, \sphinxcode{\sphinxupquote{press}}, \sphinxcode{\sphinxupquote{forceconst}}, \sphinxcode{\sphinxupquote{potential}}, \sphinxcode{\sphinxupquote{dynamic}}, \sphinxcode{\sphinxupquote{dynamicn2}},
\sphinxcode{\sphinxupquote{dynamic1p}}, \sphinxcode{\sphinxupquote{walk}}, \sphinxcode{\sphinxupquote{mm0dimer}}, \sphinxcode{\sphinxupquote{dimer}}, \sphinxcode{\sphinxupquote{orient}}, \sphinxcode{\sphinxupquote{guess}}, \sphinxcode{\sphinxupquote{fragment\_guess}}, \sphinxcode{\sphinxupquote{nss\_num}}, \sphinxcode{\sphinxupquote{steric}}, \sphinxcode{\sphinxupquote{compare}},
\sphinxcode{\sphinxupquote{isogrid}}, \sphinxcode{\sphinxupquote{functional}}, \sphinxcode{\sphinxupquote{d1fit}}, \sphinxcode{\sphinxupquote{df1\_n4}}, \sphinxcode{\sphinxupquote{df2\_n4}}, \sphinxcode{\sphinxupquote{df2\_e\_n4}}, \sphinxcode{\sphinxupquote{mbf2}}, \sphinxcode{\sphinxupquote{molbf}}, \sphinxcode{\sphinxupquote{moldf}}, \sphinxcode{\sphinxupquote{dmtype}},
\sphinxcode{\sphinxupquote{am\_integral}}, and developer tests (\sphinxcode{\sphinxupquote{mpitest}}, \sphinxcode{\sphinxupquote{multiply}}, \sphinxcode{\sphinxupquote{etest}}, \sphinxcode{\sphinxupquote{fmtest}}, \sphinxcode{\sphinxupquote{ritest}}, \sphinxcode{\sphinxupquote{riopt}}, \sphinxcode{\sphinxupquote{test\_v}}).

\end{itemize}

\sphinxAtStartPar
\sphinxcode{\sphinxupquote{check}} is not a task name. Like any unknown task name, \sphinxcode{\sphinxupquote{task=check}} sets up the run and stops quietly, which is
why it seems to work. An unknown \sphinxcode{\sphinxupquote{theory=}} name is worse: the program runs Hartree\textendash{}Fock without a message.




\section{Theory names}
\label{\detokenize{binary-keywords:theory-names}}
\sphinxAtStartPar
Besides the methods described in {\hyperref[\detokenize{reference/control:control-theory}]{\sphinxcrossref{\DUrole{std}{\DUrole{std-ref}{theory}}}}}, the program contains these method names, none of
which is documented: \sphinxcode{\sphinxupquote{OSDFT}}, \sphinxcode{\sphinxupquote{DFT\_CNO}}, \sphinxcode{\sphinxupquote{LEDFT}}, \sphinxcode{\sphinxupquote{LE1DFT}}, \sphinxcode{\sphinxupquote{DFCIS}}, \sphinxcode{\sphinxupquote{RIHFDF}}, \sphinxcode{\sphinxupquote{MP4D}}, \sphinxcode{\sphinxupquote{PT4DQ}}, \sphinxcode{\sphinxupquote{PT4SDQ}}, \sphinxcode{\sphinxupquote{PT4SDTQ}},
\sphinxcode{\sphinxupquote{CISD}}, \sphinxcode{\sphinxupquote{CCD4S}}, \sphinxcode{\sphinxupquote{CCD4ST}}, \sphinxcode{\sphinxupquote{CCSDSD}}, \sphinxcode{\sphinxupquote{CCD\&MP2}}, \sphinxcode{\sphinxupquote{LCCD+MP2}}, \sphinxcode{\sphinxupquote{SN\_CCD\&MP2}}, \sphinxcode{\sphinxupquote{SN\_CCSD\&MP2}}, \sphinxcode{\sphinxupquote{LEHF}}, \sphinxcode{\sphinxupquote{LEHF\_N3}},
\sphinxcode{\sphinxupquote{LEMP2\_N3}}, \sphinxcode{\sphinxupquote{LE2HF\_N3}}, \sphinxcode{\sphinxupquote{LE2MP2\_N3}}, \sphinxcode{\sphinxupquote{MCSCF}}, \sphinxcode{\sphinxupquote{CASSCF}}, \sphinxcode{\sphinxupquote{GSCI}}, the model Hamiltonians \sphinxcode{\sphinxupquote{QM\_N2}}, \sphinxcode{\sphinxupquote{QM\_N2D}}, \sphinxcode{\sphinxupquote{QM\_N2B}},
\sphinxcode{\sphinxupquote{QM\_N2B3}}, \sphinxcode{\sphinxupquote{QM\_N2C}}, \sphinxcode{\sphinxupquote{QM\_N2C0}}, \sphinxcode{\sphinxupquote{QM\_N2F}}, \sphinxcode{\sphinxupquote{AM\_N1}}, \sphinxcode{\sphinxupquote{MM0\_N2}}, \sphinxcode{\sphinxupquote{DM\_1}}, \sphinxcode{\sphinxupquote{BM\_1}}, \sphinxcode{\sphinxupquote{MM\_1}}, and \sphinxcode{\sphinxupquote{6\sphinxhyphen{}12}}, \sphinxcode{\sphinxupquote{FREP}}, \sphinxcode{\sphinxupquote{SCRIPT}}.



\sphinxstepscope


\chapter{Sources}
\label{\detokenize{sources:sources}}\label{\detokenize{sources:id1}}\label{\detokenize{sources::doc}}
\sphinxAtStartPar
This manual is not an official Priroda document. It was not written or reviewed by the developer of Priroda,
Dimitri Laikov. It is compiled from the sources below, which are listed roughly in the order of how much weight the
manual gives them when they disagree. Each keyword entry in the {\hyperref[\detokenize{reference/index::doc}]{\sphinxcrossref{\DUrole{doc}{input reference}}}} names the
sources it rests on.


\section{Test runs}
\label{\detokenize{sources:test-runs}}\begin{description}
\sphinxlineitem{Version 25 test runs}
\sphinxAtStartPar
Small test inputs (water, ammonia, OH) run on Priroda 25 (2025.02.28) in October 2026 for this manual. Where they
confirm or contradict an older source, the manual follows the run and says so. The inputs and outputs are kept with
the manual’s source files.

\sphinxlineitem{Version 25 executable}
\sphinxAtStartPar
The Linux build \sphinxcode{\sphinxupquote{p4\_d}} from the official site, whose banner matches the test runs, was examined (not run) to list
every group, keyword, task and theory name it reads. This settles which names exist; it says nothing about what
they do. The result is on {\hyperref[\detokenize{binary-keywords:binary-keywords}]{\sphinxcrossref{\DUrole{std}{\DUrole{std-ref}{Keywords in the program}}}}}.

\end{description}


\section{Program files and text from the developer}
\label{\detokenize{sources:program-files-and-text-from-the-developer}}\begin{description}
\sphinxlineitem{Official file index}
\sphinxAtStartPar
\sphinxurl{https://rad.chem.msu.ru/~laikov/p\_w/}, read in October 2026 (page dated June 2025). A bare listing of executables,
example inputs and outputs, basis\sphinxhyphen{}set and parameter files, and the helper scripts \sphinxcode{\sphinxupquote{ge}}, \sphinxcode{\sphinxupquote{gm}} and \sphinxcode{\sphinxupquote{grx}}. It has no
prose documentation. Example outputs from this site are the strongest evidence for defaults and output formats.
Windows builds are in the \sphinxcode{\sphinxupquote{exe64/}} subdirectory.

\sphinxlineitem{Input Summary (\sphinxcode{\sphinxupquote{input.txt}})}
\sphinxAtStartPar
An English keyword list shipped with the program, apparently written by the developer, with user annotations
appended. It is the closest thing to an official keyword reference. A few of its Russian annotations are
corrupted and cannot be read.

\sphinxlineitem{Version 5 keyword reference (KNC wiki mirror)}
\sphinxAtStartPar
The developer’s “Input Summary” for version 5, mirrored on the wiki of the Kazan Scientific Centre computing centre
together with a beginner FAQ. The wiki is offline; pages were recovered from the Internet Archive
(captures 2008\textendash{}2012). A PDF copy of the documentation page is at
\sphinxurl{https://monitor.hpc.umanitoba.ca/doc/Priroda\_Documentation\_from\_KNCWiki.pdf}.

\sphinxlineitem{\sphinxcode{\sphinxupquote{Dft.rtf}}}
\sphinxAtStartPar
A short Russian keyword list with practical advice, from an older version. Some of its defaults differ from later
versions.

\sphinxlineitem{\sphinxcode{\sphinxupquote{Dft\_ansi.txt}}}
\sphinxAtStartPar
A Russian description of the input of Priroda’s predecessor. Its group names are obsolete; it is used here only
for history.

\end{description}


\section{User\sphinxhyphen{}written material}
\label{\detokenize{sources:user-written-material}}
\sphinxAtStartPar
Forum statements attributed to the developer reach these sources only second\sphinxhyphen{}hand, through users who relayed his
e\sphinxhyphen{}mails. The manual marks such statements as relayed.
\begin{description}
\sphinxlineitem{chemport.ru forum, thread 28803 (“Priroda”)}
\sphinxAtStartPar
\sphinxhref{https://chemport.ru/forum/viewtopic.php?f=71\&amp;t=28803}{https://chemport.ru/forum/viewtopic.php?f=71\&t=28803}. The main Russian\sphinxhyphen{}language user thread, 518 posts from
2009 to 2026. The largest single source for practical advice and known problems.

\sphinxlineitem{chemport.ru forum, thread 112860 (“Difficult cases with Priroda”)}
\sphinxAtStartPar
\sphinxhref{https://chemport.ru/forum/viewtopic.php?f=71\&amp;t=112860}{https://chemport.ru/forum/viewtopic.php?f=71\&t=112860}. Six posts (2015\textendash{}2016) on SCF divergence, RI\sphinxhyphen{}MP2 memory and
parallel hangs, with complete inputs and outputs.

\sphinxlineitem{qchem.ru forum, thread “Quantum\sphinxhyphen{}chemical program Priroda”}
\sphinxAtStartPar
The forum of the Ufa Quantum Chemistry Society, 2007\textendash{}2012. The site has been offline since about 2019; posts up to
August 2012 were recovered from the Internet Archive.

\sphinxlineitem{LIMOR page, Novosibirsk Institute of Organic Chemistry}
\sphinxAtStartPar
\sphinxurl{http://limor1.nioch.nsc.ru/priroda.html}. Example inputs, tables of basis\sphinxhyphen{}set names, and Perl helper scripts
(\sphinxcode{\sphinxupquote{pri}}, \sphinxcode{\sphinxupquote{priin}}, \sphinxcode{\sphinxupquote{pri2mol}} and others) written by users. Much of what the manual says about the output format is
inferred from how these scripts parse it.

\end{description}


\section{Published papers}
\label{\detokenize{sources:published-papers}}
\sphinxAtStartPar
The papers listed under {\hyperref[\detokenize{start/about:how-to-cite}]{\sphinxcrossref{\DUrole{std}{\DUrole{std-ref}{How to cite}}}}} describe the methods behind Priroda. They are cited where the manual
explains a method, not as documentation of the input.



\renewcommand{\indexname}{Index}
\printindex
\end{document}